%\arxivnumber{2401.05492}.
%\newpage
\section{Introduction} \label{intro}
The analysis presented in this chapter is based on Ref.~\cite{Bhattacharyya:2024aeq}.
The previous chapter focused on gravitational observables for an inspiralling compact binary, where it is typically well-justified to work in a non-relativistic (post-Newtonian) expansion. In this chapter, we instead study an astrophysical \emph{scattering} event: two compact objects (black holes, or black hole--neutron star systems) approach from asymptotic infinity, interact gravitationally, and re-separate on hyperbolic trajectories. Such burst-like encounters are particularly relevant for the emerging low-frequency gravitational-wave program. Future space-based detectors promise access to the milli-Hz band, while pulsar timing arrays probe nano-Hz frequencies, potentially constraining burst signals from hyperbolic scattering. Beyond their astrophysical content, these observations provide a new arena to test Einstein’s general relativity (GR) in a strong-field, dynamical regime.

From the theory side, extracting information from scattering (or near-scattering) events demands high-precision analytical control over both the conservative dynamics (classical potential) and the emitted radiation \cite{Purrer:2019jcp}. Conceptually, this task mirrors precision computations of scattering cross-sections in particle physics: one is ultimately computing observables derived from amplitudes, organized perturbatively, and matched onto classical physics. In this spirit, perturbative quantum-gravity methods—particularly on-shell and effective-field-theory tools—have proven remarkably efficient for isolating the \emph{classical} gravitational interaction of compact objects.

There is also strong motivation to extend these high-precision predictions beyond GR. Since the first direct detections of gravitational waves, there has been renewed interest in testing GR across all accessible length scales. While GR has been extraordinarily successful, it is widely expected to be incomplete: a consistent UV-complete quantum theory of gravity should exist, and GR does not by itself explain the observed late-time acceleration of the Universe (dark energy). A conservative viewpoint is therefore that GR should emerge as an effective theory within some broader framework. Two broad routes to modifying GR are (i) adding higher-curvature operators, and (ii) introducing extra gravitational degrees of freedom (DOF) \cite{stelle1,Alexeev:1996vs,Lehebel:2018zga,Volkov:2016ehx,Kunz:2006ca}. The simplest extension of the second type is scalar--tensor gravity, where a scalar field accompanies the massless spin-2 graviton \cite{Damour:1992we,Horbatsch:2015bua,Schon:2021pcv,Rainer:1996gw,DeFelice:2011bh}. Such a scalar is phenomenologically motivated: it can appear as an effective description of dark-sector physics, and scalar dynamics can contribute to cosmic acceleration \cite{DeFelice:2011bh,Gsponer:2021obj}. Most importantly for gravitational-wave phenomenology, an additional scalar DOF generically introduces extra radiative channels and polarizations, affecting both energy loss and waveform structure. In this chapter, our primary target is precisely this: to compute the scalar polarization content of radiation emitted during a scattering process.

Theoretical interest in scalar DOF also has a long history. Early unification attempts by Weyl and Kaluza \cite{weyl} initiated a line of ideas in which extra fields arise naturally from higher-dimensional or generalized geometric frameworks. In particular, Kaluza’s five-dimensional proposal (later developed as Kaluza--Klein theory) yields additional lower-dimensional fields upon compactification, and it helped motivate scalar--tensor models such as Jordan’s theory, where the scalar couples non-minimally to the tensor sector. These motivations, together with the observational prospects, make scattering in scalar--tensor gravity a natural testing ground for modern amplitude-based methods.

Classical gravitational scattering has been extensively studied in both post-Newtonian and post-Minkowskian regimes
\cite{hyp,hyp1,hyp2,hyp3,hyp4,hyp5,hyp6,hyp7,hyp8,hyp9,hyp10,hyp11,hyp12,hyp13,hyp16,hyp14,Damour:2016gwp,Bini:2017wfr,Bini:2017xzy,Damour:2017zjx,Damour:2019lcq,Bini:2020flp,Bini:2020uiq,Damour:2020tta,Bini:2020rzn,Bini:2021gat,Bini:2022enm,Damour:2022ybd,Bini:2022wrq,Rettegno:2023ghr,Bini:2023fiz,Ceresole:2023wxg}.
Motivated by the QFT perspective and catalyzed by Damour’s proposal \cite{Damour:2017zjx}, the amplitude community has developed a powerful toolkit for the two-body problem. On the EFT side, a worldline effective-field-theory approach based on the post-Minkowskian (PM) expansion—organizing results in powers of Newton’s constant while resumming PN effects at fixed PM order—was developed for conservative dynamics in \cite{Kalin:2020mvi}, and subsequently extended to include finite-size, spin, and tidal effects \cite{Bini:2020flp,Cheung:2020sdj,Kalin:2020lmz,Haddad:2020que}. Higher-PM developments appear in \cite{Kalin:2020fhe,Dlapa:2021npj,Dlapa:2021vgp,Kalin:2022hph,Jinno:2022sbr,Dlapa:2022lmu,Dlapa:2023hsl,Riva:2021vnj}. In parallel, more direct on-shell amplitude techniques have been used to compute classical gravitational observables
\cite{PhysRevD.7.2317,Holstein:2004dn,Neill:2013wsa,Bjerrum-Bohr:2013bxa,Luna:2017dtq,Bjerrum-Bohr:2018xdl,Kosower:2018adc,Cristofoli:2021vyo,DeAngelis:2023lvf,Brandhuber:2023hhl,Aoude:2023dui,Brandhuber:2023hhy,Georgoudis:2023eke,Herderschee:2023fxh}\footnote{These references are by no means exhaustive; see the citations collected in \cite{Buonanno:2022pgc} for a broader entry point.},
including the 2PM and 3PM conservative potential \cite{Bjerrum-Bohr:2018xdl,Cheung:2018wkq,Cristofoli:2019neg,Cheung:2020gyp,Bern:2019nnu,Bern:2019crd}\footnote{In the context of soft theorems, two-body gravitational waveforms up to 3PM order have been discussed in \cite{ashoke,ashoke1,ashoke2,ashoke3,ashoke4,ashoke5,ashoke6}.}.

A natural question then arises: why do the worldline-EFT and on-shell-amplitude approaches agree so precisely? This gap was bridged by Plefka \emph{et al.} in \cite{Mogull:2020sak}. They showed that the Feynman--Schwinger representation of a graviton-dressed scalar propagator provides the key link between (i) scattering amplitudes of massive scalars and (ii) operator expectation values in \emph{Worldline Quantum Field Theory} (WQFT). WQFT is closely related to worldline EFT (WEFT), but differs in that the worldline degrees of freedom—most notably the worldline fluctuations—are quantized. This framework has enabled detailed cross-checks and new computations: first for the non-spinning case \cite{Jakobsen:2021smu}, then for spinning systems \cite{Jakobsen:2021lvp,Jakobsen:2021zvh,Jakobsen:2023ndj,Jakobsen:2023hig}, and further into higher-PM orders including tidal effects \cite{Jakobsen:2023pvx}. WQFT has also been applied to gravitational lensing \cite{Bastianelli:2021nbs}, and to aspects of the classical double copy \cite{Shi:2021qsb, Diaz-Jaramillo:2021wtl,Comberiati:2022cpm}.

\textit{Most existing scattering analyses focus on GR. In this chapter, we take a step beyond GR by computing scattering observables in a modified theory: a massive scalar--tensor model with a non-trivial cubic and quartic scalar potential. Our broader goal is to assess whether gravitational-wave observations can constrain not only the scalar mass, but also the cubic ($\lambda_3$) and quartic ($\lambda_4$) self-couplings.} As a first step, we compute two central observables in a scattering event: the \emph{impulse} and the \emph{waveform}. The impulse directly determines the scattering angle and interfaces naturally with bound-state information through the EOB map \cite{Buonanno:1998gg}. The waveform controls the radiative sector and ultimately the gravitational-wave phase. Technically, the main challenge in the WQFT computation of radiation is the evaluation of nontrivial worldline integrals. Here we make explicit progress: we compute a set of massive waveform integrals and develop analytic approximation strategies. \textit{To the best of our knowledge, these integral results and the associated methods are new.}

The chapter is organized as follows. In Section~(\ref{ch2:sec2}) we review the WQFT formalism and summarize the key results of \cite{Mogull:2020sak}. In Section~(\ref{ch2:sec3}) we describe how these results are modified by the presence of an additional massive scalar, and we derive the complete set of worldline Feynman rules for scalar and graviton interactions. In Section~(\ref{sec1}) we compute the scalar-induced corrections to the impulse, including diagrams involving scalar self-interactions and scalar--graviton couplings. In Section~(\ref{ch2:sec5}) we construct the radiation integrands and compute the time-domain waveform in the massless-scalar limit. In Section~(\ref{sec6}) we treat the massive case: we explain how the radiation integrals are complicated by more involved phase factors, list the massive radiation integrands, and evaluate the relevant integrals using the stationary-phase method (accurate in the $|x|\to\infty$ regime), discussing subtleties along the way. Appendix~(\ref{ch2:app:A}) derives the worldline effective action with an extra scalar DOF. Appendix~(\ref{ch2:app:B}) explains the relation between impulse and the scattering amplitude via the eikonal phase. Appendices~(\ref{ch2:app:C}), (\ref{ch2:app:D}), (\ref{ch2:app:E}), (\ref{3.54kk}) and (\ref{App:3.H}) provide additional details about some integrals, while Appendix~(\ref{ch2:app:F}) presents an alternative approximation based on a large-velocity expansion. 
%\newpage
\subsection*{Notations and Conventions}
\begin{itemize}[leftmargin=1.55em,itemsep=0.42em,topsep=0.45em,label=\textcolor{brandburgundy}{\small\textbullet}]
\item Metric sign convention: $(+,-,-,-)$.
\item All computations are done in units where $(c,\hbar)=1$.
\item The impact parameter $b$ is purely spacelike, $b^2=b^\mu b_\mu=-|b|^2$ and $b^{\mu}=(0,0,|b|,0)$.
\item Black-hole velocity parametrization: $v_1=(\gamma,\gamma\beta,0,0)$ and $v_2=(1,0,0,0)$.
\item Planck mass: $m_p=\frac{1}{\sqrt{8\pi G_N}}$, where $G_N$ is Newton's constant.
\item Scaled delta function: $\hat\delta^{(D)}(\cdots)\equiv (2\pi)^{D}\delta^{(D)}(\cdots)$.
\item Incomplete beta function: $B_{z}(a,b):=\rmint_0^z dt\,t^{a-1}(1-t)^{b-1}$.
\item $J_n$ and $K_n$ denote the Bessel function of the first kind and the modified Bessel function of the second kind, respectively.
\item The \textbf{Meijer G} function is defined as
\begin{align}
    G_{p q}^{m n} \left(z,r\Big|
\begin{array}{c}
 a_1\text{...}.a_p \\
 b_1\text{...} b_q \\
\end{array}\right)
=\frac{r}{2 \pi  i}\rmint_{-i\infty}^{i\infty} \frac{   \Gamma  \left(b_1+r s\right) \Gamma\left(-a_n-r s+1\right) \Gamma  \left(-a_1-r s+1\right) \Gamma\left(b_m+r s\right)}{  \Gamma  \left(a_p+r s\right) \left(-b_q-r s+1\right) \Gamma  \left(a_{n+1}+r s\right) \Gamma  \left(-b_{m+1}-r s+1\right)}z^{-s}\,ds\,.\nonumber
\end{align}
\item The Bickley–Naylor function is defined as $ \textrm{Ki}_{n}(x)=\int_0^{\pi/2}  d\theta\,e^{-\frac{x}{\cos\theta}}\cos^{n-1}\theta$.
\end{itemize}
%\newpage 
\section{A quick tour to worldline quantum field theory}\label{ch2:sec2}
In this section, we briefly summarize the results of the Worldline Quantum Field Theory (WQFT) approach to the scattering problem. In \cite{Mogull:2020sak}, precise correspondence has been derived between the scalar-graviton S-matrix element and the one-point functions of the worldline operators. In this string-inspired approach, one maps black hole scattering in worldline theory to field scattering in the field theory approach. Using the gravitationally dressed Green's functions, one can write S-matrix elements as the expectation value of operators present in the worldline theory.

Now, for spinless black holes in GR, the action can be written in the EFT framework as\footnote{To get the linearized action one should expand $\sqrt{-g}$ as,
\begin{align}
    \begin{split}
        \sqrt{-g}=1+\frac{1}{2m_p} h-\frac{1}{4m_p^2} h^{\mu\nu}h_{\mu\nu}+\frac{1}{8m_p^2}h^2,\,\, h=\text{Tr}( h_{\mu\nu}).\nonumber
    \end{split}
\end{align}},
$$S=S_{EH}+S_{gf}+\sum_i S_{pm}^i\,,$$
where $S_{EH}$ is the usual Einstein-Hilbert action, and $S_{gf}$ is the gauge-fixing action, which in the weak field approximation is given by,

$$S_{gf}=\rmint d^D x\Big(\partial_\nu h^{\mu\nu}-\frac{1}{2}\partial^\mu h^{\nu}_{\nu}\Big)^2\,.$$

This imposes the de-Donder gauge condition. Now, for an extended object, the worldline action consists of the Wilson coefficients $c_v$ and $c_R$ as follows,
\begin{align}S_{pm}=-m\rmint d\tau +c_R\rmint d\tau R(x)+c_v\rmint d\tau R_{\mu\nu}(x)\dot{x}^\mu \dot{x}^\nu+\cdots\,.\end{align}
In our paper, we will drop the second and third terms above to remove the complicacy due to the finite-size effects. In principle, there will be an infinite tower of terms consisting of higher-order derivatives of the metric represented as dots. 




Now, assuming a fixed background, we can write Green's function of the field theory as a two-point correlator : 
\begin{align}
G_i(x,x')=\mathcal{Z}_i^{-1}\rmint \mathcal{D}[\phi_i]\phi_i(x)\phi_i^{\dagger}(x')e^{i S_i},
\end{align}
where, $S_i=\rmint d^D x \sqrt{-g}(g^{\mu\nu}\partial_\mu\phi_i^{\dagger}\partial_\nu\phi_i-m^2\phi_i^{\dagger}\phi_i-\zeta R\phi_i^{\dagger}\phi_i)$.\par
%\vspace{0.3 cm}
 To this end, we have to integrate out certain field degrees of freedom with respect to a relevant scale of the problem. This will give rise to the one-loop effective action, which can be represented by a worldline path integral \cite{Strassler:1992zr,Feal:2022iyn,Ahmadiniaz:2022yam}. Then, from this path integral, we can identify the point particle action in the background of $h_{\mu\nu}$.
Now, for our interest, we evaluate the S-matrix element of two scalars with or without a final state graviton in certain \textit{classical limit} as,
\begin{align}
&\langle \Omega|T\{h_{\mu\nu}\phi_1(x_1)\phi_1^{\dagger}(x_1')\phi_2(x_2)\phi_2^{\dagger}(x_2')\}|\Omega\rangle\\&
=\mathcal{Z}_i^{-1}\rmint \mathcal{D}[h_{\mu\nu},\phi_i]h_{\mu\nu}(x)\prod_{i=1}^2\phi_i(x_i)\phi_i^{\dagger}(x_i')e^{i S'}\,.
\end{align}
Fourier transforming this relation and using LSZ reduction in the $\hbar \rightarrow 0 $ limit we get \cite{Mogull:2020sak},
\begin{align}&\langle \phi_1\phi_2(+h)|S|\phi_1\phi_2\rangle\\&
=\mathcal{Z}^{-1}\rmint d^D[x_i,x_i',x]e^{ip\cdot x-ip_i'\cdot x_i'-ik\cdot x}\times \rmint [\mathcal{D}h_{\alpha\beta}]\,\epsilon^{\mu\nu}(k)h_{\mu\nu}(x)\prod_{i=1}^2 G_i(x_i,x_i')e^{i S_{EH}+S_{gf}}\bigg{|}_{\textrm{Amputated,connected}}\,.
\end{align}
The study of the S-matrices consists of applying the LSZ reduction. Cutting the propagators on external legs, we convert correlators into S-matrices and send the external legs to infinity, where they interact weakly. We put the scalar legs on-shell, and the integration over the finite time domain extends $\tau \in (-\infty, \infty)$.
The key relation connecting the QFT form factor to the WQFT correlator is given by \cite{Mogull:2020sak},
\begin{align}
\begin{split}
&\frac{\Xi(b,v;\{\epsilon^{l},k_{l}\})}{\Xi_0}=\delta\Bigg(\sum_{l=1}^N k_l\cdot v\Bigg) \exp\Big({\sum_{l=1}^N k_l\cdot b} \Big)\mathcal{F}(p,p'|\{\epsilon^{l},k_{l}\})
\end{split}
\end{align}
\newcommand{\chthreeinlinefeynwidth}{0.13\linewidth}
\newcommand{\chthreeleadfeynwidth}{0.32\linewidth}

where,
\begin{align}
  &  \Xi(b,v;\{\epsilon^{l},k_{l}\})=\rmint \mathcal{D}[x]\rmint \mathcal{D}[a,b,c]\exp[-i\rmint d\sigma \Big[\frac{1}{4}g_{\mu\nu}(\dot{x}^\mu\dot{x}^\nu+a^\mu a^\nu+b^\mu c^\nu)\Big].
\end{align}
Having the gravitationally dressed propagator in momentum space, we put the external scalar legs on-shell to perform the LSZ reduction. Effectively, it can be represented as,


\begin{center}
\includegraphics[width=\chthreeleadfeynwidth]{fey_1.png}
\end{center}
\begin{align}
\mathcal{F}(p,p'|\{\epsilon^{l},k_{l}\})=\langle p'|\prod_{i=1}^N\epsilon_i\cdot h(k_i)|p\rangle\,.
\end{align}
For details, we refer the reader to \cite{Mogull:2020sak}.
In the following section, we will discuss how to derive the Feynman rules for the scattering events of two black holes (including extra scalar modes), which are considered to be point-like objects and, hence, spinless.

\section{Derivation of the Feynman rules}\label{ch2:sec3}
One of the key ingredients to compute the observables in a scattering problem is the partition function. 
We compute the partition function for a binary system in the Scalar-Tensor theory of gravity with a generic scalar potential. The gravitational action is given by,
\begin{align}
    \begin{split}
        S_{g}=\rmint d^4x \sqrt{- g}\Big(-\frac{m_p^2}{2} R+\frac{1}{2} {g}^{\alpha\beta}\partial_{\alpha}{\varphi}\partial_{\beta}{\varphi}-\frac{1}{2}m^2{\varphi}^2-\frac{\lambda_3}{\textcolor{black}{3!}} m_p\, \varphi^3-\frac{\lambda_4}{\textcolor{black}{4!}}{\varphi}^4\Big)\in S_{EH}+S_{scalar}.\label{3.1m}
    \end{split}
\end{align}
In this theory, the mass of the point particle which is the avatar of the scalar field $\phi_i$, is not constant but rather depends on the extra scalar degree of freedom $\varphi$ \cite{eardley}. Hence the matter action has the following form,
\begin{align}
    \begin{split}
        S_m=\sum_{i=1}^2\rmint d^4 x \sqrt{-g}(g^{\mu\nu}\partial_\mu\phi_i^{\dagger}\partial_\nu\phi_i-m_i(\varphi)^2\phi_i^{\dagger}\phi_i-\zeta R\phi_i^{\dagger}\phi_i)\,.
    \end{split}
\end{align}
Therefore, the WQFT partition function takes the form \footnote{A short derivation of the worldline effective action in the presence of extra scalar degrees of freedom appearing in the partition function is given in Appendix~(\ref{ch2:app:E}).} ,
\begin{align}
    \begin{split}
        {Z}_{\textrm{WQFT}}=\mathcal{N}\times \rmint \mathcal{D}[h_{\mu\nu},\,{\varphi}]\rmint \mathcal{D}[x_i,a_i,b_i,c_i]e^{iS_{g}}\exp\Big\{-i\sum_{i=1}^{2}\rmint d\tau_{i}\frac{m_i(\varphi)}{2}g_{\mu\nu}(\dot x_{i}^{\mu}\dot x_{i}^{\nu}+a_{i}^{\mu}a_{i}^{\nu}+b_{i}^{\mu}c_{i}^{\nu})\Big\}\label{3.2}
    \end{split}
\end{align}
where $a_i,b_i$ s are Lee-Yang Ghost and can be ignored in the classical limit. Therefore, the $n$-body partition function (non-spinning) in WQFT can be written as,
\begin{align}
    \begin{split}
        Z_{\textrm{WQFT}}^{(n)}=\mathcal{N}\rmint \mathcal{D}h_{\mu\nu}\mathcal{D}\varphi\,e^{iS_{EH}+iS_{scalar}}\Big(\prod_{i=1}^{n}\mathcal{D}x_{i}\,e^{iS_{pm}^i}\Big)\,.\label{3.3}
    \end{split}
\end{align}
It has been noticed that the partition function is related to the Eikonal phase  $\chi$ by the following exponentiation \cite{Amati:1990xe, Mogull:2020sak}.
\begin{align}
    \begin{split}
Z_{\textrm{WQFT}}:=e^{i\chi}\,.
    \end{split}
\end{align}
Now $\chi$ can be calculated by computing the connected Feynman diagrams. In order to compute the path integral in \eqref{3.3} one could start by demanding that a gravitational field is a fluctuation over Minkowski space, i.e., one can decompose the metric $g_{\mu\nu}$ as,
\begin{align}
    \begin{split}
g_{\mu\nu}=\eta_{\mu\nu}+\frac{h_{\mu\nu}}{m_p}.
    \end{split}
\end{align}
and, eventually, the worldline degree of freedom can be expanded as,
\begin{align}
    \begin{split}
        x^{\mu}(\tau)=b^{\mu}+v^{\mu}\tau+z^{\mu}(\tau)
    \end{split}
\end{align}
where, $z^{\mu}(\tau)$ is the worldline fluctuation about the straightline geodesic. Once we have the well-defined partition function, one could compute the classical observables as a correlation function in WQFT.
\begin{align}
    \begin{split}
        \mathcal{O}(b_i,v_i):=\langle\hat{\mathcal{O}}(\hat{x},\hat{h}_{\mu\nu},\hat\varphi)\rangle=\frac{1}{Z^{(n)}_{\textrm{WQFT}}}\rmint \mathcal{D}h_{\mu\nu}\mathcal{D}\varphi\,e^{iS_{EH}+iS_{scalar}}\Big(\prod_{i=1}^{n}\mathcal{D}x_{i}e^{iS_{pm}^i}\Big)\mathcal{O}(x_i,h,\varphi)
    \end{split}
\end{align}
where the worldline action can be written in Polyakov form.
The point particle action has the following form,
\begin{align}
    \begin{split}
        S_{pm}=-\sum_{i=1}^{n}\rmint_{-\infty}^{\infty}d\tau_i\frac{m_i( \varphi)}{2}(g_{\mu\nu}\dot x^{\mu}_i\dot{x}^{\nu}_i+1).\label{3.4m}
    \end{split}
\end{align}
We intend to compute the two main observables, impulse and waveform, schematically defined as,
\begin{align}
    \begin{split}
      \langle \hat h_{\mu\nu}(k),\hat\varphi(k),\hat z^{\mu}(\omega)\rangle=Z_{\textrm{WQFT}}^{-1}\rmint\mathcal{D}h_{\mu\nu}\mathcal{D}\varphi\,e^{iS_{EH}+iS_{scalar}}\Big(\prod_{i=1}^{2}\mathcal{D}x_{i}e^{iS_{pm}^i}\Big)\{h_{\mu\nu}(k),\varphi(k),z^{\mu}(\omega)\}\,.
    \end{split}
\end{align}
Before going to the computations of the partition function, we first derive all the Feynman rules. In Scalar-Tensor theory, the violation of the equivalence principle automatically implies that the mass of the binaries depends on the extra gravitational polarisation ${\varphi}$. As there is a self-interaction term so one can expand $m_{a}(\varphi)$ around one of the chosen vacuums ${\varphi}_{0}$, which for our case is zero, as,
\textcolor{black}{
\begin{align}
    \begin{split}
        m_a({\varphi})=m_a(0)\Big \{1+s_a\frac{\varphi}{m_p}+g_a\frac{\varphi^2}{m_p^2}+\mathcal{O}(\varphi^3)\Big\}.
    \end{split}
\end{align}}
In order to derive the Feynman rules it would be better to decompose the fields  ($\chi(x)\in(h_{\mu\nu},\varphi)$) in their Fourier mode assuming the black hole worldline has following the fluctuation due to the radiation reaction: $x_{(i)}^{\mu}=b_{(i)}^{\mu}+v_{(i)}^{\mu}\tau+z_{(i)}^{\mu}(\tau)$. where $z_{i}^{\mu}(\tau)$ is the worldline fluctuations about the straight line geodesic. Now, we expand the point particle action at different orders in $z$ to get the worldline vertices. In our theory we have gravitational field ($ h_{\mu\nu}$) and scalar field (${\varphi}$). One can decompose the fields in Fourier modes, when it couples with worldline, as,
\begin{align}
    \begin{split}
        &  \chi(x_i)=\rmint_{k}e^{ik\cdot(b_i+v_i\tau+z_i)} \chi(-k)\,.
    \end{split}
\end{align}
Now, expanding the exponential at different orders in $z$ and is given by,
\begin{align}
    \begin{split}
         \chi (x_i)=\rmint_{k}e^{ik\cdot(b_i+v_i \tau)}\prod_{n=0}^{\infty}\frac{i^{n}}{n!}(k\cdot z)^{n}  \chi(-k)\,.\label{3.8m}
    \end{split}
\end{align}
Again, we decompose the worldline fluctuations as,
\begin{align}
    \begin{split}
        z^{\mu}(\tau)=\rmint_{\omega}e^{i\omega \tau}z^{\mu}(-\omega)\,.
    \end{split}\end{align}
From the quadratic part of Einstein-Hilbert action, one can identify the graviton propagator as,
\begin{align}
   \langle  h_{\mu\nu}(x_1) h_{\rho\sigma}(x_2)\rangle\equiv \thesisinlinefeynman[\chthreeinlinefeynwidth]{graviton_propagator_1.png}\hspace{0.45cm}=iP_{\mu\nu;\rho\sigma}\rmint_{k}\frac{e^{-ik\cdot(x_1-x_2)}}{k^2+i\epsilon}
\end{align}
where,
\begin{align}
    \begin{split} \label{tensor}
        P_{\mu\nu,\rho\sigma}=\textcolor{black}{\frac{1}{2}}\Big(\eta_{\mu\rho}\eta_{\nu\sigma}+\eta_{\mu\sigma}\eta_{\nu\rho}-\eta_{\mu\nu}\eta_{\rho\sigma}\Big).
    \end{split}
\end{align}
and propagator for the scalar degrees of freedom has the following form,
\begin{align}
    \begin{split}
        \langle \varphi(x_1)\varphi(x_2)\rangle\equiv \thesisinlinefeynman[\chthreeinlinefeynwidth]{scalar_propagator_1.png}\hspace{0.50cm}= i\rmint_{k}\frac{e^{-ik\cdot(x_1-x_2)}}{k^2-m^2+i\epsilon} \,.
    \end{split}
\end{align}
Apart from the field degrees of freedom, we have another dynamical degree of freedom: worldline fluctuations $z^{\mu}(\omega)$. One can identify the propagator by inspecting the quadratic part of the worldline action (\ref{3.4m}).
\begin{align}
    \begin{split}
        S_{\textrm{pm}}=-m(\varphi_{0})\rmint d\tau \Big[1+\eta_{\mu\nu}v^{\mu}\dot z^{\nu}+\frac{1}{2}\eta_{\mu\nu}\dot z^{\mu}\dot z^{\nu}\Big]\label{3.12m}
    \end{split}
\end{align}
The propagator of $z^{\mu}$ can be identified from the quadratic part of (\ref{3.12m}).
\begin{align}
    \begin{split}
      \Big  \langle z^{\mu}(\tau_1)z^{\nu}(\tau_2)\Big\rangle\equiv \thesisinlinefeynman[\chthreeinlinefeynwidth]{worldline_propagator.png}\hspace{0.35cm}=-i\frac{\eta^{\mu\nu}}{m}\rmint_{\omega}\frac{e^{-i \omega (\tau_1-\tau_2)}}{(\omega\pm i\epsilon)^2}\,.
    \end{split}
\end{align}
Now, we are in a position to derive the worldline Feynman rules. The key ingredient is the worldline action $S_{\textrm{pm}}$ in (\ref{3.4m}).
\begin{align}
    \begin{split}
        S_{\textrm{pm}}\Big|_{\textrm{int.}}\in \rmint d\tau \Big[\Big(-\frac{m_a}{2}-\frac{m_a s_a}{2m_p} \varphi-\frac{m_a g_a}{2m_p^2} \varphi^2\Big)(\eta_{\mu\nu}+\frac{1}{m_p}h_{\mu\nu})(v^{\mu}+\dot z^{\mu})(v^{\nu}+\dot z^{\nu})\Big]\,.
    \end{split}
\end{align}
We can rewrite (\ref{3.8m}) by inserting the Fourier mode of the worldline fluctuation in the following way.
\begin{align}
    \begin{split}
         \chi=\sum_{m=0}^{\infty}\frac{i^m}{m!}\rmint_{k,\omega_1,...,\omega_m}e^{ik\cdot b}e^{i(k\cdot v+\sum_{i=1}^{m} \omega_{i})\tau}\prod_{i=1}^{m}[k\cdot z(-\omega_i)] \chi(-k).\label{3.15m}
    \end{split}
\end{align}
The expansion in \eqref{3.15m} helps derive the Feynman rules up to $\mathcal{O}(z^{j})$. We now list the Feynman rules coming from the scalar-worldline interaction:
\vspace{-0.4 cm}
\subsection*{\underline{Vertex for scalar field:}}
{\scriptsize
%\begin{itemize}
    \textbullet $\,\,$ $\mathcal{O}{(z^0)}:- i\frac{m_a s_a}{2m_p} \,e^{ik\cdot b_a}\,\hat\delta(k\cdot v)\rightarrow  \thesisinlinefeynman[\chthreeinlinefeynwidth]{feynmann_2_1.png}$

 \textbullet $\,\,$ $\mathcal{O}(z):\frac{m_a s_a}{2m_p}\,e^{ik\cdot b}\,\hat{\delta}(k\cdot v+\omega)(2\omega v_{\rho}+k_{\rho})\rightarrow \thesisinlinefeynman[\chthreeinlinefeynwidth]{feynmann_3_1.png}$

 \textbullet $\,\,$ $ \mathcal{O}(z): -\frac{m_a s_a}{2m_p^2}e^{i(k_1+k_2)\cdot b_a}\hat\delta(k_1\cdot v_1+k_2\cdot v_1+\omega)\{(k_{1\rho}+k_{2\rho})v^{\mu}v^{\nu}+2\omega v^{(\mu}\delta_{\rho}^{\nu)}\}\rightarrow
\thesisinlinefeynman[\chthreeinlinefeynwidth]{phi_2_z_1.png}$

 \textbullet $\,\,$ $ \mathcal{O}(z): -\frac{m_a g_a}{2m_p^2}e^{i(k_1+k_2)\cdot b_a}\hat\delta(k_1\cdot v_1+k_2\cdot v_1+\omega)\{(k_{1\rho}+k_{2\rho})v^{\mu}v^{\nu}+2\omega v_\rho\}\rightarrow
\thesisinlinefeynman[\chthreeinlinefeynwidth]{new_10.png}$

   \textbullet $\,\,$ $\mathcal{O}{(z^2)}:i\frac{m_a s_a}{2m_p} \,e^{ik\cdot b_a}\,\hat\delta(k\cdot v+\omega_1+\omega_2)  (\frac{1}{2} k_{\rho_1} k_{\rho_2} +\omega_1 k_{\rho_2}v_{\rho_1}+\omega_2 k_{\rho_1}v_{\rho_2}+\omega_1\omega_2 \eta_{\rho_1\rho_2})\rightarrow  \thesisinlinefeynman[\chthreeinlinefeynwidth]{phi_z_2.png}$
%\end{itemize}
%\vspace{-1.6 cm}
%\newpage
\subsection*{\underline{Vertex for gravitational field:}}
%\begin{itemize}
     \textbullet $\,\,$ $\mathcal{O}{(z^0)}:- i\frac{m_a }{2 m_p} \,e^{ik\cdot b_a}\,\hat\delta(k\cdot v)v_a^\mu v_a^\nu \rightarrow  \thesisinlinefeynman[\chthreeinlinefeynwidth]{h_z_0.png}$

 \textbullet $\,\,$ $\mathcal{O}(z^1):\frac{m_a }{2 m_p}\,e^{ik\cdot b}\,\hat{\delta}(k\cdot v+\omega)(2\omega v^{(\mu}\delta^{\nu)}_\rho+v^\mu v^\nu k_{\rho})\rightarrow \thesisinlinefeynman[\chthreeinlinefeynwidth]{h_z_1.png}$

 \textbullet $\,\,$ $\mathcal{O}{(z^2)}:i\frac{m_a }{m_p} \,e^{ik\cdot b_a}\,\hat\delta(k\cdot v+\omega_1+\omega_2) (\frac{1}{2} k_{\rho_1} k_{\rho_2} v^\mu v^\nu +\omega_1 k_{\rho_2}v^{(\mu}\delta^{\nu)}_{\rho_1}+\omega_2 k_{\rho_1}v^{(\mu}\delta^{\nu)}_{\rho_2}+\omega_1\omega_2 \delta^{(\mu}_{\rho_1}\delta^{\nu)}_{\rho_2})\rightarrow  \thesisinlinefeynman[\chthreeinlinefeynwidth]{h_z_2.png}$
}
\section{Computation of impulse}\label{sec1} 
In this section, we will compute the impulse in a purely conservative setting. We mainly focus on the corrections to the impulse coming from the purely scalar degree of freedom or the interaction between scalar and graviton.
The impulse is defined as,
\begin{align}
    \begin{split}
        \Delta p^{\mu}_{i}=m_i \rmint_{-\infty}^{\infty}d\tau_{i}\Big\langle\frac{d^2 z_{i}^{\mu}}{d\tau_{i}^2}\Big\rangle_{\text{WQFT}}=-m_i \omega^2\langle z_{i}^{\mu}(\omega)\rangle|_{\text{WQFT}}\Big |_{\omega=0}.\label{4.1w}
    \end{split}
\end{align}
%\begin{itemize} 
We now compute (\ref{4.1w}) order by order in Newton's constant $G_N$ up to 2PM order. %\newpage
\subsection{1PM contribution to impulse}
\textbullet $\,\,$ The simplest diagram at  1 PM order consists only of a scalar field that has the following form, 
\begin{align}
    \begin{split}
        [\Delta p_1^{\mu}]_{(a)}=\thesisinlinefeynman[\chthreeinlinefeynwidth]{z_1_1.png}
\hspace{0.35cm}&=-i\frac{m_1s_1m_2s_2}{4m_p^2}\rmint_{k_1,k_2}e^{ik_1\cdot b_1+ik_2\cdot b_2}\frac{\hat\delta(k_2\cdot v_2)\hat\delta(\omega-k_1\cdot v_1)}{k_1^2-m^2}\hat\delta^{(4)}(k_1+k_2)\\&  \hspace{0.8cm}\times(k_1^{\mu}+2\omega v_1^\mu)\Big|_{\omega=0}\,,\\ &
%=-i\frac{m_1s_1m_2s_2}{4m_p^2}\rmint_{k}e^{ik\cdot (b_1-b_2)}\frac{k^\mu \hat\delta(k_1\cdot v_1)\hat\delta(k_2\cdot v_2)}{k^2-m^2}\,,\\ &
=\frac{\partial}{\partial b_{1\mu}}\underbrace{\Big[-\frac{m_1s_1m_2s_2}{4m_p^2}\rmint_{k}e^{ik\cdot(b_1-b_2)}\frac{\delta(k\cdot v_1)\hat\delta(k\cdot v_2)}{k^2-m^2}\Big]}_{\chi}\,.
    \end{split}
\end{align}

To do this integral, we need to choose the coordinates suitably. We study the dynamics from the frame of the second black hole. Hence in the mentioned parametrization of velocities
satisfying $v_1\cdot v_2=\gamma$.
Hence \footnote{Note that $(2\pi)$ is associated with each one-dimensional delta function and $(2\pi)^4$ factor with a four-dimensional delta function. Now, for each of the momentum integrals, there is a $\frac{1}{(2\pi)^4}$ factor associated with the integration measure. One has to take into account all of these factors. Furthermore, there will be a $(2\pi)$ factor whenever we get a ${\bf K_0}$ or ${\bf J_0}$  and $(2\pi)^2$ factor whenever we get a $\boldsymbol{\arctan}$ after doing an integral. One needs to consider these to get the correct factor of $\pi$ at the end. In all the subsequent integrals, we have taken this into account. },
\begin{align}
    \begin{split}
    \chi&=-\frac{m_1m_2s_1s_2}{4m_p^2}\rmint_{k}\frac{\hat{\delta}{(\gamma k^{(0)}-\gamma\beta k^{(1)})}\hat\delta(k^{0})}{k^2-m^2}e^{ik\cdot b}\,,\\ &
    =\frac{\pi ^2 m_1m_2s_1s_2}{m_p^2\sqrt{\gamma^2-1}}
   \rmint_{\tilde k}\frac{e^{-i\tilde k \cdot b}}{\tilde k^2+m^2}=\frac{m_1m_2s_1s_2}{8\pi m_p^2\sqrt{\gamma^2-1}}K_{0}(m|b|)\,.
    \end{split}
\end{align}
Therefore,
\begin{align}
\begin{split} \label{e1:barwq2}
  [\Delta p_1^{\mu}]_{(a)}&=-\frac{m_1m_2s_1s_2}{8\pi m_p^2\sqrt{\gamma^2-1}} \frac{b^\mu}{|b|}m K_1(m|b|)\,.
\end{split}
\end{align} 
In the massless limit the impulse took the form,
\begin{align}
    \begin{split}
          [\Delta p_1^{\mu}]_{(a)}\Big|_{m\to 0}=-\frac{m_1m_2s_1s_2}{8\pi m_p^2\sqrt{\gamma^2-1}} \frac{b^\mu}{|b|^2}.
    \end{split}
\end{align}
\vspace{0.5cm}
\textbullet $\,\,$ One more diagram comes from the self-interaction vertex, which contributes at 1 PM order.
\begin{align}
    \begin{split}
        \mathcal{O}(z,\lambda_2):-\frac{\lambda_3}{\textcolor{black}{3!}} m_p\rmint d^4x\,\varphi(x)^3.
    \end{split}
\end{align}
Then\footnote{\textcolor{black}{Again note that the combinatorial factor associated with this diagram is 3!. We have multiplied it by that. For the subsequent diagrams, we will also multiply by the suitable combinatorial factors from the beginning.} },
\begin{align}
    \begin{split}
        [\Delta p_1^{\mu}]_{(b)}=\thesisinlinefeynman[\chthreeinlinefeynwidth]{z_1_4.png}&=-i\lambda_3\frac{m_1s_1}{2}\times \Big(\frac{m_2s_2}{2m_p}\Big)^2\rmint_{k_i}\hat\delta^{(4)}\Big(\sum_ik_i\Big)\frac{k_1^\mu\hat\delta(k_1\cdot v_1)\hat\delta(k_2\cdot v_2)\hat\delta(k_3\cdot v_2)}{\prod_{i=1}^3(k_i^2-m^2)}e^{ik_1\cdot b_1}e^{i(k_2+k_3)\cdot b_2}\,,\\ &
=-i\lambda_3\frac{m_1s_1}{2}\times \Big(\frac{m_2s_2}{2m_p}\Big)^2\rmint_{k_{1,2}}\frac{k_1^\mu\hat\delta(k_1\cdot v_1)\hat\delta(k_2\cdot v_2)\hat\delta(k_1\cdot v_2)}{(k_1^2-m^2)(k_2^2-m^2)[(k_1+k_2)^2-m^2]}e^{ik_1\cdot (b_1-b_2)}\,.
    \end{split}
\end{align}
Now using the proper velocity parametrization and the integral results derived in \cite{Bhattacharyya:2023kbh}, we get,
\begin{align}
    \begin{split}
         [\Delta p_1^{\mu}]_{(b)}&=-i\lambda_3\frac{m_1s_1}{2(2\pi)^3}\times \Big(\frac{m_2s_2}{2m_p}\Big)^2\frac{1}{\gamma\beta}\rmint d^2 l_1\frac{l_1^\mu e^{il_1\cdot b}}{(\vec l_1^2+m^2)|\vec l_1|}\arctan\Big(\frac{|\vec l_1|}{2m}\Big)\,,\\ &
        =\lambda_3\frac{m_1s_1}{2(2\pi)^2}\times \Big(\frac{m_2s_2}{2m_p}\Big)^2\frac{1}{\gamma\beta}\frac{\partial}{\partial b_{1\mu}}\underbrace{\rmint_{0}^{\infty}dl_1\,\frac{J_{0}(|b|l_1)}{ l_1^2+m^2}\arctan\Big(\frac{l_1}{2m}\Big)}_{I_3(m,b)}\,,\\ &
        =\lambda_3\frac{m_1s_1}{2(2\pi)^2}\times \Big(\frac{m_2s_2}{2m_p}\Big)^2\frac{1}{\sqrt{\gamma^2-1}}\frac{b^\mu}{|b|}\frac{\partial I_1(m,|b|)}{\partial |b|}\,.
    \end{split}
\end{align}  
Finally, we get, 

\begin{equation}\label{e:barwq2}
  [\Delta p_1^{\mu}]_{(b)}=\lambda_3\frac{m_1s_1}{8\pi^2}\times \Big(\frac{m_2s_2}{2m_p}\Big)^2\frac{1}{\sqrt{\gamma^2-1}}\frac{b^\mu}{|b|}\frac{\partial I_1(m,|b|)}{\partial |b|}\,.
\end{equation}\\
\textcolor{black}{To the best of our knowledge this integral does not have any closed-form expression. One can perform a numerical analysis to solve the integral. Moreover, the integral has a smooth massless limit.} In the massless limit, the contribution gives,
\begin{align}
    \begin{split}
         [\Delta p_1^{\mu}]_{(b)}\Big|_{m\to 0}=-\lambda_3 \frac{\pi m_1s_1}{8\pi}\times \Big(\frac{m_2s_2}{4m_p}\Big)^2\frac{1}{\sqrt{\gamma^2-1}}\frac{b^\mu}{|b|}\,.
    \end{split}
\end{align}
Then, the total impulse (due to the scalar field) at 1PM order is the sum of (\ref{e1:barwq2}) and (\ref{e:barwq2})  as well as the terms that come from interchanging the worldline one and two\,.
\begin{equation}
     \Delta p_1^{\mu}\Big|^{\textrm{1PM}, \textrm{Total}}_{\textrm{scalar}}=   [\Delta p_1^{\mu}]_{(a)}+ [\Delta p_1^{\mu}]_{(b)}+1\leftrightarrow 2\,.
\end{equation}
Finally, collecting all individual expressions we get, 
\textcolor{black}{
\begin{align}
\begin{split}
\Delta p_1^{\mu}\Big|^{\textrm{1PM}, \textrm{Total}}_{\textrm{scalar}}= &-\frac{m_1m_2s_1s_2}{8\pi m_p^2\sqrt{\gamma^2-1}} \frac{b^\mu}{|b|}m K_1(m|b|)+\lambda_3\frac{m_1s_1}{8\pi^2}\times \Big(\frac{m_2s_2}{2m_p}\Big)^2\frac{1}{\sqrt{\gamma^2-1}}\frac{b^\mu}{|b|}\frac{\partial I_1(m,|b|)}{\partial |b|}\,+1\leftrightarrow 2,\label{3.42b}
\end{split}
\end{align}
where, $$I_1(m,|b|)=\rmint_{0}^{\infty}dx\,\frac{J_{0}(|b|x)}{ x^2+m^2}\arctan\Big(\frac{x}{2m}\Big)\,.$$}
The expression in \eqref{3.42b} is not a convenient one; further massaging, we get the following,
\begin{align}
   \partial_{b} I_1=- \int _0^\infty\frac{x\,J_1(|b|x)}{x^2+m^2}\arctan\left(\frac{x}{2m}\right)=-2\int_0^1 ds \frac{K_1(m|b|)-\frac{2}{s}K_1\left(\frac{2m|b|}{s}\right)}{4-s^2}
\end{align}
which can be achieved by realising,
\begin{align}
    \arctan\left(\frac{x}{2m}\right)=2mx\int_0^1 \frac{ds}{x^2s^2+4m^2}.
\end{align}
In the next subsection, we extend the impulse computation to 2PM order.

\subsection{2PM contribution to the impulse}
In this subsection, we compute the 2PM contribution to impulse, mainly focusing on the scalar field contribution.

\textbullet $\,\,$  At $\mathcal{O}(z)$ simplest diagram comes from the graviton-scalar interaction vertex:
\begin{align}
    \begin{split}
        \mathcal{O}(z): -\frac{1}{2}m^2\rmint d^4x \frac{h}{2m_p}\varphi^2\,.
    \end{split}
\end{align}
%Here we consider the on-shell worldline propagator from the graviton-worldline vertex.
The corresponding contribution to the impulse is given by,
\begin{align}
    \begin{split}
        [\Delta p_1^{\mu}]_{(c)}& = \thesisinlinefeynman[\chthreeinlinefeynwidth]{z_1_2.png}\hspace{0.7 cm}\\ &
=-im^2\frac{m_1 m_2^2s_2^2}{8m_p^4}\rmint_{k_i}\hat\delta^{(4)}\Big(\sum_{i=1}^{3}k_i\Big)\frac{\hat\delta(k_1\cdot v_1)\hat\delta(k_2\cdot v_2)\hat\delta(k_3\cdot v_2)}{k_1^2(k_2^2-m^2)(k_3^2-m^2)}e^{ik_1\cdot b_1}e^{ik_2\cdot b_2}e^{ik_3\cdot b_2}k_1^\mu \underbrace{P_{\rho\rho,\sigma\delta}\,v_1^\sigma v_1^\delta}_{-1}\,,\\ &
=im^2\frac{m_1 m_2^2s_2^2}{8m_p^4}\rmint_{k_{1,2}}\frac{k_1^\mu\hat\delta(k_1\cdot v_1)\hat\delta(k_2\cdot v_2)\hat\delta(k_1\cdot v_2)}{k_1^2(k_2^2-m^2)[(k_1+k_2)^2-m^2]}e^{ik_1\cdot (b_1-b_2)}\,,\\ &
%=im^2\frac{m_1 m_2^2s_2^2}{8m_p^4}\rmint_{k_1,k_2}\frac{k_1^\mu\hat\delta(k_1\cdot v_1)\hat\delta(k_1\cdot v_2)\hat\delta(k_2\cdot v_2)}{k_1^2 (k_2^2-m^2)[(k_1+k_2)^2-m^2]}e^{ik_1 \cdot b}\,,\\ &
%=im^2\frac{m_2^2s_2^2}{16m_p^4}\rmint_{k_1,k_2}\frac{k_1^{\mu}\hat\delta(\gamma k_1^{0}-\gamma\beta k_1^{3})\hat\delta(k_1^{0})\hat\delta(k_2^0)}{k_1^2(k_2^2-m^2)[(k_1+k_2)^2-m^2]}e^{ik_1\cdot b}\\ &
%=im^2\frac{m_2^2s_2^2}{16m_p^4}\frac{1}{\gamma\beta}\rmint_{\vec k_1,k_2}\frac{\hat\delta(k_1^{3})\hat\delta(k_2^0)}{k_1^2(k_2^2-m^2)[(k_1+k_2)^2-m^2]}e^{ik_1\cdot b}\\ &
%=-im^2\frac{m_2^2s_2^2}{16m_p^4}\frac{1}{\gamma\beta}\rmint d^2 l_1\frac{l_1^{\mu}e^{il_1\cdot b}}{\vec{l}_1^2}\rmint d^3 l_2\frac{1}{(\vec {l}_2^2+m^2)[(\vec{l}_1+\vec{l}_2)^2+m^2]}\\ &
%=m^2\frac{m_2^2s_2^2}{16m_p^4}\frac{1}{\gamma\beta}\frac{\partial}{\partial b_{1\mu}}\rmint d^2 l_1\frac{e^{il_1\cdot b}}{\vec{l}_1^3}\arctan\Big(\frac{|\vec{l}_1|}{2m}\Big)\\ &
=m^2\frac{m_1m_2^2s_2^2}{32\pi^2 m_p^4}\frac{1}{\sqrt{\gamma^2-1}}\frac{b^\mu}{|b|}\partial_{|b|}\underbrace{\rmint _{0}^{\infty}dl\,\frac{J_{0}(bl)}{l^2}\arctan\Big(\frac{l}{2m}\Big)}_{I_2(m,|b|)}.\label{4.23m}
    \end{split}
\end{align}
Therefore, the contribution to the impulse from this diagram has the following form.

\begin{equation}\label{e:barwq244}
  [\Delta p_1^{\mu}]_{(c)}=m^2\frac{m_1m_2^2s_2^2}{32 \pi^2 m_p^4}\frac{1}{\sqrt{\gamma^2-1}}\frac{b^\mu}{|b|}\partial_{|b|}{I_2(m,|b|)}\,.
\end{equation}\\
Like (\ref{e:barwq2}), it also does not have any closed-form expression. Furthermore, it is evident from \eqref{4.23m} that it becomes identically zero in the massless limit as it is proportional to $m^2$.

\textbullet $\,\,$ Another contributing diagram may appear from the scalar-graviton interaction vertex where one scalar and one graviton line connect with one worldline, and the other scalar line connects with another worldline.
\begin{align}
    \begin{split}
         [\Delta p_1^{\mu}]_{(d)}& =
        \thesisinlinefeynman[\chthreeinlinefeynwidth]{hphi.png}\\ &
        =im^2 \frac{m_1 s_1 m_2^2 s_2}{8 m_p^4}\rmint_{k_{1,2}}\frac{k_1^\mu\hat\delta(k_1\cdot v_1)\hat\delta(k_1\cdot v_2)\hat\delta(k_2\cdot v_2)}{(k_1^2-m^2)(k_2^2-m^2)(k_1+k_2)^2}e^{ik_1\cdot b}\\ &
        =m^2 \frac{m_1 s_1 m_2^2 s_2}{32 \pi^2 m_p^4\sqrt{\gamma^2-1}}\frac{b^\mu}{|b|}\partial_{|b|} \underbrace{\rmint_{k_1}dk \frac{ J_{0}(|b|k_1)}{k_1^2+m^2}\arctan\Big(\frac{k_1}{m}\Big)}_{\Bar{I}_2(m,|b|)}
    \end{split}
\end{align}
Hence, the contribution to the impulse from this diagram has the following form:


\begin{equation}\label{e:barwq248}
 [\Delta p_1^{\mu}]_{(d)}=m^2 \frac{m_1 s_1 m_2^2 s_2}{32 \pi^2 m_p^4\sqrt{\gamma^2-1}}\frac{b^\mu}{|b|}\partial_{|b|} \bar{I}_2(m,|b|)
\,.
\end{equation}

%Therefore, the contribution to the impulse has the following form:
%\begin{equation}\label{e00}\begin{split}
%\Delta p_1^\mu= m^2\frac{m_2^2s_2^2}{8m_p^4\gamma\beta}\frac{b^\mu}{|b|}\frac{\partial I_1(m,|b|)}{\partial |b|}\,.\\&
%\end{split}\end{equation}
\textbullet $\,\,$  Now, we concentrate on the scalar interaction diagrams. We compute the contribution to the impulse at 2PM order coming from the scalar self-interaction vertex:
\begin{align}
    \begin{split}
        \mathcal{O}(z):-\frac{\lambda_4}{\textcolor{black}{4!}}\rmint d^4x\,\varphi^4(x)\,.
    \end{split}
\end{align}
The corresponding contribution to the impulse is shown below,
\begin{align}
    \begin{split}
          [\Delta p_1^{\mu}]_{(e)}&=\thesisinlinefeynman[\chthreeinlinefeynwidth]{z_1_3.png}\\ & =-i\lambda_4\Big(\frac{m_1s_1m_2s_2}{4m_p^2}\Big)^2\rmint_{\{k_i\}}\hat\delta^{(4)}\Big(\sum_{i=1}^{4}k_{i}\Big)\hat\delta(k_1\cdot v_1)\hat\delta(k_2\cdot v_1)\hat\delta(k_3\cdot v_2)\hat\delta(k_4\cdot v_2)\prod_{j=1}^{4}\frac{1}{k_j^2-m^2}\\ & 
\hspace{0.8cm}\times k_1^{\mu} \,e^{i(k_1+k_2)\cdot b_2}e^{i(k_3+k_4)\cdot b_2}\,,\\ &
=-i\lambda_4\Big(\frac{m_1s_1m_2s_2}{4m_p^2}\Big)^2\rmint_{k_i\ne k_4}\hat\delta(k_1\cdot v_1)\hat\delta(k_2\cdot v_1)\hat\delta(k_3\cdot v_2)\hat\delta(k_1\cdot v_2+k_2\cdot v_2)\\ &\hspace{0.8cm}\times\frac{k_1^\mu e^{i(k_1+k_2)\cdot (b_1-b_2)}}{\prod_{j=1}^{3}(k_j^2-m^2)[(k_1+k_2+k_3)^2-m^2]}\,,\\ &
\xrightarrow[]{k_1+k_2\rightarrow q}-i\lambda_4\Big(\frac{m_1s_1m_2s_2}{4m_p^2}\Big)^2\rmint_{q,k_1,k_3}\frac{k_1^\mu\hat\delta(k_1\cdot v_1)\hat\delta(q\cdot v_1)\hat\delta(k_3\cdot v_2)\hat\delta(q\cdot v_2)}{(k_1^2-m^2)(k_3^2-m^2)[(q-k_1)^2-m^2][(q+k_3)^2-m^2]}e^{iq\cdot b}\,.\\ &
%=i\lambda_4\Big(\frac{m_1s_1m_2s_2}{4m_p^2}\Big)^2\rmint k_1^\mu\, \hat\delta(\gamma k_1^0-\gamma\beta k_1^{3})\hat\delta(\gamma q^0-\gamma\beta q^3)\hat\delta(k_3^0)\hat\delta(q^0)\textcolor{black}{[\cdots]}\,.
   \label{4.26m} \end{split}
\end{align}
In \eqref{4.26m} one can see that the delta function constraints set $q$ and $k_3$ two and three-dimensional vectors, respectively and $k_1$ a three-dimensional vector with scaled variable: $
        \bar k_1^{(1)}=\frac{k_1^{(1)}}{\gamma}.
   $ 
Therefore, the integral in  (\ref{4.26m}) can be re-written as,
\begin{align}
\begin{split}
   &   [\Delta p_1^{\mu}]_{(e)}\\ &=-i\lambda_4\Big(\frac{m_1s_1m_2s_2}{4m_p^2}\Big)^2 \int e^{iq\cdot b}\,\frac{q^\mu}{q^2}\,\hat \delta(q\cdot v_1)\hat\delta(q\cdot v_2)\int_{k_3,k_1}\frac{\hat\delta(k_1\cdot v_1)\hat\delta(k_3\cdot v_2)\left((k_1^2-m^2)+q^2-((k_1-q)^2-m^2)\right)}{(k_1^2-m^2)(k_3^2-m^2)[(q-k_1)^2-m^2][(q+k_3)^2-m^2]}\,, \\ &
   =-i\lambda_4\Big(\frac{m_1s_1m_2s_2}{4m_p^2}\Big)^2\int e^{iq\cdot b} \frac{q^\mu}{q^2}\,\hat \delta(q\cdot v_1)\hat\delta(q\cdot v_2)\Bigg(\int_{k_1}\frac{\hat\delta(k_1\cdot v_1)}{(k_1-q)^2-m^2}\int_{k_3}\frac{\hat\delta(k_3\cdot v_2)}{(k_3^2-m^2)((k_3+q)^2-m^2)}\\ & \hspace{1 cm}+q^2\int_ {k_1}\frac{\hat\delta(k_1\cdot v_1)}{(k_1^2-m^2)((k_1-q)^2-m^2)}\int _{k_3}\frac{\hat\delta(k_3\cdot v_2)}{(k_3^2-m^2)((k_3+q)^2-m^2)}-\int_{k_1}\frac{\hat\delta(k_1\cdot v_1)}{k_1^2-m^2}\int_{k_3}\frac{\hat\delta(k_3\cdot v_2)}{(k_3^2-m^2)((k_3+q)^2-m^2)}\Bigg)\,,\\ &
   =-i\lambda_4\Big(\frac{m_1s_1m_2s_2}{16\pi m_p^2}\Big)^2\int e^{iq\cdot b} \frac{q^\mu}{-q^2}\,\hat \delta(q\cdot v_1)\hat\delta(q\cdot v_2)\arctan^2\left(\frac{\sqrt{-q^2}}{2m}\right)\,,\\ &
  % =\lambda_4\Big(\frac{m_1s_1m_2s_2}{16\pi m_p^2}\Big)^2\frac{b^\mu}{\sqrt{\gamma^2-1}|b|} \partial_{|b|}\int \frac{d^2\boldsymbol{q}_{\perp}}{(2\pi)^2} e^{-i \boldsymbol{q}_{\perp}\cdot b}\frac{1}{q_\perp^2} \arctan^2\left(\frac{\boldsymbol{q}_{\perp}^2}{2m}\right)\\ &
 % = \lambda_4\Big(\frac{m_1s_1m_2s_2}{16\pi m_p^2}\Big)^2\frac{b^\mu}{\sqrt{\gamma^2-1}|b|}\partial_{|b|}\int \frac{dq}{2\pi} J_{0}(q \,b)\frac{1}{q}\arctan^2\left(\frac{\boldsymbol{q}_{\perp}^2}{2m}\right)\\ &
  =- \lambda_4\Big(\frac{m_1s_1m_2s_2}{16\pi m_p^2}\Big)^2\frac{b^\mu}{\sqrt{\gamma^2-1}|b|} \int_0^\infty dq\, J_{1}(q\,b)\arctan^2\left(\frac{q^2}{2m}\right)\,,
   \label{4.19f}
    \end{split}
\end{align}
The integral in \eqref{4.19f} can be rearranged in a nicer form as,
\begin{align}
    [\Delta p_1^{\mu}]_{(e)}=-2 \lambda_4\Big(\frac{m_1s_1m_2s_2}{16\pi m_p^2}\Big)^2\frac{b^\mu}{\sqrt{\gamma^2-1}|b|^2}\int_0^1 ds\frac{K_0(2m|b|)-K_0\left(\frac{2m|b|}{s}\right)}{1-s^2}\,.\label{3.53gg}
\end{align}
Although not instructive, the integral in \eqref{3.53gg} can be written as a closed form expression in terms of Bickley–Naylor functions as,
\begin{align}
  \int_0^1 ds\frac{K_0(2m|b|)-K_0\left(\frac{2m|b|}{s}\right)}{1-s^2}  =\frac{1}{2} (\log (4m |b|)+\gamma_{E} ) K_0(2 m|b|)-\partial_\alpha \textrm{Ki}_{\alpha}(2m|b|)|_{\alpha\to0};\label{3.54kk}
\end{align}
The derivation is somewhat instructive. We put it in the Appendix~(\ref{3.54kkk}).
The massless counterpart takes the following form,
\begin{align}
    \begin{split}
        [\Delta p_1^{\mu}]_{(e)}=- \lambda_4\Big(\frac{m_1s_1m_2s_2}{32m_p^2}\Big)^2\frac{b^\mu}{\sqrt{\gamma^2-1}|b|^2}
    \end{split}
\end{align}
%\vspace{0.55cm}

\textbullet $\,\,$ Another diagram involving $\lambda_4 \varphi^4$ vertex will contribute to the impulse, where we will have three scalar lines connected with one worldline, and the other scalar line connects with another worldline.
\begin{align}
    \begin{split}
         [\Delta p_1^{\mu}]_{(f)}&=\thesisinlinefeynman[\chthreeinlinefeynwidth]{phi_4_1.png}\\ &=-i\lambda_4 \Big(\frac{m_2 s_2}{2m_p}\Big)^3\Big(\frac{m_1 s_1}{2m_p}\Big)\rmint_{\{k_i\}}\hat\delta^{(4)}\Big(\sum_{i=1}^{4}k_{i}\Big)\hat\delta(k_1\cdot v_1)\hat\delta(k_2\cdot v_2)\hat\delta(k_3\cdot v_2)\hat\delta(k_4\cdot v_2)\prod_{j=1}^{4}\frac{1}{k_j^2-m^2}\\ & 
\hspace{0.8cm}\times k_1^{\mu} \,e^{ik_1\cdot b_1}e^{i(k_2+k_3+k_4)\cdot b_2}\,,\\ &
%=-i\lambda_4 \Big(\frac{m_2 s_2}{2m_p}\Big)^3\Big(\frac{m_1 s_1}{2m_p}\Big)\rmint e^{ik_1\cdot b}\frac{k_1^\mu\hat\delta(k_1\cdot v_1)\hat\delta(k_2\cdot v_2)\hat\delta(k_3\cdot v_2)\hat\delta(k_1\cdot v_2)}{(k_1^2-m^2)(k_2^2-m^2)(k_3^2-m^2)[(k_1+k_2+k_3)^2-m^2]}\,,\\ &
=-i\lambda_4 \Big(\frac{m_2 s_2}{2m_p}\Big)^3\Big(\frac{m_1 s_1}{2m_p}\Big)\rmint e^{ik_1\cdot b}k_1^\mu \,\frac{\hat\delta(k_1\cdot v_1)\hat\delta(k_1\cdot v_2)}{k_1^2-m^2}\rmint \frac{\hat\delta(q\cdot v_2)\hat\delta(k_3\cdot v_2)}{[(q-k_1)^2-m^2](k_3^2-m^2)[(q+k_3)^2-m^2]}\,,\\ &
%=-i\lambda_4 \Big(\frac{m_2 s_2}{2m_p}\Big)^3\Big(\frac{m_1 s_1}{8\pi m_p}\Big)\rmint e^{ik_1\cdot b}k_1^\mu \,\frac{\hat\delta(k_1\cdot v_1)\hat\delta(k_1\cdot v_2)}{k_1^2-m^2}\rmint_{q}\frac{1}{|\,\vec q\,|[(\vec q-\vec k_1)^2+m^2]}\arctan\Big(\frac{|\,\vec q\,|}{2m}\Big)\,,\\ &
=\frac{\lambda_4}{(2\pi)^2} \Big(\frac{m_2 s_2}{2m_p}\Big)^3\Big(\frac{m_1 s_1}{2 m_p}\Big)\frac{1}{\sqrt{\gamma^2-1}}\frac{b^\mu}{|b|}\partial_{|b|}\rmint  d^2 k_1\,\frac{e^{-i\vec k_1\cdot \vec b}}{\vec k_1^2+m^2}\,\hat{I}(|k_1|,m)\,.
\label{4.19c}
    \end{split}
\end{align}
We can further simplify \eqref{4.19c} using Schwinger parametrization.
\begin{align}
    \begin{split}
       \hat{I}(|k_1|,m)=\frac{1}{(2\pi)^4}\rmint d^3 q\frac{1}{|\vec q\,|} \rmint_{0}^{\infty}d\alpha\, e^{-\alpha (\vec q-\vec k_1)^2}e^{-\alpha m^2}\arctan\Big(\frac{|\vec q\,|}{2m}\Big)\,.
    \end{split}
\end{align}
Now, in \eqref{4.19c}, we first perform the angular part of the $q$ integral,
\begin{align}
    \begin{split}
         [\Delta p_1^{\mu}]_{(f)} &\sim \frac{1}{(2\pi)^3}\rmint d\alpha\rmint d^3q \frac{e^{-\alpha (q^2+m^2)}}{|q|}\arctan\Big(\frac{|\vec q\,|}{2m}\Big)\rmint_{0}^\infty dk_1\frac{k_1}{ k_1^2+m^2}J_{0}(k_1|b|)e^{-\alpha k_1^2+2\alpha \vec q\cdot \vec k_1}\,,\\ &
    %    =\frac{1}{(2\pi)^2}\rmint d\alpha \,e^{-\alpha m^2}\rmint dk_1\,\frac{k_1}{ k_1^2+m^2}J_{0}(k_1|b|)e^{-\alpha k_1^2}\rmint dq \,d\lambda\,q \sin(\lambda)e^{-\alpha (q^2-2 q k_1 \cos \lambda)}\arctan\Big(\frac{q}{2m}\Big)\,,\\ &
        =\frac{1}{(2\pi)^2}\rmint d\alpha \,\frac{e^{-\alpha m^2}}{\alpha}\rmint dk_1\,\frac{1}{ k_1^2+m^2}J_{0}(k_1|b|)e^{-\alpha k_1^2}\rmint dq\,\sinh(2q k_1 \alpha)\arctan\Big(\frac{q}{2m}\Big)e^{-\alpha q^2}\,,\\ &
       % =\frac{1}{(2\pi)^2}\rmint dq\,\arctan\Big(\frac{q}{2m}\Big)\rmint dk_1\frac{1}{k_1^2+m^2}J_{0}(k_1 b)\rmint d\alpha\, e^{-\alpha(m^2+q^2+k_1^2)}\sinh(2\alpha q k_1)\,,\\ &
       % =\rmint dq\,q\arctan\Big(\frac{q}{2m}\Big)\rmint dk_1\frac{1}{k_1^2+m^2}J_{0}(k_1 b)\,\textrm{arctanh}\Big(\frac{2qk_1}{m^2+q^2+k_1^2}\Big)\\ &
        =\frac{1}{(2\pi)^2}\rmint_0^\infty dq \,dk_1\,\frac{J_{0}(k_1 b)}{k_1^2+m^2}\arctan\Big(\frac{q}{2m}\Big)\textrm{arctanh}\Big(\frac{2qk_1}{m^2+q^2+k_1^2}\Big)\,.\label{4.21cc}
    \end{split}
\end{align}
The $q$ integral can be done by taking the \textbf{logarithmic} representation of the \textbf{ArcTan} and \textbf{ArcTanh} in the following way,
\begin{align}
    \begin{split}
       & \rmint dq \arctan\Big(\frac{q}{2m}\Big)\,\textrm{arctanh}\Big(\frac{2q k_1}{q^2+k_1^2+m^2}\Big)\,,\\ &
       =\frac{i}{4}\rmint dq\, \log\Bigg[\frac{1-\frac{iq}{2m}}{1+\frac{iq}{2m}}\Bigg]\log\Big(1+\frac{4q k_1}{m^2+(q-k_1)^2}\Big)\,,\\ &
       =\frac{\pi}{2}\Big[2k_1-2m\arctan\Big(\frac{k_1}{m}\Big)-6m\arctan\Big(\frac{k_1}{3m}\Big)-k_1\log\Big(1+\frac{k_1^2}{m^2}\Big)-k_1\log\Big(1+\frac{k_1^2}{9m^2}\Big)-2k_1 \log(3m^2)\Big]\,.\label{4.25j}
    \end{split}
\end{align}
Therefore, the impulse becomes \footnote{As one can see, the impulse has a logarithmic behaviour w.r.t $|b|$. Hence, one should divide it by a UV cutoff $b_0$ to make it dimensionless. However, it does not contribute to the finite part of the impulse. The same thing can be said whenever some logarithmic terms appear. },
\begin{align}
    \begin{split}
         [\Delta p_1^{\mu}]_{(f)}&\sim \frac{1}{32\pi^3}\rmint dk_1\frac{J_{0}(k_1b)}{k_1^2+m^2}\Big[2k_1-2m\arctan\Big(\frac{k_1}{m}\Big)-6m\arctan\Big(\frac{k_1}{3m}\Big)\\ &
        \hspace{0.8cm}-k_1\log\Big(1+\frac{k_1^2}{m^2}\Big)-k_1\log\Big(1+\frac{k_1^2}{9m^2}\Big)-2k_1 \log(3m^2)\Big]\,,\\ &
        =\frac{1}{32\pi^3}\Big[2(1-\log3m^2)K_{0}(|b|m)-\rmint dk_1\frac{J_{0}(k_1 b)}{k_1^2+m^2}\Theta(k_1,m)\Big]\,,
    \end{split}
\end{align}
where,
\begin{align}
    \begin{split}
        \Theta(k_1,m)=-2m\arctan\Big(\frac{k_1}{m}\Big)-6m\arctan\Big(\frac{k_1}{3m}\Big)-k_1 \log\Big[\Big(1+\frac{k_1^2}{m^2}\Big)\Big(1+\frac{k_1^2}{9m^2}\Big)\Big]\,.
    \end{split}
\end{align}
After inserting the prefactors, we get

%\begin{align}
  %  \begin{split}
     %  \Delta p_1^\mu= \lambda_4 \Big(\frac{m_2 s_2}{2m_p}\Big)^3\Big(\frac{m_1 s_1}{8\pi m_p}\Big)\frac{b^\mu}{|b|}\partial_{|b|}\Big(2(1-\log3m^2)K_{0}(|b|m)-\rmint dk_1\frac{J_{0}(k_1 b)}{k_1^2+m^2}\Theta(k_1,m)\Big)\,.
  %  \end{split}
%
%\end{align}

\begin{equation}\label{e:barwq246}
 [\Delta p_1^{\mu}]_{(f)}=\frac{\lambda_4}{32\pi^3} \Big(\frac{m_2 s_2}{2m_p}\Big)^3\Big(\frac{m_1 s_1}{2m_p }\Big)\frac{1}{\sqrt{\gamma^2-1}}\frac{b^\mu}{|b|}\partial_{|b|}\Big(2(1-\log(3m^2/\mu^2))K_{0}(|b|m)-\rmint_0^{\infty} dk_1\frac{J_{0}(k_1 |b|)}{k_1^2+m^2}\Theta(k_1,m)\Big)\,.\,
\end{equation}\\
where $\mu$ is the UV cutoff.\textcolor{black}{ As can be seen, the integral in \eqref{4.21cc} does not have any trivial closed-form, but one can derive the massless limit}. Although the integral can be rearranged as,
\begin{align}
   \begin{split}
   & \Big(2(1-\log(3m^2/\mu^2))K_{0}(|b|m)-\rmint_0^{\infty} dk_1\frac{J_{0}(k_1 |b|)}{k_1^2+m^2}\Theta(k_1,m)\Big)\\ &
   \hspace{3 cm}=2\left(1-\log\frac{9m^2}{16\mu^2}\right)K_0(m|b|)+2\int_{0}^1 \frac{ds}{s(s+1)}K_0\left(\frac{m|b|}{s}\right)+18 \int_0^1 ds \frac{1-s}{s(9-s^2)}K_0\left(\frac{3m|b|}{s}\right)\,.
    \end{split}
\end{align}

The massless limit is given by,
\begin{align}
    \begin{split}
      [\Delta p_1^{\mu}]_{(f)}\Big|_{m\rightarrow 0}   &=\frac{\lambda_4 }{32\pi^3}\Big(\frac{m_2 s_2}{2m_p}\Big)^3\Big( \frac{m_1 s_1}{2 m_p}\Big)\frac{1}{\sqrt{\gamma^2-1}}\frac{b^\mu}{|b|}\Big[-\frac{2}{|b|}+4\partial_{|b|}\rmint_0^\infty dk_1\frac{ \,J_{0}(k_1 |b|)}{k_1}\log(k_1)\Big]\,,\\ &
    % =\lambda_4 \Big(\frac{m_2 s_2}{2m_p}\Big)^3\Big( \frac{m_1 s_1}{8\pi m_p}\Big)\frac{b^\mu}{|b|}\Big[-\frac{2}{|b|}+\frac{2 \left(\log \left(\frac{b^2}{4}\right)+2 \gamma \right)}{b}\Big]\,,\\ &
     =\frac{\lambda_4}{8\pi^3} \Big(\frac{m_2 s_2}{2m_p}\Big)^3\Big( \frac{m_1 s_1}{2 m_p}\Big)\frac{1}{\sqrt{\gamma^2-1}}\frac{b^\mu}{|b|^2}\Big[\log (|b|/b_0)+\gamma_E-\frac{1}{2}-\log 2\Big]\,.
    \end{split}
\end{align}
%\vspace{-0.4 cm}\\
\textbullet $\,\,$ Now we will deal with the derivative interactions. The simplest scalar-gravitation derivative interaction, which contributes to the impulse at 2PM order:
\begin{align}
    \begin{split}
        \mathcal{O}(z):\rmint d^4x \frac{h^{\alpha\beta}}{m_p}\,\partial_{\alpha}\varphi\partial_{\beta}\varphi\,.
    \end{split}
\end{align}
\vspace{-0.53 cm}
The vertex factor has the following form,
\begin{align}
    \begin{split}
        \mathcal{V}_{\alpha\beta}(k_1,k_2,k_3)=  -\rmint_{\{k_{i}\}}  \hat\delta\Big(\sum_{i=1}^{3}k_{i}\Big)   (k_2)_{\alpha} (k_3)_{\beta}\,.
    \end{split}
\end{align}
The contribution to the impulse is given by\footnote{Time-symmetric Feynman propagators imply an elastic scattering of the black holes. In principle, one can also begin with the retarded graviton propagators to take care of radiative effects, but we will use the first one.},
\begin{align}
    \begin{split}
         [\Delta p_1^{\mu}]_{(g)} &=\thesisinlinefeynman[\chthreeinlinefeynwidth]{z_1_5.png}\hspace{0.7 cm}\\ &
=i\Big(\frac{m_2s_2}{2m_p}\Big)^2\times \frac{m_1}{ m_p^2}\rmint_{k_i}\hat\delta^{(4)}\Big(\sum_{i}k_i\Big)k_2^\alpha k_3^\beta\,\frac{k_1^\mu P_{\sigma\delta;\alpha\beta }v_1^\sigma v_1^\delta\,\hat\delta(k_1\cdot v_1)\hat\delta(k_2\cdot v_2)\hat\delta(k_3\cdot v_2)}{k_1^2\prod_{i=2,3}(k_i^2-m^2)}e^{ik_1\cdot b_1}e^{i(k_2+k_3)\cdot b_2}\,,\\ &
%=-i\Big(\frac{m_2s_2}{2m_p}\Big)^2\times \frac{m_1}{2 m_p^2}\rmint_{k_{1,2}}k_{2\alpha}(k_1+k_2)_{\beta}\frac{k_1^\mu\hat\delta(k_1\cdot v_1)\hat\delta(k_2\cdot v_2)\hat\delta(k_1\cdot v_2)}{k_1^2(k_2^2-m^2)[(k_1+k_2)^2-m^2]}(\eta_{\sigma\alpha}\eta_{\delta\beta}+\eta_{\sigma\beta}\eta_{\delta\alpha}-\eta_{\sigma\delta}\eta_{\alpha\beta})\\ &
%\hspace{0.8cm}\times e^{ik_1\cdot b} v_1^\sigma v_1^\delta\,,\\& 
%=-i\Big(\frac{m_2s_2}{2m_p}\Big)^2\times \frac{m_1}{2m_p^2}\rmint_{k_{1,2}}\Big[2(k_2\cdot v_1)(k_1+k_2)\cdot v_1-\,k_2\cdot (k_1+k_2)\Big]\frac{k_1^\mu\hat{\delta}(k_1\cdot v_1)\hat{\delta}(k_2\cdot v_2)\hat{\delta}(k_1\cdot v_2)}{k_1^2(k_2^2-m^2)[(k_1+k_2)^2-m^2]}e^{ik_1\cdot b}\,,\\&
=\Big(\frac{m_2s_2}{2m_p}\Big)^2\times \frac{m_1}{2m_p^2}\\ &
\hspace{2 cm}\times\frac{\partial}{\partial b_{1\mu}} \underbrace{\rmint_{k_{1,2}}\Big[2(k_2\cdot v_1)(k_1+k_2)\cdot v_1-\,k_2\cdot (k_1+k_2)\Big]\frac{\hat\delta(k_1\cdot v_1)\hat\delta(k_2\cdot v_2)\hat\delta(k_1\cdot v_2)}{k_1^2(k_2^2-m^2)[(k_1+k_2)^2-m^2]}e^{ik_1\cdot b}}_{I_{4}(m,|b|)}\,.\label{4.37} 
    \end{split}
\end{align}
The computation in \eqref{4.37} is an involved one and should be done carefully. We show below the details of the integration for individual terms.
\begin{align}
    \begin{split}
        I_4(m,|b|)\Big|_{(1)}&\equiv 2\rmint_{k_{1,2}}(k_2\cdot v_1)^2\frac{\hat\delta(k_1\cdot v_1) \hat\delta(k_2\cdot v_2)\hat\delta(k_1\cdot v_2)}{k_1^2(k_2^2-m^2)[(k_1+k_2)^2-m^2]}e^{i k_1\cdot b}\,,\\ &
         =(v_1)_{\mu}(v_1)_{\nu}\rmint_{k_{1,2}}k_2^{\mu}\,k_2^{\nu}\frac{\hat\delta(\gamma k_1^{(0)}-\gamma\beta k_1^{(1)})\hat\delta( k_2^0)\hat\delta(k_1^0)}{k_1^2(k_2^2-m^2)[(k_1+k_2)^2-m^2]}e^{i k_1 \cdot b}\,,\\ &
    =\textcolor{black}{\frac{1}{(2\pi)^5}}\frac{v_{1i}v_{1j}    }{\gamma\beta}\rmint d^2 k_1\frac{e^{i k_1 \cdot b}}{\vec k_1^2}
    \underbrace{\rmint d^3 k \frac{k_i k_j}{(\vec k^2+m^2)[(\vec k+\vec k_1)^2+m^2]}}_{\mathcal{C}_{ij}}\,.
    \end{split}
\end{align}
Using the Passarino-Veltman reduction one can deduce the form of $\mathcal{C}_{ij}$,
\begin{align}
    \begin{split}
\mathcal{C}_{ij}=&\underbrace{\Big[\textcolor{black}{(2\pi)^2}\frac{3}{8 |\vec k_1|}\arctan\Big(\frac{|\vec k_1|}{2m}\Big)+\textcolor{black}{(2\pi)^2}\frac{m^2}{2 |\vec k_1|^3}\arctan\Big(\frac{|\vec k_1|}{2m}\Big)-\textcolor{black}{\frac{\pi^2}{2}}\frac{m}{|\vec k_1|^2}\Big]}_{\Xi_1(|k_1|,m)}k_1^i k_1^j\\ &
       \underbrace{- \Big[|\vec k_1|\textcolor{black}{\frac{(2\pi)^2}{8}}\arctan\Big(\frac{|\vec k_1|}{2m}\Big)+\textcolor{black}{\frac{(2\pi)^2}{2}}\frac{m^2}{ |\vec k|}\arctan\Big(\frac{|\vec k_1|}{2m}\Big)+m\textcolor{black}{\frac{\pi^2}{2}}\Big]}_{\Xi_2 (|k_1|,m)}\delta^{ij}\,,\\ &
={\Xi}_1(|k_1|,m)k_1^i k_1^j+{\Xi}_{2}(|k_1|,m)\delta^{ij}\,.  \label{new11}
    \end{split}
\end{align}
Therefore,
\begin{align}
    \begin{split} \label{eq:4.35}
        I_{4}(m,|b|)\Big|_{(1)}=\frac{2(2\pi)}{\sqrt{\gamma^2-1}}\Bigg[-v_{1i}v_{1j}\partial_{b_i}\partial_{b_j}\rmint_{0}^{\infty}dl\,\frac{J_{0}(bl)}{l}\,{\Xi}_1\Big(l,m\Big)+(\gamma^2-1)\rmint_{0}^\infty dl\,\frac{J_{0}(bl)}{l}{\Xi}_2\Big(l,m\Big)\Bigg]
    \end{split}
\end{align}
where ${\Xi}_1(l,m)$ and ${\Xi}_2(l,m)$ are defined in (\ref{new11}). 
Other parts of the integrals $I_4(m,b)$ can be computed as follows,
\begin{align}
    \begin{split}
        I_{4}(m,|b|)\Big|_{(2)}\equiv 2\rmint_{k_{1,2}}(k_2\cdot v_1)(k_1\cdot v_1)\hat\delta(k_1\cdot v_1)\cdots \rightarrow 0.
    \end{split}
\end{align}
and,
\begin{align}
    \begin{split}
        I_{4}(m,|b|)\Big|_{(3)}&\equiv- \rmint_{k_{1,2}}k_1\cdot k_2\,\frac{\hat\delta(k_1\cdot v_1) \hat\delta(k_2\cdot v_2)\hat\delta(k_1\cdot v_2)}{k_1^2(k_2^2-m^2)[(k_1+k_2)^2-m^2]}e^{i k_1\cdot b}\,,\\ &
        =\textcolor{black}{\frac{1}{(2\pi)^5}}\frac{1}{\sqrt{\gamma^2-1}}\rmint d^2 k_1 \,\frac{k_1^i e^{ik_1\cdot b}}{k_1^2}\rmint d^3 k_2 \frac{k_2^i}{(\vec k_2^2+m^2)[(\vec k_1+\vec k_2)^2+m^2]}\,,\\ &
       % =\frac{1}{2\sqrt{\gamma^2-1}}\rmint d^2 k_1 \frac{e^{-i\vec k_1\cdot \vec b}}{|\vec k_1|}\arctan\Big(\frac{|\vec k_1|}{2m}\Big)\\ &
        =\textcolor{black}{\frac{1}{(2\pi)^2}}\frac{1}{2\sqrt{\gamma^2-1}}\rmint_{0}^{\infty} dl \,J_{0}(|b|l)\arctan\Big(\frac{l}{2m}\Big) \label{eq:4.37}
    \end{split}
\end{align}
 and,
 \begin{align}
     \begin{split}
         I_{4}(m,|b|)\Big|_{(4)}&\equiv -\rmint_{k_{1,2}}k_2^2\,\frac{\hat\delta(k_1\cdot v_1) \hat\delta(k_2\cdot v_2)\hat\delta(k_1\cdot v_2)}{k_1^2(k_2^2-m^2)[(k_1+k_2)^2-m^2]}e^{i k_1\cdot b}\,,\\ &
         =-\textcolor{black}{\frac{1}{(2\pi)^5}}\frac{1}{\sqrt{\gamma^2-1}}\rmint d^2 k_1\frac{e^{-i\vec k_1\cdot \vec b }}{-\vec k_1^2}\rmint d^3 k_2 \frac{-\vec k_2^2}{(\vec k_2^2+m^2)[(\vec k_1+\vec k_2)^2+m^2]}\,,\\ &
    %     =-\frac{1}{\sqrt{\gamma^2-1}}\rmint d^2 k_1 \frac{e^{-i\vec k_1\cdot \vec b}}{\vec k_1^2}\textrm{tr}(\mathcal{C})\\ &
         =-\textcolor{black}{\frac{1}{(2\pi)^4}}\frac{1}{\sqrt{\gamma^2-1}}\rmint_0^{\infty} dl\, \frac{J_{0}(bl)}{l}\Big(l^2\,{\Xi}_1(l,m)+3\,{\Xi}_2(l,m)\Big)\,, \label{eq:4.38}
     \end{split}
 \end{align}
 where ${\Xi}_1(l,m)$ and ${\Xi}_2(l,m)$ are defined in (\ref{new11}). 
\textcolor{black}{After (\ref{eq:4.35}), (\ref{eq:4.37}) and (\ref{eq:4.38}) we get the full answer for $I_{4}(m,|b|)$ mentioned in (\ref{4.37}). Then we get, the }

\begin{equation}\label{e:barwqpp}
 [\Delta p_1^{\mu}]_{(g)}=\Big(\frac{m_2s_2}{2m_p}\Big)^2\times \frac{m_1}{2m_p^2}\frac{b^\mu}{|b|}\,\partial_{|b|} I_{4}(m,|b|)\,.
\end{equation}
Similiarly, the massless limit can be taken by using the fact that $\arctan(\infty)=\frac{\pi}{2}$ and taking the massless limit of $\Xi(l,m)$.

\textbullet $\,\,$ Another interesting diagram will contribute to the impulse from the previously mentioned derivative interaction where one scalar line connects with one worldline and the other scalar line and the graviton line connects with the other worldline.
\begin{align}
    \begin{split}
         [\Delta p_1^{\mu}]_{(h)}&=\thesisinlinefeynman[\chthreeinlinefeynwidth]{deri.png}\\ &=i\frac{m_1 s_1 m_2^2 s_2}{4m_p^4}\rmint_{k_i}\hat\delta^{(4)}(\sum_i k_i)k_2^\alpha k_3^\beta \frac{k_2^\mu P_{\alpha\beta;\sigma\delta}v_2^\sigma v_2^\delta\hat\delta(k_2\cdot v_1)\hat\delta(k_1\cdot v_2)\hat\delta(k_3\cdot v_2)}{k_1^2(k_2^2-m^2)(k_3^2-m^2)}e^{ik_2\cdot b_1}e^{i(k_1+k_3)\cdot b_2}\,,\\  &
     %   =i\frac{m_1 s_1 m_2^2 s_2}{4m_p^4}\rmint_{k_2,k_3}e^{ik_2\cdot b}\,k_2^\mu k_2^\alpha k_3^\beta\frac{\hat\delta(k_2\cdot v_1)\hat\delta(k_2\cdot v_2)\hat\delta(k_3\cdot v_2)}{(k_2+k_3)^2(k_2^2-m^2)(k_3^2-m^2)}P_{\alpha\beta;\sigma\delta}v_2^\sigma v_2^\delta\\ &
        =-i\frac{m_1 s_1 m_2^2 s_2}{8m_p^4}\rmint_{k_3,k_2}e^{ik_2\cdot b}k_2^\mu\,\hat\delta(k_2\cdot v_1)\hat\delta(k_2\cdot v_2)\frac{k_2\cdot k_3\,\hat\delta(k_3\cdot v_2)}{(k_2^2-m^2)(k_2+k_3)^2(k_3^2-m^2)}\,,\\ &
      %  =-i\frac{m_1 s_1 m_2^2 s_2}{16m_p^4}\rmint_{k_2}e^{ik_2\cdot b}\frac{k_2^\mu\, k_2^\alpha\,\hat\delta(k_2\cdot v_1)\hat\delta(k_2\cdot v_2)}{k_2^2-m^2}\rmint_{k_3}\frac{k_{3\alpha}\,\hat\delta(k_3\cdot v_2)}{(k_3+k_2)^2(k_3^2-m^2)}\\ &
      %  =-i\frac{m_1 s_1 m_2^2 s_2}{16m_p^4}\rmint_{k_2}e^{ik_2\cdot b}\frac{k_2^\mu\, k_2^\alpha\,\hat\delta(k_2\cdot v_1)\hat\delta(k_2\cdot v_2)}{k_2^2-m^2}\Bigg[\rmint_{k}\frac{k_\alpha\,\hat\delta(k\cdot v_2)}{k^2[(k-k_2)^2-m^2]}-k_{2\alpha}\rmint_{k}\frac{\hat\delta(k\cdot v_2)}{k^2[(k-k_2)^2-m^2]}\Bigg]\\ &
      %  =-\frac{m_1 s_1 m_2^2 s_2}{32m_p^4(2\pi)^3}\frac{1}{\sqrt{\gamma^2-1}}\frac{b^\mu}{|b|}\partial_{|b|}\rmint_{k_2}e^{-i\vec k_2\cdot \vec b}\frac{-\vec k_2^2}{\vec (k_2^2+m^2)|\vec k_2|}\arctan\Big(\frac{|\vec k_2|}{m}\Big)\,,\\ &
        =\frac{m_1 s_1 m_2^2 s_2}{128 \pi^2 m_p^4}\frac{1}{\sqrt{\gamma^2-1}}\frac{b^\mu}{|b|}\partial_{|b|}\underbrace{\rmint_0^\infty dk \frac{k^2}{k^2+m^2}J_{0}(k|b|)\arctan\Big(\frac{k}{m}\Big)}_{\bar{I}_4(m,|b|)}\,.
    \end{split}
\end{align}
Hence, the contribution from the diagram reads,

\begin{equation}\label{e:barwq24}
  [\Delta p_1^{\mu}]_{(h)}=\frac{m_1 s_1 m_2^2 s_2}{128 \pi^2 m_p^4}\frac{1}{\sqrt{\gamma^2-1}}\frac{b^\mu}{|b|}\partial_{|b|}\bar{I}_4(m,|b|)
\,.
\end{equation}
%\vspace{0.2 cm}
\textbullet $\,\,$ Another interesting 3-point worldline vertex, which comes from the derivative interaction,  contributing to the impulse at 2PM order, 
\begin{align}
    \begin{split} \label{vert1}
        \mathcal{O}(z): -\frac{m_1s_1}{2m_p^2}\rmint_{k_1,k_2,\omega}e^{i(k_1+k_2)\cdot b_1}\hat\delta(k_1\cdot v_1+k_2\cdot v_1+\omega)\{(k_{1\rho}+k_{2\rho})v^{\mu}v^{\nu}+2\omega v^{(\mu}\delta_{\rho}^{\nu)}\}\varphi(-k_1) h_{\mu\nu}(-k_2)z^{\rho}(-\omega)\,.
    \end{split}
\end{align}
The corresponding contribution to the impulse is given by, 
\begin{align}
    \begin{split}
         [\Delta p_1^{\mu}]_{(i)}=\thesisinlinefeynman[\chthreeinlinefeynwidth]{z_1_6.png}\hspace{0.7 cm}&=i\frac{m_1s_1}{2m_p^2}\times \frac{m_2}{2m_p}\times \frac{m_2s_2}{2m_p}\\ &
\hspace{0.35cm}\rmint_{k_{1,2}}(k_1^\mu+k_2^\mu)e^{i(k_1+k_2)\cdot b}\frac{\hat\delta(k_1\cdot v_1+k_2\cdot v_1)\hat\delta(k_1\cdot v_2)\hat\delta(k_2\cdot v_2)}{k_2^2(k_1^2-m^2)}P_{\alpha\beta;\rho\sigma}v_1^\alpha v_1^\beta v_2^\rho v_2^\sigma\,,\\ &=\frac{i m_1s_1m_2^2s_2}{8m_p^4}P_{\alpha\beta;\rho\sigma}v_1^\alpha v_1^\beta v_2^\rho v_2^\sigma\rmint_{q,k_2}\frac{q^\mu\hat\delta(q\cdot v_1)\hat\delta(q\cdot v_2)\hat\delta(k_2\cdot v_2)}{k_2^2[(q-k_2)^2-m^2]}e^{iq\cdot b}\,,\\ &
%=\frac{m_1 s_1 m_2^2 s_2}{16(2\pi)^5 m_p^4\sqrt{\gamma^2-1}}(2\gamma^2-1)\frac{b^\mu}{|b|}\partial_{b}\rmint_{q}e^{iq\cdot b}\hat\delta(q\cdot v_1)\hat\delta(q\cdot v_2)\frac{1}{|\vec q| } \arctan\Big(\frac{|\vec q|}{m}\Big)\,,\\ &
=\frac{m_1 s_1 m_2^2 s_2}{64\pi^2 m_p^4\sqrt{\gamma^2-1}}(2\gamma^2-1)\frac{b^\mu}{|b|}\partial_{|b|}\rmint_0^{\infty} dl\,J_{0}(|b|l)\arctan\Big(\frac{l}{m}\Big)\,.
    \end{split}
\end{align}
\begin{equation}\label{sa}
 [\Delta p_1^{\mu}]_{(i)}=\frac{m_1 s_1 m_2^2 s_2}{64\pi^2 m_p^4\sqrt{\gamma^2-1}}(2\gamma^2-1)\frac{b^\mu}{|b|}\partial_{|b|}\rmint_0^{\infty} dl\,J_{0}(|b|l)\arctan\Big(\frac{l}{m}\Big)\,.\vspace{0.5cm}
\end{equation}\\\
\textcolor{black}{The integral in \eqref{sa} has a closed-form expression, which can be seen by taking \textbf{logarithmic }representation of  \textbf{ArcTan} which gives,
\begin{align}
    \begin{split}
        [\Delta p_1^{\mu}]_{(i)}&=\frac{m_1 s_1 m_2^2 s_2}{64\pi m_p^4\sqrt{\gamma^2-1}}(2\gamma^2-1)\frac{b^\mu}{8\pi|b|^2}\\ &\Bigg[2 G_{3,1}^{1,3}\left(\frac{i}{{m\,b}},\frac{1}{2}\Big|
\begin{array}{c}
 1,1,\frac{3}{2} \\
 \frac{3}{2} \\
\end{array}
\right)+2 G_{3,1}^{1,3}\left(-\frac{i}{{m\,b}},\frac{1}{2}\Big|
\begin{array}{c}
 1,1,\frac{3}{2} \\
 \frac{3}{2} \\
\end{array}
\right)+G_{4,2}^{1,4}\left(\frac{i}{{m\,b}},\frac{1}{2}\Big|
\begin{array}{c}
 1,1,\frac{3}{2},\frac{3}{2} \\
 \frac{3}{2},\frac{1}{2} \\
\end{array}
\right)\\ &+G_{4,2}^{1,4}\left(-\frac{i}{{m\,b}},\frac{1}{2}\Big|
\begin{array}{c}
 1,1,\frac{3}{2},\frac{3}{2} \\
 \frac{3}{2},\frac{1}{2} \\
\end{array}
\right)\Bigg]\,.
    \end{split}
\end{align}}
Here have again used that fact that $$P_{\alpha\beta;\rho\sigma}v_1^\alpha v_1^\beta v_2^\rho v_2^\sigma=\frac{2\gamma^2-1}{2}\,.$$
Again, we have a finite massless counterpart of the diagram which gives,
\begin{align}
    \begin{split}
         [\Delta p_1^{\mu}]_{(i)}\Big|_{m\to 0}=-\frac{m_1 s_1 m_2^2 s_2}{64\pi^2 m_p^4\sqrt{\gamma^2-1}}(2\gamma^2-1)\frac{b^\mu}{|b|^3}\,.\label{4.33a}
    \end{split}
\end{align}
\textbullet $\,\,$  Lastly, there will be another diagram where $h_{\mu\nu}$ in (\ref{vert1}), is replaced by $\varphi$ in the worldline vertex. The worldline vertex under consideration is,
\begin{align}
    \begin{split}
          \mathcal{O}(z):& -\frac{m_1g_1}{2m_p^2}\rmint_{k_1,k_2,\omega}e^{i(k_1+k_2)\cdot b_1}\hat\delta(k_1\cdot v_1+k_2\cdot v_1+\omega)\{(k_{1\rho}+k_{2\rho})v^{\mu}v^{\nu}+2\omega v^{(\mu}\delta_{\rho}^{\nu)}\}\\&\hspace{0.8cm}\times\varphi(-k_1)\varphi(-k_2) \eta_{\mu\nu}z^{\rho}(-\omega)\,.
    \end{split}
\end{align}
Therefore, the impulse is given by,
\begin{align}
    \begin{split}
         [\Delta p_1^{\mu}]_{(j)}=\thesisinlinefeynman[\chthreeinlinefeynwidth]{z_1_7.png}\hspace{0.7 cm}&=i\frac{m_1g_1}{m_p^2}\times \frac{m_2^2 s_2^2}{2m_p^2}\rmint_{q,k_2}\frac{q^\mu\hat\delta(q\cdot v_1)\hat\delta(q\cdot v_2)\hat\delta(k_2\cdot v_2)}{(k_2^2-m^2)[(q-k_2)^2-m^2]}e^{iq\cdot b}\,,\\ &
=\frac{m_1g_1 m_2^2 s_2^2}{8\pi^2 m_p^4\sqrt{\gamma^2-1}}\frac{b^\mu}{|b|}\partial_{|b|}\rmint_0^\infty dl\,J_{0}(|b|l)\arctan\Big(\frac{l}{2m}\Big)\,.
    \end{split}
\end{align}
Therefore, the contribution to the impulse from the above diagram is given by,\\\
 
\begin{equation}\label{e:barwqp}
  [\Delta p_1^{\mu}]_{(j)}=\frac{m_1g_1 m_2^2 s_2^2}{8\pi^2 m_p^4\sqrt{\gamma^2-1}}\frac{b^\mu}{|b|}\partial_{|b|}\rmint_0^\infty dl\,J_{0}(|b| l)\arctan\Big(\frac{l}{2m}\Big)\,.\vspace{0.5cm}
\end{equation}
\textcolor{black}{Note that, similar to (\ref{sa}), this integral can also be recast in terms of {\bf MeijerG}. } 
The massless limit of  \eqref{4.33a} can be taken and it gives the following,
\begin{align}
    \begin{split}
         [\Delta p_1^{\mu}]_{(j)}\Big|_{m\to 0}=\frac{m_1g_1 m_2^2 s_2^2}{16 \pi m_p^4\sqrt{\gamma^2-1}}\frac{b^\mu}{|b|^3}\,.
    \end{split}
\end{align}

Finally, there will be 2PM diagrams coming from the bulk $\varphi^3$  interaction vertex. Below we will give the details of the contribution to the impulse coming from these vertices.

\textbullet $\,\,$ First, we consider the following Feynman diagram.
\begin{align}
    \begin{split}
       [\Delta p_1^{\mu}]_{(k)}&=\thesisinlinefeynman[\chthreeinlinefeynwidth]{new_6.png}\\&=-i\lambda_3 m_p\Big(\frac{m_1 g_1}{2m_p^2}\Big)\Big(\frac{m_2^3s_2^3}{8m_p^3}\Big)\rmint_{k_i}\hat\delta\Big(\sum_{i=1}^3 k_i\Big)\frac{(k_1+k_4)^\mu\hat\delta(k_1\cdot v_1+k_4\cdot v_1)\hat\delta(k_2\cdot v_2)\hat\delta(k_3\cdot v_2)\hat\delta(k_4\cdot v_2)}{\prod_{i}^4 (k_i^2-m^2)}\\ &\hspace{0.8cm}\times e^{i(k_1+k_4)\cdot b_1}e^{i(k_2+k_3)\cdot b_2}e^{-ik_4\cdot b_2}\,.\label{4.51b}
    \end{split}
\end{align}
\vspace{-0.3 cm}
Focusing only on the integral, we get,
\begin{align}
    \begin{split}
        [\Delta p_1^{\mu}]_{(k)}^\mu&\sim\rmint_{k_i}\frac{(k_1+k_4)^\mu \hat\delta(k_1\cdot v_1+k_4\cdot v_1)\hat\delta(k_2\cdot v_2)\hat\delta(k_1\cdot v_2)\hat\delta(k_4\cdot v_2)}{(k_1^2-m^2)(k_2^2-m^2)[(k_1+k_2)^2-m^2](k_4^2-m^2)} e^{i(k_1+k_4)\cdot b}\,,\\ &
        \xrightarrow[]{k_1+k_4\to q}\rmint_{k_i,q}\frac{q^\mu\hat\delta(q\cdot v_1)\hat\delta(k_2\cdot v_2)\hat\delta(k_1\cdot v_2)\hat\delta(q\cdot v_2)}{(k_1^2-m^2)(k_2^2-m^2)[(k_1+k_2)^2-m^2][(q-k_1)^2-m^2]}e^{i q\cdot b}\,,\\ &
        =\rmint _q q^\mu \hat{\delta}(q\cdot v_1)\hat\delta(q\cdot v_2)e^{i q \cdot b}\rmint_{k_1,k_2}\frac{\hat\delta(k_1\cdot v_2)\hat\delta(k_2\cdot v_2)}{(k_1^2-m^2)(k_2^2-m^2)[(k_1+k_2)^2-m^2][(q-k_1)^2-m^2]}\,,\\ &
        =\frac{1}{\sqrt{\gamma^2-1}}\frac{b^\mu}{|b|}\partial_{|b|}\rmint \hat d^2 q\, e^{-i \vec q\cdot \vec b}\, \hat L(\vec q,m)
    \end{split}
\end{align}
where,
\begin{align}
    \begin{split}
        \hat L(\vec q,m)&=\rmint_{\vec k_1,\vec k_2}\frac{1}{(\vec k_1^2+m^2)[(\vec k_1-\vec q)^2+m^2]|\vec k_1|}\arctan\Big(\frac{|\vec k_1|}{2m}\Big)\,,\\ &
        =\rmint_{0}^\infty d\alpha \,e^{-\alpha m^2}\rmint_{\vec k_1}\frac{1}{|\vec k_1|(\vec k_1^2+m^2)} e^{-\alpha(\vec k_1-\vec q)^2}\arctan\Big(\frac{|\vec k_1|}{2m}\Big)\,,\\ &
     %   =2\pi\rmint_{0}^\infty\frac{ d\alpha }{\alpha}\, e^{-\alpha (m^2+q^2)}\rmint_{0}^\infty dk_1 \frac{1}{k_1^2+m^2}e^{-\alpha k_1^2}\frac{\sinh \left(2 \alpha  k_1 q\right)}{   q}\arctan\Big(\frac{|\vec k_1|}{2m}\Big)\,,\\ &
        =\frac{2\pi}{q} \rmint dk_1 \frac{1}{k_1^2+m^2}\tanh ^{-1}\left(\frac{2 k_1 q}{k_1^2+m^2+q^2}\right)\arctan\Big(\frac{|\vec k_1|}{2m}\Big)\,.
    \end{split}
\end{align}
Therefore, the impulse is (after restoring all the prefractors),

\begin{equation}
 [\Delta p_1^{\mu}]_{(k)}=\frac{\lambda_3}{(2\pi)^4} \Big(\frac{m_1 g_1}{2m_p}\Big)\Big(\frac{m_2^3s_2^3}{8m_p^3}\Big)\frac{1}{\sqrt{\gamma^2-1}}\frac{b^\mu}{|b|}\partial_{|b|}\rmint_0^{\infty} dq \,dk_1 \,\frac{J_{0}(q|b|)}{k_1^2+m^2}\,\textrm{arctanh}\Big(\frac{2 k_1 q}{k_1^2+q^2+m^2}\Big)\arctan\Big(\frac{k_1}{2m}\Big)\label{4.54r}. \vspace{0.5cm}
\end{equation}
The integral in \eqref{4.54r} can be further simplified using the logarithmic representation of ArcTan as shown in \eqref{4.25j}. The massless limit also can be taken as \eqref{4.25j}.


%\newpage 
\textbullet $\,\,$Another Feynman diagram that we will contribute can be obtained by replacing one scalar propagator in (\ref{4.51b}) with a graviton propagator. The impulse is given by : 
\begin{align}
    \begin{split}
         [\Delta p_1^{\mu}]_{(l)}&=\thesisinlinefeynman[\chthreeinlinefeynwidth]{new_9.png}\\ &= \frac{2\gamma^2-1}{2}\frac{\lambda_3}{(2\pi)^4}\Big(\frac{m_1 s_1}{2m_p}\Big)\Big(\frac{m_2^3 s_2^2}{8m_p^3}\Big)\frac{1}{\sqrt{\gamma^2-1}}\frac{b^\mu}{|b|}\partial_{|b|}\rmint_0^\infty dq  dk_1 \frac{J_0(q|b|)}{k_1^2+m^2}\arctan\Big(\frac{k_1}{2m}\Big)\, \textrm{arctanh}\Big(\frac{2k_1 q}{k_1^2+q^2}\Big)\,.\label{4.55m}
\end{split}
\end{align}
Again this integral in \eqref{4.55m} can be further simplified using the logarithmic representation of ArcTan as \eqref{4.25j}.

%\newpage
\textbullet $\,\,$ Another Feynman topology contributing to the  impulse at 2PM order is given by: 
\begin{align}
    \begin{split}
         [\Delta p_1^{\mu}]_{(m)}&=\thesisinlinefeynman[\chthreeinlinefeynwidth]{new_8.png}\\ &= -i\lambda_3\Big(\frac{m_1^2 s_1^2}{4m_p}\Big)\Big(\frac{m_2 s_2 m_2 g_2 }{4m_p^3}\Big)\rmint_{k_i}\frac{k_1^\mu \hat\delta(k_1\cdot v_1)\hat\delta(k_2\cdot v_2)\hat\delta(k_3\cdot v_2+k_4\cdot v_2)\hat\delta(k_4\cdot v_1)}{\prod_i (k_i^2-m^2)}\hat\delta^{(4)}(k_1+k_2+k_3)\\ &
\hspace{0.8cm}\times e^{i(k_1-k_4)\cdot b_1}e^{i(k_2+k_3+k_4)\cdot b_2}\,,\\ &
        =-i\lambda_3\Big(\frac{m_1^2 s_1^2}{4m_p}\Big)\Big(\frac{m_2 s_2 m_2 g_2 }{4m_p^3}\Big)\rmint_{k_i}\frac{k_1^\mu \hat\delta(k_1\cdot v_1)\hat\delta(k_2\cdot v_2)\hat\delta(-k_1\cdot v_2-k_2\cdot v_2+k_4\cdot v_2)\hat\delta(k_4\cdot v_1)}{(k_1^2-m^2)(k_2^2-m^2)[(k_1+k_2)^2-m^2](k_4^2-m^2)}e^{i(k_1-k_4)\cdot b}\,,\\ &
        \xrightarrow[]{k_1-k_4\to q}-i\lambda_3\Big(\frac{m_1^2 s_1^2}{4m_p}\Big)\Big(\frac{m_2 s_2 m_2 g_2 }{4m_p^3}\Big)\rmint_{k_i,q}\frac{k_1^\mu\hat\delta(k_1\cdot v_1)\hat\delta(k_2\cdot v_2)\hat\delta(q\cdot v_2)\hat\delta(q\cdot v_1)}{(k_1^2-m^2)(k_2^2-m^2)[(k_1+k_2)^2-m^2][(q-k_1)^2-m^2]}e^{iq\cdot b}\,.
    \end{split}\label{4.56e}
\end{align}
Now, doing the $k_2$ integral, we will have (omitting the prefactors),
\begin{align}
    \begin{split}
         [\Delta p_1^{\mu}]_{(m)}\propto \rmint_{k_i,q}\frac{k_1^\mu\hat\delta(k_1\cdot v_1)\hat\delta(q\cdot v_2)\hat\delta(q\cdot v_1)}{(k_1^2-m^2)[(q-k_1)^2-m^2]|\vec k_1|}\arctan\Big(\frac{|\vec k_1|}{2m}\Big)e^{iq\cdot b}\,.
    \end{split}
\end{align}
\textcolor{black}{We see that due to the asymmetric structure of this diagram, it is difficult to obtain a closed-form expression, like the case of $\lambda_4 \varphi^4$ vertex. We first do the $q$ integral. It is clear from the delta function constraints that $q$ and $k_1$ are two-dimensional and three-dimensional vectors, respectively.} Therefore, after performing the $q$ integral, the result will only depend on the second and third components of $\vec k_1$. Therefore, we are left with the following, 
\begin{align}
    \begin{split}
          [\Delta p_1^{\mu}]_{(m)}&\propto 2 \pi \frac{1}{\sqrt{\gamma^2-1}}\rmint_{k_1} e^{-i ({b^{(2)}} {k_1^{(2)}}+{b^{(3)}}{k_1^{(3)}})} K_0\left(\sqrt{{b^{(2)}}^2+{b^{(3)}}^2} \sqrt{\vec k_1^2-{k_1^{(2)}}^2-{k_1^{(3)}}^2+m^2}\right)\\ &\hspace{0.8cm}\times\frac{k_1^\mu}{(k_1^2-m^2)|\vec k_1|}\arctan\Big(\frac{|\vec k_1|}{2m}\Big)\,.\label{4.57mm}
    \end{split}
\end{align}
Now, note that in our parametrisation, the impact parameter takes the form: $b^{\mu}=(0,0,1,0)$. Therefore, restoring the prefactors integral in \eqref{4.57mm} can be recasted as,
\begin{align}
    \begin{split}
         [\Delta p_1^{\mu}]_{(m)}&=\frac{\lambda_3}{(2\pi)^4}\Big(\frac{m_1^2 s_1^2}{4m_p}\Big)\Big(\frac{m_2 s_2 m_2 g_2 }{4m_p^3}\Big)\frac{1}{\sqrt{\gamma^2-1}}\rmint d^4 k_1 e^{ik_1\cdot b}\,K_{0}\Big(|b||\sqrt{k_{1(1)}^2+m^2}|\Big)\\ &
        \hspace{0.8cm}\times\frac{k_1^\mu\hat\delta(k_1\cdot v_1)}{(k_1^2-m^2)|\vec k_1|}\arctan\Big(\frac{|\vec k_1|}{2m}\Big)\,,\\ &
        =\frac{\lambda_3}{(2\pi)^4}\Big(\frac{m_1^2 s_1^2}{4m_p}\Big)\Big(\frac{m_2 s_2 m_2 g_2 }{4m_p^3}\Big)\frac{1}{\gamma\sqrt{\gamma^2-1}}\frac{b^\mu}{|b|}\partial_{|b|}\rmint d^3 k_1 e^{-i \vec k_1\cdot \vec b}K_{0}\Big(|b||\sqrt{k_{1(1)}^2+m^2}|\Big)\\ &
        \hspace{0.8cm}\times\frac{1}{(\bar k_1^2+m^2)|\vec k_1|}\arctan\Big(\frac{|\vec k_1|}{2m}\Big)\,.\label{4.58}
    \end{split}
\end{align}
Hence, the contribution has the following form,

\begin{align}
\begin{split}
  [\Delta p_1^{\mu}]_{(m)}&=\frac{\lambda_3}{(2\pi)^4}\Big(\frac{m_1^2 s_1^2}{4m_p}\Big)\Big(\frac{m_2 s_2 m_2 g_2 }{4m_p^3}\Big)\frac{1}{\gamma\sqrt{\gamma^2-1}}\frac{b^\mu}{|b|}\partial_{|b|}\rmint_{-\infty}^{\infty} d^3 k_1 e^{-i \vec k_1\cdot \vec b}K_{0}\Big(|b||\sqrt{k_{1(1)}^2+m^2}|\Big)\\ &\hspace{0.8cm}\times \frac{1}{(\bar k_1^2+m^2)|\vec k_1|}\arctan\Big(\frac{|\vec k_1|}{2m}\Big)\label{4.58o} \vspace{0.5cm}
\end{split}
\end{align}
where, $\bar k_1=(k_1^{(1)}/\gamma,k_1^{(2)},k_1^{(3)})$. \textcolor{black}{Again, we can see that the integral in \eqref{4.58} has an asymmetry in the different components of the momentum integration and therefore, to the best of our knowledge, does not have a closed-form. Further, in Appendix~\eqref{App:3.H} we progress on computation of the integral more systematically. We also comment on the convergence of the integral.} Although, for the massless case, we also do not have any closed-form due to the asymmetry, but one can take a smooth massless limit. 


%\newpage
\textbullet $\,\,$ Similar to \eqref{4.56e} we will have another contributing diagram where the $k_4$ scalar propagator will be replaced by a graviton propagator. Corresponding contribution to the impulse takes the following form,
\begin{align}
    \begin{split}
         [\Delta p_1^{\mu}]_{(n)} &= \thesisinlinefeynman[\chthreeinlinefeynwidth]{feynman-4.png} \\&=\frac{\lambda_3}{(2\pi)^4}\Big(\frac{m_1^2 s_1^2}{4m_p}\Big)\Big(\frac{m_2^2 s_2^2  }{4m_p^3}\Big)\frac{2\gamma^2-1}{2\gamma\sqrt{\gamma^2-1}} \frac{b^\mu}{|b|}\partial_{|b|}\rmint d^3 k_1 e^{-i \vec k_1\cdot \vec b}K_{0}\Big(|b||{k_{1(1)}}|\Big)\frac{1}{(\bar k_1^2+m^2)|\vec k_1|}\arctan\Big(\frac{|\vec k_1|}{2m}\Big)\,.
    \end{split}
\end{align}


\begin{equation}
  [\Delta p_1^{\mu}]_{(n)}=\frac{\lambda_3}{(2\pi)^4}\Big(\frac{m_1^2 s_1^2}{4m_p}\Big)\Big(\frac{m_2^2 s_2^2  }{4m_p^3}\Big)\frac{2\gamma^2-1}{2\gamma\sqrt{\gamma^2-1}} \frac{b^\mu}{|b|}\partial_{|b|}\rmint_{-\infty}^{\infty} d^3 k_1 e^{-i \vec k_1\cdot \vec b}K_{0}\Big(|b||{k_{1(1)}}|\Big)\frac{1}{(\bar k_1^2+m^2)|\vec k_1|}\arctan\Big(\frac{|\vec k_1|}{2m}\Big)\label{4.63pp}. \vspace{0.5cm}
\end{equation}
Again, due to the reason mentioned above, this integral does not possess any closed-form expression. One has to do it numerically componentwise. We will comment on the derivation in the Appendix.

\textbullet $\,\,$ We have another type of vertex $\lambda_3 h \varphi^3$, which also contributes to the 2PM diagrams. Calculations are similar to the computations of $\lambda_4 \varphi^4$ vertex \eqref{4.26m}.
\begin{align}
    \begin{split}
         [\Delta p_1^{\mu}]_{(o)}&=\thesisinlinefeynman[\chthreeinlinefeynwidth]{new_7.png}\\ &=-i\lambda_3\Big(\frac{m_1^2 s_1 (m_2 s_2)^2}{16m_p^4}\Big)P^{\alpha}_{\,\,\alpha,\rho\sigma}v_1^\rho v_1^\sigma\rmint_{\{k_i\}}\hat\delta^{(4)}\Big(\sum_{i=1}^{4}k_{i}\Big)\hat\delta(k_1\cdot v_1)\hat\delta(k_2\cdot v_1)\hat\delta(k_3\cdot v_2)\hat\delta(k_4\cdot v_2)\\ &\hspace{0.8cm}\times\prod_{j=1,3,4}\frac{1}{(k_j^2-m^2)k_2^2} 
 k_1^{\mu} \,e^{i(k_1+k_2)\cdot b_2}e^{i(k_3+k_4)\cdot b_2}\,.\label{4.55e}
    \end{split}
\end{align}
Now, \eqref{4.55e} has the same form as \eqref{4.26m} except that one denominator is massless. The result is given by,

\begin{align}
\begin{split}
  [\Delta p_1^{\mu}]_{(o)}&=-\frac{\lambda_3}{(2\pi)^3}\Big(\frac{m_1^2 s_1 (m_2 s_2)^2}{16m_p^4}\Big)\frac{1}{\sqrt{\gamma^2-1}}\frac{b^\mu}{|b|}\rmint_0^{\infty} dx {J_{1}(x|b|)}\arctan\Big(\frac{x}{2m}\Big)\arctan\Big(\frac{x}{m}\Big)\,,\\ &
  =-\frac{\lambda_3}{(2\pi)^3}\Big(\frac{m_1^2 s_1 (m_2 s_2)^2}{8 m_p^4}\Big)\frac{1}{\sqrt{\gamma^2-1}}\frac{b^\mu}{|b|^2}\int_0^1 ds\,\left(\frac{K_0(2m|b|)-K_0\left(\frac{m|b|}{s}\right)}{1-4s^2}+\frac{K_0(m|b|)-K_0\left(\frac{2m|b|}{s}\right)}{4-s^2}\right)\,.\label{3.99kk}
  \end{split}
\end{align}
\textcolor{black}{Like before, although this integral does not possess any closed-form expression, it has a smooth massless limit. In the massless limit, the above diagram can be calculated exactly as},
\begin{align}
    \begin{split}
         [\Delta p_1^{\mu}]_{(o)}\Big|_{m\rightarrow 0}=-{\lambda_3}\Big(\frac{m_1^2 s_1 (m_2 s_2)^2}{1024 \pi m_p^4}\Big)\frac{1}{\sqrt{\gamma^2-1}}\frac{b^\mu}{|b|^3}\,.
        \end{split}
    \end{align} 
\textbullet $\,\, $ There could be another possibility like \eqref{4.19c} where we have one line (graviton) connected with the first worldline and the other three (scalars) are connected with the second worldline. Now, the corresponding contribution to the impulse reads,
\begin{align}
    \begin{split}
         [\Delta p_1^{\mu}]_{(p)}&=\thesisinlinefeynman[\chthreeinlinefeynwidth]{new_5.png}\\ &= -i\lambda_3\Big(\frac{m_1  (m_2 s_2)^3}{16m_p^4}\Big)P^{\alpha}_{\,\,\alpha,\rho\sigma}v_1^\rho v_1^\sigma \rmint_{\{k_i\}}\hat\delta^{(4)}\Big(\sum_{i=1}^{4}k_{i}\Big)\hat\delta(k_1\cdot v_1)\hat\delta(k_2\cdot v_2)\hat\delta(k_3\cdot v_2)\hat\delta(k_4\cdot v_2)\\ &\hspace{0.8cm}\prod_{j=2}^{4} \frac{1}{(k_j^2-m^2)k_1^2}
\times k_1^{\mu} \,e^{ik_1\cdot b_1}e^{i(k_2+k_3+k_4)\cdot b_2}\,.\label{4.57a}
    \end{split}
\end{align}
Following the derivation of \eqref{4.19c} we will get,

\begin{equation}
 [\Delta p_1^{\mu}]_{(p)}=\frac{\lambda_3}{(2\pi)^3}\Big(\frac{m_1  (m_2 s_2)^3}{16m_p^4}\Big)\frac{1}{\sqrt{\gamma^2-1}}\frac{b^\mu}{|b|}\partial_{|b|}\rmint_0^\infty dq \,dk_1\,\frac{J_{0}(k_1 |b|)}{k_1^2}\arctan\Big(\frac{q}{2m}\Big)\textrm{arctanh}\Big(\frac{2qk_1}{m^2+q^2+k_1^2}\Big)\,.\label{4.58a} \vspace{0.5cm}
\end{equation}
The integral in \eqref{4.58a} can be further simplified using the logarithmic representation of ArcTan. The massless limit also can be taken as \eqref{4.25j}.


\textbullet $\,\,$ There will be another diagram topologically equivalent to \eqref{4.57a} contributing to the impulse where one line (scalar) connects with one worldline and the other three (one graviton and two scalars) connect with the other worldline. The corresponding contribution is,
\begin{align}
    \begin{split}
      [\Delta p_1^{\mu}]_{(q)} &=\thesisinlinefeynman[\chthreeinlinefeynwidth]{new_4.png}\\ &= -i\lambda_3\Big(\frac{m_1s_1 m_2 (m_2 s_2)^2}{16m_p^4}\Big)P^{\alpha}_{\,\,\alpha,\rho\sigma}v_2^\rho v_2^\sigma \rmint_{\{k_i\}}\hat\delta^{(4)}\Big(\sum_{i=1}^{4}k_{i}\Big)\hat\delta(k_1\cdot v_1)\hat\delta(k_2\cdot v_2)\hat\delta(k_3\cdot v_2)\hat\delta(k_4\cdot v_2)\\ & 
\hspace{0.8cm}\times\prod_{j=2}^{4}\frac{1}{(k_j^2-m^2)k_1^2} k_1^{\mu} \,e^{ik_1\cdot b_1}e^{i(k_2+k_3+k_4)\cdot b_2}\,,
    \end{split}
\end{align}
Similarly, following the derivation of \eqref{4.19c} we have,

\begin{equation}
[\Delta p_1^{\mu}]_{(q)} =\frac{\lambda_3}{(2\pi)^3}\Big(\frac{m_1s_1 m_2 (m_2 s_2)^2}{16m_p^4}\Big)\frac{1}{\sqrt{\gamma^2-1}}\frac{b^\mu}{|b|}\partial_{|b|}\rmint_0^\infty dq \,dk_1\,\frac{J_{0}(k_1 |b|)}{k_1^2+m^2}\arctan\Big(\frac{q}{m}\Big)\textrm{arctanh}\Big(\frac{2qk_1}{m^2+q^2+k_1^2}\Big).\label{4.66t} \vspace{0.5cm}
\end{equation}
The integral in \eqref{4.66t} can be further simplified using the logarithmic representation of ArcTan. The massless limit also can be taken as \eqref{4.25j}.

\textbullet $\,\,$ So far, the diagrams that contributed to the impulse at 2PM order require expanding the point-particle action up to $\mathcal{O}({z})\,.$ However, there is one additional diagram that must be included to complete the 2PM computation, and it requires expanding the point-particle action up to $\mathcal{O}({z}^2)$. %Having two worldline propagators and making one of them on-shell ($\omega=0$) creates some subtlety in performing the integral. 
%\vspace{-0.66 cm}
\begin{align}
    \begin{split}
       [\Delta p_1^{\mu}]_{(r)} =\thesisinlinefeynman[\chthreeinlinefeynwidth]{z_2_impulse.png}\hspace{0.7 cm}&= -im_1\Big(\frac{s_1m_2 s_2}{2\sqrt{2}m_p^2}\Big)^2\rmint_{k_i,\tilde\omega}\hat\delta(k_1\cdot v_2)\hat\delta(k_2\cdot v_2)\hat\delta(-k_1\cdot v_1+\tilde\omega)\hat\delta(-k_2\cdot v_1-\tilde\omega+\omega)\\ &\hspace{0 cm}\times (2\tilde\omega v_1^{{\rho}}-k_1^{\rho}) (\frac{1}{2}k_{2\rho}k_{2}^\mu+\tilde\omega k_{2}^\mu v_{1\rho}-\omega k_{2\rho}v_{1}^\mu-\omega\,\tilde\omega\delta^\mu_{\rho})\frac{e^{i(k_1+k_2)\cdot b}}{(k_1^2-m^2)(k_2^2-m^2)\tilde\omega^2}\Big |_{\omega=0},\\ &
  %      =-im_1\Big(\frac{s_1m_2 s_2}{2\sqrt{2}m_p^2}\Big)^2\rmint_{k_1,k}\frac{(k-k_1)^{\mu}\hat\delta(k_1\cdot v_2)\hat\delta(k\cdot v_2)\hat\delta(k\cdot v_1)e^{ik\cdot b}}{(k_1^2-m^2)[(k-k_1)^2-m^2](k_1\cdot v_1)^2}[(k_1\cdot v_1)(k-k_1)\cdot v_1\\ & \hspace{0.8cm}+2(k_1\cdot v_1)^2-\frac{1}{2}k_1\cdot(k-k_1)-(k_1\cdot v_1)^2]\,,\\ &
        =im_1\Big(\frac{s_1m_2 s_2}{4m_p^2}\Big)^2\textcolor{black}{\rmint_{k_1,k}\frac{(k-k_1)^{\mu}\hat\delta(k_1\cdot v_2)\hat\delta(k\cdot v_2)\hat\delta(k\cdot v_1)e^{ik\cdot b}}{(k_1^2-m^2)[(k-k_1)^2-m^2](k_1\cdot v_1)^2}k_1\cdot(k-k_1)}\,.\label{4.6 mm}
    \end{split}
\end{align}
One can write down the integral as ,
\begin{align}
    \begin{split}
       [\Delta p_1^{\mu}]_{(r)} &=im_1\Big(\frac{s_1m_2 s_2}{4m_p^2}\Big)^2\rmint \hat d^4k\,\hat\delta(k\cdot v_1)\hat\delta(k\cdot v_2)e^{ik\cdot b}\rmint \hat d^4 k_1\,\hat\delta(k_1\cdot v_2)\frac{k_1\cdot(k-k_1)(k-k_1)^{\mu}}{(k_1^2-m^2)[(k-k_1)^2-m^2](k_1\cdot v_1)^2}\,,\\ &
        =im_1\Big(\frac{ s_1m_2 s_2}{4m_p^2}\Big)^2\rmint \hat d^4k\,\hat\delta(k\cdot v_1)\hat\delta(k\cdot v_2)e^{ik\cdot b}\mathcal{K}^{\mu}\label{4.37a}
    \end{split}
\end{align}
where $\mathcal{K}^{\mu}(k,v_i)$ can be reduced by Passarino-Veltman reduction as,
\begin{align}
    \begin{split}
        \mathcal{K}^{\mu}=a_1 k^{\mu}+a_2 v_1^{\mu}+a_3 v_2^{\mu}\,.\label{4.38a}
    \end{split}
\end{align}
It is evident that with the support $k\cdot v_1 =0\,\text{and } k\cdot v_2=0$, one can write $ v_{2\mu}\mathcal{K}^{\mu}=0$, implying, $a_2\gamma+a_3=0$. Hence, we left with the following,
\begin{align}
    \begin{split}
\mathcal{K}^{\mu}=a_2(v_1^{\mu}-\gamma v_2^{\mu})+a_1k^\mu\,.
    \end{split}
\end{align}
From \eqref{B.22aa} it is clear that the $a_1$ coefficient is non-zero for %massive case (However it gives pure divergence in massless limit and can be ignored in classical limit.) 
and is given by,
\begin{align}
    \begin{split}
     a_1=   \frac{1}{4}(-\vec k^2-2m^2)\hat\chi_k(|\vec k|,m)
    \end{split}
\end{align}
where, $\hat\chi_{k}(k,m)$ is defined in \eqref{B.16ab}. Similarly,
$a_2$ can be fixed by contracting both side by $v_1^{\mu}$,
\begin{align}
    \begin{split}
        a_2&=\frac{1}{\gamma^2-1}\rmint \hat d^{4}k_1 \frac{\hat\delta(k_1\cdot v_2)\,k_1\cdot(k-k_1)}{(k_1^2-m^2)[(k-k_1)^2-m^2](k_1\cdot v_1+i\epsilon)}\,,\\ &
        =\frac{1}{2(\gamma^2-1)}\rmint \hat d^4 k_1\hat\delta(k_1\cdot v_2)\Big[-\frac{1}{(k_1^2-m^2)(k_1\cdot v_1+i\epsilon)}-\frac{1}{(k-k_1)^2-m^2](k_1\cdot v_1+i\epsilon)}\\\ &
        \hspace{0.8cm} \frac{k^2}{(k_1^2-m^2)[(k-k_1)^2-m^2](k_1\cdot v_1+i\epsilon)}
      -\frac{2m^2}{(k_1^2-m^2)[(k-k_1)^2-m^2](k_1\cdot v_1+i\epsilon)}  \Big]\,.\label{3.13}
    \end{split}
\end{align}
The first two terms will not contribute. Now doing a change of variable, $k_1-k\rightarrow -q$, $a_2$ can be written as,
\begin{align}
    \begin{split}
        a_2=-\frac{1}{2(\gamma^2-1)}\rmint \hat d^4q \frac{(k^2-2m^2)\hat\delta(q\cdot v_2)}{(q^2-m^2)[(q-k)^2-m^2](q\cdot v_1-i\epsilon)}\,.\label{3.14}
    \end{split}
\end{align}
Now adding (\ref{3.13}) and (\ref{3.14}) and using the fact,
\begin{align}
    \begin{split}
       - i\hat \delta(x)=\frac{1}{x+i\epsilon}-\frac{1}{x-i\epsilon}
    \end{split}
\end{align}
one gets,
\begin{align}
    \begin{split}
        a_2=\frac{-i}{4(\gamma^2-1)}\rmint \hat d^4 q \frac{(k^2-2m^2)\hat\delta(q\cdot v_1)\hat\delta(q\cdot v_2)}{(q^2-m^2)[(q-k)^2-m^2]}\,.
    \end{split}
\end{align}
%\begin{align}
    %%N_1(k,k_1,v_i)=2k_1\cdot (k-k_1)=-[(k-k_1)^2-m^2]-(k_1^2-m^2)+k^2-2m^2
   % \end{split}
%\end{align}
Thereafter, $\Delta p_1^{\mu}$ can be written as,
\begin{align}
    \begin{split}
         \hspace{0cm}[\Delta p_1^{\mu}]_{(r)} &= -m_1\Big(\frac{s_1m_2 s_2}{8\sqrt{2}\sqrt{\gamma^2-1}m_p^2}\Big)^2\rmint \hat d^4 k\,\hat d^4 q\,\hat\delta(k\cdot v_1)\hat\delta(k\cdot v_2)\hat\delta(q\cdot v_1)\hat\delta(q\cdot v_2)\,e^{i k\cdot b}\,\frac{k^2-2m^2}{(q^2-m^2)[(q-k)^2-m^2]}\\ &
\hspace{0.8cm}\times (v_1^{\mu}-\gamma v_2^{\mu})+im_1\Big(\frac{s_1m_2 s_2}{4m_p^2}\Big)^2\rmint_{k}\hat\delta(k\cdot v_1)\hat\delta(k\cdot v_2)e^{ik\cdot b}k^{\mu}a_1\,,
         \\&
        =-\frac{m_1}{(2\pi)^4}\Big(\frac{s_1m_2 s_2}{8\sqrt{2}m_p^2}\Big)^2\frac{1}{(\gamma^2-1)^3}\rmint  d^2 k\,e^{-ik\cdot b}(-k^2-2m^2)\rmint  d^2 q\,\frac{1}{(q^2+m^2)[(q-k)^2+m^2]}(v_1^{\mu}-\gamma v_2^{\mu})\\ &
      \hspace{0.8cm}  -m_1\Big(\frac{s_1m_2 s_2}{8m_p^2}\Big)^2\frac{1}{\pi^{3/2}(\gamma^2-1)^{3/2}} \frac{b^\mu}{|b|}\Big(m K_1(2 |b| m)\Big)
        \,.
    \end{split}
\end{align}
The $`q$' integral can be done using dimensional regularisation with $d=2-\epsilon$. We need to be cautious about the UV (or IR) divergences that may appear while doing the integral. The integral can be written as,
\begin{align}
    \begin{split}
      I&=  \rmint {d}^{2-\epsilon}q \frac{1}{(q^2+m^2)[(q-k)^2+m^2]}\,,
      \\ &
       \xrightarrow[]{\epsilon \rightarrow 0} (2\pi)\frac{2 \log \left(\sqrt{\frac{k^2}{4 m^2}+1}+\frac{k}{2 m}\right)}{m\,k  \sqrt{\frac{k^2}{4 m^2}+1}}\,=(2\pi)\frac{4[\log(k+\sqrt{k^2+4m^2})-\log(2m)]}{k\sqrt{k^2+4m^2}}\,.
    \end{split}
\end{align}
%The integral can be divided into two regions ($k^2<4m^2$ and $k^2>4m^2$)  and has different way of computations. For $k^2<4 m^2$,
%\begin{align}
    %%%   \xrightarrow[]{\epsilon \rightarrow 0} \frac{4m^2\,\text{arcsinh}{(\frac{k}{2m}})}{k\sqrt{k^2+4m^2}}
  %%\end{align}
%and, for $k^2>4m^2$ we will have,
%%  \begin{split}
    %    I&\sim (-1)^{1-\frac{\epsilon}{2}}\Big[\Big]\\ &
     %   I&\sim (-1)^{1-\frac{\epsilon}{2}}\Big[\Big]\\ &
      % \xrightarrow[]{\epsilon \rightarrow 0} \frac{2 \left[\log \left(\frac{4 m^2}{k^2}+1\right)+2\log \left(\frac{k}{m}\right)\right]}{k^2 \sqrt{\frac{4 m^2}{k^2}+1}}-\frac{2\times \,\,_2F_1^{(\{0,0\},\{1\},0)}\left(\left\{\frac{1}{2},1\right\},\{1\},-\frac{4 m^2}{k^2}\right)}{k^2}
   % \end{split}
%\end{align}
Now, the integral $\Delta p_1^{\mu}$ can be expressed as,
\begin{align}
\begin{split} \label{eee22}
[\Delta p_1^{\mu}]_{(r)} =&-\frac{m_1}{(2\pi)^2}\Big(\frac{s_1m_2 s_2}{8\sqrt{2}\,m_p^2}\Big)^2\frac{1}{(\gamma^2-1)^3}(v_1^\mu-\gamma v_2^\mu)\rmint_{0}^{\infty}{d}k\,k\, J_{0}(k|b|)(k^2+2m^2)
   %\frac{4m^2\,\text{arcsinh}{(\frac{k}{2m})}}{\sqrt{k^2+4m^2}}&, |k|<2m \\
      \\&\hspace{0.8cm} \Bigg( \frac{4[\log(k+\sqrt{k^2+4m^2})-\log(2m)]}{k\sqrt{k^2+4m^2}}\Bigg)\\ &
      \hspace{0 cm}-m_1\Big(\frac{s_1m_2 s_2}{8m_p^2}\Big)^2\frac{1}{\pi^{3/2}\gamma^3(\gamma^2-1)^{3/2}} \frac{b^\mu}{|b|}\Big(m K_1(2 |b| m)\Big)\,.
\end{split}
\end{align}
%\vspace{0 cm}
In the massless limit we have,
\begin{align}
    \begin{split}
     [\Delta p_1^{\mu}]_{(r)} \Big|_{m\to 0}&= -\frac{m_1}{(2\pi)^2}\Big(\frac{s_1m_2 s_2}{4\sqrt{2}\,m_p^2}\Big)^2\frac{1}{(\gamma^2-1)^3}(v_1^\mu-\gamma v_2^\mu)\Bigg(\rmint_{0}^{\infty}dk \, k\,J_{0}(k|b|) \log(2 k)-4\frac
     {\log(2\epsilon)J_{1}(|b|\Lambda)}{|b|}\Bigg)\\ &
       \hspace{0.8cm}-m_1\Big(\frac{s_1m_2 s_2}{8m_p^2}\Big)^2\frac{1}{\pi^{3/2}(\gamma^2-1)^{3/2}} \frac{b^\mu}{|b|^2}
    \end{split}
\end{align}
\vspace{-0 cm}
where $\epsilon$ and $\Lambda$ are the IR and UV cut-off, respectively. The second term is a divergent one. In fact, UV/IR  divergence is mixed, i.e., it is divergent in both the limit  ${\epsilon}\rightarrow 0 \,\& \,\Lambda\rightarrow \infty\,.$ So we ignore this in the classical limit. The finite part takes the following form,
\begin{align}
    \begin{split}
    [\Delta p_1^{\mu}]_{(r)} \Big|_{m\to 0}\sim N_1\frac{(v_1^{\mu}-\gamma v_2^{\mu})}{4(\gamma^2-1)^3} \frac{1}{|b|^2}\textcolor{black}{+N_2 \frac{b^\mu}{(\gamma^2-1)^{3/2}|b|^2}}\,.
    \end{split}
\end{align}
%\newpage
The total impulse (due to the scalar field) at 2PM order is the sum of  (\ref{e:barwq244}), (\ref{e:barwq248}), (\ref{4.19f}), (\ref{e:barwq246}), (\ref{e:barwqpp}), (\ref{e:barwq24}), (\ref{sa}), (\ref{e:barwqp}), (\ref{4.54r}), (\ref{4.55m}), (\ref{4.58o}), (\ref{4.63pp}), (\ref{3.99kk}), (\ref{4.58a}), (\ref{4.66t}) and (\ref{eee22}) as well as the terms that come from interchanging the worldline one and two. 
\begin{align}
\begin{split}
     \Delta p_1^{\mu}\Big|^{\textrm{2PM}, \textrm{Total}}_{\textrm{scalar}}= & \sum_{\upsilon=c}^{r}[\Delta p_1^{\mu}]_{(\upsilon)}\,.
     \end{split}
\end{align}
\newpage
Finally, collecting all individual expressions we get,
\begin{align}
\Delta p_1^{\mu}\Big|^{\textrm{2PM},\textrm{total}}_{\textrm{scalar}}
&=\frac{m^2\, m_1m_2^2 s_2}{32\pi^2\,m_p^4\,\sqrt{\gamma^2-1}}\frac{b^\mu}{|b|}\partial_{|b|}\Big(s_2 I_2(m,|b|)+s_1 \bar{I}_2(m,|b|)\Big)\notag\\
&\quad+\frac{m_1 m_2^2 s_2}{8\,m_p^4}\frac{b^{\mu}}{|b|}\partial_{|b|}\Big(s_2 I_{4}(m,|b|)+\frac{s_1}{16 \pi^2 \sqrt{\gamma^2-1}}\bar{I}_4(m,|b|)\Big)\notag\\
&\quad+\frac{m_1 m_2^2 s_2}{8\pi^2m_p^4\sqrt{\gamma^2-1}}\frac{b^\mu}{|b|}\partial_{|b|}\Bigg[\frac{s_1\,(2\gamma^2-1)}{8}\rmint dl\,J_{0}(|b|l)\arctan\Big(\frac{l}{m}\Big)\notag\\
&\qquad\qquad+g_1 s_2\rmint_0^\infty dl\,J_{0}(|b|l)\arctan\Big(\frac{l}{2m}\Big)\Bigg]\notag\\
&\quad-\frac{m_1}{(2\pi)^2}\Big(\frac{s_1m_2 s_2}{4\sqrt{2}\,m_p^2}\Big)^2\frac{1}{(\gamma^2-1)^3}(v_1^\mu-\gamma v_2^\mu)\rmint_{0}^{\infty}{d}k\,k\,J_{0}(k|b|)(k^2+2m^2)\notag\\
&\qquad\qquad\times\Bigg(\frac{\log(k+\sqrt{k^2+4m^2})-\log(2m)}{k\sqrt{k^2+4m^2}}\Bigg)\notag\\
&\quad-m_1\Big(\frac{s_1m_2 s_2}{8m_p^2}\Big)^2\frac{1}{\pi^{3/2}(\gamma^2-1)^{3/2}} \frac{b^\mu}{|b|}\Big(m K_1(2 |b| m)\Big)\notag\\
&\quad+\frac{\lambda_4\,m_1s_1 (m_2s_2)^2}{128\pi^3m_p^4\sqrt{\gamma^2-1}}\frac{b^\mu}{|b|}\partial_{|b|}\Bigg[m_1s_1\,I_3(m,|b|)\notag\\
&\qquad\qquad+\frac{m_2 s_2}{4\gamma}\Big(2(1-\log(3m^2))K_{0}(|b|m)-\rmint dk_1\frac{J_{0}(k_1 |b|)}{k_1^2+m^2}\Theta(k_1,m)\Big)\Bigg]\notag\\
&\quad+\frac{\lambda_3}{(2\pi)^4}\Big(\frac{m_2^3s_2^3}{16m_p^4}\Big)\frac{1}{\sqrt{\gamma^2-1}}\frac{b^\mu}{|b|}\partial_{|b|}\Bigg[m_1 g_1 I_{5}(m,|b|)+\frac{(2\gamma^2-1)\,m_1 s_1}{2}\bar{I}_{5}(m,|b|)\Bigg]\notag\\
&\quad+\frac{\lambda_3}{(2\pi)^3}\Big(\frac{m_1 (m_2 s_2)^2}{16m_p^4}\Big)\frac{1}{\sqrt{\gamma^2-1}}\frac{b^\mu}{|b|}\partial_{|b|}\Bigg[-m_1s_1 I_{6}(m,|b|)+\frac{m_2 s_2}{\gamma}\bar{I}_{6}(m,|b|)+\frac{m_2 s_1}{\gamma}\bar{I}_{6}(m,|b|)\Bigg]\notag\\
&\quad+\frac{\lambda_3}{(2\pi)^4}\Big(\frac{m_1^2 s_1^2}{4m_p}\Big)\Big(\frac{m_2 s_2}{4m_p^3}\Big)\frac{1}{\gamma\,\sqrt{\gamma^2-1}}\frac{b^\mu}{|b|}\partial_{|b|}\Bigg[m_2g_2\,I_{7}(m,|b|)+\Big(\frac{m_2 s_2\,(2\gamma^2-1)}{2}\Big)\bar{I}_{7}(m,|b|)\Bigg]\notag\\
&\quad+(1\leftrightarrow 2)\,.
\end{align}
where, \begin{align}
    \begin{split}
    &I_2(m,|b|)=\rmint _{0}^{\infty}dx\,\frac{J_{0}(|b|x)}{x^2}\arctan\Big(\frac{x}{2m}\Big)\,,\quad 
    \bar{I}_{2}(m,|b|)=\rmint_{0}^\infty dx\frac{J_{0}(b|x|)}{x^2+m^2}\arctan\Big(\frac{x}{m}\Big)\,,\\ &
    I_{3}(m,|b|)=\rmint_{0}^\infty dx\, \frac{J_{0}(x|b|)}{x}\arctan^2\Big(\frac{x}{2m}\Big),\quad
    \bar{I}_{4}(m,|b|)=\rmint_0^\infty dx \frac{x^2}{x^2+m^2}J_{0}(x|b|)\arctan\Big(\frac{x}{m}\Big)\,,\\&
    I_{5}(m,|b|)=\rmint_0^{\infty} dx \,dz \,\frac{J_{0}(x|b|)}{z^2+m^2}\,\textrm{arctanh}\Big(\frac{2 x\,z}{x^2+z^2+m^2}\Big)\arctan\Big(\frac{z}{2m}\Big)\,,\\&
    \bar{I}_{5}(m,|b|)=\rmint_0^\infty dx  dz \frac{J_0(x|b|)}{z^2+m^2}\arctan\Big(\frac{z}{2m}\Big)\, \textrm{arctanh}\Big(\frac{2x\,z}{x^2+z^2}\Big)\,,\\&
    I_{6}(m,|b|)=\rmint_0^{\infty} dx \frac{J_{0}(x|b|)}{x}\arctan\Big(\frac{x}{2m}\Big)\arctan\Big(\frac{x}{m}\Big)\,,\\&
    \bar{I}_{6}(m,|b|)=\rmint_0^\infty dx \,dz\,\frac{J_{0}(z |b|)}{z^2}\arctan\Big(\frac{x}{2m}\Big)\textrm{arctanh}\Big(\frac{2x\,z}{m^2+x^2+z^2}\Big)\,,\\&
    \bar{\bar{I}}_6(m,|b|)=\rmint_0^\infty dx\,dz\,\frac{J_{0}(z |b|)}{z^2+m^2}\arctan\Big(\frac{x}{m}\Big)\textrm{arctanh}\Big(\frac{2x\,z}{m^2+x^2+z^2}\Big)\,,\\&
   I_{7}(m,|b|)= \rmint_{-\infty}^{\infty} d^3 z e^{-i \vec z\cdot \vec b}K_{0}\Big(|b||\sqrt{z_{(1)}^2+m^2}|\Big)\frac{1}{(\bar z^2+m^2)|\vec z|}\arctan\Big(\frac{|\vec z|}{2m}\Big)\,,\\&
   \bar{I}_{7}(m,|b|)=\rmint_{-\infty}^{\infty} d^3 z e^{-i \vec z\cdot \vec b}K_{0}\Big(|b||{z_{(1)}}|\Big)\frac{1}{(\vec{z}^2+m^2)|\vec z|}\arctan\Big(\frac{|\vec z|}{2m}\Big) \label{4.86}
    \end{split}
\end{align}
and $I_{4}(m,|b|)$ is defined in (\ref{4.37}). \textcolor{black}{As discussed previously, to the best of our knowledge these integrals do not possess any closed-form expression (except $I_5(m,|b|),\bar{I}_5(m,|b|), \bar{I}_6(m,|b|)$ and $\bar{\bar{I}}_6(m,|b|)$ as they can be done to some extent using method discussed around (\ref{4.25j}). They can be evaluated numerically and possess smooth behaviour w.r.t. $|b|\,.$ They also admit smooth massless limits as discussed earlier case by case basis.} \par
Note that the expression for the total impulse also receives a contribution from the pure Einstein-Hilbert part. However, that result is already well known in the literature \cite{Mogull:2020sak}. Hence, instead of reproducing it here, we focus only on the new contributions from the various vertices involving the scalar field up to 2PM order. We also show the corresponding results for the massless case. The fall-off with respect to $|b|$ is analogous to the gravitational case.

Before ending this section, we note that although we have computed the impulse up to 2PM, we have not yet discussed its connection with the scattering amplitude. \textcolor{black}{We elaborate on this connection in Appendix~(\ref{ch2:app:B}).}
%\vspace{0.5 cm}
%Another diagram contributing to $@\mathcal{O}(z)$ having the worldline action as,
%\begin{align}
 %   \begin{split}
 %       S_{pm}=\rmint_{k_1,k_2,\omega} e^{i(k_1+k_2)\cdot b}\hat\delta(k_1\cdot v+k_2\cdot v+\omega)\{(k_{1\rho}+k_{2\rho})v^{\mu}v^{\nu}+2\omega v^{(\mu}\delta_{\rho}^{\nu)}\}\tilde\varphi(k_1)\tilde h_{\mu\nu}(k_2)z^{\rho}(\omega)
 %   \end{split}
%\end{align}
%\end{itemize}
\section{Computation of waveform}\label{ch2:sec5}
In this section, we will compute the frequency domain gravitational waveform due to different field configurations. In the wave zone, $f_{\varphi,h}(k)$ is defined by the one point function as follows,
\begin{equation}\label{e:bar}\begin{split}
&  f_{h}(k):=\frac{1}{4\pi m_p}\epsilon^{\mu}\epsilon^{\nu}k^2\langle h_{\mu\nu}(k)\rangle\Big|_{k^2\rightarrow 0}\,, \\ &
        f_{\varphi}(k) := (k^2-m^2)\frac{\langle \varphi(k)\rangle}{m_p}\Big|_{k^2\rightarrow m^2}\\
        \end{split}
\end{equation}

where $k^\mu=\Omega\, n^\mu$ describes the on-shell momentum of the graviton/scalar. As, on-shell graviton is massless then, $n_\mu n^\mu=0$  and for the scalar $n_\mu n^\mu=\frac{m^2}{\Omega^2}$. $n^\mu$ can be parameterized as, 
\begin{align}
    \begin{split}
n^\mu\Big|_{g,h}=\Big(1,\sqrt{1-\frac{m_s^2}{\Omega^2}}\boldsymbol{\hat{x}}\Big)
    \end{split}
\end{align}
where,
\begin{align}
    \begin{split}
\boldsymbol{\hat{x}}^{\mu}=e_{1}^{\mu}\cos\theta+\sin(\theta)(e_{2}^{\mu}\cos\phi+e_3^{\mu}\sin\phi),\,e_{i}^{\mu}=(0,\hat\xi_{i})\,.
    \end{split}
\end{align}
with $\hat \xi\in (\hat i,\hat j,\hat k)$. In the parametrization mentioned above, the $n$ can be written as follows,
\begin{align}
    \begin{split}
        n=(1,\cos(\theta),\sin(\theta)\cos(\phi),\sin(\theta)\sin(\phi))\,.
    \end{split}
\end{align}
First, we will compute the corrections to the waveform from the scalar degrees of freedom. For the sake of simplicity, we first take the scalar field to be massless, and in the next section, we will discuss the massive integrals and how the integrals get complicated in the massive case. It has been shown in \cite{Mogull:2020sak} that in a two-body scattering problem, the scattering amplitude with on-shell external graviton (or scalar) is related to the WQFT correlator. In general, the scattering amplitude has the following schematic form.
\begin{center}
\includegraphics[width=0.46\linewidth]{s_matrix.png}
\end{center}
\begin{align}
    \mathcal{M}[g,\varphi]\equiv \text{S-matrix topology shown above.}
\end{align}
Using the result established in \cite{Mogull:2020sak}, one can show that the connection between the S-matrix element $\langle\phi_1,\phi_2\Big |\,\mathcal{S}\,\Big |\phi_1,\phi_2,\chi\rangle$ with the WQFT correlator,
\begin{align}
    \begin{split}
     (k^2-m_{\chi}^2)   \langle \chi(k)\rangle \propto \rmint_{k_1,k_2} d\mu_{1,2}(k_i,k)\lim_{\hbar\rightarrow 0}\mathcal{M}[g,\varphi](p_i,p_i',k),\,\chi\in (h_{\mu\nu},\varphi).\label{5.5m}
    \end{split}
\end{align}
where, in \eqref{5.5m} $k_i=p_i-p_i'$ and $d\mu_{1,2}$ is the measure of integration depends on the interaction.

\subsection{Scalar waveform}
We now initiate the computation of the scalar waveform, which we evaluate up to 2PM order. Next, we list all the diagrams contributing at this order and evaluate the corresponding expressions.

\subsubsection{Scalar waveform at 1PM}
In this subsection, we list the diagrams that contribute to the 1PM scalar waveform.

\textbullet $\,\,$ We start by considering the self-interaction vertices and their contribution to the waveform. First, we consider the $\frac{\lambda_3}{\textcolor{black}{3!}} m_p\,\varphi^3$ vertex. Its contribution to the one-point function has the following form and appears at 1PM order \footnote{\textcolor{black}{Again note that the combinatorial factor associated with this diagram is 3!. We have multiplied the result by this factor. For the subsequent diagrams, we will likewise include the appropriate combinatorial factors from the outset.}}. 
\begin{align}
    \begin{split} \label{wave0}
        k^2\Big\langle\varphi(k)\Big\rangle\equiv \thesisinlinefeynman[\chthreeinlinefeynwidth]{rad_4.png}\hspace{1.2 cm}=\lambda_3\Big(\frac{m_1s_1m_2s_2}{4m_p}\Big)\rmint \frac{d\mu_{1,2}(k)}{k_1^2k_2^2}\,.
    \end{split}
\end{align}
where the integral measure is given by,
\begin{align}
    \begin{split}
        \rmint d\mu_{1,2}(k)=\rmint_{k_1,k_2}e^{ik_1\cdot b_1}e^{ik_2\cdot b_2}\hat\delta(k_1\cdot v_1)\hat\delta(k_2\cdot v_2)\hat\delta^{(4)}(k_1+k_2-k)\,.
    \end{split}
\end{align}
Now, our focus is to compute the time domain waveform, which is the Fourier transform of the frequency domain waveform. Hence, we have to do a Fourier transformation of (\ref{wave0}). This gives the following, 
\begin{align}
    \begin{split}
        f_{\varphi}(x)&= \lambda_3\Big(\frac{m_1s_1m_2s_2}{4m_p^2}\Big)\rmint_{\Omega}e^{-ik\cdot x}\rmint_{k_i}e^{ik_1\cdot b_1}e^{ik_2\cdot b_2}\frac{\hat\delta(k_1\cdot v_1)\hat\delta(k_2\cdot v_2)}{k_1^2k_2^2}\hat\delta^{(4)}(k_1+k_2-k)\,,\\ &
     =   \lambda_3\Big(\frac{m_1s_1m_2s_2}{4m_p^2}\Big) \rmint_{\Omega}e^{-ik\cdot (x-b_1)}\underbrace{\rmint_{k_2}e^{-ik_2\cdot b}\frac{\hat\delta(k\cdot v_1-k_2\cdot v_1)\hat\delta(k_2\cdot v_2)}{k_2^2(k-k_2)^2}}_{J_{(0)}(k)}\,.
     \label{5.47 m}
    \end{split}
\end{align}
The integral in \eqref{5.47 m} can be evaluated using the results derived in \eqref{A.2} and \eqref{A.8}.  


\begin{equation}\label{wave1}\begin{split}
  f_{\varphi}(x)&= \lambda_3\Big(\frac{m_1s_1m_2s_2}{4m_p^2}\Big) \rmint_{\Omega}e^{-ik\cdot (x-b_1)}\,J_{(0)}(k),\\ &
        =-\frac{\lambda_3}{(2\pi)^3}\Big(\frac{m_1s_1m_2s_2}{4m_p^2}\Big) \frac{|b|}{4\gamma}\rmint_{0}^{1}dy\,\frac{1}{\Bar\Delta(y)}\sqrt{\frac{l^2}{|b|^2\bar\Delta^2}+1},\,l:=n\cdot(x-b_1+y\,b)\\
        \end{split}
\end{equation}
\vspace{0.5 cm}

where, $\Bar\Delta$ is defined in \eqref{A.4} and \eqref{ch2:A.6}. We have restored the factors of $\pi$ that come from the integration measures and the delta functions. So at 1PM order, (\ref{wave1}) gives the entire contribution to the scalar waveform. Next, we will extend our study to a 2PM order.


\subsubsection{Scalar waveform at 2PM}
In this subsection, we intend to compute the 2PM scalar waveform coming from the purely scalar sector and the scalar-graviton interaction sector.

\textbullet $\,\,$ \textit{We start with the self-interacting vertex, namely,  $\frac{\lambda_4}{\textcolor{black}{4!}}\varphi^4$. We will show that it doesn't contribute to the waveform.} To show that, we first write down the one-point function for this case. 
\begin{align}
    \begin{split}\label{5.11 mn}
        k^2\Big\langle\varphi(k)\Big\rangle\equiv    \thesisinlinefeynman[\chthreeinlinefeynwidth]{rad_5.png}\hspace{1.4 cm}=\rmint \frac{d\mu_{1,2}(k)}{k_1^2 k_2^2 k_3^2}\,.
    \end{split}
\end{align}
where , in this case, the integral measure has the following form
\begin{align}
    \begin{split}
       \rmint d\mu_{1,2}(k)=\rmint_{k_1,k_2,k_3}e^{i k_1\cdot b_1}e^{ik_2\cdot b_2}e^{i k_3 \cdot b_2}\hat\delta(k_1\cdot v_1)\hat\delta(k_2\cdot v_2)\hat\delta(k_3\cdot v_2)\hat\delta^{(4)}(k_1+k_2+k_3-k)\,.
    \end{split}
\end{align}
Hence, the corresponding contribution to the waveform has the following form,
\begin{align}
    \begin{split}
  \label{4PHI}      f_{\varphi}(x)&=\rmint_{\Omega}e^{-ik\cdot x}\rmint_{k_i}e^{ik_1\cdot b_1}e^{i k_2\cdot b_2}e^{i k_3\cdot b_2}\frac{\hat\delta(k_1\cdot v_1)\hat\delta(k_2\cdot v_2)\hat\delta(k_3\cdot v_2)}{k_1^2 k_2^2 k_3^2}\hat\delta^{(4)}(k_1+k_2+k_3-k)\,,\\ &
=\rmint_{\Omega}e^{-ik\cdot x}\rmint_{k_2,k_3}\frac{\hat\delta((k-k_2-k_3)\cdot v_1)\hat\delta(k_2\cdot v_2)\hat\delta{(k_3\cdot v_2)}}{k_2^2 k_3^2 (k-k_2-k_3)^2}\,e^{i (k-k_2-k_3)\cdot b_1}e^{i(k_2+k_3)\cdot b_2} \,,\\ &
\xrightarrow[]{q=k_2+k_3}\rmint_{\Omega}e^{-i k\cdot x}\rmint_{k_3,q} \frac{\hat\delta(k\cdot v_1-q\cdot v_1)\hat\delta(k_3\cdot v_2)\hat\delta(k_3\cdot v_2-q\cdot v_2)}{k_3^2 (q-k)^2(q-k_3)^2}e^{i (k-q)\cdot b_1}e^{i(q-k_3)\cdot b_2}e^{ik_3\cdot b_2}\,.
    \end{split}
\end{align}
We have the following integral to be solved which can be systematically done by introducing $\alpha$ parametrization,
\begin{align}
    \begin{split}
        f_{\varphi}(x)&=\rmint_{\Omega}e^{-i k\cdot x+i k\cdot b_{1}}\rmint_{k_3,q} \frac{\hat\delta(k\cdot v_1-q\cdot v_1)\hat\delta(k_3\cdot v_2)\hat\delta(q\cdot v_2)}{k_3^2 \, (q-k_3)^2(q-k)^2}e^{-i q\cdot b}\,,\\ &
        =\int_{\Omega}e^{-i k\cdot(x-b_1)}\int_{0}^\infty d\alpha\, d\beta\, d\gamma\, \exp\Big[-i\alpha (q-k)^2-i\beta(q-k_3)^2-i\gamma k_3^2 \Big]\\ &
        \hspace{0.8cm}\times  \hat\delta(k\cdot v_1-q\cdot v_1)\hat\delta(k_3\cdot v_2)\hat\delta(q\cdot v_2) e^{-iq\cdot b}\,,\\ &
=\int_{0}^{\infty}d\hat\alpha\,d\hat\beta\,d\hat\gamma \int_{\Omega}e^{-i\Omega \,n\cdot (x-b_1)}\int_{\vec k_3,\vec q}\hat\delta(\Omega\,n\cdot v_1+\vec q\cdot \vec v_1)\,\exp\Big[i\hat\alpha(\vec q^2-2\Omega \vec q\cdot \vec n)\\ &
\hspace{0.8cm}+i\hat\beta (\vec q^2-2\vec q\cdot \vec k_3+\vec k_3^2)+{i\hat\gamma \vec k_3^2}\Big]e^{i\vec q\cdot \vec b}
    \end{split}
\end{align}
Now doing the $\Omega $ integral we have,
\begin{align}
    \begin{split}
        f_{\varphi}(x)&=\frac{1}{n\cdot v_1}\int_{0}^{\infty}d\hat\alpha\,d\hat\beta\,d\hat\gamma \int_{\vec k_3,\vec q} \exp\Big[i\frac{\vec q\cdot \vec v_1}{n\cdot v_1}n\cdot \tilde x\Big]\exp\Big[i(\hat\alpha+\hat\beta)\,\vec q^2+2i\hat\alpha \Big(\frac{\vec q\cdot \vec v_1}{n\cdot v_1}\Big)(\vec q\cdot \vec n)+i\vec q\cdot \vec b\Big]\\ &
        \hspace{0.8cm}\times \exp\Big[i(\hat\beta+\hat\gamma)\vec k_3^2-2i\hat\beta \vec k_3\cdot \vec q\Big]\,,\Tilde{x}\equiv x-b_1
    \end{split}
\end{align}
Now we do the $\vec k_3$ integral by dimensional regularisation,
\begin{align}
    \begin{split}
        f_{\varphi}(x)&=\frac{1}{n\cdot v_1}\int_{0}^{\infty}d\hat\alpha\,d\hat\beta\,d\hat\gamma \int_{\vec q}\exp\Big[i\frac{\vec q\cdot \vec v_1}{n\cdot v_1}n\cdot \tilde x\Big]\exp\Big[i(\hat\alpha+\hat\beta)\,\vec q^2+2i\hat\alpha \Big(\frac{\vec q\cdot \vec v_1}{n\cdot v_1}\Big)(\vec q\cdot \vec n)+i\vec q\cdot \vec b\Big]\\ &
    \hspace{0.8cm} \times e^{-i\frac{\pi}{4}+i\frac{\epsilon}{2}}\pi^{\frac{3}{2}-\epsilon}(\hat\beta+\hat\gamma)^{-\frac{3}{2}+\epsilon}\exp\Big(-i \frac{\hat\beta^2 \vec q^2}{\hat\beta +\hat\gamma}\Big)\\ &
    =\frac{1}{n\cdot v_1}e^{-i\frac{\pi}{4}+i\frac{\epsilon}{2}}\pi^{\frac{3}{2}-\epsilon}\int_{0}^{\infty}d\hat\alpha\,d\hat\beta\,d\hat\gamma \,(\hat\beta+\hat\gamma)^{-\frac{3}{2}+\epsilon}\int d^3q \exp\Big(i\lambda_1 \vec q^2+2i \lambda_2 (\vec q\cdot \vec v_1)(\vec q\cdot \vec n)+i\vec q\cdot \vec \lambda_3\Big)\label{5.16k}
    \end{split}
\end{align}
where,
\begin{align}
    \begin{split}
&\lambda_1(\hat\alpha,\hat\beta,\hat\gamma)=\hat\alpha+\hat\beta -\frac{\hat\beta^2}{\hat\beta +\hat\gamma}\\ &
\lambda_2(\hat\alpha,\hat\beta,\hat\gamma)=\frac{\hat\alpha }{n\cdot v_1}\\ &
\vec \lambda_3=\frac{n\cdot \tilde x}{n\cdot v_1}\vec v_1+\vec b\equiv \Upsilon(x)\vec v_1+\vec b
    \end{split}
\end{align}
The non-triviality comes in the $q$ integral due to the presence of $\vec q\cdot \vec n\,\vec q\cdot \vec v_1$ term. For the sake of simplicity, we choose $ n=(1,1,0,0)$ and the $\vec q$ integral can be done component wise,
\begin{align}
    \begin{split}
        f_{\varphi}(x)&=\frac{1}{n\cdot v_1}e^{-i\frac{\pi}{4}+i\frac{\epsilon}{2}}\pi^{\frac{3}{2}-\epsilon}\int_{0}^{\infty}d\hat\alpha\,d\hat\beta\,d\hat\gamma \,(\hat\beta+\hat\gamma)^{-\frac{3}{2}+\epsilon}\int dq_{(1)}\exp\Big(i(\lambda_1+2\lambda_2\gamma\beta\cos\epsilon)q_{(1)}^2+i q_{(1)}\Upsilon(x)\,\gamma\beta\Big)\\ &
        \hspace{0.8cm}\times \int dq_{(2)}\exp\Big(i\lambda_1 q_{(2)}^2+iq_{(2)}b\Big)\int dq_{(3)} \exp(i\lambda_1 q_{(3)}^2)\\ & =  \frac{1}{n\cdot v_1}e^{-i\frac{\pi}{4}+i\frac{\epsilon}{2}}\pi^{\frac{3}{2}-\epsilon}\int_{0}^{\infty}d\hat\alpha\,d\hat\beta\,d\hat\gamma \,(\hat\beta+\hat\gamma)^{-\frac{3}{2}+\epsilon}\Bigg(\sqrt{\frac{\pi}{\lambda_1 +2\lambda_2\gamma\beta}}\\ &
        \hspace{0.8cm}\times\exp\Big(i\frac{\pi}{4}-i\frac{\Upsilon(x)^2\gamma^2\beta^2}{4\lambda_1 +8\lambda_2\gamma\beta}\Big)\Bigg)\times \Bigg(\sqrt{\frac{\pi}{\lambda_1}}\exp\Big(i\frac{\pi}{4}-i\frac{b^2}{4\lambda_1}\Big)\Bigg)\times \sqrt{\frac{\pi}{\lambda_1}}\exp\Big(i\frac{\pi}{4}\Big)\\ &
        =\frac{\pi^3}{n\cdot v_1}\int_0^\infty d\hat\alpha \, d\hat\beta\,d\hat\gamma\,(\hat\beta+\hat\gamma)^{-3/2}\frac{1}{\lambda_1\sqrt{\lambda_1+2\lambda_2\gamma\beta}}\exp\Big(-i\frac{b^2}{4\lambda_1}-i\frac{\Upsilon(x)^2\gamma^2\beta^2}{4\lambda_1 +8\lambda_2\gamma\beta}\Big)\label{5.18k}
    \end{split}
\end{align}
The integral in \eqref{5.18k} does not admit any closed form. One has to perform the remaining integrals numerically and see whether there are any finite contribution.\par  However, we can extract some information by doing an asymptotic analysis,
\begin{align}
    \begin{split}
        \textrm{Reg.}|f_{\varphi}(x)|\le \frac{\pi^3}{n\cdot v_1}\textrm{Reg.}\int_0^\infty d\hat\alpha \, d\hat\beta\,d\hat\gamma\,(\hat\beta+\hat\gamma)^{-3/2}\frac{1}{\lambda_1\sqrt{\lambda_1+2\lambda_2\gamma\beta}}\label{5.19k}
    \end{split}
\end{align}
Now, expanding the integrand \eqref{5.19k} around two limit ($0,\infty$) we get,
\begin{align}
     \textrm{Reg.}|f_{\varphi}(x)|\le 0\implies  \textrm{Reg.}|f_{\varphi}(x)|\to 0
\end{align}
To make this statement more concrete, we make an analysis by using the method of region in Appendix~(\ref{ch2:app:D}). But a more precise numerical analysis is required to see whether there is any finite contribution from this integral, which we leave for future investigations.  In a similar fashion, one can, in principle, compute the waveform corresponding to $\lambda_3h\varphi^3$ vertex. 


\textbullet $\,\,$ The simplest 2PM contribution comes from quadratic scalar field coupling in the worldline.
\begin{align}
    \begin{split}
        k^2\Big\langle\varphi(k)\Big\rangle\equiv \thesisinlinefeynman[\chthreeinlinefeynwidth]{rad_2.png}\hspace{0.35cm}=\Big(\frac{m_1g_1}{2m_p^2}\Big)\times \Big(\frac{m_2s_2}{2m_p}\Big)\rmint d\mu_{1,2}(k)\frac{1}{k_1^2}\label{5.7mm}
    \end{split}
\end{align}
where the integral measure has the following form,
\begin{align}
    \begin{split}
        d\mu_{1,2}(k)=\rmint_{k_1}e^{-ik_1\cdot b}e^{ik\cdot b_1}\hat\delta(k_1\cdot v_2)\hat\delta(k\cdot v_1-k_1\cdot v_1)\,.
    \end{split}
\end{align}
Therefore, the  time domain waveform has the following form,
\begin{align}
    \begin{split}
        f_{\varphi}(x)&=\Big(\frac{m_1g_1}{2m_p^2}\Big)\times \Big(\frac{m_2s_2}{2m_p}\Big)\frac{1}{n\cdot v_1}\rmint_{\Omega} e^{-ik\cdot (x-b_1)}\rmint_{k_1} e^{-i k_1\cdot b}\frac{\hat\delta(k_1\cdot v_2)\hat\delta\Big(\Omega-\frac{k_1\cdot v_1}{n\cdot v_1}\Big)}{k_1^2}\,,\\ &
         =  \Big(\frac{m_1g_1}{2m_p^2}\Big)\times \Big(\frac{m_2s_2}{2m_p}\Big)\frac{1}{n\cdot v_1}\rmint_{k_1}e^{-ik_1\cdot w_1}\frac{\hat\delta(k_1\cdot v_2)}{k_1^2},\,\textrm{with}, \, w_1\equiv \frac{n\cdot (x-b_1)}{n\cdot v_1}v_1+b,\nonumber
         \end{split}
\end{align}
\begin{align}
\begin{split}
  \hspace{0cm}  & 
      = -\Big(\frac{m_1g_1}{2m_p^2}\Big)\times \Big(\frac{m_2s_2}{2m_p}\Big)\frac{1}{n\cdot v_1}\frac{1}{4\pi|\vec w_1|}\,.
    \end{split}
\end{align}
Therefore, the final result of the integrals looks,

\begin{equation}\label{e:ba}\begin{split}
 f_{\varphi}(x)&=-\Big(\frac{m_1g_1}{2m_p^2}\Big)\times \Big(\frac{m_2s_2}{2m_p}\Big)\frac{1}{n\cdot v_1}\frac{1}{4\pi|\vec w_1|}\,.
\\
        \end{split}
\end{equation}\\
\textbullet $\,\,$ Another 2PM contribution comes from the following:
\begin{align}
    \begin{split}
        k^2\Big\langle\varphi(k)\Big\rangle\Big|_{k^2\rightarrow 0}\equiv \thesisinlinefeynman[\chthreeinlinefeynwidth]{rad_1.png}\hspace{0.35cm}=m_1\Big(\frac{s_1}{2m_p}\Big)^2 \Big(\frac{m_2s_2}{2m_p}\Big)\rmint d\mu_{1,2}(k)\frac{\{2\omega (v_1)_{\rho}-(k_1)_{\rho}\}\{-2\omega (v_1)_{\rho}+ k_{\rho}\}}{\omega^2 k_1^2}\label{5.10 m}
    \end{split}
\end{align}
where, the measure $d\mu_{1,2}(k)$ has the following form,
\begin{align}
    \begin{split}
        \rmint d\mu_{1,2}(k)=\rmint_{k_1,\omega}e^{i(k-k_1)\cdot b_1}e^{i k_1\cdot b_2}\hat\delta(k_1\cdot v_1-\omega)\hat\delta(k \cdot v_1-\omega)\hat\delta(k_1\cdot v_2)\,.
    \end{split}
\end{align}
Now its contribution to the time-domain waveform is,
\begin{align}
    \begin{split}
        f_{\varphi}(x)& = -m_1\Big(\frac{s_1}{2m_p}\Big)^2 \Big(\frac{m_2s_2}{2m_p}\Big) \rmint_{\Omega} e^{-ik\cdot x} \rmint d\mu_{(a)}(k) \frac{k\cdot k_1}{{\omega}^2\,k_1^2}\,,\\ &
        =-m_1\Big(\frac{s_1}{2m_p}\Big)^2 \Big(\frac{m_2s_2}{2m_p}\Big)\rmint_{\Omega} e^{-ik\cdot x}\rmint_{k_1} e^{-ik_1\cdot b}e^{i k\cdot b_1}\frac{k\cdot k_1\hat\delta(k_1\cdot v_2)}{k_1^2\,(k_1\cdot v_1+i\epsilon)^2}\hat\delta[(k-k_1)\cdot v_1]\,.
    \end{split}
\end{align}
%The numerator has the form: $\mathcal{N}=-k\cdot k_1$.
Therefore, the whole integral can be written as,
\begin{align}
\begin{split}
      f_{\varphi}(x)&= -m_1\Big(\frac{s_1}{2m_p}\Big)^2 \Big(\frac{m_2s_2}{2m_p}\Big)\rmint_{\Omega}e^{-ik\cdot (x-b_1)}\rmint_{\omega,k_1}e^{-i k_1\cdot b}\hat{\delta}(k_1\cdot v_2)\hat{\delta}(\Omega\, n\cdot v_1-\omega)\hat\delta(k_1\cdot v_1-\omega)\frac{\Omega(n\cdot k_1)}{\omega^2 k_1^2}\,,\\ &
       = -m_1\Big(\frac{s_1}{2m_p}\Big)^2 \Big(\frac{m_2s_2}{2m_p}\Big)\frac{1}{(n\cdot v_1)^2}\rmint_{\Omega}e^{-ik\cdot (x-b_1)}\rmint_{k_1}e^{-ik_1\cdot b}\hat\delta(k_1\cdot v_2)\hat{\delta}(\Omega-\frac{k_1\cdot v_1}{n\cdot v_1})\frac{n\cdot k_1}{(k_1\cdot v_1+i\epsilon)k_1^2}\,,\\ &
     = -m_1\Big(\frac{s_1}{2m_p}\Big)^2 \Big(\frac{m_2s_2}{2m_p}\Big)\frac{n^{\mu}}{(n\cdot v_1)^2}\rmint_{k_1}e^{-i k_1\cdot w_1}\hat{\delta}(k_1\cdot v_2)\frac{k_{1\mu}}{(k_1\cdot v_1+i\epsilon)k_1^2},\,\, w_1\equiv\frac{n\cdot (x-b_1)}{n\cdot v_1}v_1+b \,.\label{7.6}
      \end{split}
\end{align}
The integral in (\ref{7.6}) can be easily done from the frame of the second particle as follows,
\begin{align}
    \begin{split}
        \mathcal{J}^{\mu}:=\rmint_{k_1}e^{-i k_1\cdot w_1}\frac{k_1^{\mu}\,\hat\delta(k_1\cdot v_2)}{(k_1\cdot v_1+i\epsilon)k_1^2}&=-i\rmint d\tau \,\theta(\tau)\,\rmint_{k_1}\hat{\delta}(k_1\cdot v_2) e^{-ik_1\cdot (w_1-\tau\, v_1)}\frac{k_1^{\mu}}{k_1^2}\,,\\ &
        =-i\rmint d\tau\,\theta(\tau)\rmint_{k_1}\hat\delta{(k_1^{0})} \exp[{-ik_1\cdot \underbrace{(w_1-\tau\, v_1)}_{\tilde{w}_1}}]\frac{k_1^{\mu}}{k_1^2}\,,\\ &
        \xrightarrow[]{k_1^0\rightarrow 0}-\rmint d\tau\, \theta(\tau)\rmint_{\vec k_1}\frac{k_1^i}{-\vec k_1^2}\,e^{i \vec{k}_1\cdot \vec {\tilde w}_1}\,,\\ &
        =-\rmint_{-\infty}^{\infty}d\tau \, \theta(\tau)\frac{(\vec w_1-\tau\,\vec v_1)^{i}}{|\vec w_1-\tau\,\vec v_1|^{3}}\,.
    \end{split}
\end{align}
Now from (\ref{7.6}), it is clear that we need to compute the following quantity,
\begin{align}
    \begin{split}
        n_{\mu}\mathcal{J}^{\mu}\rightarrow n_{i}\mathcal{J}^{i}=\rmint_{-\infty}^{\infty}d\tau\, \theta(\tau)\frac{\vec n\cdot (\vec w_1-\tau\, \vec v_1)}{|\vec w_1-\tau \vec v_1|^3}\,.
    \end{split}
\end{align}
To proceed further, we have to properly parameterize the impact parameter $b_1,\,b_2$. As we are in the frame of the second particle, it is convenient to choose $b_1=(0,0,b,0)$ and $b_2=0$.
\begin{align}
    \begin{split}
       \vec n\cdot (\vec w_1-\tau \vec v_1)&\rightarrow \Big[\frac{n\cdot (x-b_1)}{n\cdot v_1}-\tau\Big]\vec n\cdot \vec v_1+\vec n\cdot \vec b\\ &
       =(u_1-\tau)\gamma\tilde\beta+\chi,\,\,\text{with},\chi:=b\sin\theta\cos\varphi,\,\tilde\beta:=\beta \cos\theta
    \end{split}
\end{align}
and,
\begin{align}
    \begin{split}
        |\vec w_1-\tau \vec v_1|^3\rightarrow (\gamma^2-1)^{3/2}\Big[\tau^2-u_1^2-2\tau u_1+\frac{|b|^2}{\gamma^2-1}\Big]^{3/2}.
    \end{split}
\end{align}
Hence, 
\begin{align}
    \begin{split}
    n_{\mu}\mathcal{J}^{\mu}=\frac{\gamma  \tilde{\beta } \left(2 u_1^2-\frac{|b|^2}{\left(\gamma ^2-1\right)^2}\right)+\chi  \left(\sqrt{\frac{|b|^2}{\left(\gamma ^2-1\right)^2}-u_1^2}+u_1\right)}{(\gamma^2-1)^{3/2}\left(\frac{|b|^2}{\left(\gamma ^2-1\right)^2}-2 u_1^2\right) \sqrt{\frac{|b|^2}{\left(\gamma ^2-1\right)^2}-u_1^2}}
    \end{split}
\end{align}
where, we define, $u_i=\frac{n\cdot (x-b_1)}{n\cdot v_i}$.Therefore, restoring the factors of $\pi$ from delta functions and integration measures, the contribution to the waveform from the particular diagram in  (\ref{5.10 m}) has the following,
\begin{align}\begin{split}\label{5.24 mmm}
 f_{\varphi}(x)= -\frac{m_1}{(2\pi)^2}\Big(\frac{s_1}{2m_p}\Big)^2 \Big(\frac{m_2s_2}{2m_p}\Big)\frac{1}{\gamma^2(1-\tilde\beta)^2(\gamma^2-1)^{3/2}}\frac{\gamma  \tilde{\beta } \left(2 u_1^2-\frac{|b|^2}{\left(\gamma ^2-1\right)^2}\right)+\chi  \left(\sqrt{\frac{|b|^2}{\left(\gamma ^2-1\right)^2}-u_1^2}+u_1\right)}{\left(\frac{|b|^2}{\left(\gamma ^2-1\right)^2}-2 u_1^2\right) \sqrt{\frac{|b|^2}{\left(\gamma ^2-1\right)^2}-u_1^2}}\,.
\\
 \end{split}
\end{align}
%\begin{align}
    %\begin{split}
      %  (k^2-%m^2)\langle\tilde\varphi (k)\rangle\Big|_{k^2=m^2}=\rmint d\mu_{(a)}(k)\frac{\{2\omega (v_1)_{\rho}-(k_1)_{\rho}\}\{-2\omega (v_1)_{\rho}+ k_{\rho}\}}{(k_1^2-m^2)\omega^2}
  %  \end{split}
%\end{align}
%\textcolor{black}{AS FAR AS I UNDERSTAND, WE DONT NEED TO COMPUTE THE INTEGRALS IN RADIATION NOW. WHILE COMPUTING WAVEFORM IN TIME DOMAIN THERE IS ANOTHER INTEGRAL OVER `$\Omega$' AND THAT SIMPLIFIES THE INTEGRAL. IN \textbf{2101.12688} THE INTEGRALS ARE COMPUTED SYSTEMATICALLY. THE PRIMARY PAPER ON WQFT THEY DONT COMPUTE THE INTEGRALS RATHER COMPUTE THE INTEGRAND ONLY.}
%The $\delta \tilde\varphi^4$ term contributes to the $\tilde \varphi$ radiation. The contribution to the one-point function has the following form:
%\begin{align}
    %\begin{split}
      %  (k^2-m^2)\langle\tilde\varphi(k)\rangle=\rmint_{k_1,k'} d\mu_{(b)}(k)\frac{1}{(k_1^2-m^2)[(k'-k_1)^2-m^2][(k+k')^2-m^2]}
   % \end{split}
%\end{align}
%where the integral measure %%%  \rmint d\mu_{(b)}(k)=\rmint_{k_1,k_2}e^{i(k+k')\cdot b_1}e^{-ik'\cdot b_2}\hat\delta(k_1\cdot v_2)\hat\delta\{(k'-k_1)\cdot v_2\}\hat\delta\{(k+k')\cdot v_1\}
   % \end{split}
%\end{align}
%$\section{}
%space{-1 cm}
%\hrule
%\vspace{2 cm}
%\textcolor{black}{Have to be written clearly (the derivative interaction part).}\\ 
\textbullet $\,\,$ Finally, another 3-point scalar-graviton interaction vertex involving a derivative contributes to 2PM radiation. It is of the following form: $h^{\mu\nu}\partial_{\mu}\varphi\partial_{\nu} \varphi$  As we are computing the one-point function of the scalar field, it would be useful to partially integrate over the interaction Lagrangian and separate one of the scalar fields as,
\begin{align}
    \begin{split}
        S_{\textrm{int.}}=\frac{1}{m_p}\rmint d^4x \,h^{\mu\nu} \partial_{\mu}\varphi\partial_{\nu}\varphi=-\frac{1}{m_p}\rmint d^4 x \Big[\underbrace{h^{\mu\nu}\partial_{\mu}\partial_{\nu}\varphi}_{\textrm{Term I}}+\underbrace{\partial_{\mu}h^{\mu\nu}\partial_{\nu}\varphi}_{\textrm{Terrm II}}\Big]\,\varphi.\label{5.22 m}
    \end{split}
\end{align}
Now, the contribution to the scalar one-point function from the Term I is given by,
\begin{align}
    \begin{split}
        k^2\Big\langle \varphi(k)\Big\rangle\equiv \thesisinlinefeynman[\chthreeinlinefeynwidth]{rad_3.png}\hspace{0.35cm}=\Big(\frac{m_1m_2s_2}{2m_p^3}\Big)\rmint d\mu_{1,2}(k)\frac{v_1^\alpha v_1^\beta\,P_{\mu\nu;\alpha\beta}k_2^\mu k_2^\nu}{k_1^2k_2^2}\label{ch2:5.23a}
    \end{split}
\end{align}
where the integral measure has the following form,
\begin{align}
    \begin{split}
     \rmint   d\mu_{1,2}(k)=\rmint_{k_1,k_2}e^{ik_1\cdot b_1}e^{ik_2 \cdot b_2}\hat\delta(k_1\cdot v_1)\hat\delta(k_2\cdot v_2)\hat\delta^{(4)}(k_1+k_2-k)\,.
    \end{split}
\end{align}
Hence, the time domain waveform has the following form,
\begin{align}
    \begin{split}
        f_{\varphi}(x)\Big|_1=\Big(\frac{m_1m_2s_2}{2m_p^3}\Big)\rmint_{\Omega}e^{-ik\cdot x}\rmint_{k_1,k_2}e^{ik_1\cdot b_1}e^{ik_2\cdot b_2}\hat\delta(k_1\cdot v_1)v_1^{\alpha}v_1^{\beta}\hat\delta(k_2\cdot v_2)\frac{k_2^{\mu}k_2^{\nu}P_{\mu\nu,\alpha\beta}}{k_1^2k_2^2}\delta^{(4)}(k_1+k_2-k)\,.\label{5.25mmm}
    \end{split}
\end{align}
The integral in \eqref{5.25mmm} can be done using the partial fraction approach, where we can separate the denominator with undefined momentum as,
\begin{align}
    \begin{split}
        \frac{1}{k_1^2 k_2^2}=-\frac{1}{2}\frac{1}{k_1^2(k_1\cdot k)}-\frac{1}{2}\frac{1}{k_2^2(k_2\cdot k)}\,.
    \end{split}
\end{align}
Therefore, the waveform can be reduced into two independent momentum integrals,
\begin{align}
    \begin{split} \label{e:::}
        f_{\varphi}(x)\Big|_{1}=\Big(\frac{m_1m_2s_2}{4 m_p^3}\Big)\Big(f_{\varphi}(x)\Big|_{1}^{(1)}+f_{\varphi}(x)\Big|_{1}^{(2)}\Big)\,.
    \end{split}
\end{align}
where,
\begin{align}
    \begin{split} \label{5.38nm}
        f_{\varphi}(x)\Big|_{1}^{(1)}&=v_1^{\alpha}v_1^{\beta}P_{\mu\nu;\alpha\beta}\rmint_{\Omega}e^{-ik\cdot x}\rmint_{k_1,k_2}e^{ik_1\cdot b_1}e^{ik_2\cdot b_2}\frac{k_2^\mu k_2^\nu\hat\delta(k_1\cdot v_1)\hat\delta(k_2\cdot v_2)}{k_1^2(k_1\cdot k)}\delta^{(4)}(k_1+k_2-k)\,,\\ &
=v_1^{\alpha}v_1^{\beta}P_{\mu\nu;\alpha\beta}\rmint_{\Omega}e^{-ik\cdot (x-b_2)}\rmint_{k_1}e^{ik_1\cdot b}\frac{\hat\delta(k_1\cdot v_1)\hat\delta(\Omega-k_1\cdot v_2)}{k_1^2(k_1\cdot n)\Omega}(k-k_1)^{\mu}(k-k_1)^{\nu}\,,\\ &
=v_1^{\alpha}v_1^{\beta}P_{\mu\nu;\alpha\beta}\rmint_{k_1}e^{-i k_1\cdot \tilde w_1}\frac{\hat\delta(k_1\cdot v_1)}{k_1^2 (k_1\cdot n)(k_1\cdot v_2)}\Big[n^\mu n^\nu (k_1\cdot v_2)^2-2 (k_1\cdot v_2)k_1^{(\mu}n^{\nu)}+k_1^{\mu}k_1^{\nu}\Big]\,.
    \end{split}
\end{align}
The integral of concern is the following,
\begin{align}
    \begin{split}
        I^{\mu\nu}&=\rmint_{k_1}\frac{e^{-ik_1\cdot \tilde w_1}\hat\delta(k_1\cdot v_1)}{k_1^2(k_1\cdot n)(k_1\cdot v_2)}k_1^{\mu}k_1^{\nu},\, \tilde w_1=\frac{n\cdot (x-b_2)}{n\cdot v_2}v_2-b\,.\label{5.29mmm}
   %  \rightarrow -\rmint_{\tau_{1,2}}\theta(\tau_1)\theta(\tau_2)\rmint_{\bar q}\frac{\bar q^{i}\bar q^{j}}{\bar q^2}\exp[i \bar q\cdot \bar \lambda_1],\, \bar q=\Big(\frac{k_1^x} {\gamma},k_1^y,k_1^z\Big)\, \&\, \bar \lambda_1=\Big(\lambda_1^x-\beta \lambda_1^0,\lambda_1^y,\lambda_1^z\Big)\\ &
    %    \propto -8\rmint_{\tau_{1,2}}\theta(\tau_1)\theta(\tau_2)\frac{1}{|\bar \lambda_1(\tau_1,\tau_2)|}[\frac{1}{2}\delta^{ij}-\frac{3}{2}\frac{\bar\lambda_1^i(\tau_1,\tau_2) \bar\lambda_1^j(\tau_1,\tau_2)}{\bar \lambda_1(\tau_1,\tau_2)^2}]\,.
    \end{split}
    \end{align}
{\color{black}One can see that because of the delta function constraint, one can replace  $\tilde w_1$ with $\hat w_2:=-b+u_2 (v_2-\gamma v_1)$. Also, $I^{\mu\nu}$ is orthogonal to $\hat w_2$ and $v_1$. Therefore, it can be written in terms of a basis plane orthogonal to  $\hat w_2$ and $v_1$.} $I^{\mu\nu}$ is orthogonal to $\hat w_2$ and $v_1$.  Therefore,
    \begin{align}
        \begin{split}
I^{\mu\nu}&=\Pi^{\mu}_{\alpha}\,\Pi^{\nu}_{\beta}\Big[c_{vv} v_2^{\alpha}v_2^{\beta}+c_{nn} n^{\alpha}n^{\beta}+c_{nv} v_2^{(\alpha}n^{\beta)}\Big]\,.
        \end{split}
    \end{align}
One can notice that $I^{\mu\nu}$ is symmetric under the exchange of $v_2 \leftrightarrow n$ implies $c_{nn}=c_{vv}$. Now, the integral can be done using the integral reduction technique by Passarino and Veltman, redefining the basis and writing the following ansatz, 
\begin{align}
        \begin{split}
\label{Ansatz2}I^{\mu\nu}&=c_{a}\Pi_{1}^{\mu\nu}+ 2c_{b}\Big(\Pi_{1}.v_{2}\Big)^{(\mu}\Big(\Pi_{1}.n\Big)^{\nu)}
\end{split}
    \end{align}
where $c_{a}$ and $c_{b}$ are two new constants and $\Pi_{1}^{\mu\nu}= |w_{1}|^{2}\,P_{1}^{\mu\nu}\,+ w_{1}^{\mu}\,w_{1}^{\nu}$ is the projection operator orthogonal to $w_{1}^{\mu}$ and $v_{1}^{\mu}$. One can evaluate the constants from the following two equations:
\begin{align}
\label{NV2}  &  n\cdot I\cdot v_2=\rmint_{k_1}e^{-ik_1\cdot \hat w_1}\frac{\hat\delta(k_1\cdot v_1)}{k_1^2}= c_{a}\Big(n\cdot \Pi\cdot v_{2}\Big)+ c_{b}\,\Big[\Big(n\cdot \Pi\cdot v_{2}\Big)\,\Big(v_{2}\cdot \Pi\cdot n\Big)+ \Big(v_{2}\cdot \Pi\cdot v_{2}\Big)\,\Big(n\cdot \Pi\cdot n\Big)\Big]\,,\\ &
\label{NIN}  n\cdot I \cdot n= n_\mu \rmint_{k_1}e^{-ik_1\cdot \hat w_2}\frac{k_1^\mu\hat\delta(k_1\cdot v_1)}{k_1^2 (k_1\cdot v_2)}= c_{a}\Big(n\cdot \Pi\cdot n\Big)+ 2c_{b}\Big(n\cdot \Pi\cdot v_{2}\Big)\,\Big(n\cdot \Pi\cdot n\Big)\,.
\end{align}
The first equality of (\ref{NIN}) gives,
\begin{align}
\begin{split}
   n_\mu \rmint_{k_1}e^{-ik_1\cdot \hat w_2}\frac{k_1^\mu\hat\delta(k_1\cdot v_1)}{k_1^2 (k_1\cdot v_2)}&\to -i n_\mu\, \rmint d\tau\,\theta(\tau)\rmint_{ k_1} e^{-i k_1\cdot (\hat{w}_2-\tau v_2)}\frac{{k}_1^\mu}{k_1^2}\hat\delta(k_1\cdot v_1)\,,\\ &
   =-i\, \grave{n}_i\rmint d\tau\,\theta(\tau)
   \rmint_{\bar k_1}e^{i \bar k_1\cdot \grave{\underbar{$w$}}_2}\frac{\bar k_1^i}{-\bar k_1^2}\,,\\ &
   =-\grave n_i\rmint d\tau \,\theta(\tau)\frac{\grave{\underbar{$w$}}(\tau)^i}{|\grave{\underbar{$w$}}(\tau)|^3},\,\grave{\underbar{$w$}}\equiv \hat{w}_2-\tau v_2.
   \end{split}
\end{align}
where,
\begin{align}
    \begin{split}
        \grave n_i\equiv \Big\{\gamma(n^{(1)}-\beta n^{(0)}),n^{(2)},n^{(3)}\Big\}\,.
    \end{split}
\end{align}
Comparing with the second equality of (\ref{NIN}), we get,

\begin{align}\begin{split}\label{m:}
  &c_a\to -\frac{-(n\cdot \Pi\cdot v_2) (n\cdot I\cdot n) (v_2\cdot \Pi\cdot n)+2(n\cdot \Pi\cdot v_2)(n\cdot \Pi\cdot n) (n\cdot I\cdot v_2)-(v_2\cdot \Pi\cdot v_2) (n\cdot \Pi\cdot n)(n\cdot I \cdot n)}{-2 (n\cdot \Pi\cdot v_2)^2+(n\cdot \Pi\cdot v_2) (v_2\cdot \Pi\cdot n)+(v_2\cdot \Pi\cdot v_2) (n\cdot \Pi\cdot n)},\\&
 c_b\to \frac{(n\cdot \Pi\cdot n) (n\cdot I\cdot v_2)-(n\cdot \Pi\cdot v_2) (n\cdot I\cdot n)}{(n \cdot \Pi\cdot n) \left(-2 (n\cdot \Pi\cdot v_2)^2+(n\cdot \Pi\cdot v_2) (v_2\cdot \Pi\cdot n)+(v_2\cdot \Pi\cdot v_2) (n\cdot \Pi\cdot n)\right)}.\\
\end{split}
\end{align}
Next, we compute the integral of the R.H.S of the first equality of (\ref{NV2}) to get,
\begin{eqnarray}
\label{CAINT} \rmint_{k_1}e^{-ik_1\cdot \hat w_1}\frac{\hat\delta(k_1\cdot v_1)}{k_1^2} =-\frac{1}{4\pi}\frac{1}{|\grave{w}_2|}
\end{eqnarray}
where $\grave{w}_2$ is a three vector defined as,
\begin{align}
    \grave{w}_2=\Big\{\gamma(\hat{w}_2^{(1)}-\beta \hat{w}_2^{(0)}),\hat{w}_2^{(2)},\hat{w}_2^{(3)}\Big\}\,.
\end{align}
%Using the results of \eqref{CAINT} and \eqref{CB}, we compute the constant $c_{a}$ from \eqref{NV2} as
%%c_{a}= -\frac{1}{4\pi} \frac{1}{(n \cdot \Pi_{1} \cdot v_{2})|w_{1}|} - \frac{1}{2\,(\gamma^{2}-1)|b|^{2}|w_{1}|}\,\Big(\frac{(v_{2} \cdot \Pi_{1} \cdot n)}{(n \cdot \Pi_{1} \cdot n)} + \frac{(v_{2} \cdot \Pi_{1} \cdot v_{2})}{(n \cdot \Pi_{1} \cdot v_{2})}
%\Big)
%\end{eqnarray}
%Where
%\begin{eqnarray}
  %   n \cdot \Pi_{1} \cdot v_{2}&=& |b|^{2}\,(n\cdot v_{2}- \gamma)+ \gamma^{2}\beta^{2}\,u_{2}\,(n\cdot b)\\
   %  n \cdot \Pi_{1} \cdot n &=& |b|^{2} + \gamma^{2}\beta^{2}\,u_{2}^{2} + \Big(u_{2}(n\cdot v_{2}- \gamma)-(n\cdot b)\Big)^{2}\\
 % v_{2} \cdot \Pi_{1} \cdot v_{2} & = & -|b|^{2}\gamma^{2}\beta^{2}.
%\end{eqnarray}
Thus, the integral $I^{\mu\nu}$ can be finally written in a concise form from the ansatz (\ref{Ansatz2}) as,
\begin{align}
\begin{split} 
\hspace{0cm}\label{IMUNU}I^{\mu\nu} =-\Big[\frac{-(n\cdot \Pi\cdot v_2) (n\cdot I\cdot n) (v_2\cdot \Pi\cdot n)+2(n\cdot \Pi\cdot v_2)(n\cdot \Pi\cdot n) (n\cdot I\cdot v_2)-(v_2\cdot \Pi\cdot v_2) (n\cdot \Pi\cdot n)(n\cdot I \cdot n)}{-2 (n\cdot \Pi\cdot v_2)^2+(n\cdot \Pi\cdot v_2) (v_2\cdot \Pi\cdot n)+(v_2\cdot \Pi\cdot v_2) (n\cdot \Pi\cdot n)}\Big]\,\Pi^{\mu\nu} \\
+ 2\,\Big[\frac{(n\cdot \Pi\cdot n) (n\cdot I\cdot v_2)-(n\cdot \Pi\cdot v_2) (n\cdot I\cdot n)}{(n \cdot \Pi\cdot n) \left(-2 (n\cdot \Pi\cdot v_2)^2+(n\cdot \Pi\cdot v_2) (v_2\cdot \Pi\cdot n)+(v_2\cdot \Pi\cdot v_2) (n\cdot \Pi\cdot n)\right)}\Big]\Big(\Pi \cdot v_{2}\Big)^{(\mu}\Big(\Pi\cdot n\Big)^{\nu)}\,.
\end{split}
\end{align}
Once we have the closed-form expression for $I^{\mu\nu}\,,$ one can write down the exact expression for the first part of the waveform after restoring the factor of $\pi$.
\begin{align}\begin{split}\label{e:}
f_{\varphi}(x)\Big|_{1}^{(1)}=-\,\frac{v_1^\alpha v_1^\beta}{4\pi^2}P_{\mu\nu;\alpha\beta}\Big[n^{\mu}n^\nu I^{\rho\sigma}v_{2\rho}v_{2\sigma}-2v_{2\rho}I^{\rho(\mu}n^{\nu)}+I^{\mu\nu}\Big].\\
\end{split}
\end{align}
%\vspace{0.5 cm}
Now we focus on the other part.
\begin{align}
    \begin{split}
        f_{\varphi}(x)\Big|_{1}^{(2)}&=-\,v_1^\alpha v_1^\beta P_{\mu\nu;\alpha\beta}\rmint_{\Omega}e^{-ik\cdot x}\rmint_{k_1,k_2}e^{ik_1\cdot b_1}e^{ik_2\cdot b_2}\frac{k_2^\mu k_2^\nu \hat\delta(k_1\cdot v_1)\hat\delta(k_2\cdot v_2)}{k_2^2 (k_2\cdot k)}\hat\delta^{(4)}(k_1+k_2-k)\,,\\ &
        =-\,v_1^\alpha v_1^\beta P_{\mu\nu;\alpha\beta}\rmint_{\Omega}e^{-i \Omega \,n \cdot x}\rmint_{k_2}e^{i (k-k_2)\cdot b_1} e^{ik_2\cdot b_2} \frac{k_2^\mu k_2^\nu \,\hat\delta(k_2\cdot v_2)\hat\delta(\Omega\,n\cdot v_1-k_2\cdot v_1)}{k_2^2 (k_2\cdot n)\Omega}\,,\\ &
        =-\,v_1^\alpha v_1^\beta P_{\mu\nu;\alpha\beta} \underbrace{\rmint_{k_2} e^{-i k_2 \cdot w_1}k_2^\mu k_2^\nu\, \frac{\hat\delta(k_2\cdot v_2)}{k_2^2(k_2\cdot n)(k_2\cdot v_1)}}_{\bar{I}^{\mu\nu}} \textrm{with},\,w_1\equiv \frac{n\cdot (x-b_1)}{n\cdot v_1}v_1+b\,.\label{5.43mm}
    \end{split}
\end{align}
The integral in \eqref{5.43mm} can be similar to \eqref{5.29mmm} and is given by,
\begin{align}
    \begin{split} \label{e::}
f_{\varphi}(x)\Big|_{1}^{(2)}=-\,v_1^{\alpha}v_1^\beta P_{\mu\nu;\alpha\beta} \bar I^{\mu\nu}.\\
    \end{split}
\end{align}
where $\bar I^{\mu\nu}$ is defined (\ref{5.43mm}) and can be calculated as \eqref{5.29mmm}. So adding (\ref{e:}) and  (\ref{e::}) we get the full expression for $f_{\phi}(x)\Big|_1$ mentioned in (\ref{e:::}).

Now, we deal with the Term II of the interaction Lagrangian of \eqref{5.22 m}. The corresponding contribution to the waveform is given by,
\begin{align}
    \begin{split} \label{ex6}
        f_{\varphi}(x)\Big|_{2}=\Big(\frac{m_1 m_2 s_2}{2m_p^3}\Big)v_1^\alpha v_1^\beta P_{\mu\nu;\alpha\beta}\rmint_{\Omega} e^{-ik\cdot x}\rmint_{k_1,k_2}e^{ik_1\cdot b_1}e^{ik_2\cdot b_2}\frac{k_1^\mu k_2^\nu\hat\delta(k_1\cdot v_1)\hat\delta(k_2\cdot v_2)}{k_1^2 k_2^2}\hat\delta^{(4)}(k_1+k_2-k)\,.
    \end{split}
\end{align}
Now, doing the integration over $k_1$ using the delta function, we will get,
\begin{align}
    \begin{split}
         f_{\varphi}(x)\Big|_{2}=\Big(\frac{m_1 m_2 s_2}{2m_p^3}\Big)v_1^\alpha v_1^\beta P_{\mu\nu;\alpha\beta}\rmint_{\Omega}e^{-ik\cdot (x-b_1)}\rmint_{k_2}e^{-i k_2 \cdot b}\frac{(k-k_2)^\mu \,k_2^\nu}{k_2^2\,(k-k_2)^2}\hat\delta(k\cdot v_1-k_2\cdot v_1)\hat\delta(k_2\cdot v_2)\,.\label{5.46 m}
    \end{split}
\end{align}
Again, this can be split into two parts and can be dealt with separately. 
\begin{align}
    \begin{split} \label{ex3}
    f_{\varphi}(x)\Big|_{2}=\Big(\frac{m_1 m_2 s_2}{2m_p^3}\Big)\Big(f_{\varphi}(x)\Big|^{(1)}_{2}+f_{\varphi}(x)\Big|^{(2)}_{2}\Big)\,.
     \end{split}
\end{align}
The first part gives,
\begin{align}
    \begin{split}
        f_{\varphi}(x)\Big|_{2}^{(1)}=\underbrace{(v_1\cdot P\cdot v_1)_{\mu\nu}\,n^{\mu}}_{T_{\nu}}\rmint_{\Omega}e^{-ik\cdot (x-b_1)}\,\Omega\underbrace{\rmint_{k_2}\frac{k_2^\nu}{k_2^2(k-k_2)^2}\hat\delta(k\cdot v_1-k_2\cdot v_1)\hat\delta(k_2\cdot v_2)}_{J_{(1)}^\nu}\,.\label{5.57 m}
    \end{split}
\end{align}
The integral in \eqref{5.57 m} can be done using the results in \eqref{A.9 m}. We know that $J_{(1)}^{\mu}$ has two parts. We first concentrate on the first part, i.e. $\mathcal{J}^{\nu}$ in \eqref{A.9 m}. Therefore, the first term of $f_{\varphi}(x)\Big|^{(1)}_{2}$ takes the following form:
\begin{align}
    \begin{split} \label{ex}
        \widetilde {f_{\varphi}(x)}\Big|_{2}^{(1)}&=T_{\nu}\rmint_{\Omega}e^{-ik\cdot (x-b_1)}\Omega\,\mathcal{J}^{\nu}(\Omega)\,,\\ &
        =i\,\frac{\partial}{\partial l}\,\Big(f_1+f_2+f_3\Big)
    \end{split}
\end{align}
where, $f_i$'s are defined in \eqref{A.19 m}. Apart from the $\mathcal{J}^{\nu}$ part we have another part in form of $J_{(1)}^\nu$, as shown in (\ref{A.9 m}), which gives,
\begin{align}
    \begin{split} \label{ex1}
      \widetilde{ \widetilde{f_{\varphi}(x)}}\Big|_{2}^{(1)}&= \frac{|b| \,n^\nu}{\gamma}T_{\nu}\rmint_{\Omega}e^{-ik\cdot (x-b_1)}\,\Omega^2 \rmint_{0}^1dy\,y\,\,e^{-iyk\cdot b}\frac{K_1(|b|\Delta(k,y))}{4\pi\Delta(k,y)}\,, \\ &
       =\frac{|b|^2 \,T\cdot n}{4\gamma}\rmint_0^\infty dy\,y\rmint d\Omega\, e^{-i\Omega l}\,\Omega^2\, \rmint_{0}
^\infty \frac{dt}{t^2}\,\exp\Big(-t-\frac{\Omega^2|b|^2\bar\Delta(y)^2}{4t}\Big)\,,\\ &
%=\frac{b^2 \,T\cdot n}{4\gamma}\rmint_{0}^\infty dy\,y\rmint_{0}^{\infty} dt\frac{e^{-t}}{t^2}\rmint d\Omega \,\Omega^2 e^{-i\Omega l}\exp\Big(-\frac{\Omega^2 b^2\Bar\Delta(y)^2}{4t}\Big)\\ &
%=\frac{b^2 \,T\cdot n}{4\gamma}\rmint dy\,y\rmint dt\frac{e^{-t}}{t^2}\frac{4 \sqrt{\pi } t^{3/2} e^{-\frac{l^2 t}{b^2 \bar\Delta ^2}} \left(b^2 \bar\Delta ^2-2 l^2 t\right)}{b^5 \bar\Delta ^5}\\ &
=\frac{\sqrt{\pi} \,T\cdot n}{|b|^3\gamma}\rmint_{0}^1dy\,\frac{y}{\bar\Delta(y)^5}\rmint_0^\infty dt\,t^{-1/2}(|b|^2\bar\Delta^2-2l^2t)\exp\Big(-\frac{l^2 t}{|b|^2\bar\Delta^2}-t\Big)\,,\\ &
=\frac{{\pi} \,T\cdot n}{|b|^3\gamma}\rmint_0^1 dy\,\frac{y}{\Bar \Delta}\times \Bigg[\frac{b^2\Bar\Delta^2}{\sqrt{\frac{l^2}{|b|^2 \Bar\Delta ^2}+1}}-l^2\frac{1}{ \left(\frac{l^2}{|b|^2 \Bar\Delta ^2}+1\right)^{3/2}}\Bigg],\,l:=n\cdot (x-b_1-yb)\,.
\end{split}
\end{align}
Then adding (\ref{ex}) and (\ref{ex1}) we get the full expression for $f_{\phi}(x)\Big|_{2}^{(1)}\,.$
The second part in \eqref{ex3} is given by,
\begin{align}
    \begin{split}
    f_{\varphi}(x)\Big|_{2}^{(2)}=(v_1\cdot P\cdot v_1)_{\mu\nu}\rmint_{\Omega}e^{-ik\cdot (x-b_1)}\rmint_{k_2}e^{-ik_2\cdot b}\frac{k_2^\mu\,k_2^\nu}{k_2^2(k-k_2^2)}\hat\delta(k\cdot v_1-k_2\cdot v_1)\hat\delta(k_2\cdot v_2)\,.\label{5.50 m}
    \end{split}
\end{align}
Now, we will rewrite \eqref{5.50 m} to match with \eqref{5.25mmm}. To do this, we introduce an auxiliary variable $k_1$ rewrite (\ref{5.50 m}) using a four-dimensional delta function in the following way.

\begin{align}
\begin{split}f_{\varphi}(x)\Big|_{2}^{(2)}=(v_1\cdot P\cdot v_1)_{\mu\nu}\rmint_{\Omega}e^{-ik\cdot x}\rmint_{k_2}e^{ik_1\cdot b_1}e^{ik_2\cdot b_2}\frac{k_2^\mu\,k_2^\nu}{k_2^2 k_1^2}\hat\delta(k_1\cdot v_1)\hat\delta(k_2\cdot v_2)\hat\delta^{(4)}(k_1+k_2-k)\,.\end{split}
\end{align}
Then we can proceed just like the case of (\ref{5.38nm}). Finally we get,

\begin{equation}\label{ex4}\begin{split}
 f_{\varphi}(x)\Big|_{2}^{(2)}&
      =(v_1\cdot P\cdot v_1)_{\mu\nu}\Big[n^{\mu}n^\nu I^{\rho\sigma}v_{2\rho}v_{2\sigma}-2v_{2\rho}I^{\rho(\mu}n^{\nu)}+I^{\mu\nu}+\bar I^{\mu\nu}\Big].
\\
        \end{split}
\end{equation}\\
Here $I^{\mu\nu}$ and $\bar{I}^{\mu\nu}$ are defined in (\ref{IMUNU}) and (\ref{5.43mm}) respectively. Then  adding (\ref{ex4}) with (\ref{ex}) and (\ref{ex1}) we get the entire expression for $f_{\phi}(x)\Big|_2$ defined in (\ref{ex3}). Finally, adding (\ref{ex3}) with (\ref{5.25mmm}) we get the entire contribution to the waveform coming from this 3-point vertex at 2PM order.

\textbullet $\,\,$ Apart from the above-mentioned diagrams, there will be three more diagrams coming from the $\lambda_3\varphi^3$ vertex, which will contribute to the scalar wave form.
\begin{align}
    \begin{split}
        f(x) &\equiv\thesisinlinefeynman[\chthreeinlinefeynwidth]{new_2.png}\\ &=\lambda_3\frac{m_1 g_1 m_2^2 s_2^2}{8m_p^3}\rmint e^{-ik\cdot x}\rmint_{k_i}e^{i(k_1+k_2)\cdot b_1}e^{i(k_3-k_2)\cdot b_2}\frac{\hat\delta(k_1\cdot v_1+k_2\cdot v_1)\hat\delta(k_2\cdot v_2)\hat\delta(k_3\cdot v_2)}{k_1^2 k_2^2 k_3^2}\hat\delta^{(4)}(k-k_1-k_3)\,,\\ &
       % =\lambda_3\frac{m_1 g_1 m_2^2 s_2^2}{8m_p^3}\rmint e^{-i k\cdot(x-b_1)}\rmint e^{ik_2\cdot b}e^{-ik_3\cdot b}\frac{\hat\delta(k\cdot v_1-k_3\cdot v_1+k_2\cdot v_1)\hat{\delta}(k_2\cdot v_2)\hat\delta(k_3\cdot v_2)}{(k-k_3)^2 k_2^2 k_3^2}\,,\\ &
        =\lambda_3\frac{m_1 g_1 m_2^2 s_2^2}{8m_p^3}\rmint_{k}e^{-ik\cdot (x-b_1)}\rmint_{k_2} e^{ik_2\cdot b}\frac{\hat\delta(k_2\cdot v_2)}{k_2^2}\rmint_{k_3}\hat\delta(k_3\cdot v_2)\hat\delta(k_3\cdot v_1-k\cdot v_1-k_2\cdot v_1)\frac{e^{-ik_3\cdot b}}{k_3^2(k_3-k)^2}.\label{5.62x}
    \end{split}
\end{align}
The $k_3$ integral can be further simplified using Feynman parametrization.
\begin{align}
    \begin{split}
        I_{k_3}&=\rmint_{k_3}\hat\delta(k_3\cdot v_2)\hat\delta(k_3\cdot v_1-k\cdot v_1-k_2\cdot v_1)\frac{e^{-ik_3\cdot b}}{k_3^2(k_3-k)^2},\quad \textrm{with}\quad k^2=0\,,\\ &
      %  =\rmint_0^1 dy\, e^{-iy k\cdot b}\rmint_{k_3}\hat\delta(k_3\cdot v_1-(1-y)k\cdot v_1-k_2\cdot v_1)\hat\delta(k_3\cdot v_2+y k\cdot v_2)\frac{e^{-ik_3\cdot b}}{k_3^4}\,,\\ &
      %  =\rmint_0^1 dy\, e^{-iy k\cdot b}\rmint \frac{e^{i\tilde k_3\cdot b}}{[\tilde k_3^2-y^2(k\cdot v_2)^2+(\frac{y}{\beta}k\cdot v_2+\frac{1-y}{\gamma\beta}k\cdot v_1+\frac{k_2\cdot v_1}{\gamma\beta})^2]}\,,\\ &
        =\rmint_0^1 dy\, e^{-iy k\cdot b}\rmint_{0}^\infty dt\,t\,\rmint_{\tilde k_3}\exp\Big[i \tilde k_3\cdot b-t\, \tilde k_3^2-t\Sigma(y,k,k_2\cdot v_1)^2\Big]\,,\\ &
    =\frac{|b|}{\gamma}\rmint_0^1dy\,e^{-iyk\cdot b}\frac{K_1\Big[|b|\Sigma(k,y,k_2\cdot v_1)\Big]}{4\pi \Sigma(k,y,k_2\cdot v_1)}\,.
    \end{split}
\end{align}
%\vspace{-1.1 cm}
 Note that in our velocity parametrization, $\frac{k_2\cdot v_1}{\gamma\beta}=-k_2^{(1)}$, which implies that the function $\Sigma$ depends only on one component of $k_2$. Hence, we can perform the integral over the other two components of $k_2$, which reads,
 \begin{align}
     \begin{split}
         f(x)&\sim\frac{|b|}{4\pi\gamma}\rmint_0^1 dy \rmint_{k}e^{-ik\cdot (x-b_1+yb)}\rmint dz K_0(|z||b|)\frac{K_1\Big[|b|\Sigma(k,y,k_2\cdot v_1)\Big]}{\Sigma(k,y,k_2\cdot v_1)}\\ &
         =\frac{|b|}{4\pi\gamma}\rmint_{0}^1 dy\rmint_{-\infty}^{\infty} d \Omega\rmint_{-\infty}^\infty dz \,e^{-i \Omega \lambda(x,y,b)}K_{0}(|z||b|)\frac{K_1\Big[|b|\Sigma(\Omega,y,z)\Big]}{\Sigma(\Omega,y,z)},\,\lambda\equiv n\cdot(x-b_1+yb)\label{6.30d}
     \end{split}
 \end{align}
where,
\begin{align}
    \begin{split}
        \Sigma=\sqrt{\frac{\Omega^2\Big[y^2(n\cdot v_2)^2+(1-y)^2(n\cdot v_1)^2+2y(1-y)\gamma (n\cdot v_2)(n\cdot v_1)\Big]}{\gamma^2-1}+z^2-2z\,\Omega\,\Big(\frac{y}{\beta}n\cdot v_2+\frac{1-y}{\gamma\beta}n\cdot v_1\Big)}\,.
    \end{split}
\end{align}
The integral in \eqref{6.30d} can be further simplified by using the integral representation of Bessel functions,
\begin{align}
    \begin{split}
        f(x)&\sim  \frac{|b|}{4\pi\gamma}\rmint_{0}^1 dy\rmint_{-\infty}^{\infty} d \Omega\rmint_{-\infty}^\infty dz \,e^{-i \Omega \lambda(x,y,b)}K_{0}(|z||b|)\frac{K_1\Big[|b|\Sigma(\Omega,y,z)\Big]}{\Sigma(\Omega,y,z)}\,,\\ &
      =  \frac{|b|^2}{16\pi\gamma}\rmint_{0}^1 dy\rmint_{-\infty}^{\infty} d \Omega\rmint_{-\infty}^\infty dz \,e^{-i \Omega \lambda(x,y,b)}\rmint dt_1\, \frac{1}{2t_1}\exp\Big(-t_1-\frac{z^2 |b|^2}{4t_1}\Big)\rmint \frac{dt_2}{t_2^2}\exp\Big(-t_2-\frac{|b|^2\Sigma^2}{4t_2}\Big)\,,\\ &
      =\frac{|b|^2}{32\pi \gamma}\rmint_0^1 dy \rmint_{-\infty}^{\infty} d\Omega e^{-i \Omega\lambda(x,y,b)} \rmint_0^\infty \frac{dt_1 dt_2}{t_1 t_2^2}e^{-(t_1+t_2)}\rmint_{-\infty}^{\infty}dz \exp\Big(-\frac{z^2|b|^2}{4t_1}-\frac{(z^2-2z\theta_1+\theta_2)|b|^2}{4t_2}\Big)\,.
    \end{split}
\end{align}
After doing the $z$ integral we left with,
\begin{align}
    \begin{split}
        f(x)&\sim \frac{|b|^2}{32\pi \gamma}\int_0^1 dy \rmint_{-\infty}^{\infty} d\Omega e^{-i \Omega\lambda(x,y,b)} \rmint_0^\infty \frac{dt_1 dt_2}{t_1 t_2^2}e^{-(t_1+t_2)}\frac{2 \sqrt{\pi } \exp\Big({\frac{|b|^2 \theta _1^2 t_1}{4 t_2^2+4 t_1 t_2}}-\frac{\theta_2|b|^2}{4t_2}\Big)}{b \sqrt{\frac{1}{t_2}+\frac{1}{t_1}}}\,,\\ &
        =\frac{|b|\textcolor{black}{\sqrt{\pi}}}{32\pi \gamma}\rmint_0^1 dy\rmint_0^\infty \frac{dt_1 dt_2}{t_1t_2^2}\frac{e^{-(t_1+t_2)}}{\sqrt{\frac{1}{t_1}+\frac{1}{t_2}}}\rmint_{-\infty}^{\infty} d\Omega\,e^{-i\Omega \lambda}\exp\Big[-\Omega^2\Big(\frac{\hat{\theta}_2|b|^2}{4t_2}-\frac{|b|^2\hat\theta_1^2t_1}{4t_2^2+4t_1t_2}\Big)\Big]\,.\\ &
       \label{6.33d}
    \end{split}
\end{align}
In \eqref{6.33d} the $\Omega$ integral is convergent if,
$$\hat{\theta}_2>\hat{\theta}_1^2\frac{t_1}{t_1+t_2}\implies \underbrace{(\hat\theta_2-\hat\theta_1^2)}_{-y^2}t_1+\hat\theta_2 t_2>0\,,$$
which sets a bound on the $t_1,\, t_2$ integral. The region of integral lies in $t_1< \frac{\hat\theta_2}{y^2}\,t_2$. Now within this region one can do the $\Omega$ integral which reads,
\begin{align}
\begin{split}
      f(x)&   =\frac{|b|}{16 \gamma}\rmint_0^1 dy\rmint_0^\infty \frac{dt_2}{t_2^2}\rmint_0^{\frac{\hat\theta_2}{y^2}t_2}\frac{ dt_1}{t_1}\frac{e^{-(t_1+t_2)}}{\sqrt{\frac{1}{t_1}+\frac{1}{t_2}}}\frac{1}{\sqrt{\frac{\hat{\theta}_2|b|^2}{4t_2}-\frac{|b|^2\hat\theta_1^2t_1}{4t_2^2+4t_1t_2}}}\exp\Big[-\frac{\lambda^2}{4}\frac{1}{\frac{\hat{\theta}_2|b|^2}{4t_2}-\frac{|b|^2\hat\theta_1^2t_1}{4t_2^2+4t_1t_2}}\Big]\\ &\hspace{0.8cm}+ \textrm{pure divergence coming from $\Omega$ integral.}\\ &
        \xrightarrow[]{\textrm{finite part}}\frac{1}{8 \gamma}\rmint_0^1 dy\rmint_0^\infty dt_2 \rmint_{0}^{\frac{\hat\theta_2}{y^2}t_2}dt_1\frac{e^{(-t_1-t_2)}}{t_2\sqrt{t_1}\sqrt{(\hat\theta_2-\hat\theta_1^2)t_1+\hat\theta_2 t_2}}\exp\Big[-\frac{\lambda^2}{|b|^2}\frac{t_2(t_1+t_2)}{(\hat\theta_2-\hat\theta_1^2)t_1+\hat\theta_2 t_2}\Big]
    \end{split}
\end{align}
where,
\begin{align}
    \begin{split}
        &\theta_1(y)=\Omega\Big(\frac{y}{\beta}n\cdot v_2+\frac{1-y}{\gamma\beta}n\cdot v_1\Big):=\Omega\, \hat\theta_1\,,\\ &
        \theta_2(y)=\frac{\Omega^2}{{\gamma^2-1}}\Big({y^2(n\cdot v_2)^2+(1-y)^2(n\cdot v_1)^2+2y (1-y)\gamma(n\cdot v_2)(n\cdot v_1)}\Big):=\Omega^2\,\hat\theta_2\,.
    \end{split}
\end{align}
Restoring the prefactors the waveform has the following form,

\begin{align}
\begin{split}
    f(x)&=\frac{\lambda_3}{(2\pi)^5}\frac{\sqrt{\pi}m_1 g_1 m_2^2 s_2^2}{64\gamma m_p^3}\rmint_0^1 dy\rmint_0^\infty dt_2 \rmint_{0}^{\frac{\hat\theta_2}{y^2}t_2}dt_1\frac{e^{-(t_1+t_2)}}{t_2\sqrt{t_1}\sqrt{(\hat\theta_2-\hat\theta_1^2)t_1+\hat\theta_2 t_2}}\\ &\hspace{0.8cm}\times\exp\Big[-\frac{\lambda^2}{|b|^2}\frac{t_2(t_1+t_2)}{(\hat\theta_2-\hat\theta_1^2)t_1+\hat\theta_2 t_2}\Big]\,.\label{5.70l}
    \end{split}
\end{align}
The integral in \eqref{5.70l}, to the best of our knowledge, does not have a closed-form, and one can, in principle, do the integral numerically while doing further investigations.

%\newpage
\textbullet $\,\,$ Another waveform diagram has the following form:
\begin{align}
    \begin{split}
        f(x)&\sim\thesisinlinefeynman[\chthreeinlinefeynwidth]{new_3.png}\\ &=\lambda_3\frac{m_1 g_1 m_2^2 s_2^2}{8m_p^3} \rmint_{k}e^{-ik\cdot x}\rmint_{k_i}\hat\delta^{(4)}\Big(\sum_ik_i\Big)\frac{\hat\delta(k_1\cdot v_1+k\cdot v_1)\hat\delta(k_2\cdot v_2)\hat\delta(k_3\cdot v_2)}{k_1^2 k_2^2 k_3^2}e^{i(k+k_1)\cdot b_1}e^{i(k_2+k_3)\cdot b_2}\,,\\ &
        =\lambda_3\frac{m_1 g_1 m_2^2 s_2^2}{8m_p^3}\rmint_{\Omega}e^{-i\Omega\,n\cdot(x-b_1)}\rmint_{k_1,k_2}\frac{\hat\delta(k\cdot v_1+k_1\cdot v_1)\hat\delta(k_2\cdot v_2)\hat\delta(k_1\cdot v_2)}{k_1^2 k_2^2 (k_1+k_2)^2}e^{ik_1\cdot b}\,,\\ &
        =\lambda_3\frac{m_1 g_1 m_2^2 s_2^2}{8m_p^3}\frac{1}{n\cdot v_1}\rmint_{k_1,k_2}e^{-i k_1\cdot \delta_1}\frac{\hat\delta(k_1\cdot v_2)\hat\delta(k_2\cdot v_2)}{k_1^2 k_2^2 (k_1+k_2)^2}\,,\delta_1\equiv \frac{n\cdot(x-b_1)}{n\cdot v_1}v_1-b\,,\\ &
        =\lambda_3\frac{m_1 g_1 m_2^2 s_2^2}{64m_p^3(n\cdot v_1)}\rmint_{\vec k_1}\frac{e^{i \vec k_1\cdot \vec \delta_1}}{|\vec k_1|^3}\,.\label{6.35}
    \end{split}
\end{align}
The integral in \eqref{6.35} can be done using the dimensional regularisation and we get,\vspace{0.5 cm}
\begin{align}
    \begin{split}
           f(x)=\frac{\lambda_3}{\textcolor{black}{(2\pi)^3}}\frac{m_1 g_1 m_2^2 s_2^2}{8m_p^3} \frac{1}{32\pi^2(n\cdot v_1)}\Big[{-\log (\vec\delta_1 ^2/\Lambda^2)-\gamma_E +\log (4)}\Big]\,.\label{5.72l}
    \end{split}
\end{align}


\textit{Last but not the least, the total contribution to the scalar waveform at 2PM order is sum of (\ref{e:ba}), (\ref{5.24 mmm}), (\ref{5.25mmm}),(\ref{ex6}), \eqref{5.70l} and \eqref{5.72l}.} \textcolor{black}{Also, note that we need to replace worldline 1 with worldline 2 in all of our results and add them to get the total contribution, as the final result has to be symmetric under this exchange. }
\subsection{Gravitational waveform at 2PM: contribution due to extra scalar DOF}
Now, we will also spell out the correction to the gravitational waveform due to the presence of the extra scalar field. The leading order correction comes from bulk scalar graviton interaction with interacting action:
\begin{align}
    S_{int}=\frac{1}{m_p}\rmint d^4x \,h^{\mu\nu}\partial_{\mu}\varphi\partial_{\nu}\varphi\,.
\end{align}
The corresponding graviton one-point function has the following form:

\begin{align}
    \begin{split}
        k^2\Big\langle h_{\mu\nu}(k)\Big\rangle\Big|_{k^2\rightarrow 0}= \thesisinlinefeynman[\chthreeinlinefeynwidth]{feynman-1.png}\hspace{0.35cm}=\Big(\frac{m_1 s_1 m_2 s_2}{2m_p}\Big)^2\rmint d\mu_{1,2}(k) \frac{k_1^\mu k_2^\nu} {k_1^2 k_2^2}\,.\label{5.63a}
    \end{split}
\end{align}
where the integral measure takes the form,
\begin{align}
    \begin{split}
        \rmint d\mu_{1,2}(k)= \rmint_{k_1,k_2}e^{ik_1\cdot b_1} e^{i k_2\cdot b_2}\hat\delta(k_1\cdot v_1)\hat\delta(k_2\cdot v_2)\hat\delta^{(4)}(k_1+k_2-k)
    \end{split}
\end{align}
Therefore, the correction to the time domain waveform for this particular interaction is given by,
\begin{align}
    \begin{split}
        f_{h}&=\Big(\frac{m_1 s_1 m_2 s_2}{2m_p}\Big)^2\frac{1}{4\pi m_p}\epsilon^\mu \epsilon^\nu \rmint_{\Omega}e^{-ik\cdot x}\rmint d\mu_{1,2}(k)\,\frac{k_{1\mu}k_{2\nu}}{k_1^2\,k_2^2}\,,\\ &
        =\Big(\frac{m_1 s_1 m_2 s_2}{2m_p}\Big)^2\frac{1}{4\pi m_p}\epsilon^\mu \epsilon^\nu \rmint_{\Omega}e^{-ik\cdot x}\rmint_{k_2}e^{i(k-k_2)\cdot b_1}e^{ik_2\cdot b_2}\frac{(k-k_2)_{\mu}k_{2\nu}}{k_2^2(k-k_2)^2}\,\hat\delta\Big(k\cdot v_1-k_2\cdot v_1\Big)\hat\delta(k_2\cdot v_2)\,.\label{5.55mm}
    \end{split}
\end{align}
The total waveform in \eqref{5.55mm} can be separated into two parts, and the parts will be treated separately.
\begin{align}
    \begin{split}
        f_h(x)\Big|_{1}=\Big(\frac{m_1 s_1 m_2 s_2}{2m_p}\Big)^2\frac{\epsilon^\mu \epsilon^\nu}{4\pi m_p}\rmint_{\Omega}e^{-ik\cdot (x-b_1)}k_{\mu}\rmint_{k_2}e^{-ik_2\cdot b}\frac{k_{2\nu}}{k_2^2(k-k_2)^2}\,\hat\delta\Big(k\cdot v_1-k_2\cdot v_1\Big)\hat\delta(k_2\cdot v_2)\,.\label{5.56 mm}
    \end{split}
\end{align}
It is evident from \eqref{5.56 mm} that $f_h\Big|_{1}\sim \epsilon\cdot k$ and we know that $\epsilon\cdot k=0$. Therefore, the first part will not contribute to the waveform. Now, the second part of the waveform takes the following form,
\begin{align}
    \begin{split}
        f_{h}(x)\Big|_{2}=-\Big(\frac{m_1 s_1 m_2 s_2}{2m_p}\Big)^2\frac{\epsilon^\mu \epsilon^\nu}{4\pi m_p}\rmint_{\Omega}e^{-ik\cdot x}\rmint_{k_1,k_2}\frac{k_{2\mu}k_{2\nu}}{k_1^2k_2^2}\hat\delta(k_1\cdot v_1)\hat\delta(k_2\cdot v_2)\hat\delta^{(4)}(k_1+k_2-k)\,.
    \end{split}
\end{align}
Restoring the factors of $\pi$ from the delta functions and the integration  measures, this integral is the same as \eqref{5.25mmm} and results in,


\begin{equation}\label{ch2:e:grav-waveform}
\begin{split}
 f_{h}\Big|_{2}(x)=\Big(\frac{m_1 s_1 m_2 s_2}{2m_p}\Big)^2\frac{1}{2(2\pi)^2 \pi m_p }\epsilon^{\mu}\epsilon^{\nu}(I_{\mu\nu}+\bar I_{\mu\nu}).
\\
        \end{split}
\end{equation}\\
where we again used the fact that $n\cdot \epsilon=0$ and $I^{\mu}$ and $\bar{I}^{\mu\nu}$ are defined in (\ref{IMUNU}) and (\ref{5.43mm}) respectively.
%\Big(-\frac{1}{2 k_1^2(k_1\cdot k)}-\frac{1}{2k_2^2(k_2\cdot k)}\Big)
%\footnote{The $k_2$ integral can be separately done as,
%\begin{align}
 %   \begin{split}
 %       \rmint_{k_2}\frac{\hat\delta(k_2\cdot v_2)}{k_2^2[k_2-(k-\tilde q)]^2}\propto \frac{1}{|\vec k-\vec {\tilde q}|},\textrm{as, there is a constraint $\hat\delta{(k^0-\tilde q^0)}$ i.e. $ k^0=\tilde q^0$.}
 %   \end{split}
%\end{align}
%}
%\textbf{A different approach:}
%\begin{align}
 %   \begin{split}
 %       f_{\varphi}&\sim \rmint_{\Omega}e^{-ik\cdot x}\rmint_{k_i}e^{ik_1\cdot b_1}e^{i k_2\cdot b_2}e^{i k_3\cdot b_2}\frac{\hat\delta(k_1\cdot v_1)\hat\delta(k_2\cdot v_2)\hat\delta(k_3\cdot v_2)}{k_1^2 k_2^2 k_3^2}\hat\delta^{(4)}(k_1+k_2+k_3-k)\\ &
 %      = \rmint_{\Omega}e^{-ik\cdot x}\rmint_{k_i}e^{i(k-k_2-k_3)\cdot b_1}e^{i(k_2+k_3)\cdot b_2}\frac{\hat{\delta}[(k-k_2-k_3)\cdot v_1]\hat\delta(k_2\cdot v_2)\hat\delta(k_3\cdot v_2)}{(k-k_2-k_3)^2 k_2^2 k_3^2}\\ &
 %      \xrightarrow[]{q=k_2+k_3}\rmint_{\Omega}e^{-ik\cdot x}\rmint_{k_2,q}e^{i(k-q)\cdot b_1}e^{iq\cdot b_2}\frac{\hat\delta(k\cdot v_1-q\cdot v_1)\hat\delta(k_2\cdot v_2)\hat\delta(q\cdot v_2)}{(k-q)^2 k_2^2 (q-k_2)^2}\\ &  \xrightarrow[b_1=b]{b_2=0}\rmint_{\Omega}e^{-ik\cdot x}\rmint_{k_2,q}e^{i(k-q)\cdot b}\frac{\hat\delta(k\cdot v_1-q\cdot v_1)\hat\delta(k_2\cdot v_2)\hat\delta(q\cdot v_2)}{(k-q)^2 k_2^2 (q-k_2)^2}
 %   \end{split}
%\end{align}
%\vspace{-0.5 cm}
%Now do the following integral:
%\begin{align}
%    \begin{split}
   %     -\frac{\partial^2}{\partial b_i^2}f_{\varphi}&=\rmint_{\Omega}e^{-i \Omega %n\cdot x}\rmint_{k_2,q}e^{i(k-q)\cdot b}\frac{\hat\delta(k\cdot v_1-q\cdot v_1)\hat\delta(k_2\cdot v_2)\hat\delta(q\cdot v_2)}{k_2^2 (q-k_2)^2}\\ &
 %       =\rmint_{q,k_2}e^{-i \frac{q\cdot v_1}{n\cdot v_1}n\cdot (x-b)}e^{-%iq\cdot b}\frac{\hat\delta(k_2\cdot v_2)\hat\delta(q\cdot v_2)}{k_2^2(q-k_2)^2}\\ &
 %       =\rmint_{q}e^{-iq\cdot w_1}\hat\delta(q\cdot %v_2)\rmint_{k_2}\frac{\hat\delta(k_2\cdot v_2)}{k_2^2(k_2-q)^2},\,w_1\equiv \frac{n\cdot (x-b)}{n\cdot v_1}v_1+b\\ &
%      \propto \rmint_{\vec q}e^{i \vec q\cdot \vec w_1}\frac{1}{|\vec q|}=\frac{4}{\sqrt{\pi}|\vec w_1|^2}
 %   \end{split}
%\end{align}
%Therefore, $f_{\varphi}$ can be found by,
%\begin{align}
%    f_{\varphi}\propto   \mathcal{C}[\Gamma]\rmint^{b}db'\rmint^{b'}db''\frac{1}{\vec w_1(\vec b'')^2}+C_1 b+C_2.\label{7.26}
%\end{align}
%In \eqref{7.26} $\vec w_1^2$ is given by,
%\begin{align}
%    \begin{split}
%        \vec w_1^2=\kappa_0+\kappa_1 b+ \kappa_2 b^2.
%    \end{split}
%\end{align}
%where,
%\begin{align}
%    \begin{split}
%    &    \kappa_0=\gamma^2\beta^2\Big[\frac{n\cdot x}{n\cdot v_1}\Big]^2,\\ &
%    \kappa_1=-2 \sin\theta\cos\varphi \Big[\frac{n\cdot x}{(n\cdot %v_1)^2}\Big],\\ &
%    \kappa_2=1+\sin^2\theta\cos^2\varphi\frac{\gamma^2\beta^2}{(n\cdot v_1)^2}.
%    \end{split}
%\end{align}
%Therefore, the waveform can be written as,
%\begin{align}
%    \begin{split}
 %       f_{\varphi}=-\mathcal{C}[\Gamma]\frac{\log(\kappa_0+\kappa_1 b+ \kappa_2 %b^2)}{2\kappa_2}+\mathcal{C}[\Gamma]\frac{\arctan\Big(\frac{2b\kappa_2+\kappa_1}{\sqrt{4\kappa_0\kappa_2-\kappa_1^2}}\Big)}{\sqrt{4\kappa_0\kappa_2-\kappa_1^2}}\Big(\frac{\kappa_1}{\kappa_2}+2b\Big)+ C_1 b + C_2.\label{7.29}
%    \end{split}
%\end{align}
%The waveform in \eqref{7.29} has two parts one has logarithmic divergent, and another has suitable fall-off at large b. Now, demanding that the fall-off part is zero for large b, one can fix the constants $C_1$ and $C_2$.
%\begin{align}
%    \begin{split}
%    f_{\varphi}\Big|_{\textrm{fall-off}}=\mathcal{C}[\Gamma]\frac{\arctan\Big(\frac{2b\kappa_2+\kappa_1}{\sqrt{4\kappa_0\kappa_2-\kappa_1^2}}\Big)}{\sqrt{4\kappa_0\kappa_2-\kappa_1^2}}\Big(\frac{\kappa_1}{\kappa_2}+2b\Big)+ C_1 b + C_2.
 %   \end{split}
%\end{align}
%Now, for $b\rightarrow  \infty$ the waveform $f_{\varphi}\rightarrow 0$ which fixes the values of the constants.
%\begin{align}
%    \begin{split}
%        f_{\varphi}(b\rightarrow \infty)\Big|_{\textrm{fall-off}}=0=\frac{\mathcal{C}[\Gamma]}{2\kappa_2}\pi\Big(\frac{\kappa_1}{\kappa_2}+2b\Big)+ C_1 b+C_2+\mathcal{O}(1/b).
%    \end{split}
%\end{align}
%which implies,
%\begin{align}
 %   \begin{split}
 %       C_1=-\frac{\pi \mathcal{C}[\Gamma]}{\kappa_2}, C_2=-\frac{\pi \mathcal{C}[\Gamma]\kappa_1}{2\kappa_2^2}.
 %   \end{split}
%\end{align}
%Therefore, the full waveform can be written as,
%\begin{align}
%    \begin{split}
 %    \boxed{   f_{\varphi}(x)=\mathcal{C}[\Gamma]\Big[\textcolor{black}{\frac{\arctan\Big(\frac{2b\kappa_2+\kappa_1}{\sqrt{4\kappa_0\kappa_2-\kappa_1^2}}\Big)}{\sqrt{4\kappa_0\kappa_2-\kappa_1^2}}\Big(\frac{\kappa_1}{\kappa_2}+2b\Big)-\frac{\pi}{\kappa_2}b-\frac{\pi\kappa_1}{2\kappa_2^2}}-\textcolor{black}{\frac{\log(\kappa_0+\kappa_1 b+\kappa_2 b^2)}{2\kappa_2}}\Big]}
 %   \end{split}
%\end{align}
%\hrule
%\textit{Differential equation approach:}\\
%\begin{align}
%    \begin{split}
%        f_{\varphi}(x)=v_1^{\alpha}v_1^{\beta}P_{\mu\nu;\alpha\beta}\rmint_{\Omega}e^{-ik\cdot x}\rmint_{k_2}e^{i(k-k_2)\cdot b}\frac{\hat\delta(k\cdot v_1-k_2\cdot v_1)\hat\delta(k_2\cdot v_2)}{k_2^2(k-k_2)^2}
%    \end{split}
%\end{align}
%Now consider the following integral,
%\begin{align}
 %   \begin{split}
%        -\delta^{\mu\nu}\frac{\partial^2}{\partial b_\mu\partial b_\nu}f_{\varphi}&=v_1^{\alpha}v_1^{\beta}P_{\mu\nu;\alpha\beta}\rmint_{k_2}e^{-i k_2\cdot w_1}\frac{k_2^{\mu}k_2^{\nu}\hat\delta(k_2\cdot v_2)}{k_2^2}\\ &
 %       \rightarrow v_1^{\alpha}v_1^{\beta}P_{ij;\alpha\beta}\rmint_{\vec %k_2}e^{i\vec k_2\cdot \vec w_1}\frac{k_2^i k_2^j}{\vec k_2^2}\\ &
 %       =\tilde {\mathcal{C}}[\Gamma]\,v_1^{\alpha}v_1^{\beta}P_{ij;\alpha\beta}\,\frac{1}{\vec w_1^3}\,\Big[\frac{1}{2}\delta^{ij}-\frac{3}{2}\hat w_1^i\hat w_1^j\Big]
%    \end{split}
%\end{align}
%In our parametrization $b^{\mu}=(0,0,b,0)$ then the differential equation satisfied by the waveform has the following form,
%\begin{align}
 %   \begin{split}
 %       \boxed{
 %       -\frac{\partial^2}{\partial b^2}f_{\varphi}=\mathcal{K}[v_1, n, b]\frac{1}{|\vec w_1(b)|^3}
 %       }\label{7.38}
 %   \end{split}
%\end{align}
%where,
%\begin{align}
%    \begin{split}
  %      \mathcal{K}[v_1,n,b]&=v_1^\alpha v_1^\beta %P_{ij,\alpha\beta}\Big[\frac{1}{2}\delta^{ij}-\frac{3}{2}\frac{w_1^i w_1^j}{|\vec w_1|^2}\Big]\\ &
  %      =\gamma^2\beta^2+\frac{3}{2}-\frac{3}{2}\Bigg[2\Big(\frac{\vec %w_1\cdot \vec v_1}{|\vec w_1|}\Big)^2-1\Bigg]\\ &
 %       =\gamma^2\beta^2+3-3\frac{\gamma^2\beta^2\kappa_0^2-2\gamma^5(1-\beta\cos\theta)\beta^4 b+\gamma^2\beta^2(\kappa_2-1)b^2}{|\vec w_1|^2}
%    \end{split}
%\end{align}
%The solution has the following form:
%\begin{align}
%    \begin{split}
%        f_{\varphi}&=-\rmint^{b}db'\rmint^{b'}db''\mathcal{K}[v_1, n,  b'']\frac{1}{|\vec w_1(b)|^3}+c_1b+c_2\\ &
  %      =\frac{4}{\left(\kappa _1^2-4 \kappa _0 \kappa _2\right){}^2 \sqrt{b %\left(b \kappa _2+\kappa _1\right)+\kappa _0}} \Big[\kappa _0^2 \left(8 b^2 \beta ^2 \gamma ^2 \kappa _2^2+\beta ^2 \gamma ^2 \kappa _1^2+4 \kappa _2 \left(\beta ^2 \gamma ^2-3\right)-8 \beta ^2 \gamma ^2+8 b \beta ^2 \gamma ^2 \kappa _1 \kappa _2\right)\\ &+4 \beta ^2 \gamma ^2 \kappa _2 \kappa _0^3+b \kappa _1 \left(\kappa _1^2 \left(\beta ^2 \gamma ^2+3\right)- 8 b \beta ^4 \gamma ^5 \kappa _2 (\beta  \cos (\theta )-1)+\kappa _1 \left(b \kappa _2 \left(2 \beta ^2 \gamma ^2+3\right)-\\ &\beta ^2 \gamma ^2 \left(6 \beta ^3 \gamma ^3 \cos (\theta )-6 \beta ^2 \gamma ^3+b\right)\right)\right)+\kappa _0 \Big(\kappa _1^2 \left(\beta ^2 \gamma ^2+3\right)+4 \kappa _1 \left(b \kappa _2 \left(\beta ^2 \gamma ^2-3\right)-2 \beta ^2 \gamma ^2 \left(\beta ^3 \gamma ^3 \cos (\theta )-\beta ^2 \gamma ^3+b\right)\right)\\ & -4 b \kappa _2 \left(2 \beta ^5 \gamma ^5 \cos (\theta )-2 \beta ^4 \gamma ^5+b \beta ^2 \gamma ^2+3 b \kappa _2\right)\Big)\Big]+c_1 b+c_2.
  %  \end{split}
%\end{align}
%The constants $c_1,c_2$ can be fixed by demanding $f_{\varphi}\rightarrow 0$ for $b\rightarrow \infty$.
%\begin{align}
 %   \begin{split}
  %     &\textstyle{ c_1=-\frac{4 \left(8 \beta ^2 \gamma ^2 \kappa _0^2 \kappa _2^2-4 \kappa _0 \kappa _2 \left(\beta ^2 \gamma ^2+3 \kappa _2\right)+\kappa _1 \left(\kappa _1 \left(\kappa _2 \left(2 \beta ^2 \gamma ^2+3\right)-\beta ^2 \gamma ^2\right)-8 \beta ^4 \gamma ^5 \kappa _2 (\beta  \cos (\theta )-1)\right)\right)}{\sqrt{\kappa _2} \left(\kappa _1^2-4 \kappa _0 \kappa _2\right){}^2},} \\ &
  %     c_2=-\frac{2}{\kappa _2^{3/2} \left(4 \kappa _0 \kappa _2-\kappa _1^2\right){}^2} \Big[-16 \beta ^5 \gamma ^5 \kappa _0 \kappa _2^2 \cos (\theta )-4 \beta ^5 \gamma ^5 \kappa _1^2 \kappa _2 \cos (\theta )+16 \beta ^4 \gamma ^5 \kappa _0 \kappa _2^2+4 \beta ^4 \gamma ^5 \kappa _1^2 \kappa _2+\beta ^2 \gamma ^2 \kappa _1^3\\ &\hspace{0.35cm}+8 \beta ^2 \gamma ^2 \kappa _0^2 \kappa _1 \kappa _2^2+8 \beta ^2 \gamma ^2 \kappa _0 \kappa _1 \kappa _2^2-12 \beta ^2 \gamma ^2 \kappa _0 \kappa _1 \kappa _2-12 \kappa _0 \kappa _1 \kappa _2^2+3 \kappa _1^3 \kappa _2\Big].
  %  \end{split}
%\end{align}
%Therefore, the waveform looks like,
%\begin{align}
%  \boxed{  \begin{aligned}
%  f_{\varphi}&= \frac{1}{\left(4 \kappa _0 \kappa _2-\kappa _1^2\right){}^2}\Bigg(-\frac{2 \beta ^2 \gamma ^2 \kappa _1^3}{\kappa _2^{3/2} }+\frac{24 \beta ^2 \gamma ^2 \kappa _0 \kappa _1}{\sqrt{\kappa _2} }-\frac{8 \beta ^4 \gamma ^5 \kappa _1^2}{\sqrt{\kappa _2} }+\frac{8 \beta ^5 \gamma ^5 \kappa _1^2 \cos (\theta )}{\sqrt{\kappa _2}}-\\ &
%      \frac{6 \kappa _1^3}{\sqrt{\kappa _2} }-32 \beta ^4 \gamma ^5 \kappa _0 \sqrt{\kappa _2}+32 \beta ^5 \gamma ^5 \kappa _0 \sqrt{\kappa _2} \cos (\theta )+24 \kappa _0 \kappa _1 \sqrt{\kappa _2}-
 %    16 \beta ^2 \gamma ^2 \kappa _0 \kappa _1 \sqrt{\kappa _2}-16 \beta ^2 %\gamma ^2 \kappa _0^2 \kappa _1 \sqrt{\kappa _2}\Bigg)\\ &
 %\frac{1}{\left(\kappa _1^2-4 \kappa _0 \kappa _2\right){}^2 \left(b \left(b \kappa _2+\kappa _1\right)+\kappa _0\right){}^{3/2}}   \Bigg( -32 \beta ^2 \gamma ^2 \kappa _0^2+
 %    32 \beta ^4 \gamma ^5 \kappa _0 \kappa _1-32 \beta ^5 \gamma ^5 \kappa _0 \kappa _1 \cos (\theta )+\\ &
 %    12 \kappa _0 \kappa _1^2+4 \beta ^2 \gamma ^2 \kappa _0 \kappa _1^2+
 %     4 \beta ^2 \gamma ^2 \kappa _0^2 \kappa _1^2-48 \kappa _0^2 \kappa _2+
 %     16 \beta ^2 \gamma ^2 \kappa _0^2 \kappa _2+16 \beta ^2 \gamma ^2 \kappa _0^3 \kappa _2\Bigg)+\\ &
 %    \textcolor{black}{ \frac{b}{\left(\kappa _1^2-4 \kappa _0 \kappa %_2\right)^2}\left(\frac{4 \beta ^2 \gamma ^2 \kappa _1^2}{\sqrt{\kappa _2} }+16 \beta ^2 \gamma ^2 \kappa _0 \sqrt{\kappa _2}-32 \beta ^4 \gamma ^5 \kappa _1 \sqrt{\kappa _2}+32 \beta ^5 \gamma ^5 \kappa _1 \sqrt{\kappa _2} \cos (\theta )-\right.}\\ &
%    \textcolor{black}{  12 \kappa _1^2 \sqrt{\kappa _2}-8 \beta ^2 \gamma ^2 \kappa _1^2 \sqrt{\kappa _2}+48 \kappa _0 \kappa _2^{3/2}-32 \beta ^2 \gamma ^2 \kappa _0^2 \kappa _2^{3/2}-}\\ &
 %    \textcolor{black}{ \frac{32 \beta ^2 \gamma ^2 \kappa _0 \kappa _1}{\left(b \left(b \kappa _2+\kappa _1\right)+\kappa _0\right){}^{3/2}}+\frac{24 \beta ^4 \gamma ^5 \kappa _1^2}{ \left(b \left(b \kappa _2+\kappa _1\right)+\kappa _0\right){}^{3/2}}-
  %    \frac{24 \beta ^5 \gamma ^5 \kappa _1^2 \cos (\theta )}{ \left(b \left(b \kappa _2+\kappa _1\right)+\kappa _0\right){}^{3/2}}+\frac{12 \kappa _1^3}{ \left(b \left(b \kappa _2+\kappa _1\right)+\kappa _0\right){}^{3/2}}+}\\ &\textcolor{black}{
  %%    \frac{4 \beta ^2 \gamma ^2 \kappa _1^3}{ \left(b \left(b \kappa %_2+\kappa _1\right)+\kappa _0\right){}^{3/2}}+\frac{32 \beta ^4 \gamma ^5 \kappa _0 \kappa _2}{ \left(b \left(b \kappa _2+\kappa _1\right)+\kappa _0\right){}^{3/2}}-
  %    \frac{32 \beta ^5 \gamma ^5 \kappa _0 \kappa _2 \cos (\theta )}{ \left(b %\left(b \kappa _2+\kappa _1\right)+\kappa _0\right){}^{3/2}}-\frac{48 \kappa _0 \kappa _1 \kappa _2}{ \left(b \left(b \kappa _2+\kappa _1\right)+\kappa _0\right){}^{3/2}}+}\\ &
%   \textcolor{black} { \left.\frac{32 \beta ^2 \gamma ^2 \kappa _0^2 \kappa _1 \kappa _2}{ \left(b \left(b \kappa _2+\kappa _1\right)+\kappa _0\right){}^{3/2}}+\frac{16 \beta ^2 \gamma ^2 \kappa _0 \kappa _1 \kappa _2}{ \left(b \left(b \kappa _2+\kappa _1\right)+\kappa _0\right){}^{3/2}}\right)+}
%  \\ &    \textcolor{black}{\frac{b^2}{\left(\kappa _1^2-4 \kappa _0 \kappa _2\right){}^2 \left(b \left(b \kappa _2+\kappa _1\right)+\kappa _0\right){}^{3/2}}\Bigg(-4 \beta ^2 \gamma ^2 \kappa _1^2-16 \beta ^2 \gamma ^2 \kappa _0 \kappa _2+
%    32 \beta ^4 \gamma ^5 \kappa _1 \kappa _2-32 \beta ^5 \gamma ^5 \kappa _1 \kappa _2 \cos (\theta )}\\ &+
%   \textcolor{black}{ 12 \kappa _1^2 \kappa _2+8 \beta ^2 \gamma ^2 \kappa _1^2 \kappa _2-
 %     32 \beta ^2 \gamma ^2 \kappa _0^2 \kappa _2^2+48 \kappa _0 \kappa _2^2\Bigg)} \hspace{0.8cm}
%    \end{aligned}}
%\end{align}
\section{Towards massive waveform}\label{sec6}
As mentioned in the previous section, we could get exact results only when we set the mass of the scalar field to zero. In this section, we will discuss how the radiation integrals become more complicated when we make the scalar field massive. Then, we provide an approximate way to evaluate those integrals and get an analytic result. \par 
%\subsection*{Main difficulty in massive integrals}
For massive scalar, the radiated momentum has to be parametrized as follows,
\begin{align}
    k^{\mu}=\Omega n^\mu,\, n^\mu=(1,\sqrt{1-\frac{m^2}{\Omega^2}}\boldsymbol{\hat{x}})\,.\label{6.1 m}
\end{align}
\textit{From \eqref{6.1 m} it is clear that the unit vector towards the direction of the observed scalar depends on the observed frequency $\Omega$. We need to integrate over all possible frequencies $\Omega$ for computing the time domain waveform. There lies the difficulty due to the presence of complicated phase factors in the time-domain waveform integrand. }
\subsection*{Radiation integrals for massive scalar waveform}
Before proceeding further, we first list all the relevant massive integrands having a massless counterpart, which we evaluated in the previous section.
\begin{itemize}
  \item The massive integrated corresponding to $ \lambda_3 \varphi^3$ in \eqref{wave0} vertex has the following form:
    \begin{align}
        \begin{split}
          f_{\varphi}\propto \rmint_{\Omega}e^{-ik\cdot x}\rmint_{k_i}e^{ik_1\cdot b_1}e^{ik_2\cdot b_2}\frac{\hat\delta(k_1\cdot v_1)\hat\delta(k_2\cdot v_2)}{(k_1^2-m^2)(k_2^2-m^2)}\hat\delta^{(4)}(k_1+k_2-k)\,. \label{6.11 m}
        \end{split}
    \end{align}
\item The massive integrated corresponding to $\lambda_4 \varphi^4$ in \eqref{5.11 mn} vertex has the following form:
\begin{align} 
    \begin{split} \label{6.11 mm}
        f_{\varphi}\propto \rmint_{\Omega}e^{-ik\cdot x}\rmint_{k_i}e^{ik_1\cdot b_1}e^{i k_2\cdot b_2}e^{i k_3\cdot b_2}\frac{\hat\delta(k_1\cdot v_1)\hat\delta(k_2\cdot v_2)\hat\delta(k_3\cdot v_2)}{(k_1^2-m^2)(k_2^2 -m^2)(k_3^2-m^2)}\hat\delta^{(4)}(k_1+k_2+k_3-k)\,.
    \end{split} 
\end{align}
\item  
Massive integral corresponding to worldline radiation in \eqref{5.24 mmm} \footnote{Another way to approximate this integral is to take a large velocity limit. We have discussed that in detail in Appendix~(\ref{ch2:app:F}). Although this approximation helps get an approximate closed-form result for this diagram, it is unclear whether a large velocity approximation will also be useful for all other massive integrals. Hence, we focused on the stationary-phase method in the main text.}:
\begin{align}\label{6.4 mm}
f_{\varphi}(x)&\propto\rmint_{\Omega}e^{-ik\cdot (x-b_1)}\rmint_{\omega,k_1}e^{-i k_1\cdot b}\hat{\delta}(k_1\cdot v_2)\hat{\delta}[\Omega\, n(m,\Omega)\cdot v_1-\omega]\hat\delta(k_1\cdot v_1-\omega)\frac{\Omega(n\cdot k_1)}{\omega^2 (k_1^2-m^2)}\,.
\end{align}
     \item The massive integrated corresponding to the diagram \eqref{5.22 m} coming from the derivative interaction term $\rmint d^4x h^{\mu\nu}\partial_{\mu}\varphi\partial_\nu \varphi$ has the following form:
    \begin{align}
        \begin{split}
            f_{\varphi}\propto \rmint_{\Omega}e^{-ik\cdot x}\rmint_{k_1,k_2}e^{ik_1\cdot b_1}e^{ik_2\cdot b_2}\hat\delta(k_1\cdot v_1)v_1^{\alpha}v_1^{\beta}\hat\delta(k_2\cdot v_2)\frac{k_2^{\mu}k_2^{\nu}P_{\mu\nu,\alpha\beta}}{k_1^2(k_2^2-m^2)}\delta^{(4)}(k_1+k_2-k) \,.\label{6.10 m}
        \end{split}
    \end{align}
\end{itemize}
{\color{black}
\subsection{Stationary Phase (SP) approximation: the need and general procedure}
We mentioned in the previous subsection how the massive integrals get complicated. Now, we will evaluate the integrals mentioned above using the \textit{stationary phase approximation} for large observation distance (and time) $|x|\rightarrow \infty$. This will help us to get an approximate analytical result. Note that the idea behind the stationary phase approximation stems from the fact that sinusoids with rapidly varying phases interfere destructively. However, they can be added constructively if they have the same phases. Physically, in a scattering scenario, we expect a burst signal when the emitted radiation gets detected around some frequencies. Hence, we believe that using an approximation method, such as the stationary phase method, to evaluate these integrals that will give the waveform is reasonable. 
We have to evaluate the following type of integrals.
\begin{align}
\begin{split}
        f_{\varphi}\propto \rmint_{\Omega}\exp\Big({-i\underbrace{\Omega \,n(\Omega)\cdot (x-b_1)}_{l(\Omega)}}\Big)g(\Omega)\,.\label{6.13 m}
    \end{split}
\end{align}
The stationary point can be obtained by,
\begin{align}
    \begin{split}
        \frac{d}{d\Omega}\Big(\Omega \,n(\Omega)\cdot (x-b_1) \Big)\Big |_{\Omega_0}=0\,.\label{6.14 m}
    \end{split}
\end{align}
Therefore, the integral can be approximated to,
\begin{align}
    \begin{split}
        f_{\varphi}\propto g(\Omega_0) \,e^{-i l(\Omega_0)}e^{\frac{-i\pi}{4}\textrm{sgn}(l''(\Omega_0))}\sqrt{\frac{\pi}{|l''(\Omega_0)|}}+\mathcal{O}\Big(\frac{1}{|x-b_1|}\Big)
    \end{split}
\end{align}
{\color{black}
\textbf{\textit{Validity of the stationary phase approximation:}}
The stationary phase approximation is justified only under a few physical assumptions. First, the detector must be placed in the far radiation zone, namely
\begin{align}
    |x-b_1|\gg b,
\end{align}
where \(b\) denotes the characteristic length scale of the scattering process, such as the impact parameter. Second, the prefactor \(g(\Omega)\) should vary slowly in the neighbourhood of the saddle point. More precisely, if \(\Omega_0\) denotes the stationary point, then one requires
\begin{align}
    g(\Omega)\simeq g(\Omega_0)
\end{align}
within the saddle width
\begin{align}
    \Delta\Omega \sim \frac{1}{\sqrt{|l''(\Omega_0)|}} .
\end{align}
For a massive waveform, the spatial momentum satisfies  $ k(\Omega)=\sqrt{\Omega^2-m^2}.$ Thus different frequency modes propagate with different group velocities,
\begin{align}
    v_g(\Omega)
    =
    \frac{k(\Omega)}{\Omega}
    =
    \frac{\sqrt{\Omega^2-m^2}}{\Omega}.
\end{align}
Consider the phase factor in the form
\begin{align}
    l(\Omega)=\Omega T-k(\Omega)R,
    \qquad
    R=|\vec{x}-\vec{b}_1| .
\end{align}
The saddle point is determined by $   l'(\Omega_0)=0 .$ Now using $   \frac{dk}{d\Omega}
    =
    \frac{\Omega}{\sqrt{\Omega^2-m^2}}
    =
    \frac{1}{v_g(\Omega)},$ the saddle condition gives
\begin{align}
    T
    =
    R\frac{\Omega_0}{k(\Omega_0)}
    =
    \frac{R}{v_g(\Omega_0)} .
\end{align}
Therefore, the saddle chooses the frequency mode whose group velocity carries the radiation from the source to the detector in the observed time \(T\). Since a massive field has subluminal group velocity,  $  0<v_g(\Omega)<1,$, the above condition implies   $ T>R .$ Thus the massive waveform is observed inside the future light cone, rather than exactly on null infinity.
Finally, the stationary phase approximation also requires the phase evaluated at the saddle to be large. Solving the saddle condition gives
\begin{align}
    \Omega_0
    =
    \frac{mT}{\sqrt{T^2-R^2}},
    \qquad
    k(\Omega_0)
    =
    \frac{mR}{\sqrt{T^2-R^2}} .
\end{align}
Hence
\begin{align}
    l(\Omega_0)
    =
    \Omega_0 T-k(\Omega_0)R
    =
    m\sqrt{T^2-R^2}.
\end{align}
Therefore, the large-phase condition becomes   $ m\sqrt{T^2-R^2}\gg 1,$ or equivalently  $  T^2-R^2\gg \frac{1}{m^2}.$ In this regime, the integral is dominated by a small neighbourhood of the saddle point, and the stationary phase approximation is physically well justified.}}
\subsection{Dealing with the massive integrals via SP approximation}
We start with the integral mentioned in (\ref{6.11 m}) we get, 
\begin{align}
    g(\Omega_0)=\rmint_{k_2}e^{-ik_2\cdot b}\frac{\hat\delta(k_0\cdot v_1-k_2\cdot v_1)\hat\delta(k_2\cdot v_2)}{[(k_2-k_0)^2-m^2](k_2^2-m^2)},\,k_0\equiv \Omega_0 \,n({\Omega_0})\,.
\end{align}
Then the final integral over $k_2$ can be done using Feynman parametrization in the following way, 
\begin{align}
    \begin{split}
        g(\Omega_0)&=\rmint_{0}^1 dy\rmint e^{-ik_2\cdot b}\frac{\hat\delta(k_2\cdot v_1-k_0\cdot v_1)\hat\delta(k_2\cdot v_2)}{[(k_2-yk_0)^2-(1-y+y^2)m^2]^2},\,k_0^2=m^2\,,\\ &
        \xrightarrow[]{k_2-yk_0\rightarrow \bar q}\rmint_0^1 dy e^{-iy\,k_0\cdot b}\rmint \frac{e^{-i \bar q\cdot b}}{[\bar q^2-(1-y+y^2)m^2]^2}\hat\delta(\bar q\cdot v_1-(1-y)k_0\cdot v_1)\hat\delta(\bar q\cdot v_2+y k_0\cdot v_2)\,,\\ &
      %  =\frac{1}{\gamma}\rmint_0^1 \,e^{-iy\,k_0\cdot b}\rmint_0^\infty dt\,t\rmint_{\tilde q}\exp\Big[i\tilde q\cdot b-t\tilde q^2-t\Delta(k_0,y)^2-t(1-y+y^2)m^2\Big]\,,\\ &
        =\frac{1}{\gamma}\rmint_0^1\,dy\,e^{-iy\,k_0\cdot b}\rmint \frac{dt}{4\pi}\exp\Big(-\frac{|b|^2}{4t}-t\Delta(k_0,y)^2-t(1-y+y^2)m^2\Big)\,,\\ &
        =\frac{b}{4\pi \gamma}\rmint_0^1
dy\,e^{-iy\,k_0\cdot b}\frac{K_1\Big( b\sqrt{\Delta(y)^2+(1-y+y^2)m^2}\Big)}{\sqrt{\Delta(y)^2+(1-y+y^2)m^2}}\,.
\end{split}
\end{align}
Then, the contribution to the waveform takes the following form,
\begin{equation}\label{ch2:e:massive-waveform}\begin{split}
f_{\varphi}\propto \frac{b}{4\pi\gamma}\sqrt{\frac{\pi}{|l''(\Omega_0)|}}\rmint_0^1 dy\,e^{-i\sigma(b,\Omega_0,m,y)}\frac{K_1\Big( b\sqrt{\Delta(y)^2+(1-y+y^2)m^2}\Big)}{\sqrt{\Delta(y)^2+(1-y+y^2)m^2}},\,\sigma=y k_0\cdot b+l(\Omega_0)-\frac{\pi}{4}\,.
\\
        \end{split}
\end{equation}

%\textcolor{black}{Now note that the massless counterpart of (\ref{6.11 mm}) does not contribute to the waveform as shown in the previous section in (\ref{nocont}) as well as in the Appendix~(\ref{App2}). It was of a purely divergent nature. In case it has to contribute to the waveform for the massive case, the only possibility would have been that the finite part coming from (\ref{6.11 mm}) is proportional to the mass of the scalar field. It can only happen if the mass factor comes from the vertex itself. But the source of mass is the propagators and the phase factors only. Hence, by the same argument as before, it will not contribute to the waveform like its massless counterpart. An explicit check (numerical/analytical) would be good, although we have provided sufficient arguments here. We leave that for future work.}\\\\
Next, we consider the integral in \eqref{6.4 mm} and apply  the method of stationary phase to it. It gives the following,
%\begin{align}
 %   \begin{split}
 %       f_{\varphi}\propto \rmint_{\Omega}e^{-i\Omega\,n(\Omega)\cdot (x-b_1)}g(\Omega)
 %   \end{split}
%\end{align}
%where,
%\begin{align}
 %   \begin{split}
 %       g(\Omega)\equiv \rmint_{k_1} e^{-ik_1\cdot b}\frac{-k_1\cdot k}{(k_1^2-m^2)(k_1\cdot v_1+i\epsilon)^2}\hat\delta(k_1\cdot v_1-k\cdot v_1)\,.
%    \end{split}
%\end{align}
%Now, using the method of stationary phase, we will get,
\begin{align}
    \begin{split}
        f(\varphi_0)\propto g(\Omega_0)e^{-il(\Omega_0)}e^{\frac{i\pi}{4}}\sqrt{\frac{\pi}{|l''(\Omega_0)|}}
    \end{split}
\end{align}
where $g(\Omega_0)$ is given by,
\begin{align}
    \begin{split}
        g(\Omega_0)&=\frac{1}{\Omega_0^2 \,\Big(n(\Omega_0)\cdot v_1 \Big)^2}\rmint_{k_1}e^{-ik_1\cdot b}\frac{-k_1\cdot k_0}{(k_1^2-m^2)}\hat\delta(k_1\cdot v_1-k_0\cdot v_1)\hat\delta(k_1\cdot v_2)\,,\\ &
        =-\frac{k_0^\mu}{\Omega_0^2 \,\Big(n(\Omega_0)\cdot v_1 \Big)^2}\underbrace{\rmint_{k_1}e^{-ik_1\cdot b}\frac{k_1^\mu}{k_1^2-m^2}\hat\delta(k_1\cdot v_1-k_0\cdot v_1)\hat\delta(k_1\cdot v_2)}_{g^\mu(\Omega_0)}\,.
    \end{split}
\end{align}
Here, the $g_\mu(\Omega_0)$ can be evaluated using Passarino-Veltman reduction.
\begin{align}
    \begin{split}
        g^{\mu}(\Omega_0)= \lambda_b b^\mu+\lambda_1 v_1^\mu+\lambda_2 v_2^\mu\,.
    \end{split}
\end{align}
To extract the constants, one needs to contract the index structure in the RHS with the LHS, which gives the following:

\textbullet $\,\,$ Contracting with $b^\mu$ gives,
\begin{align}
    \begin{split}
        -\lambda_b \,|b|^2&=\rmint_{k_1} e^{-i k_1\cdot b} \frac{k_1\cdot b}{k_1^2-m^2}\hat\delta(k_1\cdot v_1-k_0\cdot v_1)\hat\delta(k_1\cdot v_2)\,,\\ &
        =i\lim_{\kappa\to 1}\frac{\partial}{\partial \kappa}\rmint_{0}^\infty dl\,l\,\frac{J_{0}(\kappa\,l|b|)}{l^2+\hat m^2},\,\hat m^2\equiv m^2+\Big(\frac{k_0\cdot v_1}{\sqrt{\gamma^2-1}}\Big)^2\,.\label{6.35 m}
    \end{split}
\end{align}
\textbullet $\,\,$ Contracting with $v_1^\mu$ we will get,
\begin{align}
    \begin{split}
        \lambda_1+\gamma\lambda_2&= \rmint_{k_1} e^{-i k_1\cdot b} \frac{k_1\cdot v_1}{k_1^2-m^2}\hat\delta(k_1\cdot v_1-k_0\cdot v_1)\hat\delta(k_2\cdot v_2)\,,\\ &
        =k_0\cdot v_1\rmint_{k_1}e^{-ik_1\cdot b}\frac{\hat\delta(k_1\cdot v_1-k_0\cdot v_1)\hat\delta(k_1\cdot v_2)}{k_1^2-m^2}\,,\\ &
        =k_0\cdot v_1\rmint_0^\infty dl\,l\,\frac{J_{0}(l|b|)}{l^2+\hat m^2
        }\,.\label{6.36 m}
    \end{split}
\end{align}
\textbullet $\,\,$ Contracting with $v_2^\mu$ we will get,
\begin{align}
    \begin{split}
        \gamma \lambda_1+\lambda_2=0,\,\textrm{as}\,k_1\cdot v_2=0\,.\label{6.37 m}
    \end{split}
\end{align}
\textcolor{black}{Solving \eqref{6.36 m} and \eqref{6.37 m} we will get,
\begin{align}
    \begin{split}
      &  \lambda_1=\frac{1}{1-\gamma^2}\,k_0\cdot v_1\rmint_0^\infty dl\,l\,\frac{J_{0}(l|b|)}{l^2+\hat m^2}=\frac{1}{1-\gamma^2}\,(k_0\cdot v_1)K_0\left(|b| \hat{m}\right)\,,\\ &
      \lambda_2=\frac{\gamma}{\gamma^2-1} k_0\cdot v_1\,\rmint_0^\infty dl\,l\,\frac{J_{0}(l|b|)}{l^2+\hat m^2}=\frac{\gamma}{\gamma^2-1} (k_0\cdot v_1)\,K_0\left(|b| \hat{m}\right)
    \end{split}
\end{align}
and from \eqref{6.35 m} we will get,
\begin{align}
    \lambda_b=\frac{i}{|b|}\rmint_0^\infty dl\,l^2\,\frac{J_1(l|b|)}{l^2+\hat m^2}=\frac{i}{|b|}\hat{m}\, K_1\left(|b| \hat{m}\right)\,.\label{6.19ab}
\end{align}
}
Therefore, its contribution to the waveform looks like,

\begin{align}
    \begin{split}
    f(x)\propto -\frac{1}{\Omega_0^2 \,\Big(n(\Omega_0)\cdot v_1 \Big)^2} \,e^{-il(\Omega_0)}e^{\frac{i\pi}{4}}\sqrt{\frac{\pi}{|l''(\Omega_0)|}}\,k_0\cdot[\lambda_b b+\lambda_1 v_{1}+\lambda_2 v_{2}]\,.
    \end{split}
\end{align}
Now, we analyze the integral, which comes from the derivative interaction as mentioned in \eqref{6.10 m}. It takes the following form,
\begin{align}
    \begin{split}
       f_{\varphi}\propto (v_1\cdot P\cdot v_1)_{\mu\nu} \rmint_{\Omega}\exp\Big(-i\Omega\,n(\Omega)\cdot (x-b_1)\Big)J_{(2)}^{\mu\nu}(\Omega)\,,
    \end{split}
\end{align}
where,
\begin{align}
    \begin{split}
        J_{(2)}^{\mu\nu}(\Omega)=\rmint_{q}\hat\delta(q\cdot v_1-k\cdot v_1)\hat\delta(q\cdot v_2)\frac{q^\mu q^\nu}{(q-k)^2(q^2-m^2)}e^{-iq \cdot b}\,.\label{6.20 m}
    \end{split}
\end{align}
The integral over $q$ in \eqref{6.20 m} has been done in \eqref{A.22 m}. Then using the method of the stationary phase, we get,
\begin{align}
    \begin{split}
        f_{\varphi}\propto (v_1\cdot P\cdot v_1)_{\mu\nu}\sqrt{\frac{\pi}{|l''(\Omega_0)|}}e^{-il(\Omega_0)}e^{\frac{i\pi}{4}}\,J^{\mu\nu}_{(2)}(\Omega_0)\,.
    \end{split}
\end{align}
where $l(\Omega)$ and the saddle point $\Omega_0$ is defined in \eqref{6.13 m} and \eqref{6.14 m} respectively.

\textbullet $\,\,$ Apart from the integrals listed in (\ref{6.11 m}), (\ref{6.4 mm}) and (\ref{6.10 m}), one can possibly have one term proportional to scalar mass $m$  contributing to the waveform at 2PM order. 
\textcolor{black}{$$S_{int}=-\frac{1}{2}\frac{m^2}{m_{p}}\rmint d^4x\,h\varphi^2\,.$$}
The corresponding contribution to the scalar one-point function is , 
\begin{align}
   \begin{split}
      k^2\Big\langle \varphi(k)\Big\rangle=m^2\Big(\frac{m_1m_2s_2}{8m_p^3}\Big)\rmint d\mu_{1,2}(k)\frac{v_1^\alpha v_1^\beta\,P_{\mu\nu;\alpha\beta}\,\eta^{\mu\nu}}{k_1^2(k_2^2-m^2)}\,.\label{ch2:5.23b}
    \end{split}
\end{align}
where, the integral measure has the following form,
\begin{align}
    \begin{split}
        d\mu_{1,2}(k)&=\rmint_{k_1,k_2}\hat{\delta}^{(4)}(k_1+k_2-k)e^{ik_1\cdot b_1+ik_2\cdot b_2 }\hat\delta(k_1\cdot v_1)\hat\delta(k_2\cdot v_2)\,,\\ &
        =\rmint_{k_2}e^{ik\cdot b_1}e^{-ik_2\cdot b}\hat\delta(k\cdot v_1-k_2\cdot v_1)\hat\delta(k_2\cdot v_2)\,.
    \end{split}
\end{align}
Therefore the corresponding contrubution to scalar waveform has the following form,
\begin{align}
    \begin{split}
        f_{\varphi}(x)&\propto\thesisinlinefeynman[\chthreeinlinefeynwidth]{new.png}\\ &=-m^2\Big(\frac{m_1m_2s_2}{4m_p^3}\Big)\rmint_{\Omega}e^{-ik\cdot(x-b_1)}\rmint_{k_2}e^{-ik_2\cdot b}\frac{\hat\delta(k_2\cdot v_2)\hat\delta(k\cdot v_1-k_2\cdot v_1)}{(k-k_2)^2(k_2^2-m^2)}\,,\\ &
        =-m^2\Big(\frac{m_1m_2s_2}{4m_p^3}\Big)\frac{|b|}{4\pi\gamma}\sqrt{\frac{\pi}{|l''(\Omega_0)|}}\rmint_0^1 dy\,e^{-i\sigma(|b|,\Omega_0,m,y)}\frac{K_1\Big(|b|\sqrt{\Delta(y)^2+(1-y)^2m^2}\Big)}{\sqrt{\Delta(y)^2+(1-y)^2m^2}}\,.\label{6.26f}
    \end{split}
\end{align}
One can see that the contribution to the waveform from \eqref{6.26f} identically vanishes when one takes the scalar to be massless. Hence, the contribution to the massive scalar waveform takes the form:

\begin{align}
    \begin{split} 
           f(x)\propto -m^2\Big(\frac{m_1m_2s_2}{8m_p^3}\Big)\frac{|b|}{4\pi\gamma}\sqrt{\frac{\pi}{|l''(\Omega_0)|}}\rmint_0^1 dy\,e^{-i\sigma(|b|,\Omega_0,m,y)}\frac{K_1\Big(|b|\sqrt{\Delta(y)^2+(1-y)^2m^2}\Big)}{\sqrt{\Delta(y)^2+(1-y)^2m^2}}\,.
    \end{split}
\end{align}
\textbullet $\,\,$ The massive counterpart of \eqref{5.62x} will take the form.
\begin{align}
    \begin{split}
        f(x)&\propto \rmint_{k}e^{-ik\cdot (x-b_1)}\rmint_{k_2} e^{ik_2\cdot b}\frac{\hat\delta(k_2\cdot v_2)}{k_2^2-m^2}\rmint_{k_3}\hat\delta(k_3\cdot v_2)\hat\delta(k_3\cdot v_1-k\cdot v_1-k_2\cdot v_1)\frac{e^{-ik_3\cdot b}}{(k_3^2-m^2)[(k_3-k)^2-m^2]}\,,\\ &
        =e^{-il(\Omega_0)}e^{i\frac{\pi}{4}}\sqrt{\frac{\pi}{|l''(\Omega_0)|}}g(\Omega_0)
    \end{split}
\end{align}
where,
\begin{align}
    \begin{split}
        g(\Omega_0)=\rmint_{k_2}\frac{e^{i k_2\cdot b}\,\hat\delta(k_2\cdot v_2)}{k_2^2-m^2}I_{k_3}(m,k_2,b,k_0)\,.
    \end{split}
\end{align}
Now, $I_{k_3}(m,k_2,b,k_0)$ takes the form,
\begin{align}
    \begin{split}
        I_{k_3}(m,k_2,b,k_0)&=\frac{|b|}{\gamma}\rmint_0^1dy\,e^{-iyk_0\cdot b}\frac{K_1\Big[|b|\sqrt{\Sigma(k_0,y,k_2\cdot v_1)^2+(1-y+y^2)m^2}\Big]}{4\pi\sqrt{ \Sigma(k_0,y,k_2\cdot v_1)^2+(1-y+y^2)m^2}}\,.\\ &
    \end{split}
\end{align}
Therefore, 
\begin{align}
    \begin{split}
        g(\Omega_0)=\frac{|b|}{\gamma}\rmint_0^1 dy\,e^{-iy k_0\cdot b}\rmint_{-\infty}^{\infty} dz\, K_{0}(\sqrt{z^2+m^2}|b|)\frac{K_1\Big[|b|\sqrt{\Sigma(k_0,y,z)^2+(1-y+y^2)m^2}\Big]}{4\pi\sqrt{ \Sigma(k_0,y,z)^2+(1-y+y^2)m^2}}\,.
    \end{split}
\end{align}
Therefore, the full waveform proportional to,
\begin{align}
    \begin{split}
     f(x)\propto e^{-il(\Omega_0)}e^{i\frac{\pi}{4}}\sqrt{\frac{\pi}{|l''(\Omega_0)|}}\frac{|b|}{\gamma}\rmint dy\,e^{-iy k_0\cdot b}\rmint_{-\infty}^{\infty} dz\, K_{0}(\sqrt{z^2+m^2}|b|)\frac{K_1\Big[|b|\sqrt{\Sigma(k_0,y,z)^2+(1-y+y^2)m^2}\Big]}{4\pi\sqrt{ \Sigma(k_0,y,z)^2+(1-y+y^2)m^2}}.\,
    \end{split}
\end{align}
\textbullet $\,\,$ The massive counter part of \eqref{6.35} takes the form.
\begin{align}
    \begin{split}
        f(x)&\propto e^{-il(\Omega_0)}e^{i\frac{\pi}{4}}\sqrt{\frac{\pi}{|l''(\Omega_0)|}}\rmint_{k_1,k_2}\frac{\hat\delta(k_0\cdot v_1+k_1\cdot v_1)\hat\delta(k_2\cdot v_2)\hat\delta(k_1\cdot v_2)}{(k_1^2-m^2) (k_2^2-m^2) [(k_1+k_2)^2-m^2]}e^{ik_1\cdot b}\,.
    \end{split}
\end{align}
%\subsection*{New integrals:}

%where the integral measure has the following form,
%\begin{align}
%    \begin{split}
%     \rmint   d\mu_{1,2}(k)=\rmint_{k_1,k_2}e^{ik_1\cdot b_1}e^{ik_2 \cdot b_2}\hat\delta(k_1\cdot v_1)\hat\delta(k_2\cdot v_2)\hat\delta^{(4)}(k_1+k_2-k)\,.
 %   \end{split}
%\end{align}
%Hence, the time domain waveform has the following form,
%\begin{align}
%    \begin{split}
 %       f_{\varphi}(x)&=m^2\Big(\frac{m_1m_2s_2}{8m_p^3}\Big)\rmint_{\Omega}e^{-ik\cdot %x}\rmint_{k_1,k_2}e^{ik_1\cdot b_1}e^{ik_2\cdot b_2}\hat\delta(k_1\cdot v_1)v_1^{\alpha}v_1^{\beta}\hat\delta(k_2\cdot v_2)\frac{\eta^{\mu\nu}P_{\mu\nu,\alpha\beta}}{k_1^2(k_2^2-m^2)}\delta^{(4)}(k_1+k_2-k)\,,\\&
      %  =\Big(\frac{m_1m_2s_2}{4m_p^3}\Big)\rmint_{\Omega}e^{-ik\cdot x+ik\cdot b_1}\rmint_{k_2}e^{-ik_2\cdot b_1}e^{ik_2\cdot b_2}\hat\delta((k-k_2)\cdot v_1)v_1^{\alpha}v_1^{\beta}\hat\delta(k_2\cdot v_2)\frac{\eta^{\mu\nu}P_{\mu\nu,\alpha\beta}}{(k_2-k)^2(k_2^2-m^2)}\\ &
  %      =m^2\Big(\frac{m_1m_2s_2}{8m_p^3}\Big)\eta^{\mu\nu}P_{\mu\nu,\alpha\beta}v_1^\alpha v_1^\beta\rmint_{\Omega}e^{-i k\cdot (x-b_1)}\rmint_{k_2}\frac{\hat\delta(k_2\cdot v_1-k\cdot v_1)\hat\delta(k_2\cdot v_2)}{(k_2-k)^2(k_2^2-m^2)}e^{-i k_2\cdot b}\,.
%        \label{5.25mm}
%    \end{split}
%\end{align}
%Now, using the method of stationary phase, we will get,

%\begin{equation}\label{e}\begin{split}
%f_{\varphi}(x)=m^2\Big(\frac{m_1m_2s_2}{8m_p^3}\Big) P^\mu_{\,\,\,\mu,\alpha\beta}v_1^\alpha v_1^\beta\sqrt{\frac{\pi}{|l''(\Omega_0)|}}\frac{b}{4\pi\gamma}\rmint_0^1 dy\,e^{-i\sigma(b,\Omega_0,m,y)}\frac{K_1(b\sqrt{\Delta(k_0,y)^2+(1-y)^2 m^2})}{\sqrt{\Delta(k_0,y)^2+(1-y)^2 m^2}}\,.
%\\
%        \end{split}
%\end{equation}
%\newpage
After integrating over the loop momenta, the full waveform takes the following form,

\begin{align}
    \begin{split}
    f(x)\sim e^{-il(\Omega_0)}e^{i\frac{\pi}{4}}\sqrt{\frac{\pi}{|l''(\Omega_0)|}} \rmint_0^\infty dl\,l\,\frac{J_{0}(l|b|)}{(l^2+\hat{m}^2)\sqrt{l^2+\hat{m}^2-m^2}}\arctan\Big(\frac{l^2+\hat{m}^2-m^2}{2m}\Big).\,
    \end{split}
\end{align}
\textbullet $\,\,$ The massive counterpart of \eqref{5.11 mn}, which is written in \eqref{6.11 mm} can be cast again by the stationary phase approximation and is given by,
\begin{align}
    \begin{split}
        f_{\varphi}(x)&\propto e^{-il(\Omega_0)}e^{i\frac{\pi}{4}}\sqrt{\frac{\pi}{|l''(\Omega_0)|}}\rmint_{k_3,q} \frac{\hat\delta(k\cdot v_1-q\cdot v_1)\hat\delta(k_3\cdot v_2)\hat\delta(q\cdot v_2)}{(k_3^2-m^2) \, [(q-k_3)^2-m^2][(q-k)^2-m^2]}e^{-i q\cdot b}\\ &
        = e^{-il(\Omega_0)}e^{i\frac{\pi}{4}}\sqrt{\frac{\pi}{|l''(\Omega_0)|}}\int_0^\infty d\hat\alpha\,d\hat\beta\,d\hat\gamma\int_{\vec k_3,\vec q}\hat\delta(\Omega_0\,n\cdot v_1+\vec q\cdot \vec v_1)\,\exp\Big[i\hat\alpha(\vec q^2-2\Omega_0 \vec q\cdot \vec n)\\ &
\hspace{0.8cm}+i\hat\beta (\vec q^2-2\vec q\cdot \vec k_3+\vec k_3^2)+{i\hat\gamma \vec k_3^2}+i(\hat\beta+\hat\gamma)m^2\Big]e^{i\vec q\cdot \vec b}\,.
    \end{split}
\end{align}
Now, from the delta function constraint, we can do the integral over $q_{(1)}$, and we are left with,
\begin{align}
    \begin{split}
       f_{\varphi}(x)&\sim e^{-il(\Omega_0)}e^{i\frac{\pi}{4}}\sqrt{\frac{\pi}{|l''(\Omega_0)|}}\frac{1}{\sqrt{\gamma^2-1}}\int_0^\infty d\hat\alpha\,d\hat\beta\,d\hat\gamma \,e^{i(\hat\beta+\hat\gamma) m^2}\exp\Big[i(\hat\alpha+\hat\beta)\Big(\frac{\Omega_0 n\cdot v_1}{\gamma\beta}\Big)^2+2i\hat\alpha\frac{\Omega_0^2n\cdot v_1}{\gamma\beta}\Big]\\ &
       \hspace{0.8cm}\times  e^{-i\frac{\pi}{4}}\pi^{\frac{3}{2}}(\hat\beta+\hat\gamma)^{-\frac{3}{2}}\int_{\tilde q}\exp\Big(-i \frac{\hat\beta^2 \vec q^2}{\hat\beta +\hat\gamma}\Big)\exp\Big(i(\hat\alpha +\hat\beta)\tilde q^2+i\tilde q\cdot b\Big)\\ &
       = \pi^{\frac{5}{2}}e^{-il(\Omega_0)}\sqrt{\frac{\pi}{|l''(\Omega_0)|}}\frac{1}{\sqrt{\gamma^2-1}}\int_0^\infty d\hat\alpha\,d\hat\beta\,d\hat\gamma \,e^{i(\hat\beta+\hat\gamma) m^2}\exp\Big[i(\hat\alpha+\hat\beta)\Big(\frac{\Omega_0 n\cdot v_1}{\gamma\beta}\Big)^2+2i\hat\alpha\frac{\Omega_0^2n\cdot v_1}{\gamma\beta}\\ &
       \hspace{0.8cm}-i\frac{\hat\beta^2}{\hat\beta+\hat\gamma}\frac{\Omega_0^2(n\cdot v_1)^2}{\gamma^2(\gamma^2-1)}\Big]  (\hat\beta+\hat\gamma)^{-\frac{3}{2}}\frac{e^{\frac{-ib^2}{4\lambda_1}}}{\lambda_1},\,\textrm{with},\,\lambda_1\equiv \hat\alpha+\hat\beta-\frac{\hat\beta^2}{\hat\beta+\hat\gamma}\,.\label{6.36f}
    \end{split}
\end{align}
%The integral in \eqref{6.36f} has the following form,
%\begin{align}
   % \begin{split}
    %    I\equiv \int_0^\infty d\hat\alpha\,d\hat\beta\,d\hat\gamma \,e^{i(\hat\beta+\hat\gamma) m^2}
  %  \end{split}
%\end{align}
The integral in \eqref{6.36f}, to the best of our knowledge, does not have any closed form and should be done numerically. One can analogously find the contribution from $\lambda_3 h\varphi^3$ interaction vertex.

\subsection*{Radiation integrals for gravitational waveform due to massive scalar}
Finally, we will also discuss diagrams which contribute to the gravitational waveform due to the presence of a massive scalar.

\textbullet $\,\,$ The correction to the gravitational waveform comes from bulk scalar graviton interaction with the interacting action:
\begin{align}
    S_{int}=-\frac{m^2}{2m_p}\rmint d^4x \,h\varphi^2\,.
\end{align}
Note that this diagram does not have any massless counterpart, since the vertex is proportional to the scalar mass. Its contribution to the graviton one-point function takes the following form:
\begin{align}
    \begin{split}
        k^2\Big\langle h_{\mu\nu}(k)\Big\rangle\Big|_{k^2\rightarrow 0}= \thesisinlinefeynman[\chthreeinlinefeynwidth]{feynman-12.png}\hspace{1.4 cm}=\rmint d\mu_{1,2}(k) \frac{\eta^{\mu\nu}} {(k_1^2-m^2) (k_2^2-m^2)}\,,
    \end{split}
\end{align}
where the integral measure takes the form,
\begin{align}
    \begin{split}
        \rmint d\mu_{1,2}(k)= \rmint_{k_1,k_2}e^{ik_1\cdot b_1} e^{i k_2\cdot b_2}\hat\delta(k_1\cdot v_1)\hat\delta(k_2\cdot v_2)\hat\delta^{(4)}(k_1+k_2-k)\,.
    \end{split}
\end{align}

Therefore, the correction to the time domain waveform for this particular interaction is given by,
\begin{align}
    \begin{split}
        f_{h}&=\frac{1}{4\pi m_p}\epsilon^\mu \epsilon^\nu \rmint_{\Omega}e^{-ik\cdot x}\rmint d\mu_{1,2}(k)\,\frac{\eta_{\mu\nu}}{(k_1^2-m^2)\,(k_2^2-m^2)}\,.\\ &
       \label{6.38a}
    \end{split}
\end{align}
But since $\epsilon$ is null, i.e. $\epsilon^2=0$, the index structure of the above integral ensures that its contribution to the waveform vanishes.

\textbullet $\,\,$ Next, we focus on the contribution to the gravitational waveform from a massive scalar field through derivative interaction. The massless counterpart of it is shown in \eqref{5.63a}. The time domain waveform has the following form,
\begin{align}
    \begin{split}
          f_{h}(x)&
        \propto\frac{1}{4\pi m_p}\epsilon^\mu \epsilon^\nu \rmint_{\Omega}e^{-ik\cdot x}\rmint_{k_2}e^{i(k-k_2)\cdot b_1}e^{ik_2\cdot b_2}\frac{(k-k_2)_{\mu}k_{2\nu}}{(k_2^2-m^2)[(k-k_2)^2-m^2]}\,\hat\delta\Big(k\cdot v_1-k_2\cdot v_1\Big)\hat\delta(k_2\cdot v_2)\,,\\ &
=-\frac{1}{4\pi m_p}\epsilon^\mu \epsilon^\nu \rmint_{\Omega}e^{-ik\cdot x}\rmint_{k_2}e^{i(k-k_2)\cdot b_1}e^{ik_2\cdot b_2}\frac{k_{2\mu}k_{2\nu}}{(k_2^2-m^2)[(k-k_2)^2-m^2]}\,\hat\delta\Big(k\cdot v_1-k_2\cdot v_1\Big)\hat\delta(k_2\cdot v_2)\,,\\ &
=-\frac{1}{4\pi m_p}\epsilon\cdot \hat J_{(2)}\cdot \epsilon\,.
    \end{split}
\end{align}
where,
\begin{align}
    \begin{split}
        \hat J_{\mu\nu}^{(2)}=J_{\mu\nu}^{(2)}[(1-y)^2\to (1-y+y^2)].
    \end{split}
\end{align}
where $J_{\mu\nu}^{(2)}$ is defined in (\ref{A.22 m}).

\textbullet $\,\,$
The other diagram corresponding to the graviton radiation from worldline-1 with bulk $\lambda \varphi^3$ vertex is given by,

\begin{align}
    \begin{split}
        k^2\Big\langle h_{\mu\nu}(k)\Big\rangle\Big|_{k^2\rightarrow 0}= \thesisinlinefeynman[\chthreeinlinefeynwidth]{feynman-13.png}\hspace{1.4 cm}=\rmint d\mu_{1,2}(k)\frac{1}{(k_1^2-m^2)(k_2^2-m^2)(k_3^2-m^2)}\,,
    \end{split}
\end{align}
where the measure looks like,
\begin{align}
    \begin{split}
        d\mu_{1,2}(k)=\rmint_{k_i}\hat\delta^{(4)}\Big(\sum_{i}k_i\Big)\hat\delta(k_1\cdot v_1+k\cdot v_1)\hat\delta(k_2\cdot v_2)\hat\delta(k_3\cdot v_2)e^{i(k+k_1)\cdot b_1}e^{i(k_2+k_3)\cdot b_2}\,.
    \end{split}
\end{align}
Therefore, omitting the prefactors, the corresponding waveform takes the form,
\begin{align}
    \begin{split}
        f(x)&\sim \rmint_{\Omega}e^{-ik\cdot x}\rmint_{k_i}\frac{\hat\delta^{(4)}\Big(\sum_{i}k_i\Big)\hat\delta(k_1\cdot v_1+k\cdot v_1)\hat\delta(k_2\cdot v_2)\hat\delta(k_3\cdot v_2)}{(k_1^2-m^2)(k_2^2-m^2)(k_3^2-m^2)}e^{i(k+k_1)\cdot b_1}e^{i(k_2+k_3)\cdot b_2}\,,\\ &
        =\rmint_{\Omega}e^{-i\Omega n\cdot (x-b_1)}\rmint_{k_1,k_2}\frac{\hat\delta(k_1\cdot v_1+k\cdot v_1)\hat\delta(k_2\cdot v_2)\hat\delta(k_1\cdot v_2)}{(k_1^2-m^2)(k_2^2-m^2)[(k_1+k_2)^2-m^2]}e^{ik_1\cdot b}\,,\\ &
      %  =\frac{1}{(n\cdot v_1)}\rmint_{k_1,k_2}e^{i k_1\cdot \delta_1}\frac{\hat\delta(k_1\cdot v_2)\hat\delta(k_2\cdot v_2)}{(k_1^2-m^2)(k_2^2-m^2)[(k_1+k_2)^2-m^2]}\\ &
    %    =-\frac{(2\pi)^2}{(n\cdot v_1)}\rmint_{\vec k_1}e^{-i \vec k_1\cdot \vec \delta_1}\frac{1}{(\vec k_1^2+m^2)(|\vec k_1|)}\arctan\Big(\frac{|\vec k_1|}{2m}\Big)\\ &
        =-\frac{4\pi(2\pi)^2}{(n\cdot v_1)|\vec \delta_1|}\rmint_{0}^{\infty} dk_1 \frac{\sin{(k_1|\vec \delta_1|)}}{(k_1^2+m^2)}\arctan\Big(\frac{|\vec k_1|}{2m}\Big)\,.
    \end{split}
\end{align}
Now, restoring the prefactors, the waveform can be recast as,
\begin{align}
 &   \begin{split}
    f(x)=-\frac{\lambda_3 }{2(2\pi)^3}\Big(\frac{m_1 s_1m_2^2 s_2^2}{4m_p^3}\Big)\frac{1}{(n\cdot v_1)|\vec \delta_1|}\rmint_{0}^{\infty} dy \frac{\sin{(y|\vec \delta_1|)}}{(y^2+m^2)}\arctan\Big(\frac{y}{2m}\Big)\,.    
    \end{split}
\end{align}

\section{Discussions and outlook}\label{sec8}
Inspired by the recent developments in scattering amplitude techniques in QFT, we consider a scalar-tensor theory of gravity. After briefly discussing the novel WQFT formalism \cite{Mogull:2020sak}, we compute the two main observables in a scattering event for scalar-tensor theory. We compute the impulse and waveform coming from the scalar contribution to the gravitational field by computing the one-point correlators for the fields (scalar and graviton) and worldline degrees of freedom. \textit{To the best of our knowledge, this is the first study regarding applying WQFT to compute impulse and waveform for a non-GR theory, namely the Scalar-Tensor theory. Also, note that this study helps us extend the computation of the gravitational waveform in the post-Minkwoskian regime for a theory beyond GR.} Below, we summarize the main findings of our paper.
\begin{itemize}
    \item First we compute the impulse $\Delta p^\mu$ in the massive scalar tensor theory with scalar potential $V(\varphi)\sim \lambda_3 \varphi^3+\lambda_4 \varphi^4$ up to 2PM. We mainly concentrate on the corrections to the impulse coming from the scalar degree of freedom. Also, we have provided the corresponding expressions when we take the massless limit. \textcolor{black}{Note that to the best of our knowledge some of the massive integrals as listed in (\ref{4.86}) do not possess any closed-form expression. They can be evaluated numerically and can be shown to have a smooth behaviour w.r.t. $|b|.$ Also, they possess smooth massless limits as discussed in the main text. Hence, these massless expressions serve as a consistency check of our results.}  The fall-off in that case w.r.t. $|b|$ is analogous to the gravitational case. Nevertheless, even in the one-loop case, where a fully closed-form expression is not available, the result can still be represented in terms of a $0$--$1$ parameter integral with a Bessel kernel:
\begin{align}
    I_{\text{1-loop}}
    \sim
    \int_0^1 ds\, \rho(s)\, K_0(\cdots).\nonumber
\end{align}
Moreover, we have verified explicitly in a nontrivial two-loop example that the impulse may be expressed as
\begin{align}
    I_{\text{2-loop}}
    \sim
    \int_0^1 ds \int_0^1 du\,
    \rho_1(u,s)\,\rho_2(s,u)\,
    K_0^2(\cdots)+\cdots.\nonumber
\end{align}
While establishing this form in general is technically challenging, these examples suggest that the leading $n$-loop contribution may admit a representation of the schematic form
\begin{align}
    I_{\text{$n$-loop}}
    \sim
    \prod_{i=1}^{n}\int_0^1 ds_i\;
    \rho_i(\{s_j\})\, K_0^{\,n}(\cdots)+\cdots\nonumber
\end{align}
It would be very interesting to clarify whether this pattern persists to all loop orders, to understand the underlying algebraic--geometric structure of these Bessel curve, and ultimately to identify the natural function space to which these observables belong \cite{Zhou:2018tva}.

\item Next, we compute the radiation integrals coming from the field one-point function $\langle\chi(k)\rangle$ and time domain waveform in the massless case for the scalar waveform as well as the correction due to the gravitational waveform due to the presence of bulk scalar interaction vertices. Note that the total waveform will also have contributions from the GR term, which is already known in the literature. As discussed in the introduction (\ref{intro}), the simplest way to add extra degrees of freedom is to introduce a scalar field. Further motivation was provided for considering the Scalar-Tensor theory there. Keeping this in mind, we have focused only on the correction due to the presence of extra scalar degrees of freedom. To the best of our knowledge, this is the first study of PM waveform for a Scalar-Tensor theory. 
\item Eventually, we proceed to compute the waveform where the scalar field has non-zero mass. In this case, we find that the analytical (exact) computation of the waveform becomes significantly difficult due to the presence of a complicated phase structure. \textit{We propose a procedure to handle those integrals using the stationary phase approximation.} As we expect a burst kind of signal from scattering events, this is a reasonable approximation one can make. 
%\item Interestingly, \textit{we find that $\lambda_4 \varphi^4 $ interaction vertex does not contribute to the waveform by applying the \textit{method of regions} to the relevant Feynman integral.} We also provide a more explicit illustration of our claim in Appendix~(\ref{App2}).
    \item Apart from that, we would like to emphasize that we show different approaches to computing different kinds of Feynman integrals and discuss the underlying subtleties. Some of these integrals do not appear in the corresponding GR computation. We hope that this will help while exploring WQFT methods for other non-GR theories of gravity. 
\end{itemize}
The connection to the scattering amplitude is trivial in this process. One need not compute the effective action by integrating out the different energy modes, similar to the effective field theory techniques \cite{Bhattacharyya:2023kbh}. To this end, one eventually encounters the loop integrals that one will encounter in scattering processes with intermediate loops.
To the best of our knowledge, we tried to compute diagrams of radiation and impulse with massive propagators in several places, which is new to the literature of WQFT for scalar-tensor theories. It will also be quite an interesting follow-up to incorporate the spin of each black hole by considering finite-size effects in the worldline action. In fact, WQFT of $\mathcal{N}=1 $ \textit{supersymmetric spinning particles} exist in literature \cite{Jakobsen:2021lvp}. Eventually, as a follow-up, one can also try to calculate 3PM three-body radiation. As an advantage of this procedure, one can do this without taking the classical limit ($\hbar\rightarrow 0$) explicitly in each diagram computation. Last but not least, as the connection to the scattering amplitude is apparent, it is important to investigate the Double-copy setup to make it clear in this formalism. We hope to report on some of these issues in the near future. 

%\newpage


%\newpage
\appendix 
\section{Sketching the derivation of the worldline action}\label{ch2:app:A}
In our case, the matter Lagrangian has the following form:
\begin{align}
    \begin{split}
        \mathcal{L}=g^{\mu\nu}\partial_\mu\phi_i^{\dagger}\partial_\nu\phi_i-m_i(\varphi)^2\phi_i^{\dagger}\phi_i
    \end{split}
\,.\end{align}
We start by representing the partition function (in Euclidean signature)  using the Schwinger proper time parametrization.
\begin{align}
    \Gamma[g,\varphi]=\log\Big[\rmint \mathcal{D}[\phi,\phi^{\dagger}]\,e^{-S}\Big] &=-\log\Big[\det(\nabla_{\mu}\nabla^\mu+m(\varphi)^2)\Big]\,,\\ &
    =-\textrm{Tr}\log\Big[\nabla_{\mu}\nabla^\mu+m(\varphi)^2\Big]\,,\\ &
    =\rmint_0^\infty \frac{dT}{T}\rmint \frac{d^4 k}{(2\pi)^4}\exp\Big[-\frac{1}{2}eT\Big(g_{\mu\nu}k^\mu k^\nu+m(\varphi)^2\Big)\Big]
\end{align}
where, $e$ is the einbein. Now we convert the result into the path integral over $x(\tau)$.
\begin{align}
    \begin{split}
        \Gamma[g,\phi]=\rmint_0^\infty \frac{dT}{T}\mathcal{N}(T)\rmint \mathcal{D}[x]\exp\Big[-\rmint d\tau (\frac{1}{2e}g_{\mu\nu}\dot x^\mu\dot x^\nu+\frac{e}{2}m(\varphi)^2)\Big]\,.\label{D.5}
    \end{split}
\end{align}
The result in \eqref{D.5} is a one-dimensional field theory of $x^\mu(\tau)$. For the  case of gravitational field, the integral measure becomes metric dependent which causes the existence of Lee-Yang ghost. However in classical limit the ghost fields are irrelevant and can be ignored as shown in \cite{Mogull:2020sak}.
\section{Impulse via Eikonal: a connection to scattering amplitude}\label{ch2:app:B}
In section~(\ref{sec1}) we computed the impulse up to 2PM. We now discuss its connection with the scattering amplitude. It was shown in \cite{Amati:1987wq, Amati:1990xe} that the eikonal phase $\chi$ is related to the four-point scattering amplitude as,
\begin{align}
    \begin{split}
e^{i\chi}&\equiv 1+\frac{i}{4m_1 m_2}\rmint e^{iq\cdot b}\hat\delta(q\cdot v_1)\hat\delta(q\cdot v_2)\lim_{\hbar\rightarrow 0}\mathcal{M}_4(\phi_1,\phi_2\rightarrow \phi_1,\phi_2)\,,\\ &
=1+\frac{i}{4m_1 m_2}\rmint_{\boldsymbol{q}_{\perp}}e^{i \boldsymbol{q}_{\perp}\cdot b}\lim_{\hbar\rightarrow 0}\mathcal{M}_4(\phi_1,\phi_2\rightarrow \phi_1,\phi_2)\,.
    \end{split}
\end{align}
It has been demonstrated in \cite{Bern:2020buy,Bjerrum-Bohr:2018xdl} that, in the centre-of-mass frame, the eikonal phase is related to the impulse through 2PM order and takes the following form,
\begin{align}
    \begin{split}
        \Delta p_{\perp}= \frac{\partial \chi}{\partial b}\,.
    \end{split}
\end{align}
Later, it has been shown that the result can be extended to higher PM order \cite{Maybee:2019jus,Bern:2020buy}. In WQFT, one can identify the classical part of $\chi$ to be the free energy of the WQFT at tree level and hence given by,
\begin{align}
    \begin{split}
    e^{i\chi(\hat b_i,\hat v_i)}= Z_{\textrm{WQFT}}:=\mathcal{N}\rmint \mathcal{D}h_{\mu\nu}\mathcal{D}\varphi\,\rmint\prod_{k=1}^2\mathcal{D}z_{k}\exp\Big(iS_{g}+iS^k_{pm}\Big)
    \end{split}
\end{align}
where, $\hat b$ and $\hat v$ can be related to the averaged incoming momenta $\hat p_i$ which satisfies, $\hat p_1\Delta p_1=0$. Therefore, the formula for impulse is given by,
\begin{align}
    \begin{split}
        \Delta p_{1\mu}=-\frac{\partial \chi}{\partial \hat b_1^\mu}.
    \end{split}
\end{align}
If one can compute the impulse using the Eikonal method, one may lose some extra terms in the impulse, which are proportional to $v_i^\mu$. If one blindly computes the impulse, for example, for the following diagram, 
\begin{align}
    \begin{split}
        \chi&\propto \thesisinlinefeynman[\chthreeinlinefeynwidth]{e_1.png} \\ &
= \rmint  \check{\mathcal{D}}[h_{\mu\nu}, \varphi,\{z_{i}\}]\rmint_{\{k_i,\omega_j\}}\hat{\delta}(k_1\cdot v_2)\hat{\delta}(k_2\cdot v_2) \varphi(k_1) \varphi(k_2) e^{i(k_1+k_2)\cdot b_2}\\ &
        \hspace{0.8cm}\hat\delta(k_3\cdot v_1+\omega_1)\hat\delta(k_4\cdot v_1+\omega_2)\varphi(k_3)\varphi(k_4)z^{\rho}(\omega_1)z^{\sigma}(\omega_2)\{2\omega_1 (v_1)_{\rho}+(k_{3})_{\rho}\}\\ &
        \hspace{0.8cm}
        \{2\omega_2(v_1)_{\sigma}+(k_4)_{\sigma}\}\hat\delta^{(4)}(k_2+k_4)\hat\delta^{(4)}(k_1+k_3)\delta(\omega_1+\omega_2)\,,\\&
        %=\rmint_{k_i,\omega_j}\hat{\delta}(k_1\cdot v_2)\hat{\delta}(k_2\cdot v_2)\hat\delta(k_3\cdot v_1+\omega_1)\hat\delta(k_4\cdot v_1+\omega_2)\langle \tilde \varphi(k_1)\tilde\varphi(k_3)\rangle\langle \tilde \varphi(k_2)\tilde\varphi(k_4)\rangle\\ &
      %  \hspace{0.8cm}
      %  \langle z^{\rho}(\omega_1)z^{\sigma}(\omega_2) \rangle \{2\omega_1 (v_1)_{\rho}+(k_{3})_{\rho}\}  \{2\omega_2(v_1)_{\sigma}+(k_4)_{\sigma}\}\\ &
        =\rmint_{\{k_i,\omega_j\}}\hat{\delta}(k_1\cdot v_2)\hat{\delta}(k_2\cdot v_2)\hat\delta(k_3\cdot v_1+\omega_1)\hat\delta(k_4\cdot v_1+\omega_2)\hat\delta^{(4)}(k_1+k_3)\hat\delta^{(4)}(k_2+k_4)\delta(\omega_1+\omega_2)\\ &
       \hspace{2 cm} \frac{1}{(k_1^2-m^2)(k_2^2-m^2)\omega_1^2}\{4\omega_1\omega_2+2\omega_1 v_1\cdot k_4+2 \omega_2 v_1\cdot k_3+k_3\cdot k_4\}e^{i(k_1+k_2)\cdot b_2}e^{i(k_3+k_4)\cdot b_1}\,,\\ &
=\rmint_{k,k_1,\omega_1}\frac{\hat\delta(k_1\cdot v_2)\hat\delta(k\cdot v_2)\hat\delta(k\cdot v_1)\hat\delta(\omega_1-k_1\cdot v_1)}{(k_1^2-m^2)[(k-k_1)^2-m^2]\omega_1^2}e^{ik \cdot b}\\&\hspace{0.8cm}[-4\omega_1^2+2\omega_1v_1\cdot(k_1-k)+2\omega_1v_1\cdot k_1+k_1\cdot (k-k_1)]\,.
    \end{split}
\end{align}
Integrating over $\delta(\omega_1-k_1\cdot v_1)$ and using the fact that the integral has support at $k\cdot v_1=0$, we left with the following integral,
%\begin{align}
   % \begin{split}
        %\chi\sim\rmint_{k,k_1}\frac{\hat\delta(k_1\cdot v_2)\hat\delta(k\cdot v_2)\hat\delta(k\cdot v_1)}{(k_1^2-m^2)[(k-k_1)^2-m^2](k_1\cdot v_1)^2}[k_1\cdot(k-k_1)]e^{ik\cdot b}\,.\label{7.6mm}
    %\end{split}
%\end{align}
%\begin{figure}[htb!]
  %  \centering
  %  \includegraphics[scale=0.7]{con.jpg}
  %  \caption{Contour relevant for evaluating \eqref{7.6mm}.}
  %  \label{fig:enter-label}
%\end{figure}
{%The integral is readily zero if we \textit{ignore the radiation poles} i.e. if we work in the limit $k_{1(0)}\ll |\vec k_1|$ \cite{Kalin:2020mvi}. To see this more clearly one has to do the integral from the frame of first particle,
%\begin{align}
  %  \begin{split}
     %   \chi\sim \frac{1}{2\beta}\rmint_{k} e^{-i\vec k\cdot \vec b}(-\vec k^2-2m^2)\rmint_{\vec k_1}\frac{\hat\delta(k_1^{(1)})}{(\vec k_1^2+m^2)[(\vec k-\vec k_1)^2+m^2]}\rmint dk_{1(0)}\frac{1}{(k_1^{(0)}+i\epsilon)^2}\,.
   % \end{split}
%\end{align}
}
%To evaluate this integral, we will make use of the following contour integral,
%\begin{align}
    %\begin{split}
      %  \rmint d k_1^{(0)}[\cdots]=\oint_{H^{+}}[\cdots]-\oint_{H^{-}}[\cdots]=\frac{i\pi}{2}\sum_{ k_1^{(0)}= k_1^{(0)*}\subset H^{+}}\text{Res}([\cdots]) -\sum_{ k_1^{(0)}=\bar k_1^{(0)*}\subset H^{-}}\text{Res}([\cdots])\,.
  %  \end{split}
%\end{align}
%\textcolor{black}{
%One can see that the integrand has double pole at $k_1\cdot v_1+i\epsilon=0$ which lies on the same half plane of the contour. So, we simply close the contour in the upper-half and we get zero.
% We see that the impulse corresponding to this diagram has finite non-zero value from \eqref{4.6 mm} which is completely proportional to $v_i^\mu$. Although the integral can be done in a more general setting, the procedure for recovering the result using Eikonal phase has been discussed in \cite{Kalin:2020fhe}. For completeness we show the exact computation in the next subsection.}
%\subsection*{Non-vanishing of Eikonal phase (taking into account the radiation pole):}
The integral can be done 
\begin{align}
    \begin{split}
\chi&\sim\rmint_{k,k_1}\frac{\hat\delta(k_1\cdot v_2)\hat\delta(k\cdot v_2)\hat\delta(k\cdot v_1)}{(k_1^2-m^2)[(k-k_1)^2-m^2](k_1\cdot v_1)^2}[k_1\cdot(k-k_1)]e^{ik\cdot b}\,,\\ &
        =\frac{1}{2}\rmint_{k} e^{ik\cdot b}{\hat\delta(k\cdot v_1)\hat\delta(k\cdot v_2)}(k^2-2m^2)\underbrace{\rmint_{k_1}\frac{\hat\delta (k_1\cdot v_2)}{(k_1^2-m^2)[(k-k_1)^2-m^2](k_1\cdot v_1)^2}}_{\hat\chi_{k}}\,.\label{B.8a}
    \end{split}
\end{align}
One can see that the delta function inside the $k_1$ integral reduces the four dimensional integral into a three dimensional integral which reads,
\begin{align}
    \begin{split}
        \hat\chi_{k}\sim 4\rmint_{\vec k_1} \frac{1}{(\vec k_1^2+m^2)[(\vec k-\vec k_1)^2+m^2](2\,\vec k_1\cdot \vec v_1)^2}\,.\label{B.9c}
    \end{split}
\end{align}
The integral can be done by introducing the alpha parametrization \cite{smirnov2}, 
\begin{align}
    \frac{1}{(-A)^\lambda}=\frac{i^\lambda}{\Gamma(\lambda)}\rmint_{0}^{\infty}d\alpha \, \alpha^{\lambda-1}e^{iA\alpha}
\end{align}
which gives,
\begin{align}
    \begin{split}
      (2\pi)^3  \hat\chi_{k}&=\frac{4}{\Gamma(1)^2\Gamma(2)}\rmint d^{3-2\epsilon}k_1\rmint_0^\infty  \prod_{i=1}^3 d\alpha_i\,\alpha_3 \exp\Big[i( \vec k_1^2+m^2)\alpha_1+i[(\vec k_1-\vec k)^2+m^2]\alpha_2 +i\,2(\vec k_1\cdot \vec v_1)\alpha_3\Big]\,,\\ &
        =\frac{4}{\Gamma(1)^2\Gamma(2)}\rmint_0^\infty \prod_{i=1}^3 d\alpha_i\,\alpha_3\,e^{i(\alpha_1+\alpha_2)m^2}e^{i\alpha_2\vec k^2}\rmint d^{3-2\epsilon}k_1\exp\Big[i\vec k_1^2(\alpha_1+\alpha_2)-2\vec k_1\cdot (\alpha_2 \vec k-\alpha_3 \vec v_1)\Big]\,,\\ &
      %  =\frac{4}{\Gamma(1)^2\Gamma(2)}\rmint_0^\infty \prod_{i=1}^3 d\alpha_i\,\alpha_3\,e^{i(\alpha_1+\alpha_2)m^2}e^{i\alpha_2\vec k^2}e^{i\frac{\pi}{2}(\epsilon-\frac{1}{2})}\pi^{\frac{3}{2}-\epsilon}(\alpha_1+\alpha_2)^{\epsilon-\frac{3}{2}}\exp\Big(-i\frac{\alpha_2^2 \vec k^2+\alpha_3^2\vec v_1^2}{\alpha_1+\alpha_2}\Big)\,,\\ &
        =\frac{4}{\Gamma(1)^2\Gamma(2)}e^{i\frac{\pi}{2}(\epsilon-\frac{1}{2})}\pi^{\frac{3}{2}-\epsilon}\rmint_0^\infty  \prod_{i=1}^3 d\alpha_i\,\alpha_3\,e^{i(\alpha_1+\alpha_2)m^2}(\alpha_1+\alpha_2)^{\epsilon-\frac{3}{2}}\exp\Big(i\frac{\alpha_1\alpha_2\vec k^2-\alpha_3^2\vec v_1^2}{\alpha_1+\alpha_2}\Big)\,.
    \end{split}
\end{align}
Now doing the integral over $\alpha_3$ we get,
\begin{align}
    \begin{split}
      (2\pi)^3  \hat\chi_{k}&= \frac{4}{\Gamma(2)\Gamma(1)^2}e^{i\frac{\pi}{2}(\epsilon-\frac{1}{2})}\pi^{\frac{3}{2}-\epsilon}\rmint_0^\infty  \prod_{i=1}^2 d\alpha_i \,(\alpha_1+\alpha_2)^{\epsilon-\frac{3}{2}}\,e^{i(\alpha_1+\alpha_2)m^2}\exp\Big(i\frac{\vec{k}^2\alpha_1\alpha_2-\alpha_3^2 \gamma^2\beta^2}{\alpha_1+\alpha_2}\Big)\,,\\ &
        =-\frac{4i}{\gamma^2\beta^2}e^{i\frac{\pi}{2}(\epsilon-\frac{1}{2})}\pi^{\frac{3}{2}-\epsilon}\rmint_0^\infty  \prod_{i=1}^2 d\alpha_i (\alpha_1+\alpha_2)^{\epsilon-\frac{1}{2}}\underbrace{e^{i(\alpha_1+\alpha_2)m^2}}_{\textrm{extra term coming due to non zero mass}}\exp\Big(i\frac{\vec{k}^2\alpha_1\alpha_2}{\alpha_1+\alpha_2}\Big)\,.
    \end{split}
\end{align}
Now doing the following variable change,
$$\alpha_1+\alpha_2\to \eta\,\,\,\textrm{and}\,\,\, \frac{\alpha_1}{\alpha_1+\alpha_2}\to \xi$$ we get,
\begin{align}
    \begin{split}
     (2\pi)^3   \hat\chi_{k}=&-\frac{4i}{(\gamma^2-1)} e^{i\frac{\pi}{2}(\epsilon-\frac{1}{2})}\pi^{3/2-\epsilon}\rmint_0^1 d\xi \rmint_0^\infty d\eta \,\eta^{\epsilon+\frac{1}{2}}\exp\Big[i\eta\{m^2+\vec{k}^2\xi(1-\xi)\}\Big]\,,\\ &
        =\frac{4}{(\gamma^2-1)}\lim_{\epsilon\to 0}e^{i\epsilon\pi}\pi^{5/2-\epsilon}\frac{    \sec (\pi  \epsilon ) \left(B_{z_1}\left(-\epsilon -\frac{1}{2},-\epsilon -\frac{1}{2}\right)-B_{z_2}\left(-\epsilon -\frac{1}{2},-\epsilon -\frac{1}{2}\right)\right)}{k\,\left(k^2+4 m^2\right)^{\epsilon +1}  \Gamma \left(-\epsilon -\frac{1}{2}\right)}\label{B.12}
    \end{split}
\end{align}
where,
\begin{align}
   z_1= \frac{1}{2}-\frac{|\vec{k}|}{2 \sqrt{\vec{k}^2+4 m^2}} \,\textrm{and,}\,\, z_2=\frac{1}{2}+\frac{|\vec{k}|}{2 \sqrt{\vec{k}^2+4 m^2}}\,.
\end{align}
The beta functions can be written in a simplified manner to obtain,
\begin{align}
\begin{split}
    & B_{z_1}(-\epsilon-\frac{1}{2} ,-\epsilon-\frac{1}{2} )-B_{z_2}(-\epsilon-\frac{1}{2} ,-\epsilon-\frac{1}{2} )\\ &=-\rmint_{z_1}^{z_2}dt \,t^{-\epsilon-\frac{3}{2}}(1-t)^{-\epsilon-\frac{3}{2}}\,,\\ &
      =\frac{1}{\epsilon +\frac{1}{2}}\left(\frac{2m^2}{\vec k^2+4m^2}\right)^{-\epsilon -\frac{1}{2}} \Big[\left(\frac{|\vec k|}{\sqrt{\vec k^2+4 m^2}}+1\right)^{\epsilon +\frac{1}{2}} \,\, _2F_1\left(-\epsilon -\frac{1}{2},\epsilon +\frac{3}{2};\frac{1}{2}-\epsilon ;\frac{1}{2}-\frac{|\vec k|}{2 \sqrt{ \vec k^2+4 m^2}}\right)\\ &-\left(1-\frac{|\vec k|}{\sqrt{\vec k^2+4 m^2}}\right)^{\epsilon +\frac{1}{2}} \, _2F_1\left(-\epsilon -\frac{1}{2},\epsilon +\frac{3}{2};\frac{1}{2}-\epsilon ;\frac{|\vec k|}{2 \sqrt{\vec k^2+4 m^2}}+\frac{1}{2}\right)\Big].
     \end{split}
     \end{align}
     Expanding around $\epsilon=0$ we get,
     \begin{align}
         \begin{split}
             (2\pi)^3\hat\chi_{k}&=\frac{8\,\pi^{5/2}}{(\gamma^2-1)}\frac{1}{m\,(\vec k^2+4m^2)}.\label{B.16ab}
         \end{split}
     \end{align}
     The integral $\hat\chi_k$ defined in \eqref{B.8a} can be done using Integral-by-Parts (IBP) reduction in a more sophisticated way. We have,
     \begin{align}
      \chi_k=   \int_{k_1}\frac{\hat\delta (k_1\cdot v_2)}{(k_1^2-m^2)[(k-k_1)^2-m^2](k_1\cdot v_1)^2}
     \end{align}
By unitarity cut the integral can be re written as,
\begin{align}
    \begin{split}
        \chi_k\sim \int_{k_1} \frac{1}{(k_1^2-m^2)[(k-k_1)^2-m^2](k_1\cdot v_1)^2}\Bigg[\frac{1}{k_1\cdot v_2+i\sigma}-\frac{1}{k_1\cdot v_2-i\sigma}\Bigg]\label{B.16a}
    \end{split}
    \end{align}
  The integral in \eqref{B.16a} belongs to the following family of integral,
  \begin{align}
      \begin{split}
          \mathcal{G}^{(\pm)}_{n_1,n_2,n_3,n_4}=\int_{k_1}\frac{1}{D_1^{n_1}D_2^{(\pm)n_2}D_3^{n_3}D_4^{n_4}}
      \end{split}
  \end{align}
  where,
  \begin{align}
      D_1=k_1\cdot v_1,\,D_2^{(\pm)}=k_1\cdot v_2\pm i\sigma,\,D_3=k_1^2-m^2, \,\textrm{and},\,D_4=(k_1-k)^2-m^2\,.
  \end{align}
  Hence it is evident that  the integral $\chi_k= \mathcal{G}_{2,1,1,1}$ can be evaluated  in terms of master integrals using \textbf{LiteRed} \cite{Lee:2012cn} as follows,
  \begin{align}
  \begin{split}
    \chi_k&\sim  \mathcal{G}^{(+)}_{2,1,1,1}- \mathcal{G}^{(-)}_{2,1,1,1}\,,\\ &
    =-\frac{2}{m^2(\gamma^2-1)(4m^2-k^2)}\Big[\mathcal{G}_{0,1,0,1}^{(+),\textrm{master}}-\mathcal{G}_{0,1,0,1}^{(-),\textrm{master}}\Big]+\mathcal{O}(\sigma)\,,\\ &
    =-\frac{2}{m^2(\gamma^2-1)(4m^2-k^2)}\int_{k_1}\frac{\hat\delta(k_1\cdot v_2)}{(k_1-k)^2-m^2}\,\,(\textrm{via reverse unitarity})\,,\\ &
    \propto \frac{2}{m(\gamma^2-1)(4m^2+\vec k^2)}
    \end{split}
  \end{align}
  which exactly matches with the result given in \eqref{B.16ab}. In the last line use use the fact that $\hat\delta(k\cdot v_i)=0\,,\forall i=1,2$. Finally, the Eikonal phase is given by,
     \begin{align}
         \begin{split}
             \chi&\sim\frac{1}{2} \rmint_0^\infty dk\,k\, J_{0}(k|b|)(k^2+2m^2)\chi_k(m)\,,\\ &
             =\frac{1}{m\sqrt{\pi}(\gamma^2-1)}\rmint_0^\infty dk\, k\,J_{0}(k|b|)\frac{(k^2+2m^2)}{k^2+4m^2}\,,\\ &
             =\frac{1}{m\sqrt{\pi}(\gamma^2-1)}\rmint_0^\infty dk\, k\,J_{0}(k|b|)\Big[1-\frac{2m^2}{k^2+4m^2}\Big]\,.\label{B.15}
         \end{split}
     \end{align}
     The first term in \eqref{B.15} does not have any non-zero finite value (purely UV divergent) and hence can be ignored in the classical limit. Now the second term gives,
     \begin{align}
         \begin{split}
             \chi\sim -\frac{2 m}{\sqrt{\pi}(\gamma^2-1)}K_{0}(2 m|b|)\,.
         \end{split}
     \end{align}
We can see that the massive expression in \eqref{B.15} is non-zero in general.
%\subsection*{Corrections to the \eqref{4.37a} considering the radiation poles}
 We will now show that there will be a non-vanishing. contribution to the $a_1$ %while taking into account the radiation poles. 
 start with:
\begin{align}
    \begin{split}
        \mathcal{K}^\mu= a_1k^\mu+a_2 v_1^\mu+a_3 v_2^\mu.\label{B.18b}
    \end{split}
\end{align}
$a_2$ and $a_3$ s are already derived. Now contracting both side of \eqref{B.18b} with $k^\mu$ we get,
\begin{align}
    \begin{split}
        k\cdot \mathcal{K}= a_1 k^2.
    \end{split}
\end{align}
Therefore,
\begin{align}
    \begin{split}
        a_1&=\frac{1}{k^2}\rmint_{k_1}\hat\delta(k_1\cdot v_2)\frac{k_1\cdot (k-k_1)\,k\cdot (k-k_1)}{(k_1^2-m^2)[(k_1-k)^2-m^2](k_1\cdot v_1+i\epsilon)^2}\,,\\ &
        =\frac{1}{2k^2}\rmint_{k_1}\hat\delta(k_1\cdot v_2)k_1\cdot (k-k_1)\Big[\frac{k^2}{(k_1^2-m^2)[(k_1-k)^2-m^2](k_1\cdot v_1+i\epsilon)^2}+\frac{1}{(k_1^2-m^2)(k_1\cdot v_1+i\epsilon)^2}\\ &
        \hspace{0.8cm}-\frac{1}{[(k_1-k)^2-m^2](k_1\cdot v_1+i\epsilon)^2}\Big]\,.\label{B.20a}
    \end{split}
\end{align}
The second and third terms in \eqref{B.20a} will cancel each other by virtue of relabelling the third term by $k_1-k\to -k_1$ and taking into account $\hat\delta(k\cdot v_i)$\footnote{Note that, implicitly, we take $\epsilon\to 0$ beforehand. But if we take this limit at the end, it will result in the same. For  specified $i\epsilon$ prescription the last two terms for \eqref{B.20a} takes the form,

{\footnotesize
\[
\begin{aligned}
\int_{k_1}\hat\delta(k_1\cdot v_2)\frac{k_1\cdot (k-k_1)}{k_1^2-m^2}\hat\delta'(k_1\cdot v_1)
&\to -\int_{\vec k_1}\frac{d}{dk_1^{(1)}}\hat\delta(k_1^{(1)})\\
&\quad +\int_{ k_1^{(2,3)}}(-m^2+\vec k_1\cdot \vec k)\int_{k_1^{(1)}}\frac{1}{k_1^{(1)2}+k_1^{(2)2}+k_1^{(3)2}+m^2}\frac{d}{dk_1^{(1)}}\hat\delta(k_1^{(1)})\\
&\to 0
\end{aligned}
\]
}}. Therefore, we are left with,
\begin{align}
    \begin{split}
        a_1&=\frac{1}{2}\rmint_{k_1}\hat{\delta}(k_1\cdot v_2)\frac{k_1\cdot(k-k_1)}{(k_1^2-m^2)[(k_1-k)^2-m^2](k_1\cdot v_1+ i \epsilon)^2}\,,\\ &
        =\frac{1}{4}(k^2-2m^2)\hat\chi(k,m)\,.
    \end{split}
\end{align}
Hence, the correction to the impulse takes the following form,
\begin{align}
    \begin{split}
        \Delta p_1^\mu\Big|_{\textrm{corr.}}&=im_1\Big(\frac{s_1m_2 s_2}{8m_p^2}\Big)^2\rmint_{k}e^{ik\cdot b}\hat\delta(k\cdot v_1)\hat\delta(k\cdot v_2)k^\mu(k^2-2m^2)\,\hat\chi_{k}(k,b)\,,\\ &
        =-\frac{m_1}{2\pi}\Big(\frac{s_1m_2 s_2}{8m_p^2}\Big)^2\frac{1}{\sqrt{\gamma^2-1}}\frac{b^\mu}{|b|}\partial_{|b|}\rmint_{0}^{\infty} dk\,k\,J_{0}(k|b|)(k^2+2m^2)\,\hat\chi_{k}(k,m)\,,\\ &
        % =-m_1\Big(\frac{s_1m_2 s_2}{8m_p^2}\Big)^2 -\frac{16m\pi^{5/2}}{\gamma^2(\gamma^2-1)}\frac{b^\mu}{|b|}\partial_{|b|}\rmint_{0}^{\infty} dk\,k\,J_{0}(k|b|))\,\frac{{(k^2+4m^2)}-2m^2}{m(k^2+4m^2)}\,,\\&
      %    =-m_1\Big(\frac{s_1m_2 s_2}{8m_p^2}\Big)^2 \frac{b^\mu}{|b|}\partial_{|b|}\Big(\frac{1}{m}\underbrace{\rmint_{0}^{\infty} dk\,k\,J_{0}(k|b|))}_{\text{UV Divergent, No finite part}}\,-2m\,\rmint_{0}^{\infty} dk\,k\,\frac{J_{0}(k|b|)}{(k^2+4m^2)}\Big)\,,\\&
      %  =-\frac{m_1}{2\pi}\Big(\frac{s_1m_2 s_2}{8m_p^2}\Big)^2\frac{1}{2\sqrt{\pi}\gamma^3(\gamma^2-1)^{3/2}m} \frac{b^\mu}{|b|}\partial_{|b|}\Big(-2 m K_0(2 |b| m)\Big)\,,\\ &
        =-m_1\Big(\frac{s_1m_2 s_2}{8m_p^2}\Big)^2\frac{1}{\pi^{3/2}(\gamma^2-1)^{3/2}} \frac{b^\mu}{|b|}\Big(m K_1(2 |b| m)\Big)\,.\label{B.22aa}
    \end{split}
\end{align}
\textcolor{black}{
\eqref{B.22aa} explicitly indicates that if one compute the impulse from the Eikonal phase, one may get only a part of the impulse which is proportional to $b^\mu$ and other parts (proportional to $v_i^\mu$) will be lost. Along with that \eqref{B.22aa} has a smooth massless limit which reads, $$ \Delta p_1^\mu|_{\textrm{corr.}}\sim \frac{b^\mu}{b^2}\,.$$ }
However, the result in the massless limit direct contradicts with the result in \cite{Kalin:2020mvi}, where they showed $a_1=0$ by ignoring the radiation poles. But it seems there is no such radiation region due to the presence of delta function.
%\textcolor{black}{
%\textbf{\textit{Conclusion from this appendix:}}\\
%\textbullet $\,\,$ If one ignores the radiation poles, i.e. works in the limit $k_0\ll |\vec k|$, the massive as well as the massless impulse integral becomes zero.\\
%\textbullet $\,\,$ If one do the integrals in a more general setting i.e. considering the radiation poles, one will get a non-vanishing contribution to the impulse for both the the massive and the massless  case.}
\section{Two master integrals used for the computation of  waveform}\label{ch2:app:C}
We analyze the following two integrals,
\begin{align}
    \begin{split}
     &   I^{\mu_1\cdots\mu_n}_n:=\rmint_{q}\hat\delta(q\cdot v_1-k\cdot v_1)\hat\delta(q\cdot v_2)\frac{e^{-iq\cdot b}}{q^2}q^{\mu_1}\cdots q^{\mu_n}\,,\\ &
     J^{\mu_1\cdots\mu_n}_n:=\rmint_{q}\hat\delta(q\cdot v_1-k\cdot v_1)\hat\delta(q\cdot v_2)\frac{e^{-iq\cdot b}}{q^2(k-q)^2}q^{\mu_1}\cdots q^{\mu_n}\,.
    \end{split}
\end{align}
The two particular cases that we used in the main text are: $J^{\mu}_{1}$ and $J_{0}$. We first start with $J_{0}\,.$

\begin{align}
    \begin{split}
        J_{(0)}(k)=\rmint_{q}\hat\delta(q\cdot v_1-k\cdot v_1)\hat\delta(q\cdot v_2)\frac{e^{-iq\cdot b}}{q^2(k-q)^2}\,.\label{A.2}
    \end{split}
\end{align}
We will do this integral by implementing a certain choice of frame,
\begin{align}
    v_2^\mu=\delta^\mu_0,\,v_1^\mu=(\gamma,\gamma\beta \,\hat{e}_v)\,b=(0,|b|\,\hat{e}_b)\,,
\end{align}
such that, the unit vectors $\hat{e}_v$ and $\hat{e}_b$ are mutually orthogonal i.e. $\hat{e}_b \cdot \hat{e}_v=0$.
We further proceed with this integral by using the Feynman parametrization method along with the on-shell condition: $k^2=0$. Therefore, $J_0$ can be written as,
\begin{align}
    \begin{split}
        J_{(0)}(k)&=\rmint_0^1dy\rmint_q \hat\delta(q\cdot v_1-k\cdot v_1)\hat\delta(q\cdot v_2)\frac{e^{-iq\cdot b}}{(q-yk)^4},\,k^2=0\,,\\ &
        \xrightarrow[]{\bar q\rightarrow q-yk}\rmint_0^1 dy e^{-iy k\cdot b} \rmint_{\bar q}\hat\delta(\bar q\cdot v_1-(1-y)k\cdot v_1)\hat\delta(\bar q\cdot v_2+y k\cdot v_2)\frac{e^{-i\bar q\cdot b}}{\bar q^4}\,,\\ &
      %  =\frac{1}{\gamma}\rmint_0^1 dy\, e^{-iy k\cdot b}\rmint_{\bar q}\frac{e^{i\tilde q\cdot b}}{\Big[\tilde q^2-y^2(k\cdot v_2)^2+(\frac{y}{\beta}k\cdot v_2+\frac{1-y}{\gamma\beta}k\cdot v_1)^2\Big]^2}\,,\\ &
        =\frac{1}{\gamma}\rmint_{0}^1dy\,e^{-iyk\cdot b}\rmint_0^{\infty}dt\,t\rmint_{\tilde q}\exp\Big[i\tilde q\cdot b-t\tilde q^2-t\Delta(y)^2\Big]\,,\\ &
        =\frac{1}{\gamma}\rmint_{0}^1dy\,e^{-iyk\cdot b}\rmint_0^{\infty}\frac{dt}{4\pi}\, \exp(-\frac{|b|^2}{4t}-t\Delta(k,y)^2)\,,\\ &
       = \frac{b}{\gamma}\rmint_{0}^1dy\,e^{-iyk\cdot b}\frac{ K_1\left(|b| {\Delta(k,y) }\right)}{4 \pi  {\Delta(k,y) }}\label{C.4 mm}
    \end{split}
\end{align}
where, 
\begin{align}
    \begin{split}
       \Delta(k,y)&=\sqrt{-y^2(k\cdot v_2)^2+\frac{y^2}{\beta^2}(k\cdot v_2)^2+\frac{(1-y)^2}{\gamma^2\beta^2}(k\cdot v_1)^2+2\frac{y(1-y)}{\gamma\beta^2}(k\cdot v_2)(k\cdot v_1)}\,,\\ &
        =\frac{1}{\sqrt{\gamma^2-1}}\sqrt{y^2(k\cdot v_2)^2+(1-y)^2(k\cdot v_1)^2+2y(1-y)\gamma (k\cdot v_2)(k\cdot v_1)}\,.\label{A.4}
    \end{split}
\end{align}
In thee waveform calculation, we have to further  perform a integeral over $k$, which in general has the following form,
\begin{align}
    \begin{split}
        f(x):=\rmint_{k}e^{-ik\cdot x}T_{\mu_1\cdots \mu_n}J^{\mu_1\cdots \mu_n}(k)\,.
    \end{split}
\end{align}
In the case of $\varphi^3$ interaction we have,
\begin{align}
    \begin{split}
        f(x)&=\rmint_{k}e^{-ik\cdot x}J_{0}(k), \,k=\Omega n\,,\\ &
        =\frac{b}{4\pi\gamma}\rmint_{\Omega}e^{-i\Omega n\cdot x}\rmint_{0}^1 dy\,e^{-i\,y\,\Omega(n\cdot b)}\frac{K_1[|\Omega|\,|b|\bar\Delta(y)]}{|\Omega|\bar\Delta(y)}\,\,,\bar\Delta(y)\equiv\frac{\Delta(k,y)}{|\Omega|}\,.\label{ch2:A.6}
    \end{split}
\end{align}
Now to do the integral over $\Omega$, use the following representation of modified Bessel function:
\begin{align}
    \begin{split}
        K_{\nu}(x)=\frac{1}{2}(\frac{x}{2})^{\nu}\rmint_{0}^{\infty}dt\,\exp\Big(-t-\frac{x^2}{4t}\Big)\frac{1}{t^{\nu+1}}\,.
    \end{split}
\end{align}
Therefore,
\begin{align}
    \begin{split}
        f(x)&=\frac{|b|}{4\pi\gamma}\rmint_{0}^{1}dy\, \frac{1}{\Bar\Delta(y)}\rmint_{\Omega}e^{-i\Omega n\cdot(x+yb)}\frac{|b|}{4}\Bar{\Delta(y)}\rmint \frac{dt}{t^2}\,\exp\Big(-t-\frac{\Omega^2|b|^2\bar\Delta(y)^2}{4t}\Big)\,,\\ &
        =\frac{|b|^2}{16\pi\gamma}\rmint_{0}^{1}dy\, \rmint_{0}^{\infty} \frac{dt}{t^2}e^{-t}\rmint d\Omega\,e^{-i\,\Omega\, l}\exp\Big(-\frac{\Omega^2|b|^2\bar\Delta(y)^2}{4t}\Big)\,,\\ &
         =\frac{|b|^2}{16\pi\gamma}\rmint_{0}^{1}dy\, \rmint_{0}^{\infty} \frac{dt}{t^2}e^{-t}\frac{2 \sqrt{\pi } \sqrt{t} e^{-\frac{l^2 t}{|b|^2 \bar\Delta(y) ^2}}}{|b| \bar\Delta(y) }\,,\\ &
      %   =\frac{|b|}{8\sqrt{\pi}\gamma}\rmint_{0}^1dy\frac{1}{\bar\Delta(y)}\rmint_{0}^{\infty} dt\frac{\exp\Big(-\frac{l^2}{|b|^2\Bar\Delta(y)^2}t-t\Big)}{t^{3/2}}\,,\\ &
        %  =\frac{|b|}{8\sqrt{\pi}\gamma}\rmint_{0}^1dy\frac{1}{\bar\Delta(y)}\sqrt{\frac{l^2}{|b|^2\bar\Delta^2}+1}\,\,\Gamma(-1/2)\,,\\ &
          =-\frac{|b|}{4\gamma}\rmint_{0}^{1}dy\,\frac{1}{\Bar\Delta(y)}\sqrt{\frac{l^2}{|b|^2\bar\Delta^2}+1},\,l:=n\cdot(x+yb)\,.\label{A.8}
    \end{split}
\end{align}
Now come discuss the other important vector integral:
\begin{align}
    \begin{split}
J_{(1)}^{\mu}&=\rmint_q\hat\delta(q\cdot v_1-k\cdot v_1)\hat\delta(q\cdot v_2)e^{-iq\cdot b}\frac{q^{\mu}}{q^2(k-q)^2}\,,\\ &
=\rmint_{0}^1dy\, e^{-iyk\cdot b}\rmint_{q}\hat\delta(q\cdot v_1-(1-y))\hat\delta(q\cdot v_2+yk\cdot v_2)\frac{e^{-iq\cdot b}}{q^4}(q^\mu+yk^\mu)\,,\\ &
:=\mathcal{J}^{\mu}+\frac{|b|\,k^\mu}{\gamma}\rmint_{0}^1dy\,y\,\,e^{-iyk\cdot b}\frac{K_1(|b|\Delta(k,y))}{4\pi\Delta(k,y)}\,.\label{A.9 m}
    \end{split}
\end{align}
$\mathcal{J}^{\mu}$ can be solved by reducing the vector integral in scalar integral as follows,
\begin{align}
    \begin{split}
\mathcal{J}^{\mu}=\rmint_{0}^1dy\,e^{-iyk\cdot b}\,[\lambda_b b^{\mu}+\lambda_1 v_1^{\mu}+\lambda_2v_2^\mu]\,,\label{C.11a}
    \end{split}
\end{align}
Now we need to solve the following equations:\par
    \textbullet $\,\,$ Contracting with $ b_{\mu}$ in the both side of \eqref{C.11a} we will get,
\begin{align}
    \begin{split}
 \rmint_{0}^1dy\, e^{-iyk\cdot b}\rmint_{q}\hat\delta(q\cdot v_1-(1-y))\hat\delta(q\cdot v_2+yk\cdot v_2)\frac{e^{-iq\cdot b}}{q^4}q\cdot b=\mathcal{J}\cdot b=-\rmint_0^1 dy\,e^{-iyk\cdot b}b^2\lambda_b\,.
    \end{split}
\end{align}
To carry out the $q$ integral on the left-hand side, we introduce a new parameter $\kappa$ as,
\begin{align}
    \begin{split}
        -|b|^2\lambda_b&=\lim_{\kappa \rightarrow 1}\rmint_{q}\hat\delta(q\cdot v_1-(1-y))\hat\delta(q\cdot v_2+yk\cdot v_2)\frac{e^{-i\kappa q\cdot b}}{q^4}q\cdot b\,,\\ &
        =i\lim_{\kappa \rightarrow 1}\frac{\partial }{\partial \kappa}\rmint_{q}\hat\delta(q\cdot v_1-(1-y))\hat\delta(q\cdot v_2+yk\cdot v_2)\frac{e^{-i\kappa q\cdot b}}{q^4}\,,\\ &
=i\lim_{\kappa\rightarrow 1}\frac{\partial}{\partial \kappa}\hat{J}_0(\kappa |b|),\,\hat J_{0}(|b|)\equiv \frac{|b|}{\gamma}\frac{K_{1}(|b|\Delta(k,y))}{4\pi \Delta(k,y)}\,.
    \end{split}
\end{align}
Hence, 
\begin{align}
    \begin{split}
        \lambda_b=\frac{ i\, K_0(|b| \Delta(k,y) )}{4\pi\gamma }\,.
    \end{split}
\end{align}
 \textbullet $\,\,$ Contracting with $ v_{1\mu}$ in the both side of \eqref{C.11a} we will get,
\begin{align}
    \begin{split}
  \rmint_0^1 dy\,e^{-iyk\cdot b}\rmint_{q}\hat\delta(q\cdot v_1-(1-y)k\cdot v_1)\hat\delta(q\cdot v_2+yk\cdot v_2)\frac{e^{-iq\cdot b}}{q^4}q\cdot v_1=\mathcal{J}\cdot v_1=  \rmint_0^1 dy\,e^{-iyk\cdot b}[\lambda_1+\gamma \lambda_2]\,.
    \end{split}
\end{align}
Again, the $q$ integral in LHS can be done again by introducing an auxiliary parameter $\kappa$ as,
\begin{align}
    \begin{split} \label{B.16}
        \lambda_1+\gamma \lambda_2&=\rmint_{q}\hat\delta(q\cdot v_1-(1-y)k\cdot v_1)\hat\delta(q\cdot v_2+yk\cdot v_2)\frac{e^{-iq\cdot b-i \kappa q\cdot v_1}}{q^4}q\cdot v_1\,,\\ &
      =  \rmint_{q}\hat\delta(q\cdot v_1-(1-y)k\cdot v_1)\hat\delta(q\cdot v_2+yk\cdot v_2)\frac{e^{-iq\cdot b}}{q^4}(1-y)k\cdot v_1\,,\\ &
      =(1-y)k\cdot v_1 \hat J_0\,.
    \end{split}
\end{align}
 \textbullet $\,\,$ Contracting with $ v_{2\mu}$ in the both side of \eqref{C.11a} we will get,
\begin{align}
    \begin{split}
    \rmint_0^1 dy\,e^{-iyk\cdot b}\rmint_{q}\hat\delta(q\cdot v_1-(1-y)k\cdot v_1)\hat\delta(q\cdot v_2+yk\cdot v_2)\frac{e^{-iq\cdot b}}{q^4}q\cdot v_2=\mathcal{J}\cdot v_2=  \rmint_0^1 dy\,e^{-iyk\cdot b}[\gamma\lambda_1+\lambda_2]\,.
    \end{split}
\end{align}
Hence,
\begin{align}
    \begin{split} \label{B.18}
        \gamma \lambda_1+\lambda_2&=-\rmint_{q}\hat\delta(q\cdot v_1-(1-y)k\cdot v_1)\hat\delta(q\cdot v_2+yk\cdot v_2)\frac{e^{-iq\cdot b}}{q^4}y(k\cdot v_2)\,,\\ &
        =-y k\cdot v_2\hat J_0 \,.
    \end{split}
\end{align}
Now solving (\ref{B.16}) and (\ref{B.18}) we get,
\begin{align}
    \begin{split}
      &  \lambda_2=\frac{y(k\cdot v_2)+\gamma (1-y)k\cdot v_1}{\gamma^2-1}\hat{J}_0\,,\\ &
      \lambda_1=\frac{-\gamma y k\cdot v_2-(1-y)k\cdot v_1}{\gamma^2-1}\hat{J}_0\,.
    \end{split}
\end{align}
Therefore, the contribution to the waveform takes the following form,
\begin{align}
   \begin{split}
       f(x)&=T_{\mu}\rmint_{\Omega}e^{-ik\cdot x}\mathcal{J}^{\mu}\,,\\ &
=T_{\mu}\rmint_{\Omega}e^{-ik\cdot x}\rmint_0^1dy\, e^{-iyk\cdot b}[\lambda_b b^\mu+\lambda_1 v_1^\mu+\lambda_2 v_2^\mu]\,,\\ &
=f_1+f_2+f_3\,,\label{A.19 m}
   \end{split}
\end{align}
where, 
\begin{align}
    \begin{split}
        f_1(x|l)&=\frac{-iT\cdot b}{4\pi \gamma}\rmint_{\Omega}e^{-i\Omega\,(n\cdot x)}\rmint_{0}^1dy\,K_0(|b||\Omega|\bar\Delta(y))\,,\\ &
        =\frac{-iT\cdot b}{4\pi \gamma}\rmint_{0}^1 dy\rmint_{0}^\infty \frac{dt}{2t}e^{-t}\rmint_{\Omega}d\Omega\,e^{-i\Omega l}\exp(-\frac{|b|^2\Omega^2\bar\Delta^2}{4t})\,,\\ &
    %   = \frac{-iT\cdot b}{4\pi \gamma}\rmint_{0}^1 dy\rmint_{0}^\infty \frac{dt}{2t}e^{-t}\frac{2 \sqrt{\pi } \sqrt{t} e^{-\frac{l^2 t}{b^2 \bar\Delta ^2}}}{b \bar\Delta }\,,\\ &
       =\frac{-iT\cdot b}{4 \gamma |b|}\rmint_{0}^1 dy\frac{1}{\bar\Delta(y)}\sqrt{\frac{|b|^2\bar\Delta^2}{|b|^2\bar\Delta^2+l^2}}\,,
    \end{split}
\end{align}
%and,
\begin{align}
    \begin{split}
        f_2(x|l)&=T\cdot v_1\, \bar\lambda_1\rmint_{\Omega}\Omega\, e^{-i\Omega (n\cdot x)}\rmint_0^1 dy\,e^{-i\Omega y(n\cdot b)}\frac{b}{\gamma}\frac{K_{1}(b|\Omega|\Delta(k,y))}{4\pi |\Omega|\bar\Delta(y)}\,,\\ &
        =\frac{|b|^2\, T\cdot v_1\, \bar\lambda_1}{16\pi\gamma}\rmint_0^1 dy\rmint_{0}^\infty dt\frac{e^{-t}}{t^2}\rmint_{\Omega} e^{-i\Omega l}\Omega \,\exp\Big(-\frac{\Omega^2 |b|^2 \bar\Delta(y)^2}{4t}\Big)\,,\\ &
    %    =-\frac{i b^2\, T\cdot v_1\, \bar\lambda_1}{4\sqrt{\pi}\gamma}\rmint_0^1 dy \rmint_{0}^\infty dt\frac{e^{-t}}{t^2}\frac { l e^{-\frac{l^2 t}{b^2 \Delta ^2}}}{\left(\frac{b^2 \Delta ^2}{t}\right)^{3/2}}\,,\\ &
        =-\frac{i |b|^2\, T\cdot v_1\, \bar\lambda_1}{4\gamma}\rmint_0^1 dy\, \frac{l}{|b|^3 \bar\Delta^3}\frac{1}{\sqrt{\frac{l^2}{|b|^2 \Delta ^2}+1}}\,.
    \end{split}
\end{align}
and,
\begin{align}
    \begin{split}
        f_3(x|l)= -\frac{i \, T\cdot v_2\, \bar\lambda_2}{4\gamma |b|} \rmint_0^1 dy\, \frac{l}{ \bar\Delta^3}\frac{1}{\sqrt{\frac{l^2}{|b|^2 \Delta ^2}+1}}\,,
    \end{split}
\end{align}
where, $\bar\lambda_i=\frac{\lambda_i}{\Omega}$.
%\newpage
\section{\texorpdfstring{Analysis through method of regions for $\lambda_4\varphi^4$ vertex contribution in the waveform}{Analysis through method of regions for lambda4 phi4 vertex contribution in the waveform}}\label{ch2:app:D}
We analyze the integral mentioned in (\ref{5.18k}) using the \textit{method of regions}. The method of regions \cite{Beneke:1997zp} is a universal technique for expanding Feynman integrals in various limits of momenta and masses.
We split the integration into two regions, one where $q\sim b^{-1}\gg k_{3}$ and $k_{3}\gg q\sim b^{-1}$. The integral \eqref{4PHI} for the region $q\gg k_{3}$ reduces to 
\begin{align}
    \begin{split}
    \label{FstRe1}f_{\varphi}(x)&=\rmint_{\Omega}e^{-i k\cdot x+i k\cdot b_{1}}\rmint_{k_3,q} \frac{\hat\delta(k\cdot v_1-q\cdot v_1)\hat\delta(k_3\cdot v_2)\hat\delta(k_3\cdot v_2-q\cdot v_2)}{k_3^2 \, q^2(q-k)^2}e^{i q\cdot b}\,.
    \end{split}
\end{align}
The above integration can be done easily in the rest frame of the second particle. Hence, taking $v^{\mu}_{2}=(1,0,0,0)$, the above integral simplifies to 
\begin{align}
    \begin{split}
  \label{FstRe}   f_{\varphi}(x)&=\rmint_{\Omega}e^{-i k\cdot x+i k\cdot b_{1}}\rmint_{\vec{k}_3,\vec{q}} \frac{\hat\delta(k\cdot v_1-\vec{q}\cdot \vec{v}_1)}{|\vec{k}_3^2 |\, \vec{q}^2(q-k)^2|_{q^{0}=0}}e^{i \vec{q}\cdot \vec{b}}\,.
    \end{split}
\end{align}
The integral \eqref{FstRe} is a pure divergent for the $\vec{k}_3$ integral. Similarly, for the limit $k_3\gg q$, the integral \eqref{4PHI} reduces to 
\begin{align}
    \begin{split}
    \label{2NdRe}f_{\varphi}(x)&=\rmint_{\Omega}e^{-i k\cdot x+i k\cdot b_{1}}\rmint_{k_3,q} \frac{\hat\delta(k\cdot v_1-q\cdot v_1)\hat\delta(k_3\cdot v_2)\hat\delta(k_3\cdot v_2-q\cdot v_2)}{k_3^4 \,(q-k)^2}e^{i q\cdot b}
    \end{split}
\end{align}
and solving it in the rest frame of the second particle, the $\vec{k}_{3}$ integral and the $\vec{q}$ integral factorizes and the $\vec{k}_{3}$ integral equates to zero. Thus, 
\begin{align}
    \begin{split} \label{nocont}
    f_{\varphi}(x) =&
\rmint_{\Omega}e^{-i k\cdot x}\rmint_{k_3\gg q} \frac{\hat\delta(k\cdot v_1-q\cdot v_1)\hat\delta(k_3\cdot v_2)\hat\delta(k_3\cdot v_2-q\cdot v_2)}{k_3^2 (q-k)^2(q-k_3)^2}e^{i (k-q)\cdot b_1}e^{i(k_3-q)\cdot b_2}e^{ik_3\cdot b_2}\, + \\
&\rmint_{\Omega}e^{-i k\cdot x}\rmint_{q\gg 
k_3} \frac{\hat\delta(k\cdot v_1-q\cdot v_1)\hat\delta(k_3\cdot v_2)\hat\delta(k_3\cdot v_2-q\cdot v_2)}{k_3^2 (q-k)^2(q-k_3)^2}e^{i (k-q)\cdot b_1}e^{i(k_3-q)\cdot b_2}e^{ik_3\cdot b_2}
 \end{split}
\end{align}
is purely divergent. %and the contribution to the waveform coming from this diagram must be avoided. 
%The vanishing of the contribution of the waveform from the four-point scalar vertex can also be shown alternatively using the Feynman parametrization and is done in Appendix~(\ref{App2}). Following the same argument, the 2PM radiation diagrams consisting of $\lambda_3 h \varphi^3$ bulk vertex will not contribute.
%In the main text, we have shown using the method of  region \cite{Beneke:1997zp} that the diagram, due to the self-interaction vertex $\lambda_4\varphi^4$ does not contribute to the scalar waveform as its contribution is purely divergent and has to be avoided while computing the scalar waveform. Now, we show this vanishing of the contribution of the waveform from the four-point scalar vertex alternatively using the Feynman parametrization. We start with the following integral.\\
%\begin{align}
%\begin{split}
%&\mathcal{K}\equiv\rmint_{\Omega} e^{-ik \cdot \tilde{x}}\rmint_{k_2}\frac{\hat{\delta}(k_2\cdot v_2)}{k_2^2}\rmint_{\Tilde{q}}\frac{\hat{\delta}{(\Tilde{q}\cdot v_1)}\hat{\delta}{[({k}-\Tilde{q})\cdot v_2]}}{\Tilde{q}^2[k_2-(k-\Tilde{q})]^2}e^{i \Tilde{q}\cdot b}\,,\\&
%=\rmint_0^1 dx\rmint_{\Omega} e^{-ik \cdot \tilde{x}}\rmint_{k_2}\frac{\hat{\delta}(k_2\cdot v_2)}{k_2^2} \rmint_{\Tilde{q}}\frac{\hat{\delta}{(\Tilde{q}\cdot v_1)}\hat{\delta}{[({k}-\Tilde{q})\cdot v_2]}}{(x\Tilde{q}^2+(1-x)[k_2^2+\Tilde{q}^2+2k_2\cdot \Tilde{q}-2k\cdot \Tilde{q}-2k_2\cdot k])^2}e^{i \Tilde{q}\cdot b}\,,\\&
%\xrightarrow[q'+k=\tilde{q}]{  }  \rmint_0^1 dx\rmint_{\Omega} e^{-ik \cdot \tilde{x}}\rmint_{k_2}\frac{\hat{\delta}(k_2\cdot v_2)}{k_2^2} \rmint_{q'}\frac{\hat{\delta}{\Big(({q'+k})\cdot v_1\Big)}\hat{\delta}{(q'\cdot v_2)}}{[x (q'+k)^2+(1-x)(k_2+q')^2]^2}e^{i (q'+k) \cdot b}\,.\\&
%\text{ Now replacing  $x\rightarrow (1-x)$ we achieve, }\\&
%\mathcal{K}=  \rmint_0^1 dx \, e^{ik \cdot \tilde{x}}\rmint_{q',k_2,\Omega}\frac{\hat{\delta}{(q'\cdot v_1+k\cdot v_1)}\hat{\delta}{(q'\cdot v_2)}}{k_2^2[(q'+xk_2)^2+k_2^2\,x\,(1-x)+2k\cdot q'(1-x)]^2}e^{i (q'+k)\cdot b}\,,\\&
%\xrightarrow[q'+x k_2 \rightarrow q'']{}\rmint_0^1 dx e^{-ik \cdot \Tilde{x}}\rmint_{q'',k_2,\Omega} \frac{\hat{\delta}(k_2\cdot v_2)}{k_2^2\{ (q''+k(1-x))^2+Q^2\}^2}\hat{\delta}\Big((q''-x k_2+k)\cdot v_1\Big)\\ &
%\hspace{0.8cm}\times\hat{\delta}\Big((q-x k_2)\cdot v_2\Big)\exp({i q''\cdot b- ix k_2\cdot b+i k\cdot b})\,,\\&
%=\rmint_0^1 dx e^{-ik \cdot \tilde{x}}\rmint_{q'',k_2,\Omega}  \frac{\delta(k_2\cdot v_2)}{k_2^2 \{ (q''+k(1-x))^2+Q^2\}^2}\hat{\delta}\Big(q''-x k_2+k)\cdot v_1\Big)\hat{\delta}(q''\cdot v_2)\\ &
%\hspace{0.8cm}\times\exp({i q''\cdot b- ix k_2\cdot b+ik\cdot b}),\,\,\,\,\,\, k:=\Omega \hat{n}\,,\\&
%=\rmint_0^1 dx \rmint_{\vec{k_2},\vec{q''},\Omega}\frac{1}{k_2^2 \{ (q''+\Omega \hat{n}(1-x))^2+Q^2\}^2}\hat{\delta}\Big((q''-x k_2+\Omega \hat{n})\cdot v_1\Big)\\ &
%\hspace{0.8cm}\times\exp({i q''\cdot b- ix k_2\cdot b-i{k}\cdot(\tilde{x}- b)})\,,\\&
%= \rmint_0^1 dx \rmint_{\vec{q''},\vec{k_2}}\frac{1}{\{ (q''+k(1-x))^2+\underbrace{Q^2}
%_{x(1-x)(k-k_2)^2}\}^2k_2^2 } e^{i q''\cdot l}e^{i x k_2\cdot m}\\&
% \xrightarrow[q''+\Omega \hat{n}(1-x) = q''']{ }
% \rmint_0^1 dx \rmint_{{q'''},\Omega,\vec{k_2}}\frac{1}{\{ (q''')^2+\underbrace{Q^2}
%_{x(1-x)(\Omega\hat{n}-k_2)^2}\}^2k_2^2 }\hat{\delta}(q'''-x k_2+x \,\Omega \hat{n})\cdot v_1)\textcolor{red}{\delta((q'''-\Omega\hat{n}(1-x))\cdot v_2) }\\ &
%\hspace{0.8cm}\times\exp\Big(i( q'''-\Omega \hat{n}(1-x))\cdot b\Big)\exp\Big({i x k_2\cdot b-i \Omega \hat{n}(\Tilde{x}-b)}\Big)\,,\\&
%= \rmint_0^1 dx \rmint_{\vec{q'''},\vec{k_2}}\frac{\exp({i q'''\cdot l})\exp({i x k_2\cdot m})}{\{\textcolor{red}{\alpha(q'''^{(1)}-x k_2^{(1)})^2- (\vec{q}''')^2}+
%{\textcolor{black}{x(1-x)(\frac{(q'''^{(1)}-xk_2^{(1)})}{\Gamma \, x n\cdot v_1}\vec{n}+k_2)^2+x(1-x)(\frac{(q'''^{(1)}-xk_2^{(1)})}{\Gamma \, x n\cdot v_1})^2}}\}^2k_2^2 } \,.
%\label{7.23}\end{split}
%\end{align}
%where $\vec{l}:=\vec{b}-\frac{\vec{v_1}}{ x\hat{n}\cdot v_1}\hat{n}\cdot (-\tilde{x}+x b), \vec{m}:=-\vec{b}-%\frac{\vec{v_1}}{\hat{n}\cdot v_1}\hat{n}\cdot (-\tilde{x}+x b)$.
%This change of variable can be done because as: $$\rmint_a^b f(s) ds =\rmint_a^b f(a+b-s) ds.$$

%Now to do the integral (\ref{7.23}), we use the Schwinger parametrisation. Using that (\ref{7.23}) can be written as,
%$m_2,m_3 $ are the second and third components of the $\vec{m}$ and\\
%\begin{center}
%\boxed{\begin{aligned}\Delta:=(1+\frac{\sqrt{x(1-x)}\gamma \beta}{\hat{n}\cdot v_1} ), \rho:= %\frac{\sqrt{x(1-x)}\gamma\beta \hat{n}^{1}}{(\hat{n}\cdot v_1) x}\,. \end{aligned}}
%\end{center}

%\vspace{0.5cm}
%To this end, one can easily check that the regularised value of this integral is ${\bf 0}$. Hence, we can ignore the contribution of this diagram in the waveform calculation.




%\section{Some details}

%\\
%\hspace{0cm}\textbf{Eikonal phase computation: a connection to scattering amplitudes}
%\\
%The simplest diagram has a scalar exchange between two worldline which is at $1PM$. The contribution to the Eikonal phase has the following form,
%\begin{align}
   % \begin{split}
        %\chi&=\frac{m_1m_2s_1s_2}{4m_p^2}\rmint_{k_1,k_2}e^{ik_1\cdot b_1+ik_2\cdot b_2}\hat\delta(k_1\cdot v_1)\hat{\delta}(k_2\cdot v_2)\langle\tilde \varphi(k_1)\tilde\varphi(k_2)\rangle \textcolor{black}{\delta^{4}(k_1+k_2)}\\ &
      %  =\frac{m_1m_2s_1s_2}{4m_p^2}\rmint_{k}\frac{\hat\delta(k\cdot v_1)\hat\delta(k\cdot v_2)}{k^2-m^2}e^{ik\cdot b}
   % \end{split}
%\end{align}
%To do this integral, we need to suitably choose the coordinate of the binary black holes. We study the dynamics from the frame of the second black hole. Hence we can parametrize the velocities as,
%\begin{align}
  %  \begin{split}
  %      v_2=(1,0,0,0),\,v_1=%(\gamma,0,0,\gamma\beta).
  %  \end{split}
%\end{align}
%hence,
%\begin{align}
 %   \begin{split}
 %   \chi&=\frac{m_1m_2s_1s_2}{4m_p^2}\rmint_{k}\frac{\hat{\delta}{(\gamma k^{(0)}-\gamma\beta k^{(3)})}\hat\delta(k^{0})}{k^2-m^2}e^{ik\cdot b}\\ &
   % =-\frac{m_1m_2s_1s_2}{4m_p^2\sqrt{\gamma^2-1}}
  % \rmint_{\tilde k}\frac{e^{i\tilde k \cdot b}}{\tilde k^2+m^2}=-\frac{m_1m_2s_1s_2}{4m_p^2\sqrt{\gamma^2-1}}K_{0}(m|b|).
 %   \end{split}
%\end{align}
%\\
%@1.5PM, the relevant field vertex comes from,
 %   \begin{align}
     %  \begin{split}
%\mathcal{O}(\lambda):\lambda\rmint d^4x \frac{h}{2m_p}\tilde{\varphi}^2\sim & \lambda \rmint d^4x \rmint_{k_1,k_2,k_3}h(k_1)\tilde\varphi
%(k_2)\tilde\varphi
%(k_3)   e^{-%ix\cdot(k_1+k_2+k_3)}   
%\end{split}
  %  \end{align}
%Hence the vertex factor is: 
%\begin{align}
   % \begin{split}
         %\lambda\rmint_{k_1,k_2,k_3}\hat\delta^{(4)}(k_1+k_2+k_3).
    %\end{split}
%\end{align}


%Hence, the amplitude for the relevant diagram has the form,
%\begin{align}
    %\begin{split}
        %\chi&=\lambda\Big(\frac{m_2s_2}{2m_p}\Big)^2\rmint\mathcal{D}\tilde\varphi \mathcal{D}\tilde h_{\mu\nu} \,e^{iS_{\text{quad}}}\rmint_{k_1,k_2,k_3}\hat\delta^{(4)}\Big(\sum_{i=1}^3k_i\Big)h(k_1)\tilde\varphi
%(k_2)\tilde\varphi(k_3)\rmint_{k'}\tilde\varphi(k')e^{ik'\cdot b_2}\hat\delta(k'\cdot v_2)\hat\delta(k'-k_2)\\ & \hspace{2.2 cm}
%\rmint_{k''}\tilde\varphi(k'')e^{ik''\cdot b_2}\hat\delta(k''\cdot v_2)\hat\delta(k''-k_3)\rmint_{k'''}\tilde h_{\mu\nu}(k''')e^{ik'''\cdot b_1}\hat\delta(k'''\cdot v_1)\hat\delta(k_1-k''')v_1^{\mu}v_1^{\nu},\\ &
%=\lambda\rmint_{k_1,k_2,k_3} \hat\delta(k_1+k_2+k_3)\hat\delta(k_1\cdot v_1)\hat\delta(k_2\cdot v_2)\hat\delta(k_3\cdot v_2)v_1^{\mu}v_1^{\nu}\,\langle\tilde\varphi\tilde\varphi\rangle_{k_2}\langle\tilde\varphi\tilde\varphi\rangle_{k_3}\langle h\,\tilde h_{\mu\nu}\rangle_{k_1},\\ &
%=\lambda\rmint_{k_1,k_2}\hat\delta(k_1\cdot v_1)\hat\delta(k_2\cdot v_2)\hat\delta(-k_1\cdot v_2-k_2\cdot v_2)v_1^{\mu}v_1^{\nu}\frac{e^{ik_1\cdot b}}{k_1^2(k_2^2-m^2)[(k_1+k_2)^2-m^2]}\eta_{\mu\nu}\\ &
%=\lambda\rmint_{k_1,k_2}\hat\delta(k_1\cdot v_1)\hat\delta(k_2\cdot v_2)\hat\delta(k_1\cdot v_2)v_1^{\mu}v_1^{\nu}\frac{e^{ik_1\cdot b}}{k_1^2(k_2^2-m^2)[(k_1+k_2)^2-m^2]}\eta_{\mu\nu}.    \end{split}
%\end{align}
%As, $\hat\delta(k_2\cdot v_2)$ has support at $k_2\cdot v_2=0$ we neglect the second term in the third delta function.
%So, the non-trivial integral that we have to do is the following,
%\begin{align}
   % \begin{split}
%\rmint_{k_1,k_2}\frac{\hat\delta(k_1\cdot v_1)\hat\delta(k_1\cdot v_2)\hat\delta(k_2\cdot v_2)}{k_1^2(k_2^2-m^2)[(k_1+k_2)^2-m^2]}e^{ik_1\cdot b}\label{3.11}
  %  \end{split}
%\end{align}

%In this parametrization, we can write down the integral as,
%\begin{align}
  %  \begin{split}
%\rmint_{k_1,k_2}\frac{\hat\delta(k_1\cdot v_1)\hat\delta(k_1\cdot v_2)\hat\delta(k_2\cdot v_2)}{k_1^2(k_2^2-m^2)[(k_1+k_2)^2-m^2]}e^{ik_1\cdot b}&\sim \rmint_{k_1,k_2}\frac{\hat\delta(\gamma k_1^{0}-\gamma\beta k_1^{3})\hat\delta(k_1^{0})\hat\delta(k_2^0)}{k_1^2(k_2^2-m^2)[(k_1+k_2)^2-m^2]}e^{ik_1\cdot b}\\ &
%=\frac{1}{\gamma\beta}\rmint_{\vec k_1,k_2}\frac{\hat\delta(k_1^{3})\hat\delta(k_2^0)}{k_1^2(k_2^2-m^2)[(k_1+k_2)^2-m^2]}e^{ik_1\cdot b}\\ &
%=-\frac{1}{\gamma\beta}\rmint d^2 l_1\frac{e^{il_1\cdot b}}{\vec{l}_1^2}\rmint d^3 l_2\frac{1}{(\vec {l}_2^2+m^2)[(\vec{l}_1+\vec{l}_2)^2+m^2]}\\ &
%=\frac{1}{\gamma\beta}\rmint d^2 l_1\frac{e^{il_1\cdot b}}{\vec{l}_1^3}\arctan\Big(\frac{|\vec{l}_1|}{2m}\Big)\\ &
%=\frac{1}{\gamma\beta}\rmint _{0}^{\infty}dl\,\frac{J_{0}(bl)}{l^2}\arctan\Big(\frac{l}{2m}\Big).\label{3.13}
  %  \end{split}
%\end{align}
% The integral in (\ref{3.13}) has IR divergence but for computation of impulse, we need the derivative of the integral w.r.t b. So,
 %\begin{align}
     %\begin{split}
        %  \frac{\partial \chi}{\partial b}&\sim \frac{1}{\gamma\beta}\frac{\partial}{\partial b}\rmint _{0}^{\infty}dl\,\frac{J_{0}(bl)}{l^2}\arctan\Big(\frac{l}{2m}\Big)\\ &
        %  =-\frac{1}{\gamma\beta}\rmint_{0}^{\infty}\frac{J_{1}(bl)}{l}\arctan\Big(\frac{l}{2m}\Big):=-\frac{1}{\gamma\beta}I(m,b).
   %  \end{split}
% \end{align}
%\begin{figure}
   % \centering
%\includegraphics[scale=0.25]{fp.pdf}
 %   \caption{Caption}
 %   \label{fp}
%\end{figure}
%\\
%Fig. (2) is the contour plot for $I(m,b)$ for different $m,b$. 
%Contribution to impulse:\\
%\begin{align}
   % \begin{split}
     %   \hat\delta p_2^{\mu}=\frac{\partial \chi}{\partial b_{\mu}}=\hat b^{\mu}    \frac{\partial \chi}{\partial b}\sim \frac{\hat b^{\mu}}{\gamma\beta}I_{1}(m,b).
%1  %  \end{split}
%\end{align}
%Note: If one take the massless scalar field then $\hat\delta p_2^{\mu}\sim \frac{1}{\gamma\beta}$ i.e, independent of the impact parameter b.\\ 
%Another diagram comes from the self-interaction vertex, i.e. $\delta \tilde \varphi^4$. The contribution to the Eikonal phase has the following form:
%\begin{align}
  %  \begin{split}
    %\chi&=\Big(\frac{m_1s_1m_2%s_2}{4m_p}\Big)^2\rmint_{\{k_i\}}\hat\delta^{(4)}\Big(\sum_{i=1}^{4}k_{i}\Big)\hat\delta(k_1\cdot v_2)\hat\delta(k_2\cdot v_2)\hat\delta(k_3\cdot v_1)\hat\delta(k_4\cdot v_1)\prod_{j=1}^{4}\frac{1}{k_j^2-m^2}e^{i(k_1+k_2)\cdot b_2}e^{i(k_3+k_4)\cdot b_1}\\ &
      %  \sim \rmint_{\{k_{i\ne 4}\}}\hat\delta(k_1\cdot v_2)\hat\delta(k_2\cdot v_2)\hat\delta (k_3\cdot v_1)\hat\delta(k_1\cdot v_1+k_2\cdot v_1+k_3\cdot v_1)\,\frac{e^{i(k_1+k_2)\cdot b}}{\prod_{j=1}^{3}(k_j^2-m^2)[(k_1+k_2+k_3)^2-m^2]}\\ &
        %=\rmint_{k_1,k,k_3}\hat\delta(k_1\cdot v_2)\hat\delta(k\cdot v_2)\hat\delta(k_3\cdot v_1)\hat\delta[(k+k_3)\cdot v_1]\,\frac{e^{ik\cdot b}}{(k_1^2-m^2)(k_3^2-m^2)[(k-k_1)^2-m^2][(k+k_3)^2-m^2]}\\ &
%=\rmint_{k_1,k,k_3}\hat\delta(k_1^{0})\hat\delta(k^{0})\hat\delta(\gamma k_3^{0}-\gamma\beta k_3^{3})\hat\delta(\gamma k^0-\gamma\beta k^3)[\cdots]\\ &
%=\frac{1}{\gamma^2\beta}\rmint_{\vec k_1,\Vec{k},k_3}\hat\delta( k_3^{0}-\beta k_3^{3})[\cdots]\label{3.16}
   % \end{split}
%\end{align}
%In (\ref{3.16}), one can see that ${k^{\mu}}$ and $k_1^{\mu}$ becomes completely spatial in 2-\textit{dim} and 3-\textit{dim} respectively and  $k_3^{\mu}$ has two dependent components while integrating over the delta function. To do this integral systematically, define a scaled variable as,
%\begin{align}
 %   \begin{split}
%  &     \tilde k_3^3=\frac{k_3^3}{\gamma}
 %   \end{split}
%\end{align}
%Now in terms of this redefined variable we can write down the integral as,
%\begin{align}
   % \begin{split}
      %  \chi&=\frac{1}{\gamma %\beta}\rmint d^3k_1 d^3 \tilde k_3 d^2k\frac{e^{i\vec k\cdot \vec b}}{(\vec k_1^2+m^2)(\tilde k_3^2+m^2)[(\vec k-\vec k_1)^2+m^2][(\tilde k_3+\vec k)^2+m^2]}\\ &
      %  \sim-\frac{1}{\gamma\beta}\rmint d^2 k\,\frac{e^{i\vec k\cdot \vec b}}{k^2}\arctan^2(\frac{|\vec k|}{2m})\\ &
      %  =-\frac{1}{\gamma\beta}\rmint_{0}^{\infty}dk\,\frac{J_{0}(kb)}{k}\arctan^2\Big(\frac{k}{2m}\Big)
   % \end{split}
%\end{align}Again to compute the impulse we need to take derivative with respect to the impact factor b. So,
%\begin{align}
    %\begin{split}
     %   \hat\delta p_2^{\mu}=\frac{\partial \chi}{\partial b_{\mu}}&=-\frac{\hat b^{\mu}}{\gamma\beta}\frac{\partial}{\partial b}\rmint_{0}^{\infty}dk\,\frac{J_{0}(kb)}{k}\arctan^2\Big(\frac{k}{2m}\Big)\\ &
    %    =\frac{\hat b^{\mu}}{\gamma\beta}\rmint_{0}^{\infty}J_{1}(k b)\arctan^2\Big(\frac{k}{2m}\Big):=\frac{\hat b^{\mu}}{\gamma\beta}I_{2}(m,b)
   % \end{split}
%\end{align}
%Note: we can simply take the massless limit as,
%\begin{align}
   % \begin{split}
      %   \Delta p_{2}^{\mu}(m=0)=\frac{\pi^2 \hat b^{\mu}}{4\gamma\beta}\frac{1}{b}.
   % \end{split}
%\end{align}   
%Derivative interactions also play a crucial role in the impulse.
%It comes from the interaction between the graviton and kinetic term of the scalar,
%\begin{align}
    %\begin{split}
     %   S_{int}=\rmint d^4x \,\tilde h^{\alpha\beta}\partial_{\alpha}\tilde \varphi \partial_{\beta}\tilde \varphi 
   % \end{split}
%\end{align}
%The vertex factor has the following form:
%\begin{align}
  %  \begin{split}
      %\mathcal{V}_{\alpha\beta}(k_1,k_2,k_3)=  -\rmint_{\{k_{i}\}}  \hat\delta\Big(\sum_{i=1}^{3}k_{i}\Big)   (k_2)_{\alpha} (k_3)_{\beta}
    %\end{split}
%\end{align}
%The contribution to the Eikonal phase has the following form,
%\begin{align}
  %  \begin{split}
       % \chi&\sim -\rmint_{\{k_{i}\}}  \underbrace{\hat\delta\Big(\sum_{i=1}^{3}k_{i}\Big)   (k_2)_{\alpha} (k_3)_{\beta}}_{\mathcal{V}_{\alpha\beta}}\,\hat\delta(k_1\cdot v_1)\hat\delta(k_2\cdot v_2)\hat\delta(k_3 \cdot v_2)v_1^{\mu}v_1^{\nu}\,\frac{P_{\mu\nu}^{\alpha\beta}}{k_1^2}\prod_{j=1}^{2}\frac{1}{k_j^2-m^2}e^{i k_1\cdot b_1}e^{i (k_2+k_3)\cdot b_2}\\ &
      %  =\rmint_{k_{1,2}}(k_2)_{\alpha} (k_1+k_2)_{\beta}\hat\delta(k_1\cdot v_1) \hat\delta(k_2\cdot v_2)\hat\delta(k_1\cdot v_2) v_1^{\mu}v_1^{\nu}\frac{\eta^{\alpha}_{\mu}\eta^{\beta}_{\nu}+\eta^{\beta}_{\mu}\eta^{\alpha\nu}-\eta^{\alpha\beta}\eta_{\mu\nu}}{k_1^2(k_2^2-m^2)[(k_1+k_2)^2-m^2]}e^{ik_1\cdot b}\\ &
%=\chi_{1}+\chi_{2}+\chi_{3}
  %  \end{split}
%\end{align}
%where,
%\begin{align}
    %\begin{split}
        %\chi_{1}&=\rmint_{k_{1,2}}(k_2\cdot v_1)^2\frac{\hat\delta(k_1\cdot v_1) \hat\delta(k_2\cdot v_2)\hat\delta(k_1\cdot v_2)}{k_1^2(k_2^2-m^2)[(k_1+k_2)^2-m^2]}e^{i k_1\cdot b}\\ &
       % =(v_1)_{\mu}(v_1)_{\nu}\rmint_{k_{1,2}}k_2^{\mu}\,k_2^{\nu}\frac{\hat\delta(\gamma k_1^0-\gamma\beta k_1^3)\hat\delta(\gamma k_2^0)\hat\delta(\gamma k_1^0)}{k_1^2(k_2^2-m^2)[(k_1+k_2)^2-m^2]}e^{i k_1 \cdot b}\\ &
   % =-\frac{(v_1)_{i}(v_1)_{j}    }{\gamma^{3}\beta}\rmint d^2 k_1\frac{e^{i k_1 \cdot b}}{\vec k_1^2}
   % \underbrace{\rmint d^3 k \frac{k_i k_j}{(\vec k^2+m^2)[(\vec k+\vec k_1)^2+m^2]}}_{\mathcal{C}_{ij}}
   % \end{split}
%\end{align}
%By using Passarino-Veltman reduction one can evaluate the integral %$\mathcal{C}_{ij}$ as follows,
%\begin{align}
    %\begin{split}
        %\mathcal{C}_{ij}=&\Big[\frac{3}{32\pi |\vec k_1|}\arctan\Big(\frac{|\vec k_1|}{2m}\Big)+\frac{m^2}{8 \pi |\vec k_1|^3}\arctan\Big(\frac{|\vec k_1|}{2m}\Big)-\frac{m}{16\pi|\vec k_1|^2}\Big]k_1^i k_1^j\\ &
     %  - \Big[\frac{|\vec %k_1|}{32\pi }+\frac{m^2}{8\pi |\vec k|}+\frac{m}{16\pi}\Big]\hat\delta^{ij}
  %  \end{split}
%\end{align}
%Therefore $\chi_1$ can be written as ,
%\begin{align}
  %  \begin{split}
    %    \chi_{1}=--\frac{(v_1)_{i}(v_1)_{j}    }{\gamma^{3}\beta}\rmint_{0}^{\infty} dl\,\frac{J_{0}(L|\boldsymbol{b}|)}{l}\Big[\frac{3}{32\pi\,l}\arctan\Big(\frac{l}{2m}\Big)\Big]
   % \end{split}
%\end{align}
%\newpage
%Another term coming from the Point particle action of the form $\rmint d\tau \tilde{h}_{\mu\nu} \dot{x}^{\mu} \dot{x}^{\nu} \tilde{\varphi}\tilde{\varphi}$ is given by:
%\begin{figure}
%\centering
%\scalebox{0.3}{
%\begin{feynman}
 %   \fermion[lineWidth=4]{4.00, 6.20}{8.60, 6.20}
  %  \electroweak[lineWidth=4, label=$\Tilde{\varphi}$, color=0693e3]{6.20, 6.20}{8.60, 4.00}
    %\electroweak[label=$\Tilde{\varphi}$, color=0693e3, lineWidth=4]{4.20, 4.00}{6.20, 6.20}
   % \fermion[lineWidth=4]{4.00, 4.00}{8.80, 4.00}
   % \gluon[flip=true, lineWidth=5, label=$\tilde h_{\mu\nu}$]{6.20, 6.20}{6.20, 4.00}
%\end{feynman}}
%\caption{}
%\end{figure}


%\begin{align}
   %\begin{split} 
%\chi=&\frac{g(s_2m_2)^2m_1}{m_p^6}\rmint_{k,k',k''}e^{i(k+k'+k'')\cdot b_1}v_1^{\mu} v_1^{\nu}\hat\delta((k+k'+k'')\cdot v_1)\frac{1}{k^2-m^2}\frac{1}{(k'')^2-m^2}\frac{P_{\mu\nu;\alpha\beta}}{k'^2}\\& v_2^\alpha v_2^\beta\hat\delta(k\cdot v_2)\hat\delta(k'\cdot v_2)\hat\delta(k''\cdot v_2)e^{-i(k+k'+k'')\cdot b_2}\\&
%=\frac{g(s_2m_2)^2m_1}{m_p^6}\Big\{2(v_1\cdot v_2)^2-\frac{1}{2}v_1^2v_2^2\Big\}\rmint_{k,k',k''}e^{i(k+k'+k'')\cdot b}\hat\delta((k+k'+k'')\cdot v_1)\frac{1}{k'^2}\frac{1}{k^2-m^2}\frac{1}{(k'')^2-m^2}\\&\hat\delta(k\cdot v_2)\hat\delta(k'\cdot v_2)\hat\delta(k''\cdot v_2)\\&
%=\frac{g(s_2m_2)^2m_1}{m_p^6}\Big\{2(v_1\cdot v_2)^2-\frac{1}{2}v_1^2v_2^2\Big\}\rmint_{k,k',k''}e^{i(k+k'+k'')\cdot b}\hat\delta((k+k'+k")\cdot v_1)\\&\frac{1}{k'^2}\frac{1}{k^2-m^2}\frac{1}{(k'')^2-m^2}\hat\delta(k^0)\hat\delta(k'^0)\hat\delta(k''^0)\\&
%=\frac{g(s_2m_2)^2m_1}{m_p^6}\frac{\Big\{2(v_1\cdot v_2)^2-\frac{1}{2}v_1^2v_2^2\Big\}}{\gamma \beta}\rmint_{\vec{k},\vec{k}',\vec{k''}}e^{i(\vec{k}+\vec{k'}+\vec{k''})\cdot \vec{b}}\hat\delta(k^{(3)}+k'^{(3)}+k''^{(3)})\frac{1}{k'^2}\frac{1}{k^2+m^2}\frac{1}{(k'')^2+m^2}
%\end{split}
%\end{align}
%Now to solve this, we %use the following trick:
%we know $\rmint_{0}^{\infty} dt e^{-ta^2}=\frac{1}{a^2}$\\
%So for every propagator we use the trick and we achieve an integral of the form :

%$$\rmint_{\vec{k'}} \frac{d^3k'}{(2\pi)^3} \rmint_{0}^{\infty} dt e^{-t{k'^{(1)}}^2}e^{-t{k'^{(2)}}^2}e^{-t{k'^{(3)}}^2}  e^{i{k'^{(1)}}\cdot b_1} e^{i{k'^{(2)}}\cdot b_2}e^{i{k'^{(3)}}\cdot b_3}\hat\delta((k^{(3)}+k'^{(3)}+k''^{(3)})) $$
%$$= \frac{1}{(2\pi)^2}K_0(b(k^{(3)}+k''^{(3)})e^{-ib_3(k^{(3)}+k''^{(3)})} $$

%Here we write $\sqrt{b_1^2+b_2^2}=b$ and $K_0$ is the modified bessel function. Now the exponential function cancels against the other part coming from the third component of the exponential in the initial integral. Finally we achieve again a integral of the form:
%$ \rmint_{-\infty}^{\infty} dk e^{-tk^2}e^{ik\cdot b_1}=e^{-\frac{b_1^2}{2t}}\sqrt{\frac{\pi}{t}}$.
%Hence,
%\begin{align}
%   \begin{split}
%\chi=\frac{g(s_2m_2)^2m_1}{m_p^6} \frac{\Big\{2(v_1\cdot v_2)^2-\frac{1}{2}v_1^2v_2^2\Big\}}{\gamma \beta}\frac{e^{-2m^2}}{2\pi^2}\rmint {dk \,dk''} K_0(b\sqrt 2\sqrt{k^2+(k'')^2+m^2}) K_0(b(k+k''))
 %   \end{split}
%\end{align}
%For computing the impulse we have to take the derivative of the connected Generating function with respect to the impact parameter.
%So,
%\begin{align}
   % \begin{split}
    %    \hat\delta p_2^{\mu}=\frac{\partial \chi}{\partial b_{\mu}}&=-\frac{\hat b^{\mu}}{\gamma\beta}\frac{g(s_2m_2)^2m_1}{m_p^6} \{\Big\{2(v_1\cdot v_2)^2-\frac{1}{2}v_1^2v_2^2\Big\}\}\frac{1}{2\pi^2}\frac{\partial}{\partial b}\rmint_{k,k''} K_0(b\sqrt 2\sqrt{k^2+(k'')^2+m^2}) K_0(b(k+k''))\\&
      %  =\frac{\hat b^{\mu}}{\gamma\beta}\frac{g(s_2m_2)^2m_1}{m_p^6} \{\Big\{2(v_1\cdot v_2)^2-\frac{1}{2}v_1^2v_2^2\Big\}\}\frac{1}{2\pi^2}\\&\rmint dk\, dk''\,\Bigg(\sqrt{2} \sqrt{k^2+\left(k''\right)^2+m^2} K_0\left(b \left(k+k''\right)\right) K_1\left(\sqrt{2} b \sqrt{k^2+\left(k''\right)^2+m^2}\right)+\\&\left(k''+k\right) K_0\left(\sqrt{2} b \sqrt{k^2+\left(k''\right)^2+m^2}\right) K_1\left(b \left(k+k''\right)\right)\Bigg)\\&
     %   =\frac{\hat b^{\mu}}{\gamma\beta}\frac{g(s_2m_2)^2m_1}{m_p^6} \frac{\Big\{2(v_1\cdot v_2)^2-\frac{1}{2}v_1^2v_2^2\Big\}}{\gamma \beta}\frac{1}{2\pi^2}\,\,\mathcal{I}(b)
%\end{split}
%\end{align}

%If we do series expansion of  the last integrand in b it has a fall of $\frac{1}{b}$ in leading order. The last integral over k and k'' is a one dimensional integral over only the spatial third components of the $\vec{k},\vec{k''}.$ \textcolor{black}{Is it correct?}\\
%\textbf{Another direct way  of doing the integral:}\par 
%Start withe the  relevant part of the integral,
%\begin{align}
    %\begin{split}
       %\chi&\sim\rmint_{\vec{k},\vec{k}',\vec{k''}}e^{i(\vec{k}+\vec{k'}+\vec{k''})\cdot \vec{b}}\hat\delta\Big(k^{(3)}+k'^{(3)}+k''^{(3)}\Big)\frac{1}{\vec k'^2}\frac{1}{\vec k^2+m^2}\frac{1}{\vec k''^2+m^2}\\ &
      % =\rmint_{\vec k, \vec k',\tilde k''}e^{i(\tilde{k}+\tilde{k'}+\tilde{k''})\cdot \tilde b}\frac{1}{\tilde k'^2+k'^{(3)2}}\frac{1}{\tilde k^2+\underbrace{k^{(3)2}+m^2}_{\sigma_1^2}}\frac{1}{\tilde k''^2+\underbrace{(k^3+k'^{3})^2+m^2}_{\sigma_2^2}}\\ &
    %   =\rmint d^3k \,d^3k' \frac{e^{i(\tilde k+\tilde k')\cdot b}}{(\tilde k'^2+k'^{(3)2})[\tilde k^2+\sigma_{1}^2(k^{(3)},m)]}\rmint d^2 \tilde k''\,\frac{e^{i \tilde k''\cdot b }}{\tilde k''^2+\sigma_2^2(k^{(3)},k'^{(3)},m)}\\ &
     %  =\rmint d^3k \,d^3k' \frac{e^{i(\tilde k+\tilde k')\cdot b}}{(\tilde k'^2+k'^{(3)2})[\tilde k^2+\sigma_{1}^2(k^{(3)},m)]}\underbrace{\rmint_{0}^{\infty}dl\, l\frac{J_{0}(L|\boldsymbol{b}|)}{l^2+\sigma_2^2}}_{K_{0}(b|\sigma_2|)}\\ &
      % =\rmint dk^{(3)}\,dk'^{(3)}\rmint d^2\tilde k\,\frac{e^{i\tilde k\cdot b}}{\tilde k^2+\sigma_1^2}\rmint d^2 \tilde k'\frac{e^{i\tilde k'\cdot b}}{\tilde k'^2+k'^{(3)2}}\,K_{0}(b|\sigma_2|)\\ &
     %  =\rmint _{-\infty}^{\infty}dk^{(3)}\,dk'^{(3)} K_{0}(b|k'^{(3)}|)\,K_{0}(b|\sigma_1|)\,K_{0}(b|\sigma_2|)
  %  \end{split}
%\end{align}
%here, $\tilde{()}$ denotes a 2-\textit{dim} vector lies in $x-y$ plane.
%Hence, the contribution to impulse is given by,
%\begin{align}
   % \begin{split}
     %   \hat\delta p^{\mu}&=\hat b^{\mu}\frac{\partial \chi}{\partial b}\\ &
      %  =\hat b^{\mu}\frac{g(s_2m_2)^2m_1}{m_p^6}\frac{4\gamma^2-1}{2\sqrt{\gamma^2-1}}\rmint_{-\infty}^{\infty}dx dz\,\left(-\sqrt{m^2+z^2}\right) K_0(b |x|) K_1\left(b \sqrt{m^2+z^2}\right) K_0\left(b \sqrt{m^2+(x+z)^2}\right)\\ &-K_0\left(b \sqrt{m^2+z^2}\right) \Big[|x| K_1(b |x|) K_0\left(b \sqrt{m^2+(x+z)^2}\right)+K_0(b |x|) \sqrt{m^2+(x+z)^2} K_1\left(b \sqrt{m^2+(x+z)^2}\right)\Big]
    %\end{split}
%\end{align}
%One can numerically solve the integral and which has well defines behaviour (\textcolor{black}{checked}!). One thing should be checked:
%\textcolor{black}{The two approaches gives numerically two different results but have same functional nature. May be they are just related by some numerical factor which we are missing. Need to be checked carefully!!}
%Now, we evaluate the diagrams that have world line propagators. Here, we can take the reverse unitarity techniques to evaluate the integral.
%@$\mathcal{O}(z)$ the first worldline vertex that contributes comes form the following part of the worldline action:
%\begin{align}
  %  \begin{split}
   %     S_{pm}=-\frac{m_as_a}{2m_p}\rmint_{k,\omega}e^{ik\cdot b}\hat\delta (k\cdot v+\omega)\tilde \varphi(k)z^{\rho}(\omega)\{2i\omega v_a^{(\mu}\hat\delta^{\nu)}_{\rho}+iv^{\mu}v^{\nu}k_{\rho} \}\eta_{\mu\nu}
 %   \end{split}
%\end{align}
%Therefore the vertex factor is given by,
%\begin{align}
 %   \begin{split}
  %      \mathcal{O}(z)\Big|_{1}: -\frac{m_as_a}{2m_p}\rmint_{k,\omega}e^{ik\cdot b}\hat \delta(k\cdot v+\omega)\{2\omega v_{\rho}+k_{\rho}\}
   % \end{split}
%\end{align}
%The contribution to the Eikonal phase has the following form:
%\begin{align}
    %\begin{split}
     %   \chi&=-\frac{(m_1 s_1)^2m_2s_2}{8m_p^3}\rmint  \check{\mathcal{D}}[\tilde h_{\mu\nu},\tilde \varphi,\{z_{i}\}]\rmint_{\{k_i,\textasciicaron{\omega}_j\}}\hat{\delta}(k_1\cdot v_2)\hat{\delta}(k_2\cdot v_2)\tilde \varphi(k_1)\tilde \varphi(k_2) e^{i(k_1+k_2)\cdot b_2}\\ &
     %   \hspace{0.8cm}\hat\delta(k_3\cdot v_1+\omega_1)\hat\delta(k_4\cdot v_1+\omega_2)\tilde \varphi(k_3)\tilde\varphi(k_4)z^{\rho}(\omega_1)z^{\sigma}(\omega_2)\{2\omega_1 (v_1)_{\rho}+(k_{3})_{\rho}\}\\ &
     %   \hspace{0.8cm}
      %  \{2\omega_2(v_1)_{\sigma}+(k_4)_{\sigma}\}\\&
     %   =-\frac{(m_1 s_1)^2m_2s_2}{8m_p^3}\rmint_{k_i,\omega_j}\hat{\delta}(k_1\cdot v_2)\hat{\delta}(k_2\cdot v_2)\hat\delta(k_3\cdot v_1+\omega_1)\hat\delta(k_4\cdot v_1+\omega_2)\langle \tilde \varphi(k_1)\tilde\varphi(k_3)\rangle\langle \tilde \varphi(k_2)\tilde\varphi(k_4)\rangle\\ &
     %   \hspace{0.8cm}
       % \langle z^{\rho}(\omega_1)z^{\sigma}(\omega_2) \rangle \{2\omega_1 (v_1)_{\rho}+(k_{3})_{\rho}\}  \{2\omega_2(v_1)_{\sigma}+(k_4)_{\sigma}\}\\ &
       % =\frac{(m_1 s_1)^2m_2s_2}{8m_1\,m_p^3}\rmint_{\{k_i,\omega_j\}}\hat{\delta}(k_1\cdot v_2)\hat{\delta}(k_2\cdot v_2)\hat\delta(k_3\cdot v_1+\omega_1)\hat\delta(k_4\cdot v_1+\omega_2)\delta(k_1+k_3)\delta(k_2+k_4)\delta(\omega_1+\omega_2)\\ &
     %  \hspace{0.8cm} \frac{1}{(k_1^2-m^2)(k_2^2-m^2)\omega_1^2}\{4\omega_1\omega_2+2\omega_1 v_1\cdot k_4+2 \omega_2 v_1\cdot k_3+k_3\cdot k_4\}e^{i(k_1+k_2)\cdot b_2}e^{i(k_3+k_4)\cdot b_1}\\ &
    %   =\frac{(m_1 s_1)^2m_2s_2}%{8m_1\,m_p^3}\rmint_{k,k_1,\omega_1}\frac{\hat\delta(k_1\cdot v_2)\hat\delta(k\cdot v_2)\hat\delta(k\cdot v_1)\hat\delta(\omega_1-k_1\cdot v_1)}{(k_1^2-m^2)[(k-k_1)^2-m^2]\omega_1^2}e^{ik \cdot b}\\&\hspace{0.8cm}[-4\omega_1^2+2\omega_1v_1\cdot(k_1-k)+2\omega_1v_1\cdot k_1+k_1\cdot (k-k_1)]
   % \end{split}
%\end{align}
%Integrating over $\delta(\omega_1-k_1\cdot v_1)$ and using the fact that the integral has support at $k\cdot v_1=0$, we left over with the following integral,
%\begin{align}
   % \begin{split}
    %    \chi=\frac{(m_1 s_1)^2m_2s_2}{8m_1\,m_p^3}\rmint_{k,k_1}\frac{\hat\delta(k_1\cdot v_2)\hat\delta(k\cdot v_2)\hat\delta(k\cdot v_1)}{(k_1^2-m^2)[(k-k_1)^2-m^2](k_1\cdot v_1)^2}[k_1\cdot(k-k_1)]e^{ik\cdot b}
  %  \end{split}
%\end{align}
%To evaluate this integral, we use the following,
%\begin{align}
   % \begin{split}
    %    \rmint dk_1^{(3)}[\cdots]=\oint_{H^{+}}[\cdots]-\oint_{H^{-}}[\cdots]=\frac{i\pi}{2}\sum_{k_1^{(3)}=k_1^{(3)*}\subset H^{+}}\text{Res}([\cdots]) -\sum_{k_1^{(3)}=k_1^{(3)*}\subset H^{-}}\text{Res}([\cdots]) 
   % \end{split}
%\end{align}
%One can see that the integrand has double pole at $k_1\cdot v_1+i\epsilon=0$. So, simply close the contour in the opposite side we get zero.\\


\section{Massive integral coming from derivative interaction} \label{ch2:app:E}
The integral of interest is the following,
\begin{align}
    \begin{split}
        J_{(2)}^{\mu\nu}&= \rmint_{q}\hat\delta(q\cdot v_1-k\cdot v_1)\hat\delta(q\cdot v_2)\frac{q^\mu q^\nu\,e^{-i q\cdot b}}{(q-k)^2 (q^2-m^2)}\,,\\ &
        =\rmint_0^1 dy\,e^{-iy\,k\cdot b}\rmint_q \hat\delta(q\cdot v_1-(1-y)k\cdot v_1)\hat\delta(q\cdot v_2+yk\cdot v_2)\frac{e^{-i q\cdot b}}{[q^2-(1-2y+y^2)m^2]^2}(q^{\mu}+yk^\mu)(q^{\nu}+yk^\nu)\,.\label{A.22 m}
    \end{split}
\end{align}
In \eqref{A.22 m}, the term proportional to $k^\mu k^\nu$ and $q^{\mu} k^\nu$ can be done using the same procedure as discussed in Appendix~(\ref{ch2:app:C}). We will concentrate on the term which is proportional to $q^\mu q^\nu$, which can be expanded in terms of basis vectors.
\begin{align}
    \begin{split}
        \mathcal{J}_{(2)}^{\mu\nu}(\sim q^\mu q^\nu)=\rmint_0^1 dy\, e^{-iy\,k\cdot b}[\lambda_{\eta}\eta^{\mu\nu}+\lambda_{bb}b^\mu b^\nu+\lambda_{1b}b^{(\mu}v_1^{\nu)}+\lambda_{2b}b^{(\mu}v_2^{\nu)}+\lambda_{11}v_1^\mu v_1^\nu+\lambda_{22}v_2^\mu v_2^\nu+\lambda_{12}v_1^{(\mu}v_2^{\nu)}]\,.\label{A.23 m}
    \end{split}
\end{align}
Now to extract the coefficients in \eqref{A.23 m} one needs to contract the LHS with the tensor structure of the RHS. The equations we need to solve are the following,

    \textbullet $\,\,$ Contracting with $b_\mu b_\nu$ gives:
    \begin{align}
       \rmint_q \hat\delta(q\cdot v_1-(1-y)k\cdot v_1)\hat\delta(q\cdot v_2+yk\cdot v_2)\frac{e^{-i q\cdot b}}{[q^2-(1-2y+y^2)m^2]^2} (q\cdot b)^2=\mathcal{X}_1=-\lambda_{\eta}|b|^2+\lambda_{bb}|b|^4\,.\label{A.24 m}
    \end{align}
    \textbullet $\,\,$ Contracting with $b_{(\mu}v_{1\nu)}$ gives:
    \begin{align}
    \begin{split}
        \rmint_q \hat\delta(q\cdot v_1-(1-y)k\cdot v_1)\hat\delta(q\cdot v_2+yk\cdot v_2)\frac{e^{-i q\cdot b}}{[q^2-(1-2y+y^2)m^2]^2} (q\cdot b)(q\cdot v_1)&=\mathcal{X}_2\\ &\hspace{0cm}=-\lambda_{1b}\frac{|b|^2}{2}-\lambda_{2b}\frac{\gamma |b|^2}{2}\,.
        \end{split}
    \end{align}
    \textbullet $\,\,$ Contracting with $b_{(\mu}v_{2\nu)}$ gives:
    \begin{align}
        \begin{split}
        \rmint_q \hat\delta(q\cdot v_1-(1-y)k\cdot v_1)\hat\delta(q\cdot v_2+yk\cdot v_2)\frac{e^{-i q\cdot b}}{[q^2-(1-2y+y^2)m^2]^2} (q\cdot b)(q\cdot v_2)&=\mathcal{X}_3\\ &=-\frac{|b|^2}{2}(\lambda_{2b}+\gamma \lambda_{1b})\,.
        \end{split}
    \end{align}
     \textbullet $\,\,$ Contracting with $v_{1\mu} v_{1\nu}$ gives:
    \begin{align}
        \begin{split}
             \rmint_q \hat\delta(q\cdot v_1-(1-y)k\cdot v_1)\hat\delta(q\cdot v_2+yk\cdot v_2)\frac{e^{-i q\cdot b}}{[q^2-(1-2y+y^2)m^2]^2}(q\cdot v_1)^2&=\mathcal{X}_4\\ &
           \hspace{0cm}  =\lambda_{\eta}+\lambda_{11}+\lambda_{22}\gamma^2+\lambda_{12}\gamma\,.
        \end{split}
    \end{align}
   \textbullet $\,\,$ Contracting with $v_{2\mu} v_{2\nu}$ gives:
    \begin{align}
        \begin{split}
              \rmint_q \hat\delta(q\cdot v_1-(1-y)k\cdot v_1)\hat\delta(q\cdot v_2+yk\cdot v_2)\frac{e^{-i q\cdot b}}{[q^2-(1-2y+y^2)m^2]^2}(q\cdot v_2)^2&=\mathcal{X}_5\\ &
              \hspace{0cm}=\lambda_\eta+\lambda_{22}+\lambda_{11} \gamma^2+\lambda_{12}\gamma\,.
        \end{split}
    \end{align}
    \textbullet $\,\,$ Contracting with $v_{1(\mu}v_{2\nu)}$ gives:
    \begin{align}
        \begin{split}
             \rmint_q \hat\delta(q\cdot v_1-(1-y)k\cdot v_1)\hat\delta(q\cdot v_2+yk\cdot v_2)\frac{e^{-i q\cdot b}}{[q^2-(1-2y+y^2)m^2]^2}(q\cdot v_1)(q\cdot v_2)&=\mathcal{X}_6\\ &
              \hspace{0cm}=\lambda_{\eta}\gamma+\gamma(\lambda_{11}+\lambda_{22})+\frac{\lambda_{12}}{2}(\gamma^2+1)\,.
        \end{split}
    \end{align}
      \textbullet $\,\,$ Contracting with $\eta_{\mu\nu}$ gives:
      \begin{align}
          \begin{split}
                \rmint_q \hat\delta(q\cdot v_1-(1-y)k\cdot v_1)\hat\delta(q\cdot v_2+yk\cdot v_2)\frac{e^{-i q\cdot b}}{[q^2-(1-2y+y^2)m^2]^2}q^2&=\mathcal{X}_7\\ &
                \hspace{0cm}=4\lambda_\eta-\lambda_{bb}|b|^2+\lambda_{11}+\lambda_{22}+\lambda_{12}\gamma\,.\label{A.30 m}
          \end{split}
      \end{align}
Solving \eqref{A.24 m} to \eqref{A.30 m} we will get,
\begin{align}
\lambda _{\eta }
&\to \frac{\gamma ^2 \mathcal{X}_{7}-2 \gamma  \mathcal{X}_{6}+\mathcal{X}_{4}+\mathcal{X}_{5}-\mathcal{X}_{7}}{\gamma ^2-1}+\frac{\mathcal{X}_{1}}{|b|^2},\notag\\
\lambda _{\text{bb}}
&\to \frac{2 \mathcal{X}_{1}}{|b|^4}+\frac{\gamma ^2 \mathcal{X}_{7}-2 \gamma  \mathcal{X}_{6}+\mathcal{X}_{4}+\mathcal{X}_{5}+\mathcal{X}_{7}}{|b|^2 \left(\gamma ^2-1\right)},\notag\\
\lambda _b
&\to -\frac{2 \left(\gamma  \mathcal{X}_{3}-\mathcal{X}_{2}\right)}{|b|^2 \left(\gamma ^2-1\right)},\qquad
\lambda _{2 b}\to -\frac{2 \left(\gamma  \mathcal{X}_{2}-\mathcal{X}_{3}\right)}{|b|^2 \left(\gamma ^2-1\right)},\notag\\
\lambda _{12}
&\to \frac{2 \left(-|b|^2 \left(\gamma ^3 \left(-\mathcal{X}_{7}\right)+3 \gamma ^2 \mathcal{X}_{6}-2 \gamma  \mathcal{X}_{4}-2 \gamma  \mathcal{X}_{5}+\gamma  \mathcal{X}_{7}+\mathcal{X}_{6}\right)+\gamma  \left(\gamma ^2-1\right) \mathcal{X}_{1}\right)}{|b|^2 \left(\gamma ^2-1\right)^2},\notag\\
\lambda _{11}
&\to \frac{\frac{\left(\gamma ^2-1\right) \mathcal{X}_{1}}{|b|^2}+\gamma ^2 \mathcal{X}_{7}+\left(\gamma ^2+1\right) \mathcal{X}_{5}-4 \gamma  \mathcal{X}_{6}+2 \mathcal{X}_{4}-\mathcal{X}_{7}}{\left(\gamma ^2-1\right)^2},\notag\\
\lambda _{22}
&\to \frac{\frac{\left(\gamma ^2-1\right) \mathcal{X}_{1}}{|b|^2}+\gamma ^2 \mathcal{X}_{7}+\left(\gamma ^2+1\right) \mathcal{X}_{4}-4 \gamma  \mathcal{X}_{6}+2 \mathcal{X}_{5}-\mathcal{X}_{7}}{\left(\gamma ^2-1\right)^2}\,.
\end{align}
Next, we list the values of the integrals $\mathcal{X}_i\,.$
\begin{align}
    \begin{split}
        \mathcal{X}_1&= -\lim_{\kappa\to 1}\frac{\partial^2}{\partial \kappa^2}   \rmint_q \hat\delta(q\cdot v_1-(1-y)k\cdot v_1)\hat\delta(q\cdot v_2+yk\cdot v_2)\frac{e^{-i\kappa q\cdot b}}{[q^2-(1-2y+y^2)m^2]^2}\,,\\ &
        =-\frac{\kappa |b|}{4\pi \gamma}\lim_{\kappa\to 1}\frac{\partial^2}{\partial \kappa^2}\frac{K_1(\kappa |b|\sqrt{\Delta^2+(1-2y+y^2)m^2})}{\sqrt{\Delta^2+(1-2y+y^2)m^2}}\,,\\ &
        =-\frac{|b|}{4\pi \gamma}\Big[\frac{[|b|^2 \left(\Delta ^2+m^2 (y-1)^2\right)+2]K_1\left(|b| \sqrt{m^2 (y-1)^2+\Delta ^2}\right)}{\sqrt{\Delta ^2+m^2 (y-1)^2}}-|b| K_2\left(|b| \sqrt{m^2 (y-1)^2+\Delta ^2}\right)\Big]\,.
    \end{split}
\end{align}
\begin{align}
    \begin{split}
        \mathcal{X}_2&=i(1-y)k\cdot v_1\lim_{\kappa\to 1}\frac{\partial}{\partial \kappa}\rmint_q \hat\delta(q\cdot v_1-(1-y)k\cdot v_1)\hat\delta(q\cdot v_2+yk\cdot v_2)\frac{e^{-i\kappa q\cdot b}}{[q^2-(1-2y+y^2)m^2]^2}\,,\\ &
        =-i(1-y)k\cdot v_1\frac{|b|^2 K_0\left(|b| \sqrt{m^2 (y-1)^2+\Delta ^2}\right)}{4 \pi  \gamma }\,.
    \end{split}
\end{align}
\begin{align}
    \begin{split}
        \mathcal{X}_3&=-iy\,k\cdot v_2\lim_{\kappa\to 1}\frac{\partial}{\partial \kappa}\rmint_q \hat\delta(q\cdot v_1-(1-y)k\cdot v_1)\hat\delta(q\cdot v_2+yk\cdot v_2)\frac{e^{-i \kappa\,q\cdot b}}{[q^2-(1-2y+y^2)m^2]^2}\,,\\ &
        =i y \,k\cdot v_2 \frac{|b|^2 K_0\left(|b| \sqrt{m^2 (y-1)^2+\Delta ^2}\right)}{4 \pi  \gamma }\,.
    \end{split}
\end{align}
\begin{align}
    \begin{split}
        \mathcal{X}_4&=(1-y)^2 (k\cdot v_1)^2\rmint_q \hat\delta(q\cdot v_1-(1-y)k\cdot v_1)\hat\delta(q\cdot v_2+yk\cdot v_2)\frac{e^{-i q\cdot b}}{[q^2-(1-2y+y^2)m^2]^2}\,,\\ &
        =(1-y)^2 (k\cdot v_1)^2\frac{|b|}{4\pi \gamma}\frac{K_1( |b|\sqrt{\Delta^2+(1-2y+y^2)m^2})}{\sqrt{\Delta^2+(1-2y+y^2)m^2}}\,.
    \end{split}
\end{align}
\begin{align}
    \begin{split}
         \mathcal{X}_5&=y^2 (k\cdot v_2)^2\rmint_q \hat\delta(q\cdot v_1-(1-y)k\cdot v_1)\hat\delta(q\cdot v_2+yk\cdot v_2)\frac{e^{-i q\cdot b}}{[q^2-(1-2y+y^2)m^2]^2}\,,\\ &
        =y^2 (k\cdot v_2)^2\frac{|b|}{4\pi \gamma}\frac{K_1( |b|\sqrt{\Delta^2+(1-2y+y^2)m^2})}{\sqrt{\Delta^2+(1-2y+y^2)m^2}}\,.
    \end{split}
\end{align}
\begin{align}
    \begin{split}
        \mathcal{X}_{6}&=-y(1-y)(k\cdot v_1)(k\cdot v_2)\rmint_q \hat\delta(q\cdot v_1-(1-y)k\cdot v_1)\hat\delta(q\cdot v_2+yk\cdot v_2)\frac{e^{-i q\cdot b}}{[q^2-(1-2y+y^2)m^2]^2}\,,\\ &
        =-y(1-y)(k\cdot v_1)(k\cdot v_2)\frac{|b|}{4\pi \gamma}\frac{K_1( |b|\sqrt{\Delta^2+(1-2y+y^2)m^2})}{\sqrt{\Delta^2+(1-2y+y^2)m^2}}\,.
    \end{split}
\end{align}
The integral of $\mathcal{X}_7$ is a little bit involved. One can do this integral in the following way,
\begin{align}
    \begin{split}
        \mathcal{X}_7&=\rmint_q \hat\delta(q\cdot v_1-(1-y)k\cdot v_1)\hat\delta(q\cdot v_2+yk\cdot v_2)\frac{e^{-i q\cdot b}}{[q^2-(1-2y+y^2)m^2]^2}\Big[q^2-(1-y)^2m^2+(1-y)^2 m^2\Big]\,.\label{A.38}
    \end{split}
\end{align}
The integral in \eqref{A.38} has two parts and can be evaluated separately,
\begin{align}
    \begin{split}
        \mathcal{X}_7^{(1)}&\equiv  \rmint_q \hat\delta(q\cdot v_1-(1-y)k\cdot v_1)\hat\delta(q\cdot v_2+yk\cdot v_2)\frac{e^{-i q\cdot b}}{[q^2-(1-2y+y^2)m^2]}\,,\\ &
        =\frac{1}{\gamma}\rmint_0^\infty \frac{dt}{4\pi}\,\frac{1}{t}\exp\Big(-\frac{|b|^2}{4t}-t\Delta^2-t(1-y)^2m^2\Big)\,,\\ &
        =\frac{1}{2\pi \gamma}K_{0}(|b|\sqrt{\Delta^2+(1-y)^2m^2})\,.
    \end{split}
\end{align}
and the second term takes the form,
\begin{align}
    \begin{split}
        \mathcal{X}_7^{(2)}&=(1-y)^2 m^2\,\rmint_q \hat\delta(q\cdot v_1-(1-y)k\cdot v_1)\hat\delta(q\cdot v_2+yk\cdot v_2)\frac{e^{-i q\cdot b}}{[q^2-(1-2y+y^2)m^2]^2}\,,\\ &
        =(1-y)^2 m^2\frac{|b|}{4\pi\gamma} \frac{K_1( |b|\sqrt{\Delta^2+(1-2y+y^2)m^2})}{\sqrt{\Delta^2+(1-2y+y^2)m^2}}\,.
    \end{split}
\end{align}



\section{Computation of worldline radiation diagram using large velocity approximation}\label{ch2:app:F}
We start with \eqref{6.4 mm} (omitting the overall constant, which is irrelevant to explain our point).
\begin{align}
    \begin{split}
        f_{\varphi}(x)&\propto\rmint_{\Omega}e^{-ik\cdot (x+b_1)}\rmint_{\omega,k_1}e^{-i k_1\cdot b}\hat{\delta}(k_1\cdot v_2)\hat{\delta}[\Omega\, n(m,\Omega)\cdot v_1-\omega]\hat\delta(k_1\cdot v_1-\omega)\frac{\Omega(n\cdot k_1)}{\omega^2 (k_1^2-m^2)}\,.\label{7.7}
    \end{split}
\end{align}
To do the integration over $\Omega$ in (\ref{7.7}) we first find the roots of the equation $f(\Omega):=\Omega\, n(m,\Omega)\cdot v_1-\omega=0$, which gives,
\begin{align}
    \begin{split}
       & \hat{\delta}[\Omega\, n(m,\Omega)\cdot v_1-\omega]=\frac{\hat\delta(\Omega-\Omega_1)}{|f'(\Omega_1)|}+\frac{\hat\delta(\Omega-\Omega_2)}{|f'(\Omega_2)|},\,\, \\ &
        \text{where,}\,\, \Omega_{1}(\omega)=\thesismathbox{\displaystyle\frac{\tilde\beta  \sqrt{\left(\tilde\beta ^2-1\right) \gamma ^2 m^2+\omega ^2}+\omega }{\left(1-\tilde\beta ^2\right) \gamma }},\\ &
        \text{if }\gamma  m<\omega \lor \frac{\omega }{\gamma }<m<\frac{\omega }{\sqrt{1-\tilde\beta ^2} \gamma },\\ &
        \text{and}\quad \Omega_2(\omega)=\thesismathbox{\displaystyle\frac{\tilde\beta  \sqrt{\left(\tilde\beta ^2-1\right) \gamma ^2 m^2+\omega ^2}-\omega }{\left(\tilde\beta ^2-1\right) \gamma }},\\ &
        \text{if }\frac{\omega }{\gamma }<m<\frac{\omega }{\sqrt{1-\tilde\beta ^2} \gamma }\,.
    \end{split}
\end{align}
As the observed frequency cannot be negative, only the $\Omega_1$ solution is relevant to us. Therefore, the integral in (\ref{7.7}) can be written as (with $\tilde{x}=x-b_1$),
\begin{align}
    \begin{split}
        f_{\varphi}(x)&\propto  \sum_{a=1}^2\rmint_{k_1}e^{-i \Omega_a\, n(m,\Omega_a)\cdot (x-b_1)}e^{-ik_1\cdot b}\hat\delta(k_1\cdot v_2)\frac{\Omega_a\,n(m,\Omega_a)\cdot k_1}{|f'(\Omega_a)|(k_1\cdot v_1+i\epsilon)^2(k_1^2-m^2)}\,.
    \end{split}
\end{align}
\textcolor{black}{The integral could be significantly simplified by imposing a suitable IR cutoff on the $\omega$ and $\Omega$ integrals in the regime $\frac{m^2}{\Omega^2}\ll 1$. The argument in favour of this statement is the following: $\frac{1}{m}$ is of cosmological scale, while $\Omega$ for a compact binary with an orbital period of roughly an hour is of the order of $10^{-3}$ Hz (for $1/m \sim 1$ Mpc, one has $m \sim 10^{-14}$ Hz)\,. One point to note is that placing a cutoff on the $\omega$ and $\Omega$ integrals automatically implies that there should also be a cutoff on the $k_1$ integral (coming from the condition $k_1\cdot v_1=\omega$). However, one could, in principle, argue that instead of imposing the cutoff directly on the $k_1$ integral, one should impose a suitable condition on $|v_1|$ such that $k_1\cdot v_1\gg m$, even if $|k_1|\rightarrow 0$. The quantity $\sqrt{(k_1\cdot v_1)^2-m^2\gamma^2(1-\tilde\beta^2)}$ coming from the solutions for $\Omega$ takes the following form in terms of the parametrization used here: $\sqrt{\gamma^2\beta^2 k_1^{(x)2}-m^2\gamma^2(1-\tilde\beta^2)}=\gamma\beta \sqrt{k_1^{(x)2}-\frac{m^2}{\gamma^2\beta^2}+m^2\sin^2\theta}$. Now we assume that we can place our detector at a particular direction on the celestial sphere, as shown in Fig.~(\ref{celestial}), such that $\sin^2\theta=\frac{1}{\gamma^2\beta^2}\,.$ For further simplification, we choose $\gamma\beta\rightarrow \infty$ so that $\sin \theta\rightarrow 0$.}
Hence, the integration that we have to perform is of the form as shown below,
\begin{align}
    \begin{split}
        f_{\varphi}(x)\sim \rmint_{\Omega \in| \Omega|_{\textrm{IR}}\gg m}d\Omega\rmint_{\omega \in |\omega|_{\textrm{IR}}\gg m}d\omega\rmint_{k_1,|k_1\cdot v_1|\gg m}\cdots\,.
    \end{split}
\end{align}

\begin{figure}
    \centering
    \includegraphics[scale=0.15]{celestial.jpeg}
    \caption{Sketch of the celestial sphere under consideration.}
    \label{celestial}
\end{figure}
Finally we get,
\begin{align}
    \begin{split}
        f_{\varphi}(x)\propto\frac{n^{\mu}}{(n\cdot v_1)^2}\underbrace{\rmint_{k_1}e^{-i k_1\cdot w_1}\hat{\delta}(k_1\cdot v_2)\frac{k_{1\mu}}{(k_1\cdot v_1+i\epsilon)(k_1^2-m^2)}}_{\mathcal{K}_{\mu}^{(2)}},\,\, w_1\equiv\frac{n\cdot (x-b_1)}{n\cdot v_1}v_1+b
    \end{split}
\end{align}  
where,
\begin{align}
    \begin{split}
        \mathcal{K}_{\mu}^{(2)}\to-\rmint_{-\infty}^{\infty}d\tau\,\theta(\tau)\frac{(\vec w_1-\tau \vec v_1)_{i}}{|\vec w_1-\tau \vec v_1|^3}\,[1+m |\vec w_1-\tau \vec v_1|]\,\exp(-m |\vec w_1-\tau \vec v_1|)\,.
    \end{split}
\end{align}
Then we proceed as follows:
\begin{align}
    \begin{split}
        n\cdot \mathcal{K}^{(2)}&\rightarrow -\rmint d\tau\,\theta(\tau) \frac{\vec n\cdot (\vec w_1-\tau \vec v_1)}{|\vec w_1-\tau v_1|^3}[1+m |\vec w_1-\tau v_1|]\,\exp(-m |\vec w_1-\tau \vec v_1|)\,,\\ &
        =\rmint_{-\infty}^{\infty} d\tau\, \theta(\tau)\frac{(u_1-\tau)\gamma\tilde\beta+\chi}{(\gamma^2-1)^{3/2}\Big[\tau^2-u_1^2-2\tau u_1+\frac{b^2}{\gamma^2-1}\Big]^{3/2}}\Bigg[1+m(\gamma^2-1)^{1/2}\sqrt{\tau^2-u_1^2-2\tau u_1+\frac{|b|^2}{\gamma^2-1}}\Bigg]\\ &
        \hspace{2 cm}\exp\Big[-m(\gamma^2-1)^{1/2}\sqrt{\tau^2-u_1^2-2\tau u_1+\frac{|b|^2}{\gamma^2-1}}\Big]\,.
    \end{split}
\end{align}
In the limit of $\sin\theta\rightarrow 0$, the integral has a closed-form and smooth massless limit.
\begin{equation}\label{ch2:e:collinear-waveform}\begin{split}
f_{\varphi}(x)\propto \frac{\gamma}{(n\cdot v_1)^2(\gamma^2-1)} \frac{ \exp\Big({-m\sqrt{\gamma ^2-1}  \sqrt{\frac{|b|^2}{\gamma ^2-1}-u_1^2}}\,\Big)}{\sqrt{\frac{|b|^2}{\gamma ^2-1}-u_1^2}}\,.
        \end{split}
\end{equation}
\section{Proof of the identity in \texorpdfstring{\eqref{3.54kk}}{Eq. 3.54kk}}\label{3.54kkk}
The claim we made in the main text is,
\begin{align}
     \int_0^1 ds\frac{K_0(2m|b|)-K_0\left(\frac{2m|b|}{s}\right)}{1-s^2}  =\frac{1}{2} (\log (4m |b|)+\gamma_{E} ) K_0(2 m|b|)-\partial_\alpha \textrm{Ki}_{\alpha}(2m|b|)|_{\alpha\to0};
\end{align}
\textbf{Proof:} We have,
\begin{align}
    i(a)= \int_0^1 ds\frac{K_0(a)-K_0\left(\frac{a}{s}\right)}{1-s^2}\,.
\end{align}
The idea is to first take the Mellin transformation in $a$, which gives,
\begin{align}
\begin{split}
    \mathcal{M}_{a}(p)&=\int_0^\infty da a^{p-1} \int_0^1 ds\frac{K_0(a)-K_0\left(\frac{a}{s}\right)}{1-s^2}\,,\\ &
    =2^{p-2}\Gamma^2\left(\frac{p}{2}\right)\int_0^1 ds \frac{1-s^p}{1-s^2}\,,\\ &
    =2^{p-3} \Gamma \left(\frac{p}{2}\right)^2 \left(H_{\frac{p-1}{2}}+\log (4)\right)\,.
    \end{split}    
\end{align}
Now consider the following integral,
\begin{align}
    \begin{split}
        i'(a)=\frac{1}{2}(\log(2a)+\gamma_E)K_0(a)+\int_0^\infty dt e^{-a \cosh t}\log(\cosh t)\,.
    \end{split}
\end{align}
The Mellin transformation gives,
\begin{align}
    \begin{split}
        \mathcal{M}'_a(p)&=\int_0^\infty da
        \,a^{p-1} i'(a)\,,\\ &
        =\frac{1}{4} \Gamma \left(\frac{p}{2}\right)^2 \left(\psi ^{(0)}\left(\frac{p}{2}\right)+\gamma_E +\log (2)\right)+\Gamma(p){\int_0^\infty  dt \cosh^{-p}t\, \log(\cosh t)}\,,\\ &
        =\frac{1}{4} \Gamma \left(\frac{p}{2}\right)^2 \left(\psi ^{(0)}\left(\frac{p}{2}\right)+\gamma_E +\log (2)\right)-\Gamma(p)\partial_p\underbrace{\int_0^\infty dt \cosh ^{-p}t}_{\frac{\sqrt{\pi } \Gamma \left(\frac{p}{2}\right)}{2 \Gamma \left(\frac{p+1}{2}\right)}}\,,\\ &
        =2^{p-3} \Gamma \left(\frac{p}{2}\right)^2 \left(\psi ^{(0)}\left(\frac{p}{2}\right)+\gamma +\log (4)\right)-\Gamma(p)\frac{\sqrt{\pi } \Gamma \left(\frac{p}{2}\right) \left(\psi ^{(0)}\left(\frac{p}{2}\right)-\psi ^{(0)}\left(\frac{p+1}{2}\right)\right)}{4 \Gamma \left(\frac{p+1}{2}\right)}\,,\\ &
        =2^{p-3} \Gamma \left(\frac{p}{2}\right)^2 \left(H_{\frac{p-1}{2}}+\log (4)\right)=\mathcal{M}_a(p)\,.
        \end{split}
\end{align}
Now, from the Mellin uniqueness theorem \cite{Butzer1997Mellin}, one can assert that $i(a)=i'(a)$.
Moreover, $\int_0^\infty dt e^{-a \cosh t}\log(\cosh t)$ can be written as derivative of Bickley–Naylor function $\textrm{Ki}_n$ as,
\begin{align}
    \textrm{Ki}_n(x)=\int_0^\infty dt e^{-x \cosh t} \textrm{sech}^nt\implies  -\partial_n\textrm{Ki}_n(x)|_{n=0}=\int_0^\infty dt e^{-a \cosh t}\log(\cosh t),
\end{align}
which proves our claim. One can also numerically cross-check.
\section{Comment on the integral convergence}\label{App:3.H}
The integral we would like to analyze is
\begin{align}
I^{\mu}(|b|)
=
\int_{k_i,q}
\frac{
k_1^\mu\,
\hat\delta(k_1\!\cdot\! v_1)\,
\hat\delta(k_2\!\cdot\! v_2)\,
\hat\delta(q\!\cdot\! v_2)\,
\hat\delta(q\!\cdot\! v_1)
}{
(k_1^2-m^2)(k_2^2-m^2)\bigl[(k_1+k_2)^2-m^2\bigr]\bigl[(q-k_1)^2-m^2\bigr]
}
\,e^{iq\cdot b}.
\end{align}
Our goal is to reduce this expression to a form where the integrations can be performed in a systematic way. The delta functions already impose strong kinematical constraints, and it is therefore useful to exploit them from the outset by choosing variables adapted to the geometry of the problem.
To this end, we first introduce the spacelike unit vector
\begin{align}
u_1^\mu
=
\frac{\gamma v_1^\mu-v_2^\mu}{\zeta},
\qquad
\zeta=\sqrt{\gamma^2-1}\,,
\qquad
u_1^2=-1\,,
\qquad
u_1\cdot v_1=0\,.
\end{align}
In the frame we are using, this vector takes the explicit form
\begin{align}
u_1^\mu=(\zeta,\gamma,0,0)\,.
\end{align}
The importance of $u_1^\mu$ is that, together with the directions transverse to the scattering plane, it provides a natural basis for decomposing momenta constrained by the delta functions.

Indeed, the condition $\hat\delta(k_1\!\cdot\! v_1)$ implies that $k_1^\mu$ lies entirely in the three-dimensional hyperplane orthogonal to $v_1^\mu$. It is therefore natural to write
\begin{align}
k_1^\mu=a\,u_1^\mu+\ell^\mu,
\qquad
\ell\cdot v_1=0,
\qquad
\ell\cdot v_2=0,
\end{align}
where $\ell^\mu$ is purely transverse. Similarly, the pair of delta functions
\begin{align}
\hat\delta(q\!\cdot\! v_1)\hat\delta(q\!\cdot\! v_2)
\end{align}
forces the exchanged momentum $q^\mu$ to be purely transverse as well, so that
\begin{align}
q^\mu=q_\perp^\mu.
\end{align}
Thus, both $\ell^\mu$ and $q^\mu$ live in the transverse two-dimensional subspace, which will later allow us to trade the angular integrations for Bessel functions. With this parametrization, the relevant propagators simplify considerably. Writing
\begin{align}
\ell^\mu=(0,0,\boldsymbol{\ell}_\perp),
\qquad
q^\mu=(0,0,\boldsymbol{q}_\perp),
\end{align}
we find
\begin{align}
k_1^2-m^2
&=-(a^2+\boldsymbol{\ell}_\perp^2+m^2),\\
(q-k_1)^2-m^2
&=-\bigl(\boldsymbol{q}_\perp-\boldsymbol{\ell}_\perp\bigr)^2-a^2-m^2.
\end{align}
In this form, the $q$-dependence is entirely transverse, and the Fourier transform with respect to the impact parameter can be performed in the standard way. One obtains
\begin{align}
\int e^{iq\cdot b}\,
\frac{\hat\delta(q\!\cdot\! v_1)\hat\delta(q\!\cdot\! v_2)}
{(q-k_1)^2-m^2}
\propto
-\frac{1}{\sqrt{\gamma^2-1}}\,
e^{-i\boldsymbol{\ell}_\perp\cdot\boldsymbol{b}}\,
K_0\!\left(|\boldsymbol{b}|\sqrt{a^2+m^2}\right),
\end{align}
where the modified Bessel function $K_0$ arises from the two-dimensional massive Fourier transform. At this stage, the original integral is reduced to a mixed integral over the longitudinal variable $a$ and the transverse momentum $\ell$. The next nontrivial ingredient is the bubble subintegral involving $k_2$,
\begin{align}
\mathcal{B}(a,\ell)
=
\int
\frac{\hat\delta(k_2\!\cdot\! v_2)}
{(k_2^2-m^2)\bigl[(k_1+k_2)^2-m^2\bigr]}\,.
\end{align}
Since $v_2^\mu=(1,0,0,0)$, the delta function sets $k_2^0=0$, so it is convenient to parametrize
\begin{align}
k_2^\mu=(0,c,\boldsymbol{k}).
\end{align}
The bubble integral then becomes
\begin{align}
\begin{split}
\mathcal{B}(a,\ell)
&=
\int
\frac{dc\,d^2\boldsymbol{k}}{(2\pi)^3}
\frac{1}{
(c^2+\boldsymbol{k}^2+m^2)
\left[(a\gamma+c)^2+(\boldsymbol{\ell}_\perp+\boldsymbol{k})^2+m^2-a^2\zeta^2\right]
}.
\end{split}
\end{align}
This is now a standard Euclidean-type bubble integral in three dimensions. Introducing a Feynman parameter $s$, combining denominators, and performing the Gaussian integrations gives
\begin{align}
\mathcal{B}(a,\ell)
=
\frac{1}{8\pi}\int_0^1 ds\,
\frac{1}{
\sqrt{
m^2+s(1-s)\boldsymbol{\ell}_\perp^2+(1-s)(s\gamma^2-\zeta^2)a^2+i0
}
}.
\end{align}
The parameter integral may be further evaluated in closed form, yielding
\begin{align}
\mathcal{B}(a,\ell)
=
\frac{1}{32\pi^2\sqrt{a^2\gamma^2+\boldsymbol{\ell}_\perp^2}}
\arctan\!\left(
\frac{\sqrt{a^2\gamma^2+\boldsymbol{\ell}_\perp^2}}
{m+\sqrt{m^2-a^2\zeta^2}}
\right).
\end{align}
For later convenience, it is useful to rewrite the same result in the equivalent representation
\begin{align}
\mathcal{B}(a,\ell)
=
\frac{1}{32\pi^2}\,
M_a
\int_0^1 \frac{ds}{s^2}\,
\frac{1}{\dfrac{M_a^2}{s^2}+P_\ell(a)^2},
\end{align}
where we have defined
\begin{align}
M_a=m+\sqrt{m^2-a^2\zeta^2},
\qquad
P_\ell(a)=\sqrt{a^2\gamma^2+\boldsymbol{\ell}_\perp^2}.
\end{align}
This form is particularly convenient because it separates the $\ell$-dependence in a way that combines naturally with the remaining propagator.

Substituting these ingredients back into the original expression, we are left with
\begin{align}
I^\mu
\propto
\int_{-\infty}^{\infty} da
\int d^2\boldsymbol{\ell}_\perp\,
e^{-i\boldsymbol{\ell}_\perp\cdot\boldsymbol{b}}
\frac{a\,u_1^\mu+\ell^\mu}{\boldsymbol{\ell}_\perp^2+a^2+m^2}\,
\mathcal{B}(a,\ell)\,
K_0\!\left(|\boldsymbol{b}|\sqrt{a^2+m^2}\right).
\end{align}
Now the decomposition of $k_1^\mu$ becomes especially useful. The term proportional to $a\,u_1^\mu$ is odd under $a\to -a$, while the remaining part of the integrand is even in $a$. Consequently, that term integrates to zero. Only the transverse piece survives, and we obtain
\begin{align}
\begin{split}
I^\mu
&\propto
\frac{1}{\sqrt{\gamma^2-1}}
\int_{-\infty}^{\infty} da
\int d^2\boldsymbol{\ell}_\perp\,
e^{-i\boldsymbol{\ell}_\perp\cdot\boldsymbol{b}}
\frac{\ell^\mu}{\boldsymbol{\ell}_\perp^2+a^2+m^2}\,
\mathcal{B}(a,\ell)\,
K_0\!\left(|\boldsymbol{b}|\sqrt{a^2+m^2}\right).
\end{split}
\end{align}
Rotational invariance in the transverse plane then implies that the vector structure must be proportional to $b^\mu$. Using this fact, together with the representation of $\mathcal{B}(a,\ell)$ above, we can write
\begin{align}
\begin{split}
I^\mu
&=
-\frac{i\,b^\mu}{\sqrt{\gamma^2-1}|\boldsymbol{b}|}\partial_{|\boldsymbol{b}|}
\int_0^1\frac{ds}{s^2}
\int_{-\infty}^{\infty} da\,M_a\,
K_0\!\left(|\boldsymbol{b}|\sqrt{a^2+m^2}\right)
\\
&\hspace{2.2cm}\times
\int d^2\boldsymbol{\ell}_\perp\,
e^{-i\boldsymbol{\ell}_\perp\cdot\boldsymbol{b}}
\frac{1}{
(\boldsymbol{\ell}_\perp^2+a^2+m^2)
\left(\boldsymbol{\ell}_\perp^2+a^2\gamma^2+\dfrac{M_a^2}{s^2}\right)}\propto b^\mu \partial_{|\boldsymbol{b}|
}I(b).
\end{split}
\end{align}
The remaining transverse integral is radial. Passing to polar coordinates in the transverse plane gives
\begin{align}
\begin{split}
I
&\propto
-\frac{1}{\sqrt{\gamma^2-1}}
\int_0^1\frac{ds}{s^2}
\int_{-\infty}^{\infty} da\,M_a\,
K_0\!\left(|\boldsymbol{b}|\sqrt{a^2+m^2}\right)
\int d\ell\,\ell\,
\frac{J_0(\ell|\boldsymbol{b}|)}
{(\ell^2+\Lambda_a^2)(\ell^2+\Lambda_{a,s}^2)},
\end{split}
\end{align}
where we have introduced
\begin{align}
\Lambda_a=\sqrt{a^2+m^2}\,,
\qquad
\Lambda_{a,s}=\sqrt{a^2\gamma^2+\frac{M_a^2}{s^2}}\,.
\end{align}
The appearance of the Bessel function $J_0$ is a standard consequence of the angular integration in two dimensions. The radial integral can then be evaluated using known Fourier--Bessel identities, leading to
\begin{align}
\begin{split}
I
&\propto
-2\,\frac{1}{\sqrt{\gamma^2-1}}
\int_0^1\frac{ds}{s^2}
\int_0^\infty da\,M_a\,
\frac{
K_0^2(\Lambda_a|\boldsymbol{b}|)
-
K_0(\Lambda_{a,s}|\boldsymbol{b}|)\,K_0(\Lambda_a|\boldsymbol{b}|)
}{
\Lambda_{a,s}^2-\Lambda_a^2
}.
\end{split}
\end{align}
This is already a rather compact representation: all transverse momentum integrations have been carried out, and the answer is expressed in terms of modified Bessel functions.

Finally, it is convenient to map the semi-infinite $a$-integration domain to the finite interval $u\in[0,1]$ through the change of variables
\begin{align}
a=m\,\frac{\sqrt{1-u^2}}{u}.
\end{align}
After this substitution, the integral becomes
\begin{align}
\begin{split}
I
\propto
\int_0^1
\frac{du}{u^2\sqrt{1-u^2}}
\int_0^1 ds\,
\frac{1+\rho(u)}{(1+\rho(u))^2-s^2\rho(u)^2}
\Biggl[
K_0^2\!\left(\frac{m|\boldsymbol{b}|}{u}\right)
-
K_0\!\bigl(m|\boldsymbol{b}|\,L(u,s)\bigr)
K_0\!\left(\frac{m|\boldsymbol{b}|}{u}\right)
\Biggr],
\end{split}
\end{align}
where
\begin{align}
\rho(u)=\frac{\sqrt{\gamma^2u^2-\zeta^2+i0}}{u},
\qquad
L(u,s)=
\sqrt{
\frac{\gamma^2(1-u^2)}{u^2}
+
\frac{(1+\rho(u))^2}{s^2}
}.
\end{align}
This final form is particularly useful for further analytical study and numerical evaluation, since both integrations now run over finite ranges and the entire tensor structure has been reduced to the expected overall factor $b^\mu$. Although the integral is convergent, the two-loop structure allows the propagator $(k_1+k_2)^2-m^2$ to go on shell. Consequently, the integral acquires both real and imaginary parts. The conservative part of the impulse is extracted from the real part alone. The imaginary part contributes to the radiative dynamics. The dependency on the impact parameter is shown in Fig~\eqref{fig:3.2}. In a similar fashion other two loop asymmetric integrals can be done.
\begin{figure}
    \centering
    \includegraphics[width=0.5\linewidth]{I_b.pdf}
    \caption{Figure showing dependency of the integral $I(b)$ as a function of $|b|$ for $m=1\,.$}
    \label{fig:3.2}
\end{figure}
\section{Comments on the algebraic properties of the Bessel curve}
\paragraph{One-loop integrals, the Bessel curve, and topological recursion.}
The massive one-loop integrals encountered in our analysis have the schematic form
\begin{align}
    I_{\mathrm{1\mbox{-}loop}}
    \sim
    \sum_r \int_0^1 ds\, \rho_r(s)\, K_0\!\big(\alpha_r(s)\,b\big),
\end{align}
where \(\rho_r(s)\) is a rational function of \(s\). Thus the special-function content is entirely carried by the modified Bessel kernel \(K_0\), which may be viewed as the wavefunction associated with the quantization of the Bessel curve.

We define the classical Bessel curve as the rational spectral curve \cite{DoNorbury2018},
\begin{align}
    x(z)=\frac{z^2}{2},
    \qquad
    y(z)=\frac{1}{z},
    \qquad\Longrightarrow\qquad
    2xy^2-1=0.
\end{align}
We quantize this curve by promoting the classical variables \(x,y\) to non-commuting operators
\begin{align}
    \hat x=x,
    \qquad
    \hat y=\hbar\,\partial_x,
    \qquad
    [\hat x,\hat y]=\hbar.
\end{align}
With the symmetric ordering prescription, the quantum curve associated with
\(
H(x,y)=2xy^2-1=0
\)
is
\begin{align}
    \hat H\,\psi=(2\hat y\,\hat x\,\hat y-1)\psi=0.
\end{align}
In the \(x\)-polarization this becomes
\begin{align}
    \left(2(\hbar\partial_x)\,x\,(\hbar\partial_x)-1\right)\psi(x)=0.
\end{align}
Rewriting the equation in terms of the spectral parameter \(z\), with \(x=z^2/2\), one finds
\begin{align}
    \hbar^2 z^2\psi_0''(z)+\hbar^2 z\psi_0'(z)-z^2\psi_0(z)=0.
\end{align}
This is precisely the modified Bessel equation of order zero, and its decaying solution is
\begin{align}
    \psi_0(z)\propto K_0\!\left(\frac{z}{\hbar}\right).
\end{align}

The same wavefunction can be reconstructed from the Eynard--Orantin topological recursion on the Bessel curve. In the quantum-curve/WKB formalism one writes
\begin{align}
    \psi_0(z,\hbar)
    =
    \exp\!\left(
        \frac{1}{\hbar}S_0(z)+S_1(z)+\hbar S_2(z)+\hbar^2 S_3(z)+\cdots
    \right),
\end{align}
where
\begin{align}
    S_0(z)=\int^z y(z')\,dx(z'),
    \qquad
    S_1(z)=-\frac{1}{2}\log\frac{dx}{dz},
\end{align}
and, for \(k\ge 2\),
\begin{align}
    S_k(z)
    =
    \sum_{2g-1+n=k}
    \frac{(-1)^n}{n!}
    \int_{\infty}^{z}\cdots\int_{\infty}^{z}
    \omega_{g,n}(z_1,\ldots,z_n).
\end{align}
Here \(\omega_{g,n}\) are the symmetric meromorphic multidifferentials on \((\mathbb{CP}^1)^n\) produced recursively by the topological recursion. For the Bessel curve one finds, on the decaying branch,
\begin{align}
    S_0(z)=-z,
    \qquad
    S_1(z)=-\frac{1}{2}\log z,
    \qquad
    S_2(z)=-\frac{1}{8z},
    \qquad
    S_3(z)=\frac{1}{16z^2},
\end{align}
up to additive constants that only affect the overall normalization. Therefore
\begin{align}
    \psi_0(z,\hbar)
    \sim
    \exp\!\left(
        -\frac{z}{\hbar}
        -\frac{1}{2}\log z
        -\hbar\frac{1}{8z}
        +\hbar^2\frac{1}{16z^2}
        +\cdots
    \right).
\end{align}
Equivalently,
\begin{align}
    \psi_0(z,\hbar)
    \sim
    z^{-1/2}e^{-z/\hbar}
    \left(
        1-\frac{\hbar}{8z}
        +\frac{9\hbar^2}{128z^2}
        +O\!\left(\frac{\hbar^3}{z^3}\right)
    \right),
\end{align}
which matches the standard large-argument asymptotic expansion of the modified Bessel function,
\begin{align}
    K_0\!\left(\frac{z}{\hbar}\right)
    \sim
    \sqrt{\frac{\pi\hbar}{2z}}\,
    e^{-z/\hbar}
    \left(
        1-\frac{\hbar}{8z}
        +\frac{9\hbar^2}{128z^2}
        +O\!\left(\frac{\hbar^3}{z^3}\right)
    \right).
\end{align}
Thus the asymptotic expansion of the Bessel kernel may be recovered directly from the topological-recursion data of the Bessel curve.

In our one-loop integrals, the kernel appears in the combination
\begin{align}
    K_0\!\big(\alpha_r(s)\,b\big),
\end{align}
so each term may be interpreted as the Bessel-curve wavefunction evaluated at the parameter-dependent point
\begin{align}
    z=\alpha_r(s)\,b,
    \qquad
    x=\frac{z^2}{2}=\frac{\alpha_r(s)^2b^2}{2}.
\end{align}
Hence the one-loop amplitude is naturally viewed as an integral over parameter space of wavefunctions of the quantized Bessel curve.

At this point it is important to identify the correct semiclassical parameter. The integration variable \(s\) should not itself be identified with \(\hbar\). Rather, the natural semiclassical regime of the full integral is the large-impact-parameter limit \(b\to\infty\) with \(m\) fixed. In this regime one may introduce the effective small parameter
\begin{align}
    \hbar_{\mathrm{eff}}:=\frac{1}{mb},
\end{align}
so that, schematically,
\begin{align}
    K_0\!\big(\alpha_r(s)\,b\big)
    =
    K_0\!\left(\frac{\alpha_r(s)}{m}\,\frac{1}{\hbar_{\mathrm{eff}}}\right).
\end{align}
Thus the one-loop integral is semiclassical in the limit \(\hbar_{\mathrm{eff}}\to0\), namely for \(mb\gg1\). In particular, for a kernel of the form
\begin{align}
    K_0\!\left(\frac{2mb}{s}\right),
\end{align}
the argument is large throughout the large-\(b\) regime and becomes even larger near the endpoint \(s\to0\). Therefore the WKB expansion of the Bessel wavefunction yields a controlled large-\(b\) asymptotic expansion of the integral:
\begin{align}
    K_0\!\left(\frac{2mb}{s}\right)
    \sim
    \sqrt{\frac{\pi s}{4mb}}\,
    \exp\!\left(-\frac{2mb}{s}\right)
    \left(
        1-\frac{s}{16mb}
        +\frac{9s^2}{512m^2b^2}
        +O\!\left(\frac{s^3}{(mb)^3}\right)
    \right).
\end{align}
This makes clear that the role of the Bessel curve is to supply the local semiclassical structure of the kernel, while the large-impact-parameter limit controls the asymptotics of the full one-loop integral.\\
